% Generated by scripts/build_guide.py from reviewed-section paragraphs.
\def\EditionMode{3}
\def\EditionLabel{A guide for the reader}
\def\EditionStatus{Complete guide to the checked typescript and appendix}
\documentclass[10pt,oneside]{book}
\usepackage{fontspec}
\setmainfont{Latin Modern Roman}
\setsansfont{Latin Modern Sans}
\usepackage{amsmath,amssymb,mathtools}
\usepackage{unicode-math}
\setmathfont{Latin Modern Math}
\usepackage{xcolor,tikz-cd,lua-ul,marginnote,pdflscape}
\usetikzlibrary{arrows.meta,decorations.pathmorphing,calc,shapes.geometric}
\tikzcdset{ampersand replacement=\&}
\tikzcdset{every label/.append style={text=black}}
\usepackage[a4paper,left=18mm,right=45mm,top=18mm,bottom=22mm,marginparwidth=28mm,marginparsep=7mm]{geometry}
\usepackage[unicode,colorlinks=true,urlcolor=black,linkcolor=black]{hyperref}
\usepackage{bookmark}
\setlength{\parindent}{0pt}
\setlength{\parskip}{4pt plus 1pt minus 1pt}
\setlength{\emergencystretch}{1.5em}
\raggedbottom
% Include final-line descenders inside the declared text-height rectangle.
\maxdepth=0pt
\definecolor{PSHand}{RGB}{25,58,132}
\definecolor{PSDelete}{RGB}{146,37,37}
\definecolor{PSTypedOld}{RGB}{95,95,95}
\ifnum\EditionMode=2\colorlet{blue}{black}\else\colorlet{blue}{PSHand}\fi
\RenewDocumentCommand{\textcircled}{m}{\tikz[baseline=(PSCircle.base)]\node[draw,circle,inner sep=.5pt,line width=.4pt](PSCircle){#1};}
\newif\ifPSInApparatus
\newcounter{psnote}
\newcounter{psvariant}
\newcounter{pslinkmark}
\newif\ifPSInLink
\newsavebox{\PSNoteBox}
\newdimen\PSNoteLimit
{\footnotesize\global\PSNoteLimit=3\baselineskip}
\newcommand{\CurrentScanURL}{}
\newcommand{\CurrentKey}{}
\newcommand{\CurrentScan}{}
\newcommand{\pssource}[3]{\psflushmargins\clearpage\gdef\CurrentKey{#1}\gdef\CurrentScan{#2}\gdef\CurrentScanURL{#3}\setcounter{psvariant}{0}\typeout{PS-SOURCE: #1}}
\newcommand{\psdivision}[2]{\pdfbookmark[0]{#1}{#2}\par\noindent{\large\bfseries #1\par}\medskip}
\newcommand{\pslinkmarker}[2]{\stepcounter{pslinkmark}\latelua{texio.write_nl('log','PS-LINK-SHIP|#1|\CurrentKey|\arabic{pslinkmark}|#2')}}
\newcommand{\pslink}[2]{\ifPSInLink\pslinkmarker{#1}{nested}#2\else\begingroup\PSInLinktrue\href{\CurrentScanURL}{\pslinkmarker{#1}{outer}#2}\endgroup\fi}
\newcommand{\scanpage}[1]{\ifnum\EditionMode=2\leavevmode\marginpar{\scriptsize\sffamily\pslink{scanpage}{scan #1}}\else\par\noindent{\small\sffamily\pslink{scanpage}{[scan #1]}}\par\medskip\fi}
\newcommand{\mspage}[1]{\ifnum\EditionMode=2\leavevmode\marginpar{\scriptsize\sffamily p.\,#1}\else{\small\sffamily[p.\,#1]}\fi}
\newcommand{\hand}[1]{{\ifnum\EditionMode=2\color{black}\else\color{PSHand}\fi #1}}
\let\over\relax
\NewDocumentCommand{\over}{O{} m m}{{\textcolor{PSTypedOld}{#2}\ensuremath{^{\!\nearrow}}\if\relax\detokenize{#1}\relax\hand{#3}\else #3\fi}}
\NewDocumentCommand{\inserted}{O{} m}{{\if\relax\detokenize{#1}\relax\hand{\ensuremath{\langle}#2\ensuremath{\rangle}}\else\ensuremath{\langle}#2\ensuremath{\rangle}\fi}}
\newcommand{\deleted}[1]{{\color{PSDelete}\ifmmode\text{[del.: }#1\text{]}\else[\textsf{del.:} #1]\fi}}
\newcommand{\xout}[1]{{\color{PSTypedOld}\ifmmode\text{[x: }#1\text{]}\else[\textsf{x:} #1]\fi}}
\let\PSMathUnderline\underline
\RenewDocumentCommand{\underline}{+m}{\ifmmode\PSMathUnderline{#1}\else\underLine{#1}\fi}
\NewDocumentCommand{\uline}{+m}{\underline{#1}}
\NewDocumentCommand{\uuline}{+m}{\ifmmode\PSMathUnderline{\PSMathUnderline{#1}}\else\underLine{\underLine{#1}}\fi}
\newunderlinetype\PSWaveLine{\leaders\hbox{\lower.7ex\hbox{\tikz[x=1pt,y=1pt,baseline=0pt]{\draw[line width=.3pt] (0,0) .. controls (1,1) and (2,1) .. (3,0) .. controls (4,-1) and (5,-1) .. (6,0);}}}}
\NewDocumentCommand{\uwave}{+m}{{\PSWaveLine #1}}
\newcommand{\unsure}[1]{\pslink{unsure}{[#1\ifmmode\,?\else\,?\fi]}}
\newcommand{\illegible}[1]{\ifmmode\text{\pslink{illegible}{[illeg.\ #1]}}\else\pslink{illegible}{[illeg.\ #1]}\fi}
\NewDocumentCommand{\diagram}{+m}{\ifmmode\vcenter{\hbox{#1}}\else\par\medskip{\centering #1\par}\medskip\psflushmargins\fi}
\newenvironment{diagramblock}{\par\begin{center}}{\end{center}\par}
\newcommand{\PSMathNotes}{}
\makeatletter
\NewDocumentCommand{\psplacednote}{m +m}{%
      \sbox{\PSNoteBox}{\parbox{\marginparwidth}{\footnotesize\raggedright\hand{\textsf{m#1.} #2}}}%
      \ifdim\dimexpr\ht\PSNoteBox+\dp\PSNoteBox\relax<\PSNoteLimit
        \leavevmode\marginpar{\usebox{\PSNoteBox}}%
      \else
        \par\nobreak\smallskip\hrule\nobreak\smallskip
        {\footnotesize\textsf{Marginal note m#1.}\quad\hand{#2}\par}
        \nobreak\smallskip\hrule\smallskip
      \fi}
\newcommand{\psflushmargins}{\ifmmode\else\ifinner\else\ifx\PSMathNotes\@empty\else\let\PSTempNotes\PSMathNotes\global\let\PSMathNotes\@empty\PSTempNotes\fi\fi\fi}
\NewDocumentCommand{\margin}{+m}{%
  \ifPSInApparatus\ifmmode\hand{\text{[marg.: }#1\text{]}}\else\hand{\textsf{[marg.: }#1\textsf{]}}\fi
  \else\stepcounter{psnote}%
    \ifmmode
      \text{\textsuperscript{m\arabic{psnote}}}%
      \begingroup\edef\PSQueue{\endgroup\noexpand\g@addto@macro\noexpand\PSMathNotes{\noexpand\psplacednote{\arabic{psnote}}{\unexpanded{$#1$}}}}\PSQueue
    \else
      \textsuperscript{m\arabic{psnote}}\psplacednote{\arabic{psnote}}{#1}%
    \fi
  \fi}
\AtBeginDocument{\everymath\expandafter{\the\everymath\aftergroup\psflushmargins}\everydisplay\expandafter{\the\everydisplay\aftergroup\psflushmargins}}
\makeatother
\newcommand{\psnumber}[2]{\expandafter\gdef\csname psnum@#1\endcsname{#2}}
\NewDocumentCommand{\var}{+m +m m m}{\ifnum\EditionMode=0 #1\else\ifnum\EditionMode=1\ifmmode\text{\textsuperscript{\tiny\csname psnum@#4\endcsname}}\else\textsuperscript{\tiny\csname psnum@#4\endcsname}\fi\fi #2\fi}
\newcommand{\psapparatusstart}{\par\bigskip\hrule\smallskip\PSInApparatustrue\begingroup\footnotesize\raggedright\textsf{Apparatus for scan \CurrentScan}\par}
\newcommand{\psapparatusend}{\par\endgroup\PSInApparatusfalse}
\NewDocumentCommand{\psentry}{m m m +m +m m}{\par\noindent\textbf{#2. #3.}\quad\textit{Typescript:} #4\quad\textit{Corrected:} \if\relax\detokenize{#5}\relax\textnormal{[omitted]}\else #5\fi\quad #6\par}
\hyphenation{Gro-then-dieck abel-ian-iza-tion}

\begin{document}
\hypersetup{pageanchor=false}
\begin{titlepage}\raggedright
{\Large Alexander Grothendieck\par}\vspace{12mm}
{\Huge Pursuing Stacks\par}\medskip{\LARGE 1983\par}\vspace{12mm}
{\Large A guide for the reader\par}\vspace{8mm}
{\large\EditionStatus\par}\vfill
Université de Montpellier, Grothendieck archive, cotes 134-2--134-8.\par
Including the appendix translation by R. Brown, cote 134-2, sites 115--143.\par
\vspace*{2mm}\end{titlepage}\hypersetup{pageanchor=true}
\chapter*{Editorial statement}

AI-transcribed by Claude Opus 5.5 (Anthropic, effort high); checked against the scans, and this edition built, by GPT-6 Astra (OpenAI, ChatGPT Pro mode); prompts and records by Claude (Anthropic).

The witness consulted and not reproduced is the LaTeX edition of Mateo Carmona with the collaboration of Ulrik Buchholtz (arXiv 2111.01000 v2, 2021), extended from the 2015 edition by ``the scrivener'', CC0 1.0.

The translation in the appendix is R. Brown's, prepared with assistance from L. Breen and J-L. Loday, as his note says; the rights in it belong to its translator or his heirs.

The rights in the original work belong to its author or the author's representatives or heirs. Where we hold rights, in the AI-produced transcription, translation, corrections and typesetting, they are dedicated to the public domain under CC0 1.0. This edition is published for educational, non-profit purposes, and we claim no rights we do not hold. If you hold rights in this work and object to its publication here, please open a GitHub issue and we will take it down.


\section*{Using this guide}
Each paragraph describes a section already checked against both supplied scan halves. It identifies the section's task, terminology, dependencies, revisions and consequential witness differences. Mathematical cautions are editorial observations, distinguished from the author's statements; they do not alter the readers. Archive sites and internal typescript folios are kept distinct, especially where the initial letter has been reordered or renumbered. All 706 source pages have completed their recorded image review. Statements that remain conjectural or provisional in the typescript are identified as such.

The section paragraphs follow the original order, including the separately numbered section~16-prime. Additional paragraphs cover the front matter, the complete letter to Quillen, the associated letters and notes, and Brown's translation. The supplied scans govern the account; the arXiv edition is a comparison witness.
\chapter*{Section guide}
\subsection*{Front matter: plans, contents, and dedication}
The front matter (cote 134-2, sites 1--13) should be read as evidence of a developing book: the title page and chapter-dependency diagram precede several layers of contents, including canceled headings, renumbered sections, and an entirely handwritten replacement page. The diagram uses the author's chapter labels and places derivators within the projected sequence; its heavily retraced connections should not be treated as a settled logical dependency chart. The canceled chapter-IV block on site 7 and the replacement handwritten contents on site 8 show the reorganization of the asphericity discussion. Sites 9--12 then project the progression through homology and abelianization, schematization, and linearization: the proposed comparison between schematic and discrete homotopy theory leads to the ``breakdown of an idyllic picture'' of section 119, then to revised programs and questions about the six operations in sections 124 and 134--139. These titles announce later developments; they do not themselves establish the claims. In particular, the pair of arrows listed for section 121 points from homotopy types to abelian homotopy types above and back below, with handwritten arrows and accents over typed objects and letters. Several continuations on site 11---\emph{topos}, \emph{bundles}, and \emph{valences}---are typing guided by handwritten curves, a material distinction lost in the initial transcription. The doubled \emph{ff} in the title for section 131 is visibly typed. The French dedication on site 13 retains a canceled address to former pupils above the dedication to Dodek Shapiro, together with its canceled and unresolved words. It remains French in all readers. The contents' cancellations and the final handwritten ``to be completed'' are particularly useful reminders that the source records successive intentions, including abandoned formulations, rather than a retrospectively regularized exposition (Grothendieck, 1983, cote 134-2, sites 1--13).

\subsection*{The letter to Quillen as a whole (cote 134-2, sites 14--36)}
The letter and its dated continuations supply the motivating problem for the subsequent notes: describe homotopy types by higher groupoid structures, relativize those structures over topoi to obtain stacks, and organize their operations through a universal categorical pattern. The progression is from the need for coherent higher composition, through globular source and target data and categories of standard amalgamations, to successive freely adjoined fillers and their comparison with topological or model-category realizations. The name ``coherator'' arrives only after substantial revisions to the proposed construction. The letter's mathematical force lies partly in those revisions: strict formulas give way to specified homotopies; the supposed universality of a category of actual disc maps becomes doubtful; economical choices of operations are treated as secondary to higher comparison; and unrestricted closure under fiber products is retracted in the postscript. The main text later returns to the program in sections 18--19 and develops a different, extensive investigation through test categories, asphericity, localizers, and integration. The earlier Breen correspondence reproduced in the appendix explains the topos and coefficient background invoked at the start and near the end of the letter. The archive places centered folios 4 and 3 in the reverse order at sites 16--17; the readers retain that archive order, while the section paragraphs identify the mathematical continuation. The witness's expanded marginal notes, silent index corrections, omitted signatures, and retained canceled passages must therefore be distinguished from the source's own stages. The guide records the local formula cautions in the relevant section paragraphs, without replacing the author's provisional mathematics by a reconstructed finished theory.

\subsection*{1. Foundational questions and tame topology (letter pp. 1--1')}
The letter dated 19 February 1983 starts from the papers sent by Ronnie Brown and recalls the 1975 correspondence with Larry Breen as a program joining homotopical and homological algebra over topoi. Before pursuing that program, Grothendieck argues for an even more urgent foundation for ``tame'' topology and stratified tame spaces, contrasting the conceptual difficulty with the inertia of established habits. The next section returns to the homotopical program; section 14 later resumes the criticism of the category of topological spaces and its inadequacy for isotopy and stratification. The opening is therefore a statement of priorities, not a theorem about a proposed category. Its cancellation history matters: on cote 134-2, site 14, the addressee is crossed out, and the long passage about the recipient's axiomatic homotopical algebra and lecture notes is canceled and replaced by a shorter handwritten phrase. The arXiv witness prints that canceled continuation as live text. It also substitutes ``as for schemes'' for the clearly handwritten ``or for schemes''. The section's descriptive title in the witness is absent from the scanned opening, which gives only its number. The continuation on site 15 concludes the comparison with programs of homological, homotopical, or topological algebra.

\subsection*{2. Higher groupoids and coherent composition (letter, sites 15--18)}
Returning to the Breen program recalled in section 1, Grothendieck proposes that $n$-truncated homotopy types be described by $n$-groupoids, with a fundamental $n$-groupoid, geometric realization, and an $n$-equivalence comparison; relativization over a topos should yield $n$-gerbs and $n$-stacks as noncommutative cohomological coefficients. The assertion that every $n$-groupoid should be $n$-equivalent to a fundamental one is present on site 15 but omitted from the arXiv witness, which joins the surrounding clauses. The attempted globular description then introduces sets $F_i$, source, target, identity and inversion maps, and composition at different levels. Strict associativity is rejected as insufficient for the intended homotopy types; equalities are to become specified homotopies, whose own coherence leads from ``secondary'' to ``ternary'' and further structure, illustrated by the pentagon at the start of site 18. Section 3 asks whether this can support a model-category description; sections 18--19 later revise the coherator and realization proposals. The internal reading order here is site 15, then site 17 (centered folio 3), then site 16 (folio 4), then site 18 (folio 5); the readers preserve the archive's different order. Several formulas require caution. On site 17 the first factor of the composition domain has an overprinted $t$, with the underlying letter still doubtful, where the witness gives $s$; the marginal identity notation is $id_a$, not $d_a$. The handwritten target formula ends with $s^i_1(u)$ in its second right-hand factor, an apparent source/target slip. On site 16 the inversion formula places separate inversion marks on both left-hand factors and appears tautological, although the first right-hand letter remains overwritten and doubtful; the witness substitutes a question mark for that right side. The handwritten interchange range also permits equal levels, $1\leq\alpha\leq\nu$: for a one-object group at equal levels, the displayed interchange would force the two middle factors to commute, so this range is too broad for arbitrary fundamental groups. These source readings remain unchanged. The author replaces ``implying'' by ``expressing'' beneath the compatibility conditions, a correction the witness does not apply. The resulting discussion is a developing proposal for coherence, and these local defects should not be treated as established definitions or silently repaired.

\subsection*{3. The model-category question (letter, sites 18--19)}
Grothendieck asks Quillen whether an appropriate notion of higher category already exists, recalling a description by multiple semisimplicial sets but distinguishing strict associativity from the coherent structure sought in section 2. He proposes two connections with homotopical algebra: the conjectural category of infinity-groupoids should model ordinary homotopy types, and a fundamental infinity-groupoid construction should be available in any model category from a cylinder object of its final object. The latter is presented as informative chiefly in the nonpointed case, and the handwritten ending explicitly identifies the cylinder as playing the role of the unit segment. Section 4 immediately reopens the skeletal data of that proposed construction; section 14 returns to the limitations of topological-space models, and the later interval discussion develops the cylinder intuition. These are proposals and questions, without a completed definition or asserted model-structure theorem. On site 18, the handwritten replacement of ``your'' by ``Quillen's'' in the reference to axiomatics is not applied by the arXiv witness, and its ``contents'' misreads the typed ``comments''. Site 19 closes the first letter with a handwritten signature and is numbered 5-prime on the scan, where the witness labels it 3-prime. The ensuing dated continuation on site 20 begins a substantive revision rather than simply continuing an established definition.

\subsection*{4. Skeletal data and successive structure (letter, sites 20--21)}
In the continuation dated 20 February, Grothendieck revises the undifferentiated ``primary'' structure proposed in section 2: the skeletal datum is only the family of sets with source and target maps, while inversion, degeneration and first composition belong to structure added subsequently. Equivalences of ordinary groupoids motivate replacing strict involutions by pseudoinvolutions, and choices of interval maps explain why degeneracies and path composition should also be nonprimitive. Godement's next composition is then secondary, and the further composition levels successively ternary and higher. This ordering prepares the tower of universal categories in section 5 and makes explicit a change of mind: the text says that the preceding day's operations were hastily put into one pot. Two details deserve care. On site 20 the formal maps $s_o,t_o$ are described as coming from $F_0$ to $F_{-1}$ but are assigned ``$i=1$''; the evident index mismatch is retained. On site 21 the handwritten relocation yields ``are rather'', while the witness retains ``rather are''. The replacement ``the real'' on that page is typed above the line, not handwritten as the initial transcription classified it. The source says ``all these operations'', where the witness shortens this to ``all the operations'', and retains the two handwritten marginal composition symbols absent from that witness. The overwritten folio at the head of site 21 remains uncertain between 7 and 4; the archive order and all scan links are unaffected.

\subsection*{5. A tower of universal categories (letter, sites 21--24)}
Grothendieck proposes a sequence of categories $C_0,C_1,\ldots$ with common objects and progressively added arrows, whose inductive limit should classify the proposed infinity-groupoid structures by functors preserving specified limits. The discussion first compares universal structures in categories with finite products, finite projective limits, and the stronger exactness available in topoi; it then discards the separate finite-product requirement while writing site 23, restricting the common objects to ``standard'' iterated fiber products determined by source and target. This is the categorical organizing device promised in section 4. Section 6 gives it a topological model; the later coherator discussion revises how the successive operations should be adjoined. The text expressly presents the strict increase at every stage as a conviction still needing proof. On site 23 the letter in ``exist in F'' is visibly F, although the ambient category is C and the witness prints C; the source letter is retained as an apparent local slip. The iterated-source indices on that page contain typed 1-shaped bases with handwritten ell/chi-like corrections and remain marked, rather than being settled solely by the witness's ell. The live insertion says that the word ``iterated'' will henceforth be dropped, while the canceled marginal version says ``qualification''; the witness combines the latter wording with the live text. On site 24 the marginal note assigns order ell to the composition bearing lower index ell, where the witness gives that composition lower index 2. Some of this note's lettering and cancellation status remain doubtful, but the difference of indices should not be silently harmonized. The section's revisions distinguish a proposed universal construction from a finished coherence theorem.

\subsection*{6. Discs and a proposed universal co-groupoid (letter, sites 24--25)}
The tower of section 5 receives its motivating topological interpretation: a space $X$ supplies sets of maps from discs $D_i$, and the source and target maps are induced contravariantly by the two boundary embeddings. Compositions and further operations are to arise from maps between discs and their prescribed amalgamated sums, producing a co-infinity-groupoid whose transform under $h_X$ is the desired infinity-groupoid. Grothendieck then proposes that the category $B_\infty$ generated by these cells, sums and arrows should be dual to the universal category $C_\infty$. The expected universality would make the category independent of the successive arbitrary choices up to unique isomorphism, while he insists that the formal construction of $C_n$ remains meaningful without that expectation. Section 7 begins to specify the disc maps and their coglobular identities. Site 25 records a significant change of mind: the confident parenthesis ``I believe it is'' acquires handwritten exclamation marks, and a marginal addition dated 23 February says ``But no longer now'' and refers to page 10, section 11. The crossing pencil hatching leaves the note's cancellation status uncertain; the readers retain that doubt. The arXiv witness paraphrases this retraction, omits its explicit page reference, and prints the canceled article before ``various'' amalgamated cells. On site 24 the unit-segment formula retains the typed interlinear $I=$ and the live $D_1$, an identification shortened in the witness. These distinctions preserve both the proposed model and the author's later withdrawal of confidence in its universality.

\subsection*{7. Coglobular discs and standard amalgamations (letter, sites 25--27)}
The two hemisphere maps between successive discs give the cellular skeleton needed in section 6. Grothendieck describes the boundary identifications and the decomposition into open cells, then interprets composition at level ell as a choice of extension from a boundary sphere into an amalgamated sum. Contractibility supplies the proposed existence argument; its inductive dependence on lower-level operations explains the hierarchy introduced in section 4. He next seeks a combinatorial account of ``standard'' amalgamations by attaching cells along prescribed subcells, dualized into relations among iterated sources and targets in a product of the objects $F_i$. Section 8 recasts these data as finite ordered sets. Two scan readings materially alter the comparison witness: on site 26 the range is strictly $0\leq j<i$, and the open-cell rings are present on $D_n$, both $S_j$ symbols, and the argument of the positive map. The corresponding negative-map argument has no visible ring; this apparent local omission is retained. The same page draws undirected circle examples, where the witness adds arrowheads. On site 27 the proposed restriction on pairwise cell intersections receives a marginal counterexample. Its three regions are cut out by two arcs sharing their endpoints, so $f\cap f''=\{s,s'\}=S_0$; a central three-sector diagram would destroy that counterexample and has been corrected in the checked transcription. The note says ``see this'', with its first word and hatching status still uncertain, whereas the witness substitutes explanatory prose. The product formula uses $F_{n(i)}$, not the witness's $D_{n(i)}$, and the handwritten sign variables are epsilon and epsilon-prime throughout. These qualifications keep the proposed combinatorial scheme and its own marginal objection visible together.

\subsection*{8. The finite ordered sets underlying $B_o$ and $C_o$ (letter, site 28)}
The dated continuation of 21 February replaces each amalgamated cell complex by its finite ordered set of cells, with distinguished subposets isomorphic to the standard $2n+1$-element posets $J_n$. The incidence relation and the specified identifications encode the combinatorial structure sought in section 7; their opposite category is proposed as $C_o$, the starting point for the inductive construction in section 9. The intersection $K_i\cap K_{i'}$ is expressly a closed subposet isomorphic to some $J_n$; the arXiv witness misprints its intersection sign as an arrow. Grothendieck then excludes the empty structure and claims stability under binary amalgamated sums, extending the claim to all $C_n$ and describing the desired functors by preservation of those sums. The page itself objects: the all-categories phrase has a wavering underline and a marginal note saying ``This is obviously false, see P.S. p.12'', with repeated calls to note 13. The pencil hatching leaves that note's status uncertain, but its wording is clear. The distinction between amalgamated sums in $B_o$ and limits in its opposite is also essential to reading this developing assertion; the readers preserve the claim rather than repairing its variance. The witness removes the all-categories phrase and rewrites the note with later qualifications referring to sections 13 and 18. The guide therefore treats the closure claim as a proposal already challenged on its own page, pending the later revisions.

\subsection*{9. Universal coherence and comparison between definitions (letter, sites 28--30)}
Conditions (A) and (B) propose that suitable parallel arrows in $C_n$ acquire a filler at the next stage and that $C_{n+1}$ be freely generated from such fillers subject to their prescribed boundaries. This turns the informal coherence program of sections 2 and 4 into a proposed universal construction, while section 10 will consider how economically to choose the new arrows. The phrase ``an isomorphism'' is visibly canceled beneath the typed replacement ``a homotopy'' on site 29: the intended coherence is itself higher-dimensional, rather than an identification of all operations. Grothendieck nevertheless says that these conditions do not characterize $C_\infty$ up to equivalence. His sequence of seven mathematicians on site 30 explains the intended weaker independence: competing definitions should have comparison functors, the composites should be infinity-equivalent to the originals, different comparison choices should themselves admit equivalences, and comparisons between those equivalences should continue indefinitely. The promised finite mathematical statement of this story remains explicitly unwritten. Two local cautions concern the construction itself. The filler has target $F_{i+1}$ on site 28 but $F_i$ on site 29; the source mismatch is preserved. The marginal correction on site 29 says that amalgamated sums exist only in $C_o$, not the initial transcription's $C_\infty$, and another note calls the asserted stability of all $C_n$ inaccurate. The underlined demonstrative and abbreviated ``should'' in the first note, as well as its hatching status, remain uncertain. These source objections continue the revisions noted in section 8 and prevent the asserted universal construction from being read as a settled theorem. The arXiv witness retains the canceled ``can'' instead of the live ``may'', omits the source article in ``at all a sense'', and supplies articles before ``same'' that the scan does not contain. None of those changes is imported into the readers.

\subsection*{10. Redundant operations and choices of coherence data (letter, sites 30--31)}
Grothendieck considers progressively narrower choices of fillers for the construction of section 9: one for every admissible pair, one for every isomorphism class of such situations, one only when no previous filler exists, and finally only for pairs that do not factor nontrivially through a smaller object. Even these restrictions leave many operations serving the same purpose. In particular, binary composition and all simultaneous $n$-ary compositions appear at the first stage; the second stage relates them by homotopies and the third relates composites of those homotopies. The comparison with two descriptions of a category with associative tensor composition gives a finite-dimensional analogue of the seven-mathematician story. Its claim about nonequivalence of the ordinary categories of algebraic structures is explicitly qualified by the typed insertion ``I believe''; it should not be read as denying the higher comparison the preceding section seeks. The marginal qualification ``even not in the grou-poid context'' remains uncertain in status because of pencil hatching, and the arXiv witness expands it into a different footnote mentioning Gr-categories. Grothendieck concludes that economy of the presentation is secondary and leaves the notion partly indeterminate, subject to (A), (B), and the description of $C_0$. Section 11 tests the alternative of selecting actual maps of topological cells. On site 31 the first variable is visibly capital $U_1$, while the witness prints $u_1$; this local notational mismatch is preserved. The initial transcription also mistook the live overprinted letters of ``operation'' for a deletion, which would have produced ``opeion'' in a corrected reader; the image check restores their proper typed-overprint status.

\subsection*{11. Allowable maps and the faithfulness question (letter, site 32)}
Grothendieck returns to the topological realization proposed in section 6, asking whether the universal category can be described as the opposite of a category of actual maps between amalgamated cells. He proposes subsets of allowable maps from discs, stable under composition and admitting extensions from their boundary spheres. These conditions model closure and the filler requirement, and elementary tameness conditions seem capable of enforcing them. They do not by themselves secure the freeness required by condition (B): actual composition of continuous maps might impose extra relations not forced by the universal algebra. Thus the stepwise universal construction gives a contravariant realization functor, but its faithfulness is neither evident nor, in the closing judgment, essential. This is the substantive withdrawal of confidence anticipated by the dated marginal retraction on site 25. The later comparison of competing realizations can therefore proceed without identifying the abstract category with its image in topological spaces. The arXiv witness drops K from ``every K a subset'' and changes the source index $C_{n-1}$ to $C_{n+1}$ when describing the universal properties of $C_n$; the latter source reading is clearly minus 1 and suits the stated inductive construction. The overprinted A/a in the reference following condition b) remains visually doubtful, but the reference concerns the earlier filler requirement; it should not be silently interpreted as the local closure condition a). The superscript mark on $S_n$ also remains unresolved, and the handwritten subscript in the phrase about a piecewise-linear structure is uncertain between i and n. These uncertainties are preserved in the readers rather than decided by the witness.

\subsection*{12. A co-groupoid in a model category (letter, site 33)}
The proposed disc construction is transferred from topological spaces to an arbitrary model category, as anticipated in section 3. Grothendieck chooses trivial fibrant objects $D_i$ and cofibrations from the successive boundaries, then seeks to extend the resulting realization functor from $B_o$ through the universal stages. He expects standard gluings of those objects to remain fibrant and trivial and proposes localization of the associated infinity-groupoids to recover the homotopy category. For semisimplicial sheaves over a topos, this is related to the operation called ``integration'' or ``sections'' of gerbs and stacks in the earlier Breen program. These expectations are explicitly presented as things still to be checked. In particular, the asserted stability of fibrancy under the indicated pushouts is not a consequence of the model axioms: in simplicial sets, glue the nerves of two contractible two-object groupoids along a common vertex. All three inputs are contractible Kan complexes and the inclusions are cofibrations, but a composable pair crossing the shared vertex lacks a filling 2-simplex in the pushout, so that pushout is not Kan. The mathematical assertion therefore remains as written and should not be silently converted into a cofibrancy assertion. The subsequent terminology changes in section 13 retain the conceptual program while abandoning an unqualified name for its choice-dependent objects. On this page the witness changes the source name (Ho) to Hot and the handwritten lower-case pi to capital Pi. The final two interlinear replacements of ``groupoids'' by ``gerbs'' are typed and are retained as such. The overprinted disc indices and the supposed equals correction after $S_o$ remain marked where the image does not securely separate their layers.

\subsection*{13. Coherators, intrinsic structure, and retractions (letter, sites 34--36)}
Grothendieck replaces the potentially misleading name ``infinity-groupoid'' for the auxiliary algebraic pattern by ``coherator'', reserving ``stacks'' for structures modelled on such a pattern. Building on the universal constructions of sections 9--12, he asks how much of the chosen globular objects, source and target maps, and filtration can be recovered from the abstract category. The successive answers on site 35 must be read with their retractions: the proposed automorphism groups $(\pm1)^i$ are challenged by the observation that the inversion operations need be neither involutions nor automorphisms; an irredundancy condition is then offered to recover rigidity of the whole globular system. The marginal objection dated 25 February says that the proposed intrinsic construction can give ``two large a subcategory'' $C_{n+1}$, and the closing marginal note on site 36 suggests supplying the primitive subcategory $C_o$ as extra structure. These are changes within a developing definition, not an established characterization theorem. The continuation in sections 18--19 revisits this program after the intervening reflections on topology. The final postscript, visibly dated 23.2 rather than the witness's 25.2., explicitly withdraws the claim that all the categories $C_n$ admit the fiber products previously asserted: the surviving statement concerns the fiber products already present in $C_o$, their preservation in $C_n$, and functors whose restrictions to $C_o$ preserve them. This resolves the objections on sites 28--29 without deleting the earlier claims. Three local mathematical slips remain in the text. On site 35, the intrinsic characterization of the source--target pair refers to chosen $F_i$ and $F_{i+1}$, although the preceding arrows have codomain $F_{i-1}$; the witness silently substitutes $i-1$. Later the second boundary composite uses an ununderlined $t$, and the final endomorphism is named $h$ while its boundary equalities use $f$. Neither notation is silently harmonized. The source's marginal notes have shorter, sometimes differently located wording than the witness footnotes: call 13 follows ``amalgamation'', call 19 rejects the automorphism claim, and call 20 qualifies the filtration. Two of those notes were hidden in comments in the initial LaTeX and have been restored to printable anchors. Their readable wording remains visible with uncertainty about crossing pencil hatching. The extra cluster after ``Maybe'' on site 36 and the second numeral in its body call remain unresolved, while both calls in the postscript read 13 followed by an exclamation mark. The typed interlinear ``later'' gives ``a later elaboration'', where the witness rearranges the phrase. The letter thus ends by stating the intended role of stacks as noncommutative cohomological coefficients while expressly leaving the coherator axioms open to further testing.

\subsection*{Introduction to the Breen letters (cote 134-2, site 37)}
The appendix introduces letters dated 5 February, 17 February, and 17--19 July 1975 as the earlier program recalled in the Quillen letter. Grothendieck explains that the first two are reproduced in French and the third in Ronnie Brown's English translation, prepared with assistance from Larry Breen and J.-L. Loday; he identifies his present interventions as comments, numbered subtitles, and corrections chiefly occasioned by the handwriting of the original. The first two letters develop the proposed connection between higher categories, stacks, and truncated homotopy types. The third broadens the program to duality, geometric local and global class field theory, tame topology, and the inverse system of indexed homotopy types associated with equisingular stratifications. These themes supply the historical setting for the later investigations of integration and tame topology without asserting that the 1983 formalism has already been established. Three witness differences materially affect that setting: the typed copy is dated 1975 rather than 1945; the intended phrase is ``fine homotopy type'', not ``finite homotopy type''; and the handwritten replacement of ``ordered set'' is ``inverse system'', not ``inverted system''. The source also specifies notes 22--31 and subtitles App.1--App.18, details shortened in the witness. The stepped pen cue begins a new paragraph for the third letter, while the retained cancellations and insertions distinguish the 1983 editorial introduction from the letters it presents.

\subsection*{Appendix 1: Picard categories and truncated complexes (cote 134-2, sites 38--40)}
The first part of the letter of 5 February 1975 asks what geometric objects, rather than their isomorphism classes alone, realize cohomology groups. Its four examples move from gerbes and their higher torsors to Picard stacks, stacks of $(\mathrm{Gr})$-categories, and Azumaya algebras. For a band $L$ with center $Z$, the $Z$-gerbes form a strict Picard 2-category governed by the degree-two truncation of $R\Gamma_X(Z)$; the $L$-gerbes form a pseudo-2-torsor, whose nonemptiness has an obstruction in $H^3(X,Z)$. Passing from global objects to the corresponding stacks shifts the description toward 2-gerbes and a degree-three truncation, anticipating Appendix 2's proposed interpretation of noncommutative complexes by higher categories. Picard stacks with invariants $M,N$ introduce the unusual $\mathrm{Ext}^2$ attached to $R\mathrm{Hom}(X(M),N)$, distinguished from the ordinary lower Ext groups and from the strict case governed by $R\mathrm{Hom}(M,N)$; the source range for those lower groups is $0\leq i\leq1$, whereas the witness extends it to 2. Likewise, the source repeatedly has $Z$-1-gerbes or $Z$-1-champ, not the witness's $Z$-$L$ expressions. For a group $G$ acting on $N$, the proposed invariants are relative cohomology groups in degrees 3 and 2 together with $Z^1(G,N)$, explicitly the cocycle group, which must not be replaced by $H^1$. The Brauer example assigns to an Azumaya algebra its gerbe of trivializations and distinguishes the Brauer subgroup of $H^2(X,\underline G_m)$ from the full cohomology group; this leads Grothendieck to question effective 2-descent for the proposed fibered 2-category. The middle displayed invariant on site 40 is literally $H^1(X,\underline H_m)$: the surrounding $\underline G_m$-gerbe construction and the adjacent degree-two and degree-zero terms indicate the intended coefficient $\underline G_m$, which the witness silently supplies, but the source $\underline H_m$ is retained in every reader. The geometric trivialization variable on that page is handwritten $\phi$, not an empty-set symbol. The change of emphasis from classes to structured objects, and then from global objects to stacks, provides the dependency needed for the ensuing Lefschetz discussion; the later notes 22--23 qualify the band and the complex $X(M)$ already cited here.

\subsection*{Appendix 2: Lefschetz theorems as a test for higher stacks (cote 134-2, sites 40--42)}
The preceding Picard examples suggest interpreting truncated complexes, explicitly as objects of derived categories, by strict Picard $n$-categories and then forgetting the additive structure to obtain candidates for noncommutative complexes. Grothendieck motivates this proposal through weak Lefschetz theorems in étale cohomology: depth conditions and the strongest available statements in degrees zero and one are naturally formulated for stacks, rather than merely for cohomology sets or fundamental groups. Site 41 distinguishes the prime-to-characteristic restrictions accompanying duality and smooth or generic base change from the unrestricted proper-case statements, and identifies proper base change, smooth base change, and generic cohomological properness as the indispensable tests of a useful formalism. The witness omits the constructibility hypothesis in the generic case, together with the transition explaining that smooth and generic base change serve the generic Lefschetz version; the omission materially weakens the stated assumptions. The handwritten étale-cohomology phrase belongs after ``coéfficients discrets'', and the canceled phrase before ``les résultats les meilleurs possibles'' reads ``en fait'', not the initial transcription's tentative ``aussi''. Site 42 extends the proposed framework to higher stacks and complex analytic variants, while insisting that serious geometric theorems, rather than formal manipulation alone, must justify it. The concluding dissatisfaction with the definition of triangulated category prepares note 24 and the later concern with retaining structure lost in localization. Appendix 3 then supplies the promised heuristic link with simplicial homotopy theory; the Quillen letter and the main body return to that link through test categories and higher groupoids. The present passage states a program and its tests, not a constructed general theory of noncommutative derived categories.

\subsection*{Appendix 3: Higher groupoids and truncated homotopy types (cote 134-2, sites 42--45)}
The postscript proposes the homotopy hypothesis in its finite-dimensional form: an $n$-category in groupoids, up to $n$-equivalence, should encode a simplicial homotopy type with homotopy groups above degree $n$ discarded. Paths, homotopies of paths, and successive higher homotopies motivate $C_n(K.)$, with the top-dimensional arrows taken modulo homotopy; the failure of strict associativity is explicitly a reason to seek weak higher categories and coherent higher functors. The source inequality on site 42 is $i\geq n+1$, while the witness reverses it, thereby changing which homotopy groups are discarded. Site 43 defines the invariants through equivalence classes of iterated endomorphisms; the witness omits the full clause separating the set of components from the next endomorphism-class construction, and also misbrackets the displayed functor $C_n(K.)\to C_n(K'.)$. The capital $N$ in one sentence and the omitted basepoint in $\pi_i(C_n(K.))$ remain literal source features; the retraced subscript in the earlier definition is genuinely doubtful. Site 44 considers localization, truncation, and adjoining identity arrows, then passes toward infinity and relates strict Picard objects to the Dold--Puppe picture. His suggestion that localization might additionally require surjectivity on $\pi_{n+1}$ is a hesitation, not a needed condition for the category of $n$-types: for nonzero $A$, the map from a point to $K(A,n+1)$ is already an isomorphism on all groups through degree $n$, although it is not surjective in degree $n+1$. The higher categories in this comparison must retain the groupoid qualification: ordinary categories localized only at categorical equivalences do not identify the contractible interval category with the terminal category. Site 45 tests the proposal on loop objects, mapping objects, homotopy fibers, Eilenberg--Mac Lane objects, and Postnikov decomposition, leaving suspension without a satisfactory categorical candidate. These constructions feed Appendix 4's relative version and the later Quillen-letter program. The proposed classifying topos of a higher groupoid is expressly reconsidered in the following section, so the optimistic Postnikov description here should be read together with that later retraction.

\subsection*{Appendix 4: Relativization and the classifying-topos difficulty (cote 134-2, sites 45--46)}
The absolute comparison is relativized over a topos $X$, pairing simplicial sheaves with higher stacks in groupoids. The full source specification of the simplicial-sheaf side is omitted in the witness. More seriously, the source compares the evident inverse-image operation for simplicial sheaves with the subtler inverse image for higher stacks, whereas the witness substitutes direct image at that point; the next sentence then reverses the perspective, comparing direct image of stacks with a derived direct image of simplicial sheaves. Grothendieck explicitly hesitates between $Lf_*$ and $Rf_*$, and the readers retain both symbols and the hesitation. The Postnikov analogy is linked to the decomposition of a stack in groupoids into its component sheaf and a gerbe over the induced topos. On site 46, the paragraph headed ``Remords'' retracts the preceding expectation that extending classifying topoi to higher stacks should be straightforward: he cannot construct the requisite object in higher degree, and reduces the difficulty to obtaining a functorial topos of homotopy type $K(\pi,n)$ for an abelian group $\pi$. The proposed nerve construction, compatibility of homotopy and cohomology invariants, and the possibility of needing higher topoi are presented as unresolved requirements. His stated low-dimensional exception and the surrounding inequalities remain as written. This correction of perspective explains why Appendix 5 treats the proposed theory as a coordinated program involving several kinds of objects, rather than a completed identification of them; it also anticipates the later use of simplicial and categorical models to construct classifying spaces and homotopy types.

\subsection*{Appendix 5: The proposed topological algebra (cote 134-2, sites 46--47)}
The closing section of the first letter gathers four families of objects into one proposed theory: topological spaces and topoi; simplicial sets and simplicial sheaves; higher categories, especially higher groupoids, and higher stacks; and complexes of abelian groups or abelian sheaves. Their relations and possible additional structures, rather than one favored presentation alone, define the program. This synthesis depends on the Picard examples, the Lefschetz test, the homotopy hypothesis, and the relative difficulties developed in Appendices 1--4. Grothendieck proposes ``algèbre topologique'' as a shorter name than noncommutative homological algebra, argues that the familiar collection of topological rings, fields and groups does not constitute a sufficiently unified doctrine to monopolize that name, and invokes Serre's usage of analytic geometry as an analogy for terminological reassignment. The witness repeats a phrase in the enumeration on site 46 and omits the handwritten word ``vocable''; the checked source retains the actual single enumeration and the recovered word. The closing initials on site 47 are also present. The unusual extra typed acute above the m of the final analytic-geometry phrase is preserved as a source feature, not converted into a standard spelling. The next letter, beginning Appendix 6, takes up the promised relationship between locally constant higher stacks and a topos's homotopy type rather than claiming that this proposed unified theory already exists.

\subsection*{Appendix 6: Local systems, fundamental higher groupoids, and classifying topoi (cote 134-2, sites 48--51)}
The second letter begins by comparing two characterizations of an $n$-equivalence $f:X\to Y$: invariance of cohomology with locally constant coefficients through degree $n$, including the appropriate noncommutative low degrees, and equivalence on categories of sections of locally constant $(n-1)$-stacks. The proposed equivalence between conditions A and B turns the earlier homotopy hypothesis into a reconstruction problem for the fundamental $n$-groupoid $\Pi_n(X)$: locally constant $(n-1)$-stacks should be $n$-functors from that groupoid to $(n-1)$-categories. When the local connectivity hypothesis fails, a pro-$n$-groupoid is proposed instead, with the classification of tori cited as a degree-one precedent. Site 49 connects the fundamental groupoid to the homotopy groups through degree $n$ and to Postnikov invariants; its clear range $0\leq i\leq n$ is reversed twice in the witness. The earlier heavily retraced dimensional inequality on that page remains doubtful in the checked text and should not be silently settled from the later clear occurrences. Recovering a groupoid from its local systems leads on site 50 to the prospective higher topos $\mathcal T_n=(n\text{-}\underline{\mathrm{Hom}})(\mathcal C_n,((n-1)\text{-Cat}))$, whose suitably exact fiber functors should recover $\mathcal C_n$. The source distinguishes these two symbols throughout, whereas the witness repeatedly substitutes $C_n$ for $\mathcal T_n$ and once removes the $n$ from $n$-topos. It also reverses the two limit arrows: the source correctly asks an inverse-image functor to preserve arbitrary inductive limits and finite projective limits. Grothendieck then considers the weaker hope of an ordinary classifying topos that 2-represents $T\mapsto\tau_1\underline{\mathrm{Hom}}(\Pi_n(T),\mathcal C_n)$, and immediately doubts its existence already for $n=2$. The mixed typed and handwritten marginal objection ``cela semble suspect!'' is absent from the witness but is integral to this reconsideration. His proposed obstruction is failure of a sufficiently strong van Kampen principle for higher groupoids; at the top of site 51 he asks whether locally constant 1-stacks on a union of two closed subsets can be recovered from their restrictions and an equivalence on the intersection. The intersection sign is visibly $\cap$, which the witness changes to $\cup$ and thereby changes the gluing datum. This passage records an unresolved reconstruction and descent problem, rather than a proof of the suggested failure. It continues Appendix 4's retraction, supplies the local-system description used immediately in Appendix 7, and anticipates the later main-text investigation of classifying spaces, weak equivalences, and the information preserved by categorical models.

\subsection*{Appendix 7: Integration of higher stacks and the proposed cohomology formulas (cote 134-2, sites 51--54)}
Using Appendix 6's local-system interpretation, Grothendieck identifies integration over $X$ with an internal Hom from the final higher stack, then translates this operation into cohomology. A strict Picard $n$-stack associated with a complex in degrees zero through $n$ gives $H^i(X,L)=\pi_{n-i}\Gamma_X(\mathcal F)$ in the stated range; iterating the classifying-object operation extends the definition to higher degrees. For a general stack in groupoids only the initial range is asserted, while a group structure permits one further delooping and degree. The proposed necessity of strict Picard structure for a second delooping on site 53 is explicitly tentative and should not be promoted to a theorem: extra coherent multiplication is the issue, and the text has not supplied its classification. The explicit correction concerning highly iterated loop spaces in Appendix 9 (site 56) must be read alongside this earlier expectation. The source phrase ``à isomorphisme unique près'' on site 51 also needs care, since an object of a derived category can have nontrivial automorphisms; uniqueness requires additional specified data, not merely passage to that category. On site 52 the comparison from the classifying object of global sections to the global sections of the classifying stack is literally called a morphism of $(n-1)$-groupoids, although the construction and the following range $1\leq i\leq n+1$ indicate the higher dimension $n+1$. That index remains unchanged in every reader. Conversely, the opening parenthesis in the boxed formula is actually visible in the scan: its restoration corrects the initial transcription and is not an emendation of Grothendieck's mathematics. The small handwritten 4 following the bracketed explanation is also present. Site 53 transports the formulas to an infinity-groupoid and a local system, raises the canonical comparison in degree $n+1$ as a question, and restates A and B as conditions A$'$ and B$'$ for a truncated functor. The first cohomology formula has $\pi_{n-i}$, not the witness's $\pi_{n-1}$. The right-hand final object on site 54 has exponent $\mathcal C_{n+1}$, not the witness's $\mathcal D_{n+1}$; both handwritten $\Gamma$ abbreviations take $\mathcal F$ as their argument, whereas the witness puts it into the subscript. Condition B$'$ needs a substantive qualification as written. Already for $n=0$, the functor from the terminal groupoid to the one-object groupoid of a nontrivial group $G$ induces an equivalence on zero-truncations, but a $G$-set $S$ gives the section map $S^G\to S$, which need not be a bijection; a nonempty free $G$-set makes this failure explicit. Thus the asserted equivalence of zero-truncations alone does not imply the stated section comparison for every local system. This elementary check identifies a hypothesis or truncation problem without altering the text. Appendix 8 then compares this proposed categorical construction with derived-category and semi-simplicial constructions of cohomology, while the later main text's work on integration and local equivalences develops the broader program announced here.

\subsection*{Appendix 8: Three constructions of cohomology (cote 134-2, sites 54--55)}
The second letter closes by asking whether the higher-stack construction of cohomology can be developed without complexes of abelian sheaves, injective resolutions, or any semi-simplicial machinery. Grothendieck notes its apparent kinship with \v Cech methods through hypercoverings, but the sentence explicitly questioning whether semi-simplicial algebra can be avoided is omitted by the witness; preserving it is essential to the conditional status of the proposed independence. He distinguishes three approaches: commutative homological algebra through complexes, resolutions and derived categories; homotopical algebra through \v Cech or semi-simplicial constructions; and categorical or noncommutative homological algebra through higher stacks. Site 55 summarizes their differing objects of resolution: coefficients in the first approach, the base space or topos in the second, and apparently neither in the third. This is a conceptual comparison dependent on Appendix 7's proposed integration formulas, not an assertion that the third construction has already been supplied. The underlined expression crosses the source-page boundary, and the handwritten commas closing the first two items remain commas rather than the witness's semicolons. The source's repeated spelling ``résoud'', unpunctuated closing salutation, and typed signature are retained; the witness omits the signature. Appendix 9 resumes the constructive problem after Breen's response, reconsidering both the required higher-category axioms and the relationship between strict Picard objects and topological commutativity.

\subsection*{Appendix 9: Breen's construction and a noncommutative Dold--Puppe problem (cote 134-2, sites 56--58)}
The letter dated 17/19 July 1975 responds to Breen's proposed non-strict higher categories and nerve construction, welcoming a precise beginning while criticizing its technical character and the absence of a geometrically transparent functor from simplicial sets to higher groupoids. Grothendieck proposes an explicit list of desired objects, operations and comparisons that might determine a construction or an existence-and-uniqueness theorem; the handwritten note 25 qualifies that uniqueness expectation. He asks strict Picard $n$-categories to recover the appropriate category of chain complexes, but expressly withdraws his earlier belief that sufficiently associative and commutative iterated loop spaces correspond to commutative topological groups. This correction on site 56 belongs with the earlier Picard and delooping speculations, especially Appendix 7. Sites 57--58 ask for the noncommutative analogue of the Dold--Puppe comparison and work out the first categorical dimension through a homomorphism $d:L_1\to L_0$, an action $\theta$, equivariance of $d$, and the identity relating $\theta(d(x_1))$ to inner conjugation. These are the crossed-module conditions, giving the component group, the central kernel, and the Postnikoff--Sinh class in $H^3(\pi_0,\pi_1)$. The source's second axiom lacks one closing parenthesis before the equals sign; its intended identity is determined by the preceding action and the conclusions, but the missing parenthesis remains in the readers. At the opening of site 58, the action factors through $\pi_0$ on the invariant subgroup $\pi_1$, not generally on all of $L_1$: the preceding conjugation axiom itself shows why the latter interpretation would be too strong. The witness also changes several explicit multiplication dots, plural categories, and the singular word structure, which the checked text preserves. An application to formal groups motivates asking for noncommutative chain complexes with the higher analogues of the action and its coherence conditions. The concluding handwritten replacement specifies \emph{pointed connected} spaces as the corresponding delooped objects, an essential qualification. This section revisits the homotopy hypothesis from Appendix 3 in algebraic terms and anticipates the main text's later insistence on weak higher groupoids and coherent structures; Appendix 10 then turns to Breen's separate questions about duality and homology.

\subsection*{Appendix 10: The six operations, homology, and autoduality (cote 134-2, sites 58--60)}
Breen's questions about biduality lead Grothendieck from higher-stack constructions to the six-operation formalism and its interpretation of homology. Part A surveys the settings already available to him and those still proposed, then takes a separated finite-type scheme $f:X\to e=\operatorname{Spec}k$ with torsion coefficients prime to the characteristic. The exceptional inverse image $f^!\Lambda_e$ plays the role of a complex of singular chains, and its compactly supported direct image supplies the corresponding relative homology object. The argument depends on global duality and, under the stated constructibility and coefficient conditions, biduality with respect to a dualizing complex. The example $\Lambda=\mathbb Z/n\mathbb Z$ makes the base coefficient ring itself a possible dualizing object; the text then mentions the analogous $\ell$-adic setting and a relative extension under cohomological local triviality, with smooth morphisms as an example. The pen correction on site 59 reads ``no need to truncate'': the initial transcription had failed to apply the replacement e and cancellation of on, so this is a repaired transcription, not a grammatical emendation. The typed hat over the f in the structural-morphism parenthesis is retained, and the tiny s/S coefficient subscript on site 60 remains doubtful. Part B recalls Artin--Mazur's autoduality work for relative Jacobians and its use in a component-group duality problem, then traces Grothendieck's broader geometric class-field-theory program to Serre's treatment of Rosenlicht--Lang theory and pro-algebraic groups. His remarks on possible proper-flat hypotheses are explicitly provisional; they should not be read as a theorem with fully stated assumptions. The source's paired operations are $f_!$ and $f^!$, whereas the witness prints the exceptional inverse image twice. The source also says ``flat of relative'', retains the final t of Rosenlicht, and says ``residual field''. The handwritten replacement of ``he suggested'' by ``I sketched'' changes the attribution of the proposed programs and must be retained. Their arithmetic forms are to follow by descent from geometric statements over an algebraic closure. This prepares part C's proposed complex of generalized Jacobians and the subsequent global and local duality sections. The partial pen cancellation in ``residuefields'' leaves a literal terminal ds outside its strokes; the readers preserve that physical reading, while the witness silently supplies the cleaned word.

\subsection*{Appendix 11: A homological generalized Jacobian (part C)}
Cote 134-2, sites 60--61 (folios 45--46), recalls a complex $J_{*X/k}$ whose degree-zero part contains the generalized Jacobian and whose proposed universal role is expressed by $H^i(X,G_X)\simeq\mathbb E xt^i(J_{*X/k},G)$. The residual-complex construction and local generalized Jacobians are meant to produce a homology object representing cohomology with commutative algebraic-group coefficients, continuing part B's class-field-theoretic motivation and preparing part D's continuous duality. The author explicitly revises the old Zariski-topology construction toward étale and finally f.p.p.f. cohomology, and changes the anticipated range from dimension $n$ to $2n$, allowing zero-dimensional components beyond $n$; these are a research program and recollections, not an established representability theorem. The base-change claim is deliberately only derived. At site 61 the literal ``homeomorphism'' $X\to\underline H_0$ cannot be read as an asserted topological equivalence establishing the universal algebraic mapping property: its intended morphism and topology require specification. Neither that word nor the heuristic bounds are repaired in the readers. The witness materially changes the notation by making the argument of Spec a subscript, printing capital O for the homology index 0, moving the typed i inside the handwritten parentheses around J, dropping the added coefficient G before the question mark, and replacing the plural pen word ``localizations'' by a singular spelling. The edition preserves the visible formulas, the abbreviation ``comm.'', and the original punctuation, including a closing parenthesis after Alb; the two unrecovered pen traces remain marked.

\subsection*{Appendix 12: Global duality and geometric class field theory (part D)}
Cote 134-2, sites 62--65 (folios 47--50), begins with prime-to-characteristic duality for a smooth separated finite-type scheme of dimension$d$, then asks for its flat-cohomological continuation in characteristic$p$. Curves are the motivating case: cohomology groups with finite coefficients acquire proalgebraic structure, and connected unipotent groups are dualized by Serre's Ext-one construction, while finite étale groups use Pontrjagin duality in degree zero. Combining the two in a suitable derived category suggests the continuous Serre-duality operator$\Delta$; Cartier duality$D$ upstairs and Serre duality downstairs should exchange proper and arbitrary supports, with shift[2] for a curve. This continues parts B--C's geometric class-field program, and the author says explicitly that the coefficient categories and functors must still be constructed. Site 62's example literally contains$\mu_m^{\otimes}$ instead of the preceding$\mu_n^{\otimes d}$; neither the changed index nor the absent exponent is repaired in the readers. Likewise, the opening assertion for an arbitrary sheaf requires the appropriate finiteness and dualizability hypotheses before its claimed perfect cohomology pairings follow. On site 63, the displayed shift[2] must govern the surrounding schematic identification of$f^!(\mathbb Q/\mathbb Z)$ with$\mu_\infty$: omitting that shift is not a separate degree-zero identification. Site 64 applies Lang's relation between constant-kernel isogenies and rational points over a finite field to obtain the additional shift[1], hence arithmetic duality with total shift[3]. Its$\pi_1^{ar}(J)$ must be read in that stated isogeny sense, not silently identified with the full arithmetic étale fundamental group. The printed$H_!$ has no degree star, and the duplicate closing parenthesis remains literal. The final warning about non-étale coefficients such as$\alpha_p$ introduces formal completion at the boundary in the compact-support triangle on site 65; the coefficient formalism and its formal direct-image notation remain part of the proposed construction. The witness changes the live plural ``categories'' to singular on site 63, drops the distinctions of italic Ext and the underlined RHom, and changes the lower-case base-field k to capital K twice on site 64. These differences matter to the interpretation of the functors and are preserved in the source-based edition. Part E then makes the boundary's local duality problem explicit.

\subsection*{Appendix 13: Local duality and the Greenberg construction (part E)}
Cote 134-2, sites 65--66 (folios 50--51), starts with a complete discrete valuation ring$V$, fraction field$K$, and residue field$k$, assuming either a coefficient-field lifting or a perfect residue field of positive characteristic. The sought object$R\underline\Gamma_K(F)$ should carry the continuous$k$-algebraic structure suggested by elementary examples and satisfy$R\Gamma_K(F)\simeq R\Gamma_k(R\underline\Gamma_K(F))$. Witt vectors and a cohomological Greenberg functor are proposed as the construction, with$Rj_*$ supplying the first step. This develops the formal-boundary term of part D and predicts a geometric duality shift[1], increased to[2] by Lang's construction when the residue field is finite. Three literal slips at site 66 require attention without emendation: the opening formula has$R\Gamma_k(F)$ on the left even though$F$ is on$K$; the inverse-image duality identity uses$D_S$ on both sides although$i^*F$ is on$s$; and the closed point is then called$p=\operatorname{Spec}k$ after the morphism$i:s\to S$ has been introduced. The corresponding types are recovered by reading$K$, $D_s$, and the closed point$s$ in the intended roles, but all three printed readings remain in every reader. The prime-to-residue-characteristic case is presented as already understood, with the mixed-characteristic primary case singled out as subtle. Remark (c) explicitly limits the author's confidence to finite coefficients; the proposed extension to abelian varieties and Néron models, including$H^1(K,F)\simeq\operatorname{Ext}^1(G',\mathbb Q/\mathbb Z)$, is marked with a question mark and ``without any guarantee''. The handwritten prime on the dual abelian scheme$F'$ is visible and restored, whereas the witness drops it; its blackened predecessor stays unreadable. The source also distinguishes the doubled R in the first displayed functor from plain R later, retains the lower-case k omitted by the witness in the comparison formula, and abbreviates ``char.''. The fully handwritten final request about Deligne closes this program; part F compares fppf and crystalline cohomology. No conjectural assertion here is promoted to a theorem by the corrected reader.

\subsection*{Appendix 14: The scope of flat cohomology (part F)}
Cote 134-2, site 67 (folio 52), revisits the hope that finite and continuous coefficients in fppf cohomology might recover a p-adic Weil cohomology. Serre's degree-one constructions and the generalized-Jacobian complex of part C motivate the hope, although the literal backward reference is to (B). Crystalline cohomology, whose properties Grothendieck attributes here to Berthelot, now provides his comparison standard. He argues that the part obtainable from algebraic-group coefficients stays within Frobenius slopes at most one, explaining the success for degree one and the difficulty already in degree two, then explicitly labels recovery of the entire slope-at-most-one part a conjecture and admits that he has lost contact with the current work. This is a retrospective assessment of a proposed program, not a proof of a precise comparison theorem. For mathematical interpretation, Frobenius slopes belong to the rationalized crystalline object over a suitable perfect base; torsion in the integral Witt-vector modules is not described by a slope decomposition. The source preserves its stated hypotheses and its closed bound at one, without importing a modern comparison statement. The witness changes the decisive degree-one expression from the visible $H^1$ to $H^i$, drops the closing quotation after the slope bound, expands the pen abbreviation ``char.'', and misspells Katz. The original handwritten superscript i, correction of ``parts of points'' to ``part of slope'', and Frobenius arrow are retained, while the blackened old exponent in the coefficient $\mu_{p^n}$ remains doubtful. This closes the duality and cohomology sequence of parts A--F; Appendix 15 returns to the information retained by a homotopy type and its locally constant higher stacks.

\subsection*{Appendix 15: Homotopy types and higher local systems}
Cote 134-2, sites 67--69 (folios 52--54), distinguishes the homotopical algebra performed inside a rich topos from the much poorer homotopy type obtained by localization. The retained information is expressed through locally constant higher stacks, fundamental higher groupoids, and ``special'' or ``n-galois'' categories, arranged in a triangle of proposed mutual reconstructions. This resumes Appendix 6's higher Galois reconstruction: exact functors should recover the fundamental groupoid from its category of local systems, with finite projective limits and arbitrary inductive limits again explicitly distinguished. The handwritten replacement of ``must'' by ``should'' at site 68 preserves the conjectural status of the comparison. Its dimensions also require care: the definition uses locally constant $(n-1)$-stacks but then writes functors into $(n\text{-}\mathrm{Cat})$, followed by lowercase $(n\text{-}\mathrm{cat})$ in the reconstruction formulas. With the usual meaning of n-Cat these targets are one categorical level higher than the preceding local-system description; already at $n=1$, category-valued systems differ from set-valued systems. The text does not settle that convention here, and the printed indices and capitalization remain unchanged. Site 69 makes the important qualification that a locally homotopically trivial situation can have an ordinary homotopy type, while étale topoi generally require a pro-n-groupoid and an accompanying ind-category of local systems. Its emphatic negative assertion about the étale topos should not be read as excluding every special scheme: the étale topos of an algebraically closed field already gives the elementary constant case. The witness loses the underline emphasizing pro and normalizes the parenthesized Cat/cat distinctions; the visible final d alone is canceled in interpreted, leaving the literal interprete. The isolated typed famliliar is corrected with apparatus because familiar occurs correctly later in the same paragraph. Appendix 16 then addresses the dissatisfaction motivating this section by retaining stratified constructible data instead of only locally constant data.

\subsection*{Appendix 16: Constructible data, gluing, and tubular neighborhoods}
Cote 134-2, sites 69--72 (folios 54--57), responds to the loss of structure lamented in Appendix 15 by retaining sheaves, complexes, or higher stacks constructible for a fixed stratification P. The initial quasi-compact and quasi-separated hypothesis is materially changed to quasi-coherent by the witness. Grothendieck distinguishes ordinary constructibility under operations from the stronger stability for one fixed stratification, proposing sufficiently fine equi-singular stratifications for the latter. He then analyzes two strata, a closed part $X_0$ and its open complement $X_1$, by gluing locally constant sheaves with a map $F_0\to\phi(F_1)$. The literal formula at site 70 contains an extra superscript star on the second inclusion symbol, in addition to its subscript star; the required gluing functor is $i_0^*i_{1*}$, and the later formula $p_*j^*$ confirms its variance. The witness introduces a further error by raising that lower star onto the index 1. The intended étale tubular neighborhood is motivated on site 71 by a conic neighborhood over $X_0$, with boundary fibration p and an inclusion j into the cut-open complementary piece; it yields $\phi\simeq p_*j^*$ and, by adjunction, a gluing map $p^*F_0\to j^*\widetilde F_1$. Two geometric conventions need clarification in the literal text: the conic bundle is called a bundle over X rather than over $X_0$, and the definition $\widetilde X_1=X_1\setminus(T\setminus X_0)$ removes the boundary if the earlier closed-cone convention for T is used, although the diagram requires that boundary to lie in $\widetilde X_1$. Deleting the appropriate interior or choosing compatible open and closed neighborhoods repairs the construction conceptually; no such repair is inserted into the readers. Site 72 offers the less symmetric but more canonical diagram using the punctured neighborhood and then T itself, distinguishing an actual inclusion from a map regarded homotopically as a fibration. The qualification about compact finite-dimensional smooth fibers is part of this geometric proposal, not a property of every replacement by an arbitrary fibration. Note 28 later comments on the étale-neighborhood construction. The witness also substitutes tilde-T for the visible tilde-X and changes two diagram relation signs; the scans show a tilde above two bars. The circled transposition of topological intuition is now represented faithfully, and the new section's circled spacing note reads espacer, not the initial tentative à part. Appendix 17 generalizes this gluing diagram to indexed homotopy types over many strata.

\subsection*{Appendix 17: Stratified diagrams and the fine homotopy type}
Cote 134-2, sites 72--75 (folios 57--60), generalizes the two-stratum construction to a geometric description of constructible sheaves and higher stacks by diagrams of homotopy types. Site 73 insists on taking the homotopy category of the category of diagrams, rather than merely forming diagrams in a homotopy category; note 29 supplies the later categorical perspective. For tame compact spaces with equi-singular strata, the proposed diagram consists of manifolds with boundary, fibrations with compact fibers, and inclusions of boundary pieces, from which the space should be reconstructed by gluing. The witness's singular space and its singular inclusion of the boundary fail to preserve the visible plural data. Site 74 extends the proposal to a morphism and to the six operations, then proposes a corresponding construction for excellent schemes through étale tubular neighborhoods. Its displayed list mixes $Rf_*,Rf_!$ with $Lg^*,g^!$, although $g$ is subsequently introduced as the restriction to locally closed unions of strata; the variance of the four operations must be read with a consistent choice of the ambient or restricted map. The baseline hyphens in $P_X$- and $P_Y$-constructible are not superscripts on the indices, as the witness makes them. The handwritten word replacing some natural is literally mutal, not natural. Refinement of the stratification leads to a fine homotopy type, intended to retain the whole system of constructible data rather than the ordinary homotopy type alone. The assertion at site 75 that passage to an inductive limit gives all of them requires the preceding constructibility restriction: refinement through finite suitable stratifications does not make every arbitrary sheaf locally constant on strata. The proposed strongest finiteness theorem says that the constituent ordinary types, or more cautiously their prime-to-characteristic profinite completions, should come from finite polyhedra or compact manifolds with boundary. This is a programmatic proposal here, with the characteristic-zero route contrasted explicitly with the difficulty of varieties resisting even birational lifting in positive characteristic. Appendix 18 turns from these geometric diagrams to the algebraic structures that could recover them; the later notes revisit both the diagram category and the gluing construction. The call after gluing is handwritten 30, not the initial transcription's50.

\subsection*{Appendix 18: Linearization and recovery of a homotopy type}
Cote 134-2, sites 75--77 (folios 60--61bis), asks which algebraic data could recover the geometric constructions of Appendices 16--17. Grothendieck first considers cohomology with variable twisted constant coefficients, its exact-sequence structure, cup products, and further operations, assuming the fundamental group and its coefficient category given. He then proposes the triangulated category of complexes with locally constant cohomology, together with derived tensor product and internal Hom, as a more satisfactory candidate, while explicitly changing have little idea to the cautious do'nt really know. Note 31 revisits whether these abelian data suffice. By contrast, he expects the full non-abelian system of categories of locally constant higher stacks to recover the fundamental infinity-groupoid; the claim depends on the earlier higher-categorical program, not merely on the list of ordinary cohomology groups. For a fixed stratification, the analogous proposal uses complexes whose cohomology is locally constant on each stratum and, as needed, the four operations for maps between locally closed unions of strata. The question about knowing what a triangulated category is concerns the adequate formalism and its non-commutative analogue; note 32 adds the derivator perspective. The final line of site 76 changes fibrations and cofibrations to fibers and cofibers in the author's spelling: pen es overwrites at, the remaining ions and old closing quotes are canceled, and new quotes close the shortened words. This physical correction is more specific than the initial transcription's deletion of the whole ending. The witness changes more simply to more simple. Site 77 cancels an N.B. crediting an early Quillen approach, ends the letter with a question about the Bloch--Quillen formula, and retains the typed signature. The subscription is literally your's, with its old comma overwritten and no surviving terminal comma; the readers preserve it. The following notes qualify several earlier claims and are essential to distinguishing the 1975 letters from their later annotated reprinting.

\subsection*{Notes 1--11: Retrospective qualifications of the Quillen letter}
Cote 134-2, site 78 (folio 62), opens the notes to ChapterI and the appendix. A long pen arrow and the renumbering move the notice about reproducing the Breen letters before the account of the unanswered Quillen letter; the diplomatic form retains the physical order and the corrected forms apply the move. Note 3 points toward tame topology in the Esquisse, while notes 4--5 qualify the earlier judgment of the Bangor group's relation to geometry. Note 6 restricts the earlier assertion to strict associativity. Note 7 calls the claim about existence and uniqueness of the functor rash and points to section 12 and note 17; its canceled final word is goes. Note 8 gives torsors and the classifying topos as another example, with an ordinary typed B rather than the witness's blackboard-style symbol. Note 9 recalls the Buffalo seminar. The initial transcription of note 10 as a fraction was materially wrong: the scan has a star with typed upper i and a looped handwritten lower ell. The witness instead gives a lower 2, while its following order is ell. In the source the latter position contains a clear typed 1 with a tall pen correction plausibly read as ell; that live reading remains doubtful, so the readers do not settle it by appeal to the intended composition order. Note 11 explicitly withdraws an earlier belief and refers to section 11. These retrospective statements should be read with the earlier tentative construction of successive layers of structure, not as part of its original uninterrupted argument.

\subsection*{Notes 12--21: Products, test functors, and primitive structure}
Cote 134-2, site 79 (folio 63, initially misread 53), collects substantial retractions. Note 12 gives an overlapping-circle counterexample; note 13 rejects the earlier assertion for positive dimensions and adds that even the postscript to section 13 remains inaccurate, referring onward to section 18. Note 14 confines the required fiber products to the category marked by a lower bullet and requires the relevant functors to preserve those products; the witness omits that bullet, thereby losing a mathematical distinction. Note 15 doubts the proposed equivalence even for groupoids or Gr-categories. Notes 16--17 revisit fibrancy and the functors into a homotopy model: even if every object is fibrant, a suitable uniqueness statement is needed before the resulting functors to (Ho) can be said to be canonically isomorphic. The discussion links this issue to existence and uniqueness of test functors out of the category of standard hemispheres and to section 90, then records the shift toward asphericity and contractibility structures. The hemisphere symbol has a visible curved meridian, and the witness's missing target M changes the displayed functor. Note 18 explains gerb as a rendering of gerbe and proposes stack of groupoids or Gr-stack. Note 19 explicitly rejects the earlier claim; note 20 warns that the suggested construction may give overly large categories; note 21 recommends retaining the primitive subcategory as additional structure rather than expecting to recover it from the ambient category alone. These notes qualify the Quillen letter's statements about coherators and their generators. The isolated it for if in note 16 is corrected only in the corrected forms and is recorded in the apparatus.

\subsection*{Notes 22--24: Picard stacks, resolutions, and derived structure}
Cote 134-2, site 80 (folio 64), supplies the background for the Breen letters. Note 22 explains lien or tie and Picard categories and stacks, referring to Giraud and to Deligne's account in SGA 4, XVIII,1.4. The source visibly changes plural categories to singular category and its verb phrase to the typed is a groupoi; the missing final d is an isolated slip confirmed by groupoid later in the paragraph. The stated range of the presumed higher-stack formalism is changed from n=2 to $n\leq3$, not to the witness's n=3. The proposed connection with truncated complexes requires the strict commutativity condition that the symmetry on $L\otimes L$ be the identity; the note assumes without further explanation that the appropriate condition extends to higher Picard stacks. That assumption should not be read as a supplied higher-categorical proof. Note 23 recalls the Mac-Lane resolution by sums of free abelian sheaves on finite products of M, relating Ext invariants to the cohomology of the slice topos. The word left is a typed interlinear insertion; the enclosing curves around the resolution's name are placement marks rather than parentheses. Note 24 makes the search for an adequate version of Verdier's triangulated-category formalism part of the projected later work and adds a handwritten pointer to section 69 and derivators. The spelling developped is clearly typed and remains unchanged; the cut-off rathe at the scan edge is not completed from the witness.

\subsection*{Notes 25--32: Coherence, homotopy invariants, and later references}
Cote 134-2, sites 80--81 (folios 64--65), finishes the retrospective notes. Note 25 broadens uniqueness to account for the choice of a coherator, recalling section 13. Note 26 corrects the theorem's name to Dold--Kan; note 27 connects the earlier structure with Whitehead's crossed modules and with a then forthcoming generalization by Conduché. Note 28 records Cox's work on algebraic tubular neighborhoods while explicitly withholding judgment on whether it supplies the precise stratified dévissage machinery proposed in Appendix 16. Notes 29--30 connect the diagram construction with a derivator and point to tame topology in the Esquisse. Note 31 reports a negative answer to the reconstruction question in Appendix 18 using simply connected rational homotopy types: the rational cohomology ring alone does not determine such a type. The argument's asserted identification of that ring with all the contemplated data must be distinguished from a statement about any richer, fully specified enhancement of the derived structure; the note does not provide such an enhancement or a comparison theorem. A handwritten addition qualifies the assertion by finite-dimensionality of each cohomology group and asks whether a counterexample remains for a homotopy type of finite type, without defining that latter condition here. The addition begins At least, not the witness At last. Its long connector, canceled false start, and reference-marker layers are retained. The apostrophe in it's is an authorial pen addition, not a deletion, so the witness's its fails to preserve his final spelling. Note 32 substitutes the derivator formulation for an earlier description of the non-commutative problem and recalls note 24. These qualifications provide the author's later perspective on the letters and should accompany their claims about reconstruction, uniqueness and homological algebra.

\subsection*{14. Homotopical algebra and the limits of present topology (pp. 1--2)}
Grothendieck opens the continuation of the letter to Quillen by distinguishing the study of spaces of maps and their successive homotopies from the study of spaces of homeomorphisms, immersions, embeddings and fibrations. His terms ``wildness'', ``isotopy'', ``stratified structures'' and ``tubular neighbourhoods'' delimit problems for which the familiar category of topological spaces seems poorly adapted; the contrast with piecewise linear structures is that their maps are too sparse, notably because real-valued piecewise linear functions are not closed under multiplication. This develops the complaint recalled explicitly from the letter, p.~1, and supplies the contrast on which section~15 builds its more favorable account of localization and homotopy types. The interlinear revisions narrow the opening description to ``maps between spaces'' and replace ``numerical'' by $\mathbb R$-valued (p.~2). The witness differs in two mathematically relevant places: the typescript says ``pro-topoi'' on p.~1, where arXiv has ``pretopoi'', and arXiv retains ``numerical'' without applying its handwritten replacement on p.~2. These readings are established by the scan; the edition does not identify the two notions or silently repair either authorial formulation.

\subsection*{15. Models of homotopy types and the place of stacks (pp. 2--4)}
The argument passes from the failure of the Hauptvermutung as a literal assertion about triangulations to its much weaker homotopical analogue: appropriate localizations of topological spaces and semi-simplicial sets yield the same homotopy category. Grothendieck compares CW-spaces, Kan complexes and Serre's cubical complexes, then proposes stacks, previously called ``$\infty$-groupoïdes'', as a third description particularly close to the homotopy groups themselves. His analogy with chain complexes and the Dold--Puppe equivalence makes the intended role of noncommutative structure explicit; section~16 immediately generalizes the motivation from spaces to topoi. The distinction between describing and computing invariants matters: on p.~3 he twice strikes ``description'' and writes ``computation'', corrections which the arXiv witness does not apply. On p.~4 a marginal note dated 3.5 begins ``This is an error'' and retracts the claim that stack cohomology has no immediate description from the primitive structure, referring instead to the nerve and to forgetting structure. Broad pencil strokes across this note make its final cancellation status uncertain, and one verb remains doubtful; both uncertainties are preserved. The arXiv footnote is substantially longer than this marginal wording and must not be substituted for it as though it were a literal transcription.

\subsection*{16. Stacks over topoi and noncommutative cohomological algebra (pp. 4--6)}
The motivation for stacks is enlarged from section~15's models of homotopy types to coefficients for cohomology over an arbitrary topos. Starting with Giraud's $0$-, $1$- and $2$-dimensional theory, Grothendieck treats $n$-stacks as coefficients whose ``integration'' is evaluation at the final object, with the homotopy groups of the resulting stack supplying cohomology after a dimension shift. His letters to Breen and Mme Raynaud's work on Lefschetz comparisons provide the geometric motivation: higher noncommutative cohomology should unify homological and homotopical comparison statements. The passage to infinite-order stacks is presented as a qualitative enlargement of perspective, and the following section asks how much of the homotopy-theoretic side ordinary categories can already model. The scans matter for the formulas: p.~4 has $H^i$, not the witness's $K^i$; p.~5 has $R\Gamma F_*$ and complexes $F_*$, where the supplied arXiv comparison reports a missing $F$ and a $j$ in place of the star. The handwritten qualification ``[non commutative]'' on p.~5 must also remain visible. These are witness differences, not authorial mathematical corrections introduced by this edition.

\subsection*{16$'$. Categories and the ``impressive bunch'' of models (pp. 7--9)}
Grothendieck associates a homotopy type to a small category both through its presheaf topos and through its nerve, comparing the first construction with Verdier and Artin--Mazur pro-homotopy types. He recalls and rejects his former guess that passage to the enveloping groupoid preserves this type: Quillen had pointed out that ordinary categories realize arbitrary homotopy types. This corrects the initial groupoid intuition while continuing sections~15--16's search for models, and motivates the later study of the ``impressive bunch'' and of coherators. A second proposed shortcut is itself explicitly retracted: the margin of p.~7, dated 5.3, says that the asserted weak equivalence $K\longrightarrow NN'(K)$ is false and directs the reader to section~24, pp.~21--22. The note is shorter than its arXiv counterpart; its call number is (18), and the section has been renumbered $16'$ rather than left as $16$bis. Page~8 proposes cofibrations injective on objects and arrows with a further call (19), then contrasts general $n$-stacks with $n$-Gr-stacks and their truncated homotopy types. Its wording ``left derived functors'' for the projective-limit cohomology construction should be treated as an authorial formulation requiring care, not silently changed to the usual right-derived terminology. The final paragraph on p.~9 introduces the dependence on a coherator and the universal-algebra question, despite professing little inclination to pursue it.

\subsection*{17. The Artin--Mazur cohomological criterion (p. 9)}
This short section supplements the unification program by recalling that weak equivalences can be detected through cohomology with twisted coefficients, including noncommutative $H^0$ and $H^1$ as well as the commutative higher groups. Grothendieck explains how this permits étale cohomological information to constrain homotopy types, while wondering whether its success discouraged the development of higher noncommutative cohomology. The reference back to section~15 is the author's own handwritten cross-reference and is retained even though the immediate unification discussion occupies section~16 in the present numbering. The section bridges that motivation and the concrete corrections of section~18. Its handwritten insertion reads ``(namely, unification)'', not the arXiv ``messy unification''; the final sentence deletes ``that'' and says ``their observation'', distinctions that are obscured by the witness. The section number itself is corrected from 18 to 17 on the scan.

\subsection*{18. Corrections to the coherator and higher coherence (pp. 10--11)}
The return to the letter (L) is explicitly corrective. Grothendieck withdraws the existence of arbitrary binary amalgamated sums in $\underline B_o$: the geometric intuition of gluing contractible cells along embedded subcells had concealed the presence of nonmonomorphic maps, including the codiagonal. He also revises the construction of $\underline C_{n+1}$ from $\underline C_n$ by insisting on the total set of new arrows, and rejects the assertion that intersections of maximal subcells must be subcells. The marginal drawing on p.~10 makes the latter failure concrete: three regions $f,f',f''$ meet so that $f\cap f''=\{s,t\}$. The scan therefore supplies a two-point intersection, not the quite different formula in the initial transcription. A separate correction explains why associativity beyond primary composition cannot be represented by a single homotopy arrow: the sources and targets are themselves only homotopic, so squares, cubes and higher coherence data intervene. Page~11 names the dual categories left and right coherators, then connects the proposed non-groupoidal stacks to Benabou's bicategories and their cellular, topological interpretation. This prepares section~19's realization and singular-stack adjunction. The source wording $L\to K$ and $M\to K$ beside $K\amalg_L M$ on p.~10 is retained as an apparent authorial slip; it does not give the usual span for that pushout. Calls (11) and (14) on p.~11 are present without the explanatory footnotes printed in arXiv.

\subsection*{19. A provisional itinerary and the sheaf temptation (pp. 12--13)}
The first two items of the itinerary are the definition of $\underline B_o$, its dual and the coherators, followed by a realization/singular-stack adjunction with topological spaces. Unlike arbitrary presheaves on a left coherator, Gr-stacks must carry the designated amalgamated sums to fiber products. The realization functor nevertheless extends to all presheaves by preserving inductive limits; its values on basic cells are prescribed, while extensions to the coherator require choices. The anticipated independence of those choices is therefore tied to an as-yet undeveloped notion of equivalence between Gr-stacks. On p.~13 Grothendieck asks whether the exactness conditions could instead be sheaf conditions for a topology, which would make Gr-stacks a topos, and proposes weak equivalences detected by realization, monomorphisms as cofibrations and a lifting-property definition of fibrations. These are programmatic proposals rather than established results. The ``quick glance'' suggesting a topology is retracted immediately in section~20; even the concluding sentence about the inclusion, revised to composite, functor is already struck on p.~13. The source has ``one reward'', not the witness's ``one ward'', and handwritten ``as cofibrations'', not ``for cofibrations''. Its call (13) contains no footnote text on this page.

\subsection*{20. Model categories, sites and a retracted canonical topology (pp. 14--15)}
The section opens by correcting section~19's ``quick glance'': the proposed topology on the coherator would make the component functors fiber functors and force componentwise inductive limits, which the author recognizes as false. The category of categories supplies a warning and a new candidate construction. Grothendieck considers families of functors whose nerves cover every finite composable string, proposes identifying them with universally effectively epimorphic families, and associates to a category $A$ the topos of sheaves on $(\mathrm{Cat})_{/A}$. The handwritten margin immediately marks the asserted equivalence as false and points to section~24; the digits are readable, while pencil hatching leaves the note's cancellation status uncertain. Page~15 broadens the proposal to a natural site structure on model categories, with the homotopy type of $X$ identified with that of the induced site or topos over $X$. Its further presumption that the topology should be canonical is itself contradicted by a dated margin note absent from the arXiv comparison. These local corrections are essential to reading the section as an evolving investigation. The four arrows in the diagram of effective, universal and ordinary epimorphisms are handwritten, as are the topos tildes in $(M_{/X})^{\sim}=M^{\sim}_{/X}$; the object labels are typed. The inquiry leads into section~21's wider inventory of models rather than establishing a general theorem about canonical sites.

\subsection*{21. Multicomplexes, symmetry and restricted coherence (pp. 15--16)}
The inventory now includes symmetric variants of semi-simplicial and cubical complexes and their multiple-complex analogues. A diagonal complex already describes the intended homotopy type, but Grothendieck explicitly pursues the variety of algebraic presentations, not a unique most convenient model. He compares the passage from complexes to multicomplexes with that from ordinary categories to higher categories or stacks, recalling Quillen's suggestion of defining higher categories through multiple semi-simplicial complexes. The return to the coherator is concrete: the coherence diagrams examined in Benabou's bicategory axioms use restricted amalgamated sums describable through multisimplices, suggesting a smaller starting category $\underline B_o$ than the generous one in the original letter. The handwritten cross-references on p.~16 revise the old page references to sections~8 and~9, and the scan preserves the tense and spelling ``This I ckecked'', where arXiv reads ``This I check''. The suggested nerve-like functor from $n$-stacks to $n$-semi-simplicial complexes is a goal for later investigation. Section~22 next reverses direction, seeking simpler special cases of categories, beginning with preorders and ordered sets. The uncertain final singular/plural in ``connection(s)'' on p.~15 remains an image-edge problem and has no mathematical consequence for this comparison.

\subsection*{22. Ordered sets, realizations and maquettes (pp. 16--19)}
Preorders and ordered sets offer simpler categorical presentations of homotopy types. He relates ordered strata in conically stratified spaces to a realization whose barycentric simplices are finite totally ordered subsets, or flags. The resulting realization is canonically compared with that of the nerve. On p.~17 the corrected chain is $(\mathrm{Ord})\hookrightarrow(\mathrm{Preord})\hookrightarrow(\mathrm{Cat})\hookrightarrow(\mathrm{Ss\ sets})$, not the reversed first pair in arXiv. He cautions that the adjoint passage from a category to its object preorder can lose homotopy information, as the one-object case shows; the handwritten qualification ``but not Fl(A)'' is indispensable to that example. Page~18 gives a second realization through a nonseparated topology on the underlying ordered set and identifies sheaves with covariant functors, hence presheaves on the opposite order. Reversing the order preserves the flags and yields homeomorphic geometric realizations, while the two associated presheaf topoi need not be equivalent. The scan confirms hats in the latter claim and the author's rounded inclusion sign for the initial order comparison. Finally, maquettes are described by vertices and finite simplices and embedded in ordered sets through their face posets. The body claims that this category lacks finite products, while a dated marginal note adds that the notion is not an algebraic structure species; the hatching and call number of that note remain uncertain. Section~23 consequently shifts attention from satisfying Quillen's full package to constructing and localizing a natural functor into $(\mathrm{Hot})$. The edition preserves these claims and reservations as stages of the investigation, without treating the later witness's explanatory footnotes as text written on these pages.

\subsection*{23. The functor to homotopy types and a cellular core (pp. 19--21)}
Grothendieck separates the initial construction of a natural functor $M\to(\mathrm{Hot})$ from the further demand that $M$ carry a Quillen model structure. The functor would define weak equivalences by inverse image of isomorphisms; only then would one choose monomorphisms as cofibrations and test the resulting lifting-property definition of fibrations. The proposed site-theoretic construction uses the homotopy type of the induced site over each object, provided its pro-homotopy type is essentially constant. The canonical-topology suggestion is already qualified by a handwritten retraction pointing to section~24. Page~20 adds a structural requirement, universal and disjoint coproducts, and presents an algebraic species by a category $\underline C$ with finite projective limits, its dual $\underline B$, and the anticipated equivalence $\mathrm{Ind}(\underline B)\simeq M$. This connects localization to finite presentations and to the question of which exactness properties are already visible in $\underline B$. Page~21 then isolates a smaller category $\underline B_{\mathrm{oo}}$ of basic cells and monomorphic boundary maps: the ordered set of proper subobjects should realize the boundary sphere of the corresponding cell. That construction uses section~22's ordered-set realizations and develops section~18's correction of the initial gluing intuition. The distinction between the broad species-level proposal and the geometric cell data remains heuristic. In the scan, the long p.~20 insertion includes ``can be interpreted'', omitted by arXiv, and the identity-functor arrow is drawn with two shafts; p.~21 preserves the printed slip ``of cubical'' and the typed replacements ``define'', ``family'' and ``sphere''.

\subsection*{24. Truncated nerve topologies and the failure of uniqueness (pp. 21--24)}
This section supplies the correction repeatedly anticipated in the earlier margins. Covering all composable pairs does not cover arbitrary strings: Grothendieck defines a descending sequence $T_N$ by epimorphic families of $N$-truncated nerves, identifying the first two levels with universal epimorphy and universal effective epimorphy. The comparison lemma relates the associated sheaf topoi to presheaves on truncated simplex categories. The finite levels retain only skeletal information, so the full nerve topology $T_\infty$ is needed for the proposed homotopy-type construction. The claim on p.~22 that an $N$-skeleton has the same homotopy and cohomology invariants through dimension $N$ requires caution at the endpoint: the $N$-skeleton of an $(N+1)$-simplex is an $N$-sphere, although the full simplex is contractible. The edition preserves the author's inequality rather than silently strengthening or repairing this intermediate assertion. The same loss of information explains the failure of $K\to NN'(K)$ discussed in section~16$'$: $N'$ factors through the two-dimensional data, while the nerve itself can still induce the desired equivalence after localization. Page~23 then replaces simplicial diagrams by commutative cubical diagrams, producing a strictly coarser topology with a different associated topos but the same intended homotopy interpretation. Multiple complexes suggest infinitely many further suitable topologies; the conclusion is that no unique topology is privileged. This leads directly to section~25's search for an equivalence relation among them. The scan's slice notation on p.~22 places $/A$ outside the parentheses ending after $T_N$, followed by an additional handwritten closing parenthesis; the earlier transcription's placement inside the subscript was corrected, while the extra parenthesis remains as written.

\subsection*{25. Hot-equivalence of topologies and global asphericity (pp. 24--25)}
The lack of a privileged suitable topology leads to a proposed equivalence relation. For comparable topologies $T\geq T'$ on $M$, Grothendieck compares the induced morphism of sheaf topoi and asks that every induced morphism over an object $A$ preserve homotopy type. He calls the resulting property ``globally aspheric'', strengthening the familiar cohomological notion of global acyclicity by allowing noncommutative coefficients. The primitive relation between comparable topologies is not yet an equivalence relation; taking the relation it generates is a first proposal, and he asks whether a coarser useful notion could do better. The intended invariant is a canonically associated functor $M\to(\mathrm{Hot})$, perhaps strong enough to recover the topology class itself. The analogy with Giraud's recovery of a topology from its associated subtopos motivates, rather than proves, a hoped-for correspondence. Page~25 narrows the aspiration to a single natural Hot-equivalence class for suitable species, admitting that even $M=(\mathrm{Cat})$ has no intrinsic characterization yet. Section~26 then reformulates the problem through small subcategories. Two revisions on p.~24 are essential: ``inverse image'' becomes ``direct image'', and the prime is moved so that the inclusion is from $T$-sheaves into $T'$-sheaves. The vertical margin note actually says ``just means'' and ``namely contractible'' with the qualification ``in the Čech sense''; the initial transcription's ``weakly contractible'' was not a literal reading. Its pencil shading leaves cancellation status uncertain. The body ends with the typed slip ``he'' for ``be'', and its final ``this'' is visibly canceled before the continuation on p.~25.

\subsection*{26. Test categories, test functors and the passage back to $M$ (pp. 25--29)}
The proposed intrinsic characterization of suitable topologies is reformulated through small categories $B$ inside a model category $M$. The cellular cores of section~23 and the simplex and cube categories of section~24 suggest that the presheaf category $\widehat B$ should itself describe homotopy types: its natural functor to $(\mathrm{Hot})$, computed from categories of objects over a presheaf, should be a localization. Grothendieck separates this internal requirement on $B$ from the eligibility of a functor $B\to M$; marginal terminology on p.~27 calls them ``test categories'' and ``test functors''. The comparison lemma gives an equivalence between the induced sheaf topos and $\widehat B$ under either of two fullness hypotheses. The decisive qualification comes on p.~28: localization of that sheaf category is not yet localization of $M$ itself. Generation by strict epimorphisms is considered as a way of making $M\to\widehat B$ fully faithful, but the author then proposes dropping both fullness requirements and directly demanding that $M\to\widetilde M\to\widehat B\to(\mathrm{Hot})$ be a localization. For categories $\underline B$ of finite presentations he further hopes for a localization onto homotopy types of finite type. The questions whether equal weak equivalences determine the functor, and whether suitable choices form one natural class, feed sections~27--28. The margin of p.~27 actually says ``same question arises'' for Hot-equivalence of topologies; ``arises'' was absent from the initial transcription. Its wording is now readable, but pencil hatching still leaves the note's cancellation status uncertain. The call after objects ``over a'' on p.~26 is (23), not (3).

\subsection*{27. Geometric realization as a third construction (p. 29)}
Calling the preceding discussion ``thin air conjecturing'', Grothendieck warns that both test categories and test functors may need stronger conditions. He places geometric or ``topologic'' realization beside the two constructions through topologies and through test functors. Semi-simplicial and cubical complexes offer comparatively direct realizations, hence so do the categories embedded in semi-simplicial sets. For Gr-stacks, the prescribed realization of the primitive $\underline B_{\mathrm{oo}}$ structure has to extend to a left coherator, and that extension involves the arbitrary choices already encountered in sections~18--19. Rather than immediately eliminating them, he compares this indeterminacy with the variety of test topologies and functors developed in sections~24--26. The paragraph thus states the next intended problem, relating realization to the two test constructions; section~28 immediately postpones it in favor of rigidity of $(\mathrm{Hot})$. The scan preserves ``topologic'', where the witness expands it to ``topological'', and the word ``is'' in ``is that a suitable'', where the witness has ``it''. Its call (28) supplies no text of the witness's retrospective footnote. The typed insertions ``localization'', ``left'' and ``between'' sharpen the objects of the comparison; the two round letters in $\underline B_{\mathrm{oo}}$ and the primitive graph designation remain literal rather than silently replaced by a joined infinity glyph.

\subsection*{28. Rigidity of $(\mathrm{Hot})$ and the first modelizers (pp. 30--33)}
The intended return to geometric realization is suspended by an ``Assumption'': the category of self-equivalences of $(\mathrm{Hot})$ and their natural isomorphisms should be equivalent to a point. Grothendieck separates its two demands, every self-equivalence is isomorphic to the identity and the identity has no nontrivial automorphism, while wondering whether arbitrary natural endomorphisms could also be admitted. This is asserted as a plausible assumption, not proved; it supplies the desired uniqueness up to unique isomorphism of a localization functor once its weak equivalences are known. A ``modelizing category'' or ``modelizer'' is then a category $M$ with a saturated class $W$ whose localization is equivalent to $(\mathrm{Hot})$. A model-preserving functor carries weak equivalences to weak equivalences and induces an equivalence of localizations, equivalently commuting up to isomorphism with the type functors. This formalizes sections~26--27's questions without yet specifying cofibrations and fibrations. The diagram on p.~31 compares ordered sets, preorders, categories, semi-simplicial sets and cubical sets; the category-of-elements functors $i_A$ lead to right adjoints $j_A$. Page~32 carefully distinguishes those right adjoints from the familiar nerve functors, contrasting infinite categories of simplices with finite ordered face structures when degeneracies are omitted. Mixed multicomplexes then enlarge the diagram and motivate the test-category definition in section~29. The mixed-product display on p.~32 visibly has $\Delta$ in both factors, although the adjacent prose says $p$ times simplicial and $q$ times cubical; this apparent authorial symbol slip is preserved. The same page changes ``modelators'' to ``modelizers'' by hand, a correction the witness does not apply, and contains a small branching diagram omitted there. The p.~30 cross-reference is to section~26; calls (28$'$), (29) and (30) are present without the witness's expanded footnotes.

\subsection*{29. The basic modelizer and a stronger test condition (pp. 33--36)}
The category of small categories now becomes the foundational modelizer: $(\mathrm{Hot})$ is its localization, and weak equivalences are described through the induced morphisms of presheaf topoi and the Artin--Mazur cohomological criterion recalled from section~17. A small category $A$ is provisionally a test category when $i_A\colon\widehat A\to(\mathrm{Cat})$, sending a presheaf to its category of elements, is model preserving; its presheaf category is an ``elementary modelizer''. This replaces a special preference for semi-simplicial sets by a comparison among all such presentations, while acknowledging the computational convenience of the nerve and the simplicial construction of pro-homotopy types. The general adjunction calculation of pp.~34--35 asks when both adjoints induce inverse equivalences after localization. Applied to $i_A\dashv j_A$, it leads to condition (T): every counit $i_Aj_A(B)\to B$ should be a weak equivalence. Grothendieck explicitly leaves open whether this strengthens the preceding definition, and proposes strengthening it if necessary; the call (33) does not contain the witness's retrospective answer. Products of test categories and direct comparison of elementary models become the next questions. The closing heuristic on p.~36 would test (T) by a fibration with contractible fibers, either after taking nerves or directly in $(\mathrm{Cat})$; section~30 immediately attacks the model-category assumption behind the latter route. On p.~35 the source twice writes that $M'$ or $M$ is saturated, although the discussion concerns the arrow classes $W'$ and $W$; these apparent symbol slips remain unchanged. The typed ``elementary normalizer'' on that page likewise remains source text. The scan's baseline slash in $M/_{b}$ on p.~34 is preserved, and its reference reads ``section 28'', not the witness's page reference. The broad emphasis of the witness also differs from the selective underlining and marginal double bar of p.~33.

\subsection*{30. From the model-category obstruction to asphericity (pp. 36--38)}
The attempt to treat every monomorphism of categories as a cofibration now runs into the two elementary horns drawn on p.~37. Lifting against both would force the arrows of a fibrant category to have left and right inverses, leaving groupoids, whose homotopy types cannot exhaust those of arbitrary categories. The objection concerns that proposed choice of cofibrations; the source does not prove the nonexistence of every possible model structure on categories. Grothendieck nevertheless keeps the expectation that elementary modelizers should admit closed model structures, and turns instead to the asphericity criterion introduced in section~25. A morphism of topoi is required to remain a weak equivalence after every localization over the target; when the topoi come from categories, generators reduce the condition to the asphericity of each slice $C'_{/a}$. Applying this to the counit $i_Aj_A(B)\to B$ from section~29 gives three equivalent conditions: the strengthened test condition, asphericity of every counit, and asphericity of $i_Aj_A(B)$ whenever $B$ has a final element. That last formulation opens the explicit interval argument of section~31. The scan's four bracketed ``[small]'' restrictions on p.~38 are genuine hand additions, including two absent from the witness. Its caret under the $A$ in $\operatorname{Hom}(A_{/a},B)$ is retained in its actual position without asserting a mathematical interpretation. The apparent inline hat on $f$ reported by the initial comparison on p.~37 belongs instead to the display beneath it; the inline symbol has been corrected to plain $f$. The word ``closed'' on p.~37 and the replacement ``aspheric'' on p.~38 are interlinear typing, although their placing strokes are handwritten.

\subsection*{31. An interval proof of the test-category criterion (pp. 38--41)}
The search for a direct proof that the simplex category is a test category yields three sufficient conditions. T1 asks that $A$ be aspheric; T2 asks that $A_{/a\mathrm{x}b}$ be aspheric for any two objects, with their product interpreted in the presheaf topos when necessary; T3 supplies an aspheric interval object $I$ with disjoint final subobjects $e_o,e_1$. For $F=j_A(B)$ with $B$ admitting a final element, the aim is a homotopy from the identity to the constant section. Diagram (D) on p.~39 expresses the cylinder $F\mathrm{x}I$ and its two endpoint inclusions. Given $u:a\to I$, the inverse image of $e_o$ determines a sieve $C_u$ in $A_{/a}$; a functor $f$ is extended by its own values on $C_u$ and by the final object of $B$ on the complementary objects. The sieve property makes the remaining arrows uniquely definable. This constructs the homotopy without T1, T2 or the asphericity of $I$. Those hypotheses enter the second step: the slice criterion of section~30 reduces the asphericity of $F\mathrm{x}I\to F$ first to $a\mathrm{x}I$ and then to $a\mathrm{x}b$, precisely T2. T1 then converts the weak equivalence to the final presheaf into absolute asphericity. Page~40 separately justifies the retract step through the Artin--Mazur criterion rather than claiming it is formal for an arbitrary arrow class. The scan gives genuinely hooked tails to both arrows into the cylinder in (D), contrary to the initial transcription; all diagram arrows are handwritten and its object strings are typed. The p.~39 reference is ``cf théorème p 46'', while the marginal comparison involving T1 and T2 retains an uncertain short connector. The live punctuation after ``checked'' on p.~41 is a colon, and the source's ``asperic'' remains uncorrected.

\subsection*{32. A revised program after the interval criterion (p. 42)}
Dated 15.3, the new planning page treats T1--T3 as convincing candidates rather than a finished characterization. Its first three tasks are to examine necessity, construct the familiar simplicial and cubical examples, and prove stability under products so that different elementary models can be compared. The fourth returns to section~30's unresolved closed-model question; the fifth restores the larger program of modelizers attached to algebraic structure species and the possible uniqueness of their weak equivalences. The author explicitly notices that stacks have receded into the background. A sixth task broadens the inquiry to a satisfactory replacement for Verdier's provisional triangulated-category notion, including the exactness, generator and cogenerator questions for $(\mathrm{Hot})$. Section~33 begins carrying out the first task and introduces the distinction between test and strict test categories. The scan matters to the direction of the program: item 5 says categories $M$, not the witness's $\widehat A$. Its call (9) is present without the witness's explanatory footnote. ``Gave to this notion'' and ``general'' are typed interlinear revisions, while ``ways of'' and the displaced ``of'' are handwritten. The canceled ``topos'' in item 4 is legible beneath both machine and pen deletion; the cropped completion of ``which'' and the old word after ``constructing'' remain doubtful.

\subsection*{33. Strict test categories, the Lawvere interval and a qualified theorem (pp. 43--47)}
The provisional terminology first separates test categories from ``strict'' test categories satisfying T1--T3. T1 is necessary by the final-category case of the criterion on p.~38. T3 admits a canonical candidate: applying $j_A$ to the two-object ordered interval produces a presheaf whose sections are sieves, hence the ``Lawvere element'' representing subobjects. This ties the explicit interval construction of section~31 to an intrinsic object of the presheaf topos. A proposed concluding statement about an intrinsic T3$'$ is completely canceled on p.~43. The remaining issue is T2. Pages~44--45 connect it with preservation of binary products by the localization functor: the comparison $A_{/F\mathrm{x}G}\to A_{/F}\mathrm{x}A_{/G}$ is a weak equivalence for all $F,G$ exactly when the representable-product condition holds, and the slice criterion in fact supplies asphericity. The blanket claim about finite inductive and projective limits in $(\mathrm{Hot})$ on p.~44 is stronger than the finite sums and products used here; its marginal (!35) must not be expanded into an unqualified theorem in the guide. The result assembled on p.~46 compares product-preserving localization, the localized adjunction $i_A\dashv j_A$, and T1--T3. Its later margin dated 9.4 is essential: condition (i) is said to be ``not strong enough'', and the left continuation proposes adding that $j_A=i_A^*$ is model preserving, with a reference to remark~3 on p.~174. The left note is readable but crossed by broad pencil hatching, so its cancellation status remains uncertain. The source says ``model preserving'', not the witness's ``equally modelizing''. The author also explains how to avoid assuming that $(\mathrm{Cat})\to(\mathrm{Hot})$ preserves binary products: formulate the condition directly for $i_A$ up to weak equivalence. Page~47 provisionally reserves ``test category'' for the strict case, with call (38), and characterizes elementary modelizers through presheaf-topos structure, localization and finite products; a wider modelizing-topos notion is deferred. The scan retains ``asphéric'' on p.~43, three rounded intersection caps rather than the transcription's wedges, typed interlinear additions to the theorem, and the actual oval marginal annotations. The two opposite bent arrows on p.~47 remain graphic annotations, without an invented reordering instruction.

\subsection*{34. Product-stable examples and the homotopy circularity (pp. 47--50)}
The first examples strengthen the sufficient hypotheses again. A final object gives T1$'$, and an actual object $I$ of $A$ with two disjoint sections gives T3$'$; if $A$ is also stable under binary products, T2 follows because each product is representable and its slice has a final object. This constructs test categories inside any suitable ambient category, even one that is not small, by selecting a small product-stable subcategory. The economical construction uses powers of an interval and only projections and constant endpoint maps as components. Its cubical category includes permutations and diagonals, so p.~49 distinguishes it from the more familiar semicubical category. Nonempty finite sets give a corresponding full simplicial construction, with affine simplices as its geometric interpretation. The notation on p.~48 literally says to include $I$ and $\Delta_o,\Delta_1$; these triangular symbols are retained even though the preceding discussion specifies endpoint maps $\delta_o,\delta_1$. Returning to the ordinary simplex and cube categories removes the convenient product-closure hypothesis. Grothendieck therefore tries to prove T2 by making every representable deform to the final presheaf. The decisive revision on p.~50 is his own detection of a circular argument: section~31's homotopy proof used $F\mathrm{x}I\to F$ as a weak equivalence, and that was established from T2. A separate condition, called (2), that every $a\mathrm{x}I$ be aspheric must therefore accompany the deformation-retract condition TH 1. Section~35 postpones this unfinished verification for the intrinsic theory of modelizing topoi. The formula $u_a=p_a c_a$ on p.~50 is retained in that visible order; with $c_a:e\to a$ and $p_a:a\to e$, the surrounding right-to-left composition convention would instead require $c_a p_a$ for an endomorphism of $a$. The scan also confirms the literal spellings ``axists'', ``oddurs'' and ``consequance''. The extra presheaf caret on the full simplicial category on p.~49 is present despite the witness omission, and the stepped mark on p.~48 is retained graphically in the diplomatic reading and applied as a paragraph break only in the master reading.

\subsection*{35. Local and total asphericity of topoi (pp. 50--55)}
The detour from section~34's unfinished simplex proof translates its conditions into the language of topoi. An aspheric topos is required to be $0$-connected and to have trivial cohomology with constant coefficients in every defined positive degree, including the noncommutative first degree. For an object or arrow, the induced topos supplies the corresponding notion. A locally aspheric topos has a generating family of aspheric objects; these generators suffice for testing whether an arrow remains a weak equivalence after every localization (p.~51). Proposition~1 identifies product stability of aspheric objects with their being aspheric over the final object. Total asphericity combines that property with local asphericity, and the interlinear ``non empty'' inserted on p.~52 is essential: without it, the preceding assertion that T2 already implies T1 would admit the empty case. The ensuing \v{C}ech argument gives absolute asphericity. For topological spaces the much stronger total condition becomes irreducibility, so a Hausdorff example must be a point. This sharpens section~33's distinction between a local supply of cells and the global product condition. Proposition~2 then separates the interval requirement from total asphericity. Its original ``totally aspheric'' hypothesis is visibly canceled, and p.~53 explicitly says that no local-asphericity assumption is being made. The remaining equivalence is between the Lawvere element being aspheric over the final object and the existence of an object with that property and two disjoint sections. The homotopy argument of section~31 supplies the converse, and localization upgrades the resulting weak equivalence to asphericity. A connected topological space provides a useful obstruction: the proof with direct summands of its Lawvere element shows that this element is disconnected precisely in the irreducible case, where it has two connected components. Consequently no topological space can simultaneously meet total asphericity and the interval requirement (pp.~53--54). The proposed term ``modelizing topos'' packages both, while a second, still conjectural characterization asks for a saturated localization equivalent to $(\mathrm{Hot})$, compatible with finite sums and products and with the Verdier--Artin--Mazur construction. The closing inquiry about morphisms $X\to\operatorname{Top}(B)$ seeks a categorical description closer to the given indexing category than a large inverse system of semi-simplicial objects (pp.~54--55). The source retains ``0-connexity'', ``independant'', ``Corollay'' and ``monomorphisme''. Its upper margin on p.~54 reads ``cf (34)'' beside a body call (39); the lower strictness-related note remains partly doubtful. These are retained source features, not expanded into the arXiv witness's retrospective commentary.

\subsection*{36. Shuffles, nerves and comparison with categories (pp. 55--58)}
The simplex proof postponed in section~34 is completed by separating its two tasks. The ordered map $\Delta_n\mathrm{x}\Delta_1\to\Delta_n$ contracts toward the last element, giving TH 1. For TH 2, the maximal flags of the product order---the ``shuffles''---are attached successively so that each new intersection is a subflag. A Mayer--Vietoris-type lemma then propagates asphericity from two subobjects and their intersection to their union (p.~55). The author's inserted ``and covering e'' is indispensable to that lemma. He immediately warns against replacing this argument by the observation that the product ordered category has a final element: that shortcut assumes the still-unproved compatibility between the category's homotopy type and that of its nerve. The literal typed replacement ``apshericity'' remains uncorrected in the reader. The adjacent handwritten call is (41), not the asterisk proposed by the first transcription. The comparison proposition of pp.~56--57 addresses the missing compatibility in a more general form. A functor $a\mapsto\Delta_a$ from $A$ to categories defines $\alpha(C)(a)=\operatorname{Hom}(\Delta_a,C)$. If each $\Delta_a$ has a final element, $A$ is aspheric, and $\alpha(\Delta_1)$ is aspheric over the final presheaf, choosing final elements defines an aspheric functor $\varphi:A_{/C}\to C$, natural in $C$. The choice of final rather than initial elements makes the extension to arrows possible. Its slices identify with the corresponding construction over $C_{/x}$, so it suffices to treat $C$ with a final element; applying $\alpha$ to the ordinary categorical contraction supplies the required interval homotopy. Thus $i_A\alpha$ preserves weak equivalences and induces the identity on $(\mathrm{Hot})$. The simplex case now needs only the verified TH 2, while the category of nonordered simplices is sent to the ordered set of nonempty faces of each finite set. A proposed sieve criterion for the interval condition follows on p.~58, but the entire corollary and its proof are diagonally canceled before the new date 21.3. Its uncancelled presentation in arXiv must not conceal that revision. The source diagram on p.~56 literally labels its upper objects $a',a$ rather than $\Delta_{a'},\Delta_a$; the reader preserves this abbreviation or slip. The canceled formula on p.~58 has $C_u\cap C_o$, not the witness's composition sign, and its first $F\to e$ occurrence is actually drawn with a mapsto tail. Section~37 turns the repeated argument into a new systematic account.

\subsection*{37. Homotopy intervals and the key comparison theorem (pp. 58--62)}
The new account dated 21.3 begins by identifying the gap in section~36's canceled sieve criterion: the choices $C_u$ must be functorial in $u$, so that they actually define a map to the proposed interval. The replacement of ``automatic'' by a claim restricted to the intended examples makes that qualification explicit (p.~59). Grothendieck now formulates the argument in a category with finite products and a class $W$ satisfying isomorphism, two-out-of-three and retraction conditions. The last condition is stronger than the earlier mild saturation convention. An interval has two disjoint sections and an aspheric projection, with the required base changes understood; the margin calls the terminology ``provisional'' and adds ``separating''. The first homotopy lemma contracts an object to a section and concludes universal asphericity, while explicitly observing that disjointness was not used. The comparison lemma then transfers this property along an endpoint-preserving map to an object equipped with a unit-and-zero multiplication (p.~60). The revised formulas consistently make endpoint 1 the unit and endpoint 0 the zero, although the earlier ``respectively'' sentence is not fully adjusted. The Lawvere element supplies the universal example. In its proof on p.~61, the pullback along $e_0\to L$ is still called a subobject of $e_1$; the edition preserves this apparent index slip. The resulting theorem separates two strengths of comparison for a functor $i:A\to(\mathrm{Cat})$ whose values have chosen final objects. The choices need not be preserved by the arrows of $A$. Part (a) equates asphericity of the evaluation $A_{/C}\to C$, its being a weak equivalence, and the resulting identity on localized homotopy types, with equivalent tests on categories having a final element. Part (b) asks for asphericity relative to $A\times C$, equivalently for the appropriate representable products, and reduces this stronger requirement to universal asphericity of $i^*(\Delta_1)$. Together with asphericity of $A$, part (b) implies part (a); the first remark explicitly doubts the converse (pp.~61--62). The proof again transports the canonical categorical contraction through the left-exact functor $i^*$, while nonemptiness of each $i(a)$ ensures separation of the two endpoints. A superseded proposed condition in part (b) is canceled by both machine and pen. The marginal reference to the sequel in volume~2 has a still-uncertain first page numeral, tentatively 64--65. Section~38 extracts the consequences for the Lawvere criterion and asks how far the final-object hypothesis can be removed.

\subsection*{38. The Lawvere criterion and generalized nerves (pp. 63--66)}
The theorem's main gain is the reduction to one interval condition: universal asphericity of $i^*(\Delta_1)$, together with asphericity of $A$, makes $i_Ai^*$ induce the identity on $(\mathrm{Hot})$. Both that interval condition and the localized conclusion make sense without chosen final objects in the categories $i(a)$, and p.~63 therefore anticipates a more general comparison theorem, while acknowledging that a direct evaluation map would then be unavailable. For the canonical choice $a\mapsto A_{/a}$, the right adjoint is $j_A$ and $j_A(\Delta_1)$ is the Lawvere element $L_A$. Corollary~1 consequently gives a sufficient criterion for a test category in the wide sense: $L_A$ is aspheric over the final presheaf and $A$ is aspheric (p.~64). The scan actually has a separate handwritten caret in ``unit element \ldots\ of $\widehat A$'' on p.~63; the initial transcription's bare $A$ was not a source slip. For general $i$, a crible in $i(a)$ determines the crible of arrows $b\to a$ whose images factor through it. This defines an endpoint-preserving map $i^*(\Delta_1)\to L_A$, compatible with intersection. Section~37's comparison lemma transfers the interval condition to $L_A$. Corollary~2 then makes both $i^*$ and $i_A$ modelizing and identifies the functor induced by $i^*$ with the canonical quasi-inverse induced by $j_A$ (pp.~64--65). Thus the localized functor does not depend on the admissible choice of $i$, although the resulting concrete presheaf models can be much simpler than those obtained from the often infinite slice categories $A_{/a}$. The ordinary nerve is the motivating example of this ``generalized nerve functor''. Corollary~3 makes the remaining verification practical: if $A$ has a final object $e$, $i(e)=\Delta_0$, and a homotopy interval $I$ admits an endpoint-preserving map $i_!(I)\to\Delta_1$, adjunction and the same comparison lemma give the desired asphericity of $i^*(\Delta_1)$ (p.~66). In the intended examples $I$ belongs to $A$ and $i(I)=\Delta_1$, so only asphericity of all products $I\mathrm{x}a$ remains to check. The reference to the canceled corollary is revised to the end of section~36. Page~64 literally says ``for a in J'' where the construction indexes by $a\in A$, and retains a stray typed $u$ beneath a final-object caret; these visible irregularities remain in the readers. Section~39 returns to total asphericity and refines the terminology accordingly.

\subsection*{39. Strict, local and wide-sense test categories (pp. 66--69)}
The return to total asphericity first reuses the homotopy lemma: if every representable admits a contraction to a chosen section through a fixed homotopy interval, then every representable is aspheric over the final object and $\widehat A$ is totally aspheric (p.~67). The proof does not use separation of the endpoints. The following assertion that separation follows from the other assumptions is too broad if those assumptions are taken alone: the terminal category with $I=e$ supplies a counterexample. The boxed and diagonally canceled paragraph had discussed this exceptional case; its removal is preserved, together with the still-live preceding assertion. This is one of the places where the revision marks matter to the logical reading. The terminology is now reorganized around the Lawvere element. A local test category has $L_A$ aspheric over the final presheaf; a test category is additionally aspheric; a strict test category additionally has a totally aspheric presheaf topos. Their presheaf categories are correspondingly local elementary modelizers, elementary modelizers and elementary strict modelizers. A general strict modelizer is characterized by its weak-equivalence class and a localization equivalent to $(\mathrm{Hot})$ that preserves finite products; further requirements concerning sums and infinite products are still being considered (pp.~67--68). The weaker adjunction criterion $i_Aj_A(C)\to C$ remains the definition of a test category ``in the wide sense''. The source's dismissal of the likely usefulness of that condition is an uncancelled parenthetical, while an adjacent faint marginal note is more favourable; its exact wording and final Delta superscript remain uncertain. The page's note call is 45. The later shorthand ``T1 + T3'' must be read with care: the earlier T3 on p.~39 asked only for an aspheric object with disjoint endpoints, and its margin used T2 to pass to asphericity over the final object. The operative interval condition in the present definition is explicitly relative. The terminology for general modelizing topoi is proposed with an express reservation that its connection to actual models has not yet been established (p.~68). The distinction between test and strict test categories is then made concrete. A first proposed small category is generated by interval powers together with $a_o\mathrm{x}I^n$; Grothendieck expects the missing product $a_o\mathrm{x}a_o$ in the presheaf category to fail asphericity. A more decisive construction takes any test category $A$ and a homotopy interval $I$: the slice $A_{/I}$ is again local test and aspheric, hence a test category, but is never strict because the two endpoint subobjects of its final object are nonempty and disjoint (p.~69). Total asphericity would force their intersection to be aspheric, so it cannot hold. The first construction remains an expectation, whereas the slice construction supplies the stated universal counterexample. The scan's opening insertion on p.~69 reads ``as'', giving ``and as maps \ldots\ all those'', rather than the witness-based ``those'' of the initial transcription. Section~40 next supplies the product-stability lemma that the preceding arguments have been using.

\subsection*{40. Products and relative weak equivalences (pp. 69--71)}
The product argument supplies a fact already used in section~37: if $A$ is aspheric, the projection $C\mathrm{x}A\to C$ is aspheric for every $C$. The slice criterion reduces first to $C$ with a final element and then to products $C\mathrm{x}A_{/a}$, which themselves have final elements (p.~69). Grothendieck strengthens this to stability of weak equivalences under cartesian products. It is enough to treat $f\mathrm{x}\mathrm{id}_C$ with $f:A'\to A$ a weak equivalence. Viewing this as a morphism over $C$ motivates weak equivalence relative to a base topos $Y$: the property must survive each localization by an object of $Y$, and testing on generators suffices. For categories over $C$, the test is the induced functor over every $c\in C$. Thus the product statement reduces to $C$ with a final element, where the projection square and two-out-of-three prove it (p.~70). The displayed localization arrow literally has target $U$ in $Y_{/U}\to U$; its described role requires target $Y$, so the reader preserves the visible letter and flags the mismatch here. The final square's upper handwritten label has lowercase $c$, another source detail retained as drawn. Universal weak equivalences strengthen this relative notion by requiring stability under every base change in $(\mathrm{Cat})$, not only localizations (p.~71). Their test on base-change categories with a final element suggests an analogy with trivial fibrations. Grothendieck proposes that fibering or cofibering functors with aspheric fibers, and more generally such functors whose transition functors are weak equivalences, might supply the categorical counterparts of Quillen's fibrations. The inserted qualification ``in the purely categorical sense'' limits the proposal. Cofibrations remain unidentified, but monomorphisms alone appear insufficient, prompting the revised references to section~30 and (9). These are prospective comparisons rather than an established model structure on $(\mathrm{Cat})$. The margin labels the terminology ``awkward'' and carries call (46). The section then deliberately returns to test categories, leaving section~41 to clarify exactly which assumptions the preceding proofs have used.

\subsection*{41. The guiding assumption and weakened saturation (pp. 71--72)}
The conjectural triviality of the category of autoequivalences of $(\mathrm{Hot})$ had encouraged defining a modelizer through its weak-equivalence class, with the functor to homotopy types following essentially uniquely. Grothendieck now emphasizes, by changing the final period to an exclamation mark, that none of the proved statements has used that guiding assumption (p.~71). The operative definition of modelizer has also changed: the initial requirement of strong saturation, meaning that $W$ consists of precisely the arrows inverted by localization, has been replaced by the three weaker conditions written in section~37. The authorial transposition under ``wrote explicitly down'' supplies ``wrote down explicitly'' in the applied reading; the diplomatic text keeps the original order and the pen curve. The source spelling ``drapping'' is left visible where the earlier transcription had conjectured an $o$ correction. Strong saturation can nevertheless be transferred between $(\mathrm{Cat})$ and an elementary modelizer through the established comparison functors, so one example suffices for all (p.~72). Grothendieck invokes the known closed model structure on simplicial sets and his cited proposition~1 in Quillen's I~5.5 for that example. His remaining question concerns the proof of that model structure: can the present methods replace the specialized simplicial argument by one applicable to general test categories, or at least strict ones? He explicitly postpones this question and returns to comparing test categories and completing the inventory of standard examples. The handwritten call following ``strongly saturated'' is (47), while the separate margin has an asterisk and a reference to underlined VIII whose short abbreviation remains uncertain. This transition motivates section~42's matching revision of the definition of a model-preserving functor.

\subsection*{42. Model-preserving functors and full subcategories (pp. 72--74)}
The weakened saturation convention of section~41 requires a corresponding change to the definition of model-preserving functor: in addition to inducing an equivalence after localization, $f:M\to M'$ must satisfy the equality $W=f^{-1}(W')$, so that it reflects as well as preserves the specified weak equivalences (p.~72). The revised text first returns to the unresolved modelizing status of ordered sets and preorders. It then distinguishes the equivalence of localized categories from a canonical comparison with the chosen category of homotopy types. The earlier rigidity assumption would supply such a comparison, but has not been established. A subsidiary question asks whether an endofunctor of $(\mathrm{Hot})$ that is the identity on objects must be naturally isomorphic to the identity (p.~73). These questions remain conjectural; they do not become consequences merely from the existence of the concrete comparison functors already constructed. The constructive result uses a full subcategory containing the models produced by a test category. In the preliminary abstract argument an endofunctor $h$ of a modelizer factors through $M'\subset M$ and induces an autoequivalence on the localization; the text then describes the induced inclusion and factor as quasi-inverse (pp.~73--74). That exact formulation is stronger than its stated hypothesis: their composite is the localized $h$, and an autoequivalence need not be isomorphic to the identity. For example, taking $M'=M$ and inclusion equal to the identity still leaves that distinction. In the intended application $h=i_Ai^*$, the key comparison theorem does provide the required natural comparison to the identity. The resulting proposition says that a full subcategory of $(\mathrm{Cat})$ containing all categories of elements $A_{/F}$ for a test category $A$ inherits a modelizer structure and has model-preserving inclusion (p.~74). Categories locally isomorphic to $A$ supply a useful reformulation of the required objects. This reopens the ordered-set problem: a suitable test category defined by an ordered set would put all its categories of elements inside $(\mathrm{Ord})$. Such a category could not be strict, by the earlier obstruction for topoi attached to topological spaces in section~35. The last remarks contrast localization stability for test categories with the stricter total-asphericity condition, and lead directly to the standard examples of section~43. The source retains its literal ``T1 + T3'' shorthand. Its interlinear ``inclusion'', ``modelizer'', ``the test category'', ``algebraic'' and ``a suitable'' are typed revisions; the canceled locally-ordered qualification is a genuine pen deletion.

\subsection*{43. Ordered, unordered and face-only simplices (pp. 75--78)}
The return to standard examples applies section~38's test-functor criterion directly. For ordered simplices, the inclusion $i:\Delta\to(\mathrm{Cat})$ is fully faithful and carries the interval to the ordinary two-object ordered category. Thus $i^*(\Delta_1)$ is the representable interval, and the previously proved asphericity of its products with simplices gives the test condition. Contractions to final vertices then give total asphericity and hence strictness (p.~75). The canceled preliminary proof is replaced by this more precise argument. For unordered simplices, a small skeleton of nonempty finite sets has finite products and a two-point interval. Sending a finite set to its ordered set of nonempty subsets supplies the test functor. On the two-point set this produces the span of the two singleton subsets into the full set; the endpoint-preserving map to $\Delta_1$ sends $\{0\}$ to 0 and both other subsets to 1 (pp.~75--76). The comparison square with ordered sets identifies the resulting functor on ordered simplices with the flag, or barycentric-subdivision, construction. The single untilded $\Delta$ in the unordered-category conclusion on p.~75 is a visible source omission; the later functor-domain notation has an additional uncertain wavy stroke, rather than a securely readable presheaf hat. Removing degeneracies gives a substantially different answer. Grothendieck first emphasizes that face operators already determine homology, cohomology and the fundamental groupoid, and hence the cohomological test for a weak equivalence of the associated complexes. Nevertheless, in the category $\Delta^f$ of strictly increasing simplex maps, the product of the zero-simplex with an $n$-simplex has a discrete category of elements with $n+1$ objects. It is therefore not aspheric when $n>0$ (p.~76). Removing the zero-simplex does not supply a representable interval with sections, since arbitrarily large simplices cannot map into a fixed smaller one. The more flexible attempt using the presheaf $i^*(\Delta_1)$ also fails the required relative interval test. The canonical Lawvere element gives the same two-point obstruction over the zero-simplex (p.~77). Thus $\Delta^f$ is not a local test category, while the possibility that it is a test category in the wide sense remains open at this point. The added name ``weak test category'' anticipates the next section's systematic terminology. The remaining weak condition asks for absolute asphericity of $i_Aj_A(C)$ for every category $C$ with a final object. Its first tests are asphericity of $A$ itself and of the Lawvere element, expressly distinguished from asphericity over the final presheaf (p.~77). A positive answer would let section~42's proposition construct further full-subcategory modelizers. To pursue the ordered-set case, Grothendieck considers the poset $J$ of nonempty finite subsets of an infinite vertex set. It has the desired local simplex structure and, for countable vertices, is a category of elements of an infinite face simplex. Its filtering property suggests absolute asphericity. The local test condition still fails because a minimal vertex has a one-point slice, which cannot be a test category (p.~78). The closing calculation identifies the category of elements of $L_J$ with pairs $K\subset T$, with $T$ finite and nonempty, ordered by $T'\subset T$ and $K'=K\cap T'$. It begins an explicit test of the surviving weak condition rather than asserting its conclusion. Section~44 opens the following day by announcing that the face-only simplex category does satisfy that weaker notion. The readers retain the source's shifted apostrophes, the canceled product/intersection argument on p.~78, and the distinction between the two inline typed round-letter indices and the displayed handwritten infinity.

\subsection*{44. Weak, local, strict test categories and test functors (pp. 79--88)}
The discovery that the face-only simplex category is a weak test category rehabilitates the weakest member of the three notions and prompts a systematic comparison. Part~a recalls the adjunction $i_A:\widehat A\rightleftarrows(\mathrm{Cat}):j_A$, with $i_A(F)=A_{/F}$, and compares eight formulations of the weak test condition: equivalence of localized categories, control of the adjunction maps, asphericity of $A_{/j_A(C)}$ for $C$ with a final object, and preservation of weak equivalences by $i_Aj_A$ together with asphericity of $A$ (pp.~79--81). The local asphericity argument identifies the stronger counit conditions; the remaining implication from the initial equivalence condition is revisited in the later marginal note dated 27.3, which observes that every autoequivalence preserves a final object. Thus the extra final-object hypothesis is unnecessary in that formulation. The notes numbered 48 and 49 have broad pencil hatching whose cancellation status remains doubtful. A literal composition on p.~80 is reversed: condition~(v) prints $j_Ai_A(C)$ where both its domain and the displayed category of elements require $i_Aj_A(C)$. This source slip is preserved. Part~b identifies local test categories through weak test slices, a separated interval aspheric over the final presheaf, the corresponding property of the Lawvere element, and asphericity of the comparison $i_Aj_A(C)\to A\times C$ (pp.~81--83). Its proof uses the interval comparison and homotopy argument of section~37. At the reduction on p.~82 the text says that $A$ has a final element, although replacing $A$ by $A_{/F}$ for aspheric $F$ supplies only asphericity in general; the argument needs that weaker condition. The same page's final reference says ``in same \S.31'', despite the preceding reference to section~37. These formulations remain in the readers. A local test category that is also aspheric is a test category. Grothendieck calls the corresponding presheaf topoi locally modelizing and modelizing, and asks whether the weak condition can also be made intrinsic to the topos rather than to its chosen generating category. Part~c characterizes strictness by preservation of binary products after localization, equivalently by total asphericity together with the local test property (pp.~83--85). The concise conditions T~2 and T~3 require the added hypothesis $A\ne\varnothing$. Their equivalence with the product formulation uses the asphericity test for $A_{/F\times G}\to A_{/F}\times A_{/G}$. The source's reference call is 36, and the interlinear words ``in A'' with their added hat do not cancel the following word ``of'' on p.~84. Part~d introduces weak test functors $i:A\to(\mathrm{Cat})$ through the model-preserving property of $i^*$, including reflection of weak equivalences; equivalently the composite $i_Ai^*$ is model preserving (pp.~85--86). An induced autoequivalence alone supplies no chosen comparison with the identity. With final objects in the $i(a)$, the concrete evaluation functor $A_{/i^*(C)}\to C$ gives that comparison by the earlier homotopy argument. The revised proposition on p.~86 states only $(ii)\Leftrightarrow(iii)\Rightarrow(i)$: relative asphericity of $i^*(\Delta_1)$ is equivalent to an endpoint-compatible interval map $i_!(I)\to\Delta_1$, with asphericity of $A$ explicitly added. It does not assert the converse from weak test functor to these stronger conditions. The latter define a test functor on p.~87. Whether every weak test functor on a test category is of this stronger sort remains open here; strictness of $A$ removes the absolute-versus-relative asphericity obstruction. The closing remarks then correct a further logical shortcut: recovering reflection of weak equivalences from an induced equivalence uses strong saturation of $W$, and the necessary Quillen-model argument is not yet checked in this presentation (p.~88). The left adjoint $i_!$ need not itself be model preserving, even for a strict test category. Grothendieck also leaves open whether model preservation of $i_A$, or of a more general $i^*$, already forces the weak test property; the canonical $j_A$ with asphericity is the known case. The deleted fourth remark returns to the final-object observation already incorporated in the earlier margin. The arXiv witness omits the live ``even'' in Remark~2 and rewrites the ending of Remark~3, losing ``assuming'', ``is'', and its comma; these omissions obscure the successive qualifications preserved here. The next main section uses this review to examine the interplay of weak equivalences, aspheric objects, homotopy intervals and test functors; the intervening archive leaves contain later notes to this chapter.

\subsection*{Later notes to Chapter II (cote 134-3, sites 89--93; leaves II i--v)}
These five leaves follow the main chapter in archive order and annotate it retrospectively; their lower foliation 89--93 must not be confused with the next chapter's main pages 89--93 in cote 134-4. Notes~1--3 revisit the Hauptvermutung, the semisimplicial construction of the Leray spectral sequence, and the attribution of the Dold--Kan theorem. Note~3 expressly calls ``Dold-Puppe'' an inaccurate name, contrary to the witness's ``accurate''. Note~4 retracts the earlier obstruction to expressing stack cohomology through the primitive source-and-target structure, introduces hemispherical or globular complexes, and identifies the category of standard hemispheres with the one used in David W.~Jones's thesis. Its face operations have indices $i\geq1$, whereas its degeneracies begin at $i\geq0$; the arXiv witness changes the latter bound to 1. The handwritten category symbol has an inner curved arc, retained in the readers. The note compares hemispheres, simplices and cubes as test categories and contractors, then asks for explicit test functors and suitable existence and uniqueness results, pointing to section~90 (sites 89--90). Notes~5--6 refer to the Breen appendix. Note~7 recalls the nerve construction in a Bourbaki talk but leaves the bibliographical slot blank, marked by a marginal X; that unfinished reference is not supplied editorially. Notes~8--9 retract the proposed topology characterization and the early cofibration description, directing the reader to sections~24 and~30 and to the later discussion of Thomason in section~87. Notes~10--12 record the expanding scope of the project, the preference for hemispheres, and the connection with Benabou's work (site~90). Note~13 makes the larger methodological change explicit: Quillen's axioms recede as asphericity and contractibility structures become the more immediate tools for constructing modelizers and their canonical homotopy-type functors. Notes~15 and~19 reject the unrestricted use of the canonical topology, especially for $(\mathrm{Cat})$. Note~16 bases the suitability of mixed multiple complexes on the test-category criterion and stability under products, referring to sections~31 and~74. Note~17 corrects the earlier compactness description of cells: a finiteness condition on the ordered set is needed. Notes~18 and~20 leave the relevant modelizer assertion unproved and set finite-presentation restrictions aside (site~91). Notes~21--22-prime describe the displacement of cellular structures by integrators and cointegrators, and of homotopy-equivalent topologies by test notions; the proper-map correction in note~22 must be read with the condition it qualifies in the earlier body. Notes~23--28 identify the slice-topos equivalence, the evolving test terminology, canonical modelizers, strictly generating families, and the eventual abandonment of topological realization as the main route. The fully handwritten note~28-prime points to the weakened condition of section~41. Note~29 still treats the modelizing status of ordered sets and preorders as unproved. Notes~30--32 connect product test categories and their natural comparison functors with sections~74 and~82, prefer hemispheres again, and question the efficiency of the semisimplicial route to the Cech--Cartier--Verdier functor (site~92). The last leaf continues that question toward a projected return to topoi, revises a section reference from 38 to 39, and cancels an entire example about failure of total asphericity in slices. Note~35 explicitly retracts the earlier confusion between ordinary limits and homotopy limits. Its claim that only finite products and arbitrary sums exist is itself too restrictive on products: arbitrary products exist in the homotopy category, represented by products of Kan complexes, since maps and homotopies into a product are determined componentwise. The source wording is preserved. Note~36 visibly has a corrected numeral XIII, with an uncertain older layer. The dated note~37 strengthens a modelizing condition by also requiring $j_A=i_A^*$ to be modelizing and refers to section~65, D), Remark~3. Note~38 finally sends the reader back to section~44 for the revised weak, ordinary and strict test notions. Throughout these leaves the checked text distinguishes typed interlinear amendments from handwriting; in particular the live final i and typed ve of ``primitive'' must survive, while the apparent commas attached to several insertion hooks are not textual punctuation.

\subsection*{45. Comparing modelizers by test functors (pp. 89--92)}
The review of section~44 leads back to the organizing problem of section~26: how weak equivalences, aspheric objects, homotopy intervals and test functors determine one another in a general modelizer (p.~89). Test categories are compared with spectra of discrete valuation rings, which can be sent into other objects to test their structure. A proposed test functor $i:A\to M$ should induce a modelizing functor $i^*:M\to\widehat A$, whereas the canonical $i_A:\widehat A\to(\mathrm{Cat})$ sends an elementary modelizer toward the common categorical model (p.~90). Their exactness properties run in opposite directions: $i^*$ preserves projective limits, while $i_A$ preserves inductive limits. The left adjoint $i_!$ is generally not expected to be model preserving; the canonical pair $(i_A,j_A)$ is exceptional in having both adjoints model preserving. Composing gives $i_Ai^*:M\to(\mathrm{Cat})$, but this composite generally loses both exactness properties. For a strict modelizer its preservation of finite products survives up to weak equivalence. The discussion also points to preservation of sums, with a faint pencil margin adding fiber products; that addition is absent from the witness. In the reverse direction, finding $(\mathrm{Cat})\to M$ is equivalent, through the comparison functors already available, to finding a modelizing functor $\widehat A\to M$. Neither existence nor useful exactness is asserted for arbitrary $M$. Topological spaces provide a familiar special case: a suitable realization $r:A\to(\mathrm{Spaces})$ extends to $r_!$, and the adjunction $(r_!,r^*)$ can have both members modelizing (p.~91). The resulting categorical realization is explicitly $r_!i^*:(\mathrm{Cat})\to(\mathrm{Spaces})$ in formula~(6), not the witness's $r_!r^*$, whose right-hand factor would have the wrong source here. This analogy does not make spaces an equally satisfactory basic modelizer: their structure is less algebraic, the realization choices are less canonical, and the existence of such adjoint comparisons for broad classes of test categories remains unsettled (p.~92). The conclusion therefore singles out the two constructions $i^*$ and $i_A$ as the main ones and prepares the definitions in section~46. The witness omits the intervening words ``test functor (still to be defined !) of a test category into an arbitrary'' from the first construction's description on p.~92; retaining them is necessary to see which functor remains to be defined. The checked source also retains the author's paragraph instruction before the final reflection and the joining curve in ``def ined'', without converting either into manuscript prose.

\subsection*{46. Weak test functors into an arbitrary modelizer (pp. 93--95)}
A modelizer here consists of a category $M$ and a class $W_M$ satisfying the mild saturation conditions of p.~59, with localization equivalent to $(\mathrm{Hot})$; no particular equivalence is included in the structure. The ``inspiring assumption'' supplies motivation but is expressly withheld from the technical argument (p.~93). For a weak test category $A$, the proposed weak test condition on $i:A\to M$ is that $i^*:M\to\widehat A$ be modelizing: it must reflect as well as preserve weak equivalences and induce an equivalence of localizations. The source uses $W_A$ in this formula, without a hat on its subscript. Aspheric objects are defined by their maps to a final object $e_M$. Their images are final in $H_M$ if $e_M$ remains final there; the converse additionally uses strong saturation. The dated marginal note 29.3 observes that existence of a weak test functor verifies the final-object condition. In the special categorical case, the striking extra feature is an equivalence between preservation of weak equivalences, preservation of aspheric objects and the full weak test condition (p.~94). The explanation explicitly restores a hypothesis omitted from the displayed observation: every category $i(a)$ must have a final object. With this hypothesis, section~37's key result on p.~61 proves the nontrivial implication, and it even suffices to test categories $C$ having final objects. The inserted reference is $(a\,iv)\Rightarrow(a\,ii)$, with typed clauses and a handwritten arrow. The resulting question asks when a functor between general modelizers that preserves weak equivalences, or merely aspheric objects, must already be model preserving. Possible assumptions include compatibility with projective limits or existence of a left adjoint, but no general theorem is asserted (pp.~94--95). The first proposed test cases are $i^*:M\to\widehat A$ with $M=(\mathrm{Cat})$ or another elementary modelizer $\widehat B$. They lead directly to the deferred task of dropping the final-object hypothesis on $i(a)$. The following section begins on 30.3 by announcing a satisfactory answer under strictness of $A$; the explicit words ``test category'' in that announcement are omitted by the witness. Several dark marginal or punctuation-like marks remain unresolved in the checked source, including the raised mark after $M$ in the general question; they do not supply additional hypotheses.

\subsection*{47. The categorical comparison through functor categories (pp. 95--98)}
The proposed extension starts with a distinction that the notation must preserve. The presheaf $i^*(C)$ has sets $\mathrm{Hom}(i(a),C)$ as values, whereas the category $A_{//C}$ is fibered over $A$ with the full functor categories $\underline{\mathrm{Hom}}(i(a),C)$ as fibers (p.~96). The set of objects, regarded as a discrete category, gives a canonical comparison $A_{/C}\to A_{//C}$, usually far from an equivalence on individual fibers. Constant functors supply the other comparison, giving the zigzag $A_{/C}\to A_{//C}\leftarrow A\times C\to C$. This is the first explicit use of the bicategory structure of $(\mathrm{Cat})$ in the argument: natural transformations now provide morphisms inside the fibers. The statement that $C$ is identified with the full subcategory of constant functors needs qualification before assumptions on $i(a)$ are imposed. It is valid for nonempty connected $i(a)$; with several connected components, transformations between constant functors can vary by component, and an empty domain collapses the objects. The later contractibility condition supplies the needed connectedness and nonemptiness. If all three maps in the zigzag are weak equivalences, they give a canonical isomorphism between $i_Ai^*(C)$ and $C$ after localization. Thus $i_Ai^*$ preserves weak equivalences and induces the identity up to isomorphism. When $A$ is weak test, the known equivalence induced by $i_A$ then transfers the equivalence property to $i^*$ (p.~97). The source calls the latter a ``test-functor'' at this point, although the earlier definition applies that name to $i$ and calls $i^*$ modelizing. The three arrows lead to different hypotheses. The projection $A\times C\to C$ is a weak equivalence for all $C$ exactly when $A$ is aspheric, by section~40. For the middle arrow, Grothendieck admits a general fiberwise comparison proposition: a cartesian functor between fibered categories, or a cocartesian functor between cofibered categories, is a weak equivalence if it is so on every fiber. Its second alternative literally prints ``i is cartesian'', although the proposition quantifies the functor $u$; that variable slip remains in the source. This reduces the middle arrow to the maps $C\to\underline{\mathrm{Hom}}(i(a),C)$. Their being weak equivalences for every $C$ is announced as a characterization of contractibility of $i(a)$, a stronger notion than mere asphericity. The marginal correction dated 9.4 postpones the promised proof to p.~121 and its converse to p.~143. The first arrow is then expected to yield to the asphericity criterion, producing a generalized key result with twice as many equivalent formulations (p.~98). This section states the program and its reductions; it does not yet prove that final equivalence theorem. The closing terminology correction explains why earlier uses of ``contractible'' for ``aspheric'' must be read cautiously outside the CW-space setting. Categories with final objects remain the evident contractible examples. The next section develops homotopies and contractibility in a general category with weak equivalences. Its opening retains the explicit refusal to assume a strict initial object until needed, contrary to the witness's omission of ``not'', and adds binary products as a precaution.

\subsection*{48. Interval homotopies, quotient categories and weak intervals (pp. 98--107)}
The generalities begun on 31.3 supply the machinery deferred in section~47. The starting data are a category $M$, weak equivalences $W$ satisfying the earlier mild saturation assumptions, binary products and an interval $\underline I=(I,\delta_0,\delta_1)$. Strictness of the initial object is deliberately withheld (p.~98). A homotopy from $f$ to $g$ is a map from the product of their source with $I$, whose restrictions at the two endpoints are $f$ and $g$. The reformulated homotopy lemma makes three consequences explicit: homotopic maps have the same image in $W^{-1}M$, weak equivalence is invariant under homotopy, and a homotopy from an isomorphism to a constant map makes the relevant objects aspheric over the final object (p.~99). The last assertion is a revision of a canceled, broader claim about an arbitrary aspheric map. A single interval homotopy need be neither symmetric nor transitive. For $M=(\mathrm{Cat})$, the ordered two-object interval gives natural transformations, the equivalent-to-a-point two-object groupoid gives natural isomorphisms, and the subdivided interval gives a common natural-transformation target. These examples separate the directed, symmetric and transitive properties. Passing to the equivalence relation generated by strict interval homotopies produces the category $M_{\underline I}$ with the same objects as $M$; finite zigzags, allowing the opposite interval, are essential to this wide sense of homotopy (pp.~99--100). An $\underline I$-homotopism is a map that becomes invertible in this quotient, and an $\underline I$-contractible object becomes isomorphic to the final object. Grothendieck distinguishes the stronger saturation condition involving both composites $gf$ and $fg$ from the earlier retract condition, and from strong saturation defined by the inverse image of isomorphisms under localization (pp.~100--101). The repeated display of condition $c')$ on p.~101 visibly has $gf$ twice, despite the earlier distinct $gf,fg$ formulation; the live repeated $gf$ has been checked independently and is preserved as an apparent authorial formula slip. The quotient functor preserves finite products, and finite sums when sums distribute over the required products. The restrictions to finite operations matter: ``arbitrary'' is canceled in this discussion, and the later paragraph on p.~106 proposing an extension to infinite operations from a characteristic interval is entirely canceled. A fixed interval by itself does not provide a uniform bound on the zigzags needed in infinitely many factors. To turn wide homotopies into single homotopies, the author considers amalgamation of intervals along successive endpoints. He requires the needed amalgamated sums, stability of the two endpoint maps and their intersection, asphericity, and compatibility of the construction with products. Under those additional hypotheses, finite strings of $\underline I$ and its opposite form $\operatorname{Comp}(\underline I)$, and wide $\underline I$-homotopy is realized by a strict homotopy for an interval in that family (pp.~102--103). The formula for the opposite of an amalgamation on p.~102 keeps the order $I,J$; with the stated endpoint convention, reversing the endpoints should reverse the two factors. That visible order is retained and flagged. The ensuing topos argument is exploratory: a Mayer--Vietoris statement would supply the required asphericity, but the proposed passage to $(\mathrm{Cat})$ does not make a point a closed subtopos in general. The source leaves its lemma reference blank and abandons this digression; it does not establish an unconditional amalgamation theorem. Its reference to condition $c')$ on p.~103 points contextually to the amalgamation condition $c)$, not to the distinct saturation axiom. An endpoint-preserving map $\underline J\to\underline I$ makes $I$-homotopy imply $J$-homotopy and therefore gives a quotient functor $M_{\underline I}\to M_{\underline J}$. The direction of this functor is important. A weaker sufficient hypothesis is that the two endpoints of $I$ already be $J$-homotopic; more generally, any congruence compatible with composition and products identifies maps joined by an interval whose endpoints it identifies (pp.~103--104). A fundamental or characteristic interval $\underline I_0$ is one whose wide homotopy relation contains all the others. For categories, $\Delta_1$ has this property because its wide homotopy relation on maps from the final category detects membership in the same connected component; connectedness of an aspheric interval then joins its endpoints. This assertion concerns the generated relation, not existence of a single directed natural transformation. General existence of characteristic intervals remains a question (pp.~104--105). Binary products of intervals supply a common broader relation for any pair, so the union of all interval relations is an equivalence relation $\sim_W$. Its quotient $M_W$ is the filtering limit of the $M_{\underline I}$, with the preorder on intervals opposite to the existence of an interval map. The sentence describing the underlying object of the product interval literally says ``of M'' where the construction concerns $\underline K$; the source wording is preserved. The canonical chain $M\to M_{\underline I}\to M_W\to H_M$ has a common further localization at weak equivalences. The quotient $M_W$ itself is not thereby identified with $H_M$ (pp.~105--106). The resulting $W$-homotopisms and $W$-contractible objects generalize the fixed-interval definitions. Homotopisms are weak equivalences, but merely having weak equivalences in both directions does not give a homotopy inverse. The section closes with a significant relaxation: the arguments have used the interval's aspheric projection, not the disjointness of its endpoints. A ``weak homotopy interval'' therefore requires only that projection to be aspheric; even the final object with its repeated section supplies one (p.~106). Agreement between homotopy generated by weak intervals and by actual intervals is a further issue. It is clear in the categorical case from connectedness, but is not asserted for arbitrary $M$. Semisimplicial face complexes without degeneracies are cited as a modelizer without a homotopy interval, motivating the weaker language before the question is postponed (p.~107). This qualification is what allows section~49 to use base change without an unnecessary hypothesis about strict initial objects. The checked revisions distinguish typed interlinear additions from pen words throughout. In particular ``the one and unique'' on p.~104 leaves one live under a placing curve, and the corresponding witness disagreement arose from the initial transcription's mistaken deletion.

\subsection*{49. Contractibility, asphericity and perfect intervals (pp. 107--109)}
The section isolates the relation between contractibility and universal asphericity. An object $X$ is contractible when its identity is homotopic, in the wide sense, to a constant map through a section of the final-object projection. The section is visibly named $c_o$ on p.~107, not the initial transcription's $e_o$. The displayed constant map is nevertheless written $p_Xc_o:X\to e\to X$, reversing the order required by the surrounding right-to-left composition convention, as in the analogous formula on p.~50. The source formula is retained. Homotopy invariance and mild saturation first make $p_X$ a weak equivalence. Applying the same reasoning in every slice $M_{/S}$ then makes it universally so: base change preserves the products used in a homotopy and preserves the interval's aspheric projection. The disjoint-endpoint condition is unnecessary to this argument; insisting on it would introduce an extraneous assumption about preservation of initial objects. Thus contractible objects are aspheric over $e$. The two typed insertions ``one would expect it'' and ``that px is'' record the author's movement from expectation to proof. Their placing curves leave ``to be'' and ``an'' live, contrary to the initial transcription's deletions and the witness's attempted repair. The converse fails in the usual simplicial and topological settings, so an aspheric interval with two disjoint endpoints need not itself be contractible. The source's phrase ``in the sense of singular cohomology'' should not be expanded into a general assertion that acyclicity implies weak contractibility; ordinary acyclicity alone does not eliminate the fundamental group. Grothendieck recognizes that he had tacitly conflated all homotopy intervals with contractible ones (pp.~107--108). His first proposed term ``strict'' is repeatedly replaced by ``contractible'', leaving room for a stronger condition. That stronger structure is a multiplication $h:I\times I\to I$ with $h(e_0,x)=x$ and $h(e_1,x)=e_1$: endpoint~0 is a left unit and endpoint~1 a left zero. This is a single elementary self-contraction, not merely a zigzag. It reinterprets the comparison lemma of p.~60: a map from an aspheric interval $J$ to such an $I$ transfers the contraction to $J$ and makes $I$ aspheric. With this multiplication, asphericity over $e$, contractibility and $W$-homotopy of the two endpoints become equivalent (p.~108). The explicit $W$ prefix in the last formulation is on the manuscript but omitted by the witness. The endpoint convention should be followed as printed here, where the unit and zero are numbered 0 and 1. A Lawvere element carries the required idempotent multiplication from intersection of subobjects. With a strict initial object, its full and empty sections are disjoint. The earlier comparison corollary therefore upgrades existence of any homotopy interval to existence of a contractible one and even of an interval with this multiplication, now named ``strict'' (p.~109). The displayed implications point from strict to contractible to ordinary; the surrounding Lawvere argument supplies the converse existence implication under its hypotheses. The ordered categorical interval $\Delta_1$, representing sieves, and the two-object category equivalent to a point, representing full subcategories, both obtain intersection laws. Their homotopy relations differ: $\Delta_1$ is characteristic, whereas the groupoid interval detects only natural isomorphism. A ``perfect'' interval is defined to be both strict and characteristic. This closes the section's refinement of the earlier interval vocabulary and leads into section~50's reflection on the four interacting homotopy notions. The right-edge ending of contractible on p.~109 and the scope of the pen S over the terminal cells of represents remain doubtful; neither is silently regularized. The new section begins with a handwritten date 2.4 and the number 50 alone, without the descriptive title added in arXiv.

\subsection*{50. The proposed intrinsic structure of canonical modelizers (pp. 109--110)}
The reflections dated 2.4 draw together four notions developed in sections~48--49: homotopy relations between maps, homotopy intervals, homotopisms and contractible objects. Their dependence on a preassigned class $W$ now appears weaker than their mutual dependence. A fixed interval already defines its homotopies and quotient category without $W$; even arbitrary binary products can be replaced by the ``squarability'' of that interval, meaning existence of $X\times I$ for every $X$. In the $W$-based theory, asphericity of the interval supplied the homotopy lemma, but the preceding congruence argument permits the weaker condition that the two endpoints be $W$-homotopic. Such a relation is expressed by a finite chain of aspheric intervals mapping into the given object (p.~110). Grothendieck therefore proposes that, in the intended examples, intrinsic $0$-connectedness may provide the appropriate interval condition. This is a program for recovering homotopy from the category, not a proved equivalence for all categories with products. Recalling the observation and question of pp.~94--95, he next hopes that contractible objects themselves determine weak equivalences in the nicest modelizers, including $(\mathrm{Cat})$ and elementary modelizers. If both expectations hold, $W$ is uniquely determined by the underlying category and any equivalence of the underlying categories automatically preserves the modelizing structure. ``Canonical modelizers'' names the anticipated class with this property; an overall description analogous to elementary modelizers through test categories is left as the next ambition. Section~51 accordingly starts again without a pregiven $W$. The scan's additions of in and W are typed interlinear insertions, and the corresponding is placed after the live of near the close. Two short canceled rows remain unresolved; the longer canceled phrase has visible quotation marks around joined. None of those doubtful older layers supplies an extra hypothesis to the proposed intrinsic theory.

\subsection*{51. Four homotopy structures and their provisional hierarchy (pp. 111--124)}
Starting without a pregiven class of weak equivalences, Grothendieck compares four kinds of data: a congruence $R_h$ on arrows with fixed endpoints, a strongly saturated class $W_h$ whose localization is a quotient on arrows, a saturated family of weak intervals, and a class of contractible objects (pp.~111--113). Their proposed relations must be read together with the corrections made in the margin and in section~52. In particular, an arbitrary quotient congruence need not be the kernel of localization at the arrows it makes invertible. The margin on p.~114 denies the proposed map $r_{21}$ and literally refers to section~51, whereas the next page's note refers to section~52 and calls the preceding assertion ``rash''. The initial diagram and its proposed retractions consequently describe a developing framework. The sufficient condition discussed on p.~115 is more substantial: homotopic maps factor through two sections of a common homotopism $p:X'\to X$. Inverting $p$ then identifies the sections and hence the maps. This is the cylinder mechanism familiar from the earlier homotopy lemma; the source first names the maps $f,g$ and then writes $f_o,f_1$, a visible notation shift retained in the readers. For interval structures, saturation means that the endpoints of an admitted interval can be joined by a finite chain of homotopies using intervals from the family and their opposites. The generated relation is compatible with composition and finite products, and products of intervals give common enlargements of the corresponding relations (pp.~116--117). The claimed passage through $\mathrm{Hom}_3$, $\mathrm{Hom}_2$ and $\mathrm{Hom}_1$ remains subject to the subsequent correction. Page~116 switches from $\Sigma_h$ to $R_h$ while speaking of a subset of $\operatorname{Int}(M)$, and the first displayed $r_{23}$ on p.~117 has target $\mathrm{Hom}_1$ although the next occurrence has target $\mathrm{Hom}_2$; neither mismatch is silently repaired. The fourth structure requires more than declaring a class $C_h$ of pointed objects. Its basic condition says that every object of $C_h$ is contractible through finite chains parameterized by objects of $C_h$; saturation requires conversely that every object so contracted belong to the class. Thus $C_h$ equals the class $\overline C_h$ generated from its own parameter objects. An arbitrary interval structure need not have its generating intervals contractible, so a retraction from interval structures to these contractibility structures is explicitly left doubtful (pp.~117--119). The prose on p.~117 mistakenly says that the latter class contains $\overline C_h$ where the displayed inclusion concerns $C_h$, and the contraction formulas on p.~118 write $M\to M$ and $\mathrm{id}_M$ although their object is $X$. The live corrected wording there is ``homotopism structure'', not the initially transcribed ``homotopy''. The next aim is a small generating family of intervals, preferably a single contractible generator. A finite family can be replaced by its product; an interval generating its own homotopy relation is contractible precisely when its identity is joined to a constant map in that relation. A stronger, elementary contraction is a multiplication $I\times I\to I$ with endpoint~0 a left unit and endpoint~1 a left zero. Such an interval is now called a ``contractor'', replacing an earlier term in the draft. A contractor becomes contractible as soon as its endpoints are homotopic, whether or not it generates the whole structure. The standard categorical interval $\Delta_1$ motivates the expectation that the intended modelizers possess a single generating contractor (pp.~119--120). A Lawvere element carries the intersection multiplication, with full and empty sections as unit and zero, and strictness of the initial object separates them. The claim on p.~121 that it is final among separated intervals, with a unique endpoint-preserving map from every such interval, is too strong in the stated generality: in sets, a three-element interval has two endpoint-preserving maps to the two-element Lawvere interval, according to the image of its third element. Existence of suitable maps is enough for the transfer of endpoint homotopy used here, but it does not establish that the Lawvere interval generates every homotopy relation. The manuscript's uniqueness claim remains literal. The alternative construction takes the internal endomorphism object of $X$, with identity as unit and a constant endomorphism as left zero. The proposition on p.~121 restores, by dashed marks and a circled ``biffé par erreur'', its previously canceled hypothesis that the homotopy structure be of type $\mathrm{Hom}_3$. That authorial restoration is applied in the corrected forms. Under the interval-generated hypothesis, currying turns homotopies of maps $X\to Y$ into homotopies between sections of the internal Hom object. Consequently contractibility of $X$, contractibility of its endomorphism object, and homotopy of the latter's distinguished endpoints are equivalent; moreover $\underline{\mathrm{Hom}}(Y,X)$ is contractible and the constant-map arrow $Y\to\underline{\mathrm{Hom}}(X,Y)$ is a homotopism (pp.~121--122). The first corollary weakens a proposed contractor assumption to a weak interval acting on an object $H$, with endpoint actions the identity and an idempotent $p$. The image of $p$ is then a homotopy retract. Its displayed identity literally reads $pi_o=\mathrm{id}_{H_o}$; the preceding definitions instead require $p_oi=\mathrm{id}_{H_o}$, so the index placement is flagged here and preserved in the text. Finally, part~F lists the ambient hypotheses: finite products, arbitrary disjoint universal sums, decomposition into nonempty connected objects, connectedness of the final object, and, as a separate stronger condition, connectedness of products of connected objects (pp.~122--124). Component decomposition gives a functor $\pi_0$ left adjoint to the constant-object functor; the last condition makes it preserve finite products. The source's term is ``O-connected'', with one slashed round glyph whose interpretation remains doubtful. The phrase ``set of subobjects'' on p.~123 is live beneath a placing curve, contrary to the initial transcription's proposed cancellation. On p.~124 the author says that condition \emph{c}, concerning component decomposition, might be weakened by allowing $\pi_0$ to take values in strict pro-sets; the witness substitutes \emph{e}, changing the mathematical reference. Connectedness of the final object is treated as a convenience for avoiding a product decomposition of the ambient category. The stronger product condition, informally related to total O-asphericity, will be used only when stated. Its proposed equivalence with total asphericity for a topos is expressly a suspicion, and the locally connected topological example makes its restrictiveness apparent. The new date 3.4 then opens section~52 with an explicit correction of the first two structures, resolving the difficulty already marked in this section's margins.

\subsection*{52. Correcting the quotient--localization correspondence (pp. 124--129)}
The notes of 3.4 correct the first two notions of section~51 by taking the basic object to be a functor $M\to M'$ that is simultaneously a quotient and a localization: it is bijective on objects, surjective on arrows, and localizing (pp.~124--126). A relation is now a homotopy relation only when its quotient functor is localizing; a strongly saturated class is a homotopism structure only when its localization is surjective on arrows. With those restrictions, the two descriptions are equivalent, giving the corrected bijection $\mathrm{Hom}_1(M)\leftrightarrow\mathrm{Hom}_2(M)$. The distinction is strict enough that an equivalence under $M$ becomes an isomorphism under $M$ in the stated situation. For arbitrary congruences and strongly saturated classes, write their ordered sets as $Q$ and $L$. The map $r:Q\to L$ takes the inverse image of the isomorphisms in $M/R$, while $s:L\to Q$ takes the equality relation induced by $M\to W^{-1}M$. The natural comparison $\alpha_R:M_{r(R)}\to M_R$ is bijective on objects and surjective on arrows, so its faithfulness tests whether $R$ is a homotopy relation. Dually, $\beta_W:M_{s(W)}\to M_W$ is bijective on objects and injective on arrows, so its fullness tests whether $W$ is a homotopism structure (pp.~126--127). These are different tests. The first monoid example cites noninjectivity of $\overline G\to G'$ immediately after discussing whether $r(R)$ is a homotopism structure; that noninjectivity directly tests $R$, not $r(R)$. Taking $G=\mathbb N$ and its unit quotient supplies the intended failure of the latter property, but the relevant obstruction is that $\mathbb N\to\mathbb Z$ is not surjective. In the converse paragraph, nonsurjectivity of a localization likewise tests $W$, and alone does not show that its kernel relation is outside $Q_0$: for this same example the kernel relation is equality, whose quotient is already a localization. The source examples are retained with this qualification. The corrected homotopy subsets are $Q_0=\{q:sr(q)=q\}$ and $L_0=\{\ell:rs(\ell)=\ell\}$. Monotonicity and the inequalities $sr(q)\le q$, $rs(\ell)\le\ell$ characterize associated pairs by the simultaneous reverse inequalities $\ell\le r(q)$ and $q\le s(\ell)$ (pp.~127--128). They also explain exactly what the cylinder condition $C_{1,2}$ of p.~115 accomplishes: it gives $R\le s(W)$, and with $W=r(R)$ forces $R=sr(R)$. Thus the earlier sufficient condition survives, while the claim about arbitrary relations does not. The interval construction offers a different route from unrestricted data to a homotopy structure. The new diagram has a retraction from $Q(M)$ to weak-interval structures, and composing it with $s$ gives a construction from arbitrary $W$. The comparison back to the original quotient or localization requires the accessory finite-product condition. Grothendieck emphasizes that the interval method is not self-dual, unlike Quillen's left and right homotopies, and connects this asymmetry with the disjoint-sum assumption adopted in section~51 (pp.~128--129). The connected pen rules on p.~128 instruct a paragraph before this return to the diagram; the corrected forms apply that instruction. The revised inventory therefore has three essentially different levels: general homotopy structures, weak homotopy interval structures, and contractibility structures. The hoped-for modelizer examples still lie at the strongest level and should even possess a single generating contractor with commutative idempotent multiplication. This remains an expectation, setting up section~53's analysis of functors compatible with the structures. Several source details matter to the reading: p.~126 defines a functor $i$ but then writes $u(f)$, its live rewritten cross-reference is (2), and a separate (*) remains unrevised; p.~127 names $r$ even while describing the direction from $L$ to $Q$, where $s$ is required; p.~128 writes ``R and L'' for the pair $(R,W)$ and closes the displayed set with a parenthesis instead of a brace. The readers preserve these slips. The margin on p.~129 asks that the crowded diagram be spread out, while small pen graphics after the caption's ``or'' remain distinct from any secure cancellation. The typed overprints giving final capital H in ``THe'' and w over t in ``which'' are preserved as observed; the witness's substituted ``that'' is not adopted.

\subsection*{53. Functors compatible with a homotopy structure (pp. 130--131)}
For an initially unrestricted functor $u:M\to N$, the section compares four tests suggested by section~51: equality of the images of homotopic arrows, equality of the two images of each weak interval's endpoints, inversion of every homotopism, and passage of every contractible object to a final object (p.~130). The last formulation assumes that $u$ already preserves the final object, and the interval and contractibility language assumes finite products in $M$. Section~52's identification of the homotopy quotient with the appropriate localization makes the first and third tests equivalent: both say that $u$ factors through $M\to M_h$. They define compatibility of $u$ with the homotopy structure. The stronger endpoint and contractibility tests recover this compatibility under additional hypotheses: for a weak-interval structure, preservation of finite products lets endpoint equality kill each elementary homotopy, while for a contractibility structure passage of contractible objects to a final object first gives endpoint equality. The diagram explicitly marks the $\mathrm{Hom}_3$, $\mathrm{Hom}_4$ and product-preservation qualifications. The witness omits the whole typed line that finishes the contractibility implication and starts the implication from endpoint equality to equality of homotopic maps; the checked reader retains that essential distinction. Its raised pair in condition~(i) literally reads $(f;g)$ with a semicolon. The small dark sign following the diagram's lower-right node remains a doubtful star/1-like mark, and the heavily retraced prime in the prose implication retains its doubt; neither supplies a new hypothesis. The opening of p.~131 states the practical conclusion: a product-preserving functor may be checked on intervals or on contractible objects when the homotopy structure admits the respective description. Section~54 then applies these tests to a localization functor and separates compatibility from the stronger requirement that homotopisms belong to the specified class itself.

\subsection*{54. Compatibility with weak equivalences and the maximal interval structure (pp. 131--134)}
The compatibility tests of section~53 are now applied to the localization $u:M\to W^{-1}M$. Compatibility means that $u$ kills the given homotopy relation, while \emph{strict} compatibility asks that every homotopism already belong to $W$. These coincide when $W$ is strongly saturated; otherwise strict compatibility retains information about the specified class before strong saturation (p.~131). The first displayed definition literally calls $N=W^{-1}M$ a localization ``of N'', an apparent source-variable slip preserved in the reader. For a contractibility structure, the initial test using product preservation and the projections of contractible objects is substantially rewritten: the product-preservation hypothesis is deleted, and the projections are required to be universally in $W$. The resulting proposition applies to a generating family of weak intervals whose projections are universally in $W$. Mild saturation gives invariance of membership in $W$ under the generated homotopy relation; the stronger saturation condition then gives $W_h\subset W$. The proof invokes the earlier homotopy lemma and the reduction of a homotopism to endomorphisms homotopic to identities (pp.~131--132). This returns to section~48's construction from a class of weak equivalences, but deliberately abandons the old requirement that the endpoints be disjoint. The family of all intervals whose underlying objects are universally in $W$ generates the structure now denoted $h_W$. Its resulting weak intervals can be more numerous than those generators: being a weak interval for $h_W$ does not itself force universal $W$-asphericity. A later margin dated 11.4 expressly warns that ``W-aspheric'' here means that $I\to e$ is universally in $W$, calls the terminology ambiguous, and refers forward to p.~181 and following. That qualification should accompany the argument, especially when comparing later uses of the same term. The question whether $h_W$ is the widest structure strictly compatible with $W$ must be restricted to structures generated by weak intervals. A category with a zero object has only coincident interval endpoints and therefore only the trivial interval-generated relation, although other homotopy quotients can be nontrivial. The example mentions projective complexes of modules and quasi-isomorphisms; the category of projective complexes need not be abelian, despite the preceding wording, but the zero-object observation suffices for the intended obstruction (pp.~132--133). The ensuing proposition separates two conditions: maximality of $h_W$ among strictly compatible interval structures, and the assertion that two sections of any object which become equal under $u$ are already $h_W$-homotopic. The latter implies the former. Conversely, strong saturation of $W$ and preservation of finite products by $u$ make the endpoint-equality construction a strictly compatible interval structure, so maximality yields the endpoint condition. Grothendieck expects later connectedness hypotheses to verify this criterion for reasonable modelizers. He nevertheless keeps the construction from universally $W$-aspheric intervals as the practical definition: strong saturation and product preservation are harder to check, product preservation fails for non-strict elementary modelizers, and equality of endpoint images in a localization is less directly testable than the projection condition (pp.~133--134). The revisions are essential to this progression. The formerly proposed Kan-complex example is canceled in favor of a general homotopism structure; ``natural'' is typed above the order comparison with ``of the'' left live beneath its placing curve; the inclusion on p.~132 has the literal parenthesized subscript $W_{(h_W)}$, whereas the later strict-compatibility formula there has $W_h$. The blackened final s of homotopy intervals and the overlapping middle letters of Clearly remain doubtful. On p.~133 the placement curve puts Clearly before the new ``for such an h'' sentence, and the final comma after strongly saturated is visibly penned. The isolated typing slip W-asperic at the close is corrected only in the corrected forms, with the diplomatic spelling retained.

\subsection*{55. Morphisms of homotopy structures and tests on generators (pp. 134--135)}
A morphism $(M_1,h_1)\to(M_2,h_2)$ is a functor $u$ descending to the corresponding homotopy quotients. Its descended functor is unique, and the commutative square expresses either preservation of homotopy between arrows or preservation of homotopisms (p.~134). This uses section~52's corrected identification of quotient and localization, rather than the unrestricted earlier definitions. If final objects are preserved, weak intervals are carried to weak intervals. Conversely, preservation of finite products lets one test a morphism on weak intervals when $h_1$ is interval-generated, or on contractible objects when it is a contractibility structure. The test further reduces to any chosen generating family: the images of the two endpoints of each generator must be $h_2$-homotopic. If the generators are themselves contractible, their images being contractible is equivalent to that endpoint condition. A single generating contractor is the intended especially economical case: products transport its multiplication, while compatibility supplies the endpoint homotopy needed for contractibility of its image (pp.~134--135). The last formula on p.~134 literally writes $u(\underline h)$ where the preceding quantified interval is $\underline I$; this apparent notation slip is preserved. Page~135 then explicitly supplies tests omitted from the two preceding sections. For a product-preserving functor $M\to N$, compatibility may be checked by equality of the images of the endpoints of the generating intervals; for contractible generators, their images must be final objects of $N$. For a class $W$ whose localization preserves finite products, membership of each generating interval projection $I\to e$ in $W$ suffices for strict compatibility and is necessary when the generators are contractible. Without product preservation, the sufficient condition becomes universal membership of these projections in $W$, and remains necessary for contractible generators. This is the generator form of section~54's proposition, not a claim that an arbitrary endpoint-equality test works without the finite-product hypothesis. The overwritten paragraph on p.~135 makes that distinction visible: the old necessary-and-sufficient formulation is canceled and replaced by the qualification about contractible generators, followed by the separate universal-asphericity statement. The placing curves before ``for strict compatibility'' and the typed ``for $h_2$'' are not textual opening parentheses; the strike scope at the remaining opening parenthesis before the rewritten condition is itself doubtful and remains marked. The word contractor, the phrase factor through in the following section, and under the beneath its later insertion all remain live. The section ends before the handwritten date 4.4 and the number 56 corrected from 55; the descriptive section titles supplied by the witness are absent from the manuscript.

\subsection*{56. Recovering weak equivalences and the provisional work plan (pp. 135--139)}
The notes dated 4.4 return to the intuition behind the preceding generalities: useful test functors should take their values in the full subcategory $M_c$ of contractible objects. With a contractible generating interval, an asphericity test on its image under $i^*$ is expected to characterize a test functor $i:A\to M$; more generally one should test a small full subcategory $M_c^o$ of contractible objects generating the homotopy structure (pp.~135--136). The handwritten interval symbol here is $\Delta$ in both formulas, whereas the witness substitutes an underlined $I$. Fully faithful inclusions containing the chosen generators would supply the most immediate examples, and their smallness explains why small generating interval families matter. The reciprocal ambition is to reconstruct weak equivalences from contractibility: for an inclusion $i:A\hookrightarrow M$ of a small full generating subcategory, define $W(A)$ as the inverse image under $i^*$ of the weak equivalences in the presheaf category $\widehat A$. Independence of $A$ is expressly a question to be checked, possibly with further restrictions. Closure of $A$ under finite products is proposed as a temporary convenient hypothesis giving a strict test category, while the author insists that eventual test functors must be more general than fully faithful inclusions (p.~136). The accompanying claim that $\Delta$ cannot embed faithfully in spaces is too strong for ordinary topological spaces: the usual geometric-simplex functor is faithful because distinct order-preserving maps have different effects on vertices. The source really says faithfully, with no fully, and is preserved; this qualification does not alter the preceding demand to allow general test functors. Combining recovery of $W$ from $h_W$ with an intrinsic description of $h_W$ using $0$-connectedness would make the modelizer structure canonical in terms of its underlying category. Both steps remain expectations here. The suggested relation between test categories and the unique coarse contractibility structure in which all objects are contractible then exposes a limitation of the recent definitions: requiring finite products in the category of test objects is too restrictive. Grothendieck compares this with the early restriction of sites to those admitting fiber products, crediting Giraud with recognizing the need to remove that condition (p.~137). This is an internal recollection, not a theorem supplied by the current section. The resulting work plan has four parts. First, relate homotopy structures to $\pi_0$ and define canonical homotopy structures. Second, write the key result on test functors to $(\mathrm{Cat})$, expected also to characterize its weak equivalences axiomatically. Third, extend that result to test functors into presheaf modelizers and study maps between test categories, then use these results to identify canonical modelizers. Fourth, pursue the alternative direct construction of weak equivalences from a contractibility structure, related to test categories carrying the coarse structure (pp.~137--138). Parts a) and b) are chosen immediately; precedence between c) and d) is left open, with a provisional preference for d). The older six-question list on p.~42 is not declared wholly solved: four questions are settled or expected to be settled by this program, while the closed-model-category question for presheaves and the homotopy structure of $(\mathrm{Cat})$ remain. The unwritten proof that the face-only simplex category $\Delta^{\mathrm f}$ is a weak test category is another deferred item. The opening of p.~139 confirms the decision to follow the canonical-modelizer problem without unrelated digressions before section~57 starts work on a). The checked revisions distinguish typed replacements such as reasonable, carefully, outline and homotopy from pen additions, preserve the canceled indeed omitted from the corrected argument, and apply the small pen connector joining canonical only in the corrected forms. Unreadable canceled words and uncertain overwritten cells remain marked; they do not supply additional hypotheses or completed claims.

\subsection*{57. Connectedness, canonical homotopy\newline and geometric weak equivalences (pp. 139--143)}
The first task of section~56 begins under the standing conditions a)--d) of section~51, while the extra product condition e) is initially withheld. A homotopy structure is $0$-admissible when $\pi_0$ factors through its homotopy quotient, equivalently when homotopic arrows induce the same component map or homotopisms induce bijections on components (p.~139). Its weak interval endpoints then lie in one connected component, and its contractible objects are connected. When $\pi_0$ preserves finite products, these necessary conditions are also sufficient for interval-generated structures and contractibility structures respectively. Without that product hypothesis, the relevant components must instead remain $0$-connected after product with every $0$-connected object. This universal condition is necessary as well in the contractibility case. The proposed higher $n$-admissibility, expressed through truncated homotopy or prohomotopy types, remains a suggestion: the expectation that geometric interval structures are already controlled by $0$-admissibility is not proved here (p.~140). Suprema supply a widest admissible structure in each of the three classes $\mathrm{Hom}_0$, $\mathrm{Hom}_3$ and $\mathrm{Hom}_4$. The chosen canonical structure $h_M$ is the one in $\mathrm{Hom}_3$, generated by weak intervals. For totally $0$-connected $M$, all connected intervals form a generating family. Without total connectedness, equality of the endpoint components must be tested after product with every $0$-connected object; universal connectedness of their common component is sufficient, but its necessity is expressly left open (pp.~140--141). A saturated class $W$ supplies its own structure $h_W$ through intervals universally $W$-aspheric over the final object. Requiring $W$ to preserve $\pi_0$ is stronger than merely requiring $h_W\leq h_M$. The proposition on p.~142 characterizes equality by two conditions: pointed objects universally $W$-aspheric over the final object are connected, and any two sections of a connected object can be joined by a finite chain of intervals of this kind. The typed first endpoint equation is $s_i=\delta_o$, rather than the expected $s_0=\delta_o$; that apparent index slip remains literal. Grothendieck then names $W$ \emph{geometric} after strengthening the first condition to $0$-admissibility of $W$ itself. Thus geometricity means this stronger admissibility together with $h_W=h_M$, not the equality alone. Recovering $W$ from its homotopy structure is deferred again, with the handwritten margin explicitly returning to p.~135. The immediate gain is that a connected contractor is canonically contractible, hence $h_W$-contractible and universally $W$-aspheric; the converse connectedness implication follows from admissibility. Assuming all internal Hom objects exist, the final proposition uses this observation to give equivalent tests on an object $X$: connectedness, ordinary $W$-asphericity, universal $W$-asphericity and $h_W$-contractibility of every $\underline{\mathrm{Hom}}(Y,X)$, or only of the endomorphism object; equivalently, the canonical maps $Y\to\underline{\mathrm{Hom}}(X,Y)$ are in $W$, universally in $W$, or exhibit deformation retracts (p.~143). The equivalence with contractibility of $X$ itself is explicitly qualified by the added assumption that $X$ has a section. The handwritten revision both inserts ordinary asphericity as a separate condition and moves it, together with connectedness, before the stronger two conditions. The diplomatic reader preserves the original physical order; the corrected forms apply the bracket and arrow. The additional explanation that c-prime means universal membership in $W$ is absent from the arXiv witness and is retained here. For $(\mathrm{Cat})$ with its usual weak equivalences, the proposition returns to the result following p.~97 and explains the long intervening study of homotopy notions. Testing the same framework on an elementary presheaf modelizer is the next intention, postponed at the late-night close. The image collation also distinguishes the typed interlinear finite, on M, main and somewhat from pen insertions: the placing curve under finite leaves the preceding a live. The heavily overwritten letters of retract and check, the final clipped letter of modelizer, and several older canceled readings remain explicitly doubtful rather than borrowing certainty from the witness.

\subsection*{58. Sections, functor homotopies and the partial localization of categories (pp. 144--146)}
The notes dated 5.4 first strengthen the connectedness setting. For the canonical structure, two sections of an object are homotopic exactly when they lie in the same connected component; hence homotopy classes of sections inject into $\pi_0$. Surjectivity is equivalent to the new condition f), requiring every $0$-connected object to have a section. Under a)--d), this implies the product condition e): chosen sections in two components of $X\times Y$ can be joined through $(s_1,t_1)$, $(s_2,t_1)$ and $(s_2,t_2)$, using the connected copies of $X$ and $Y$ (p.~144). The margin calls d) a consequence of f), but this needs a nondegeneracy qualification: in the category with one object and one arrow, initial and final objects coincide, there are no $0$-connected objects, and f) is vacuous while d) fails. The printed proposition retains d), so this marginal shortcut does not compromise its stated proof. The extra section condition is strictly stronger than total connectedness. The example is revised from an ordered-set formulation to the category of sheaves on an irreducible space; its nonempty proper open subobjects of the final sheaf obstruct sections, so f) forces the one-point topos in this example. Grothendieck calls the resulting condition \emph{strict total $0$-connectedness}, with connected rather than the witness's disconnected. It identifies homotopy classes of sections with components and, whenever internal Hom exists, gives $\overline{\mathrm{Hom}}(X,Y)\simeq\pi_0(\underline{\mathrm{Hom}}(X,Y))$ (pp.~144--145). For $(\mathrm{Cat})$, he now supplies the details omitted in the previous night's claim that the usual weak equivalences are geometric: the cohomological definition preserves components, and the walking-arrow category $\Delta_1$ both generates the canonical structure and is universally weakly equivalent to the point. Thus elementary homotopies of functors are natural transformations, while the generated equivalence relation allows finite zigzags of such transformations. These homotopy classes are the connected components of the functor category, a relation broader than natural isomorphism. The point of tolerating this unfamiliar relation is its compatibility with the map to $(\mathrm{Hot})$. The quotient $\overline{(\mathrm{Cat})}$ is simultaneously localization at canonical homotopisms, giving $(\mathrm{Cat})\to\overline{(\mathrm{Cat})}\to(\mathrm{Hot})$ (pp.~145--146). The saturated weak-equivalence class descends to this partial localization and is recovered as its inverse image. A dated margin on p.~146, added 6.4, explicitly removes the previously listed requirement that weak equivalences preserve $\pi_0$ from this factorization argument: the interval projections being in $W$, together with saturation, suffice for that step. This is a local correction to the factorization proof, not a withdrawal of admissibility from the definition of geometricity. A calculus of right fractions, perhaps of either handedness, is suggested as a possible route to strong saturation without the closed-model machinery, then deliberately postponed. The proof is not supplied here. The scan also preserves an unresolved raised definition label on p.~144, including a marked extra cell before def that cannot safely be discarded. Both Gamma symbols there are formed from a typed I with an added pen bar. On p.~145 the actual cancellation leaves the physical continuation star/ting; the original transcription had duplicated the canceled t. The isolated typed felationship is corrected only in the corrected forms, while the doubtful overprints, clipped letters and the nonstandard apostrophe in the dated margin remain explicit.

\subsection*{59. Replacing intrinsic components by a given functor (pp. 146--148)}
Topological spaces interrupt the effort to make everything intrinsic to the underlying category. General spaces do not decompose as coproducts of their ordinary connected components unless those components are open; requiring this for open subspaces as well leads to local connectedness. Homotopy theory instead uses path components. Their functor still preserves coproducts, but it is not generally left adjoint to the discrete-space functor: an arbitrary function on path components need not define a continuous map into a discrete space. Locally path-connected spaces provide a sufficient setting for that adjunction, because their path components are open. The typed interlinear pathwise is inserted before locally and does not cancel it (pp.~146--147). Grothendieck therefore proposes supplying a functor $\pi_0:M\to(\mathrm{Sets})$ as additional structure rather than deriving it from coproduct decompositions. Preservation of sums is the first possible axiom, but he observes that the recent uses more directly involved finite products and their preservation. For a prospective modelizer, a comparison $M\to(\mathrm{Hot})$ would already supply such a functor by composition with the usual component functor. The ambition is then reversed: perhaps the comparison with homotopy types can be recovered from $\pi_0$, in the manner of a total left derived functor. The text explicitly calls this speculation ``thin air''; no derived-functor construction or universal property is established (p.~147). What is immediately definable is relative $0$-connectedness, meaning singleton $\pi_0(X)$, admissibility with respect to that functor, and the widest admissible weak interval structure $h_{\pi_0}$. The pair $(M,\pi_0)$ now replaces $M$ alone. The page further claims that, when $\pi_0$ preserves finite products, all relatively $0$-connected intervals generate this widest structure, and that the preceding comparisons with $h_W$ carry over verbatim (p.~148). That generation claim needs an additional realization assumption in this new generality. For example, take the free category with finite products generated by an object $I$ and two distinct sections $\delta_0,\delta_1:e\to I$. Define a product-preserving functor to sets by sending $I$ to a two-element set and both sections to the same element. Its only objects with singleton image are copies of $e$, so its $0$-connected intervals generate the trivial relation. Nevertheless, the interval $I$ itself generates a compatible nontrivial relation identifying the two sections, since their images agree and products are preserved. The earlier decomposition into actual connected components supplied the missing ability to realize the common endpoint component by an object; an arbitrary supplied functor does not. The source assertion and its literal overwritten phrase ``by terms'' are preserved, with this limitation confined to the guide. The final reflection makes the methodological retreat explicit: the geometrical language of actual components has given way to a formal functor, prompted by the modelizer of spaces. Section~60 then returns to topoi, where concrete decomposition conditions can again be examined.

\subsection*{60. Connected topoi, geometric weak equivalences and initial axioms (pp.~148--151)}
In returning to test categories, Grothendieck distinguishes local connectedness, connectedness of the final object, and closure of connected objects under products. Local connectedness is expressed through a generating family of 0-connected objects; total 0-connectedness adds the product condition, while strict total 0-connectedness requires every nonempty object to have a section. The product test may be checked on generators. This applies the interval machinery of sections 54--59 to presheaf topoi, while the semisimplicial category without degeneracies provides a warning: its representables have no global sections, and on p.~149 Grothendieck deletes the proposed exception for the vertex. He asks, rather than proves, whether modelizing together with total 0-connectedness forces the stronger conditions. Under strict total asphericity he does prove that weak equivalences are geometric and that $h_W$ agrees with the canonical homotopy structure; the handwritten strengthening says that any connected generating family supplies the testing intervals, and the marginal reference points forward to p.~268. This equality reduces the modelizing question to the existence of an object with two distinct sections, or equivalently in the stated situation to connectedness of the Lawvere contractor. The topos of sets fails this condition, and p.~150 conjectures that it is the only counterexample under these hypotheses. The section then replaces cohomologically defined weak equivalences by a general saturated class $\underline W$ of functors, pulls it back through $i_A(F)=A_{/F}$, and lists 0-admissibility, invariance under product with the interval, and the final-object axiom. On p.~151 the final-object axiom is recognized as a consequence of the interval axiom; this shortens the initial list and leads directly to section 61's comparison theorem. A caution is needed for the intervening assertion that a category has a final object if and only if its identity admits an interval homotopy to a constant functor. The forward implication is valid, but the converse is false: the one-object category with endomorphism monoid $\{1,0\}$, where 0 is absorbing, has a natural transformation from its identity to its constant functor with component 0, yet its sole object is not final because it has two endomorphisms. Thus contractibility by such a transformation does not establish finality. This does not invalidate the deduction of the final-object axiom from interval invariance, which uses only the forward implication. The readers preserve the false converse, the literal superscript in $\Delta^1$, and the author's tentative questions; the guide separates their logical roles from the proved geometricity statement. The later argument in section 96 should also be read with the actual section-existence hypothesis: having only two subobjects of the final object alone does not imply that every nonempty object has a section.

\subsection*{61. The promised key result and a first list of axioms (pp.~151--156)}
Grothendieck resumes the comparison of a functor $i:A\to(Cat)$ with its induced $i^*:(Cat)\to\widehat A$, temporarily allowing two saturated arrow classes $\underline W$ and $\underline W'$. The desired equivalence of localizations is strengthened to the exact inverse-image condition $\underline W'=f_i^{-1}(\underline W)$, where $f_i=i_Ai^*$, so that successful sufficient conditions might also characterize the usual weak equivalences. The comparison diagram from pp.~96--98 is rebuilt using relative asphericity of $A$, contractibility of the categories $i(a)$, preservation of weak equivalences by fibered categories under condition e), and a slice criterion f). The decisive computation identifies the slice over $(a,p)$ of $A_{/C}\to A_{//C}$ with the category of elements of a pullback of $i^*(\underline{Fl}(C))\to i^*(C)$; thus the problem becomes universal membership of this last map in $W_A$. This develops the earlier homotopy and internal-Hom machinery and supplies the lemma reformulated in section 62. The text records its own weakening on p.~154: 0-admissibility b) is unnecessary for the intermediate comparison, since only $h_{(Cat)}\leq h_{\underline W}$ is needed. Two source slips require care. Diagram (11), p.~153, ends in $A$, although the intended comparison and relative-asphericity argument require the projection $A\times C\to C$, as the earlier diagram on p.~96 explicitly shows. The same paragraph changes the section reference to 57 while retaining pp.~96--98; readers should follow those pages for the diagram. Neither is repaired in the readers. Page 155's literal $f_Y$ and lower mark by $Y$ are retained, and the visible label $p$ on the pullback square, omitted by the witness, is restored. At p.~156 Grothendieck immediately withdraws the independent choice of the two arrow classes: the next section assumes $\underline W=\underline W'$ rather than claiming that this first axiomatization succeeded.

\subsection*{62. General localizers and the failure of the first characterization (pp.~156--159)}
The new entry, visibly dated 6.3. rather than the witness's 6.4., retracts the attempt to choose two independent classes of arrows. With one saturated $\underline W$, the comparison diagram proves both $\underline W=f_i^{-1}(\underline W)$ and a canonical isomorphism of the induced $\overline f_i$ with the identity on $\underline W^{-1}(Cat)$; the middle arrow in the resulting chain is used through its inverse after localization. The hypotheses combine relative asphericity of $A$, contractibility of $i(a)$, the interval condition, the fibered-category condition, and the slice criterion developed in section 61. They require no restrictive 0-admissibility assumption, so even the class of all arrows satisfies them. This forces the main conceptual revision: the argument supplies a general modelizing theory rather than a unique characterization of ordinary weak equivalences. Grothendieck sketches coefficient-dependent, twisted and truncated cohomological variants, adopts $H_{\underline W}$ or $(Hot)_{\underline W}$ for the localized category, and asks whether ordinary weak equivalences form the smallest acceptable class, so that other examples would be localizations of the usual homotopy category. He also entertains finer notions, including simple homotopy types, as possibilities rather than results. These reflections motivate section 63's terminology and the later work with fundamental localizers. The marginal note on p.~159 is consequential: it explicitly says that the impression is incorrect and suggests instead the biggest class under an extra requirement; the two words crossed by body typing remain doubtful in the checked text. The witness omits this qualification, as well as the inserted possibly before refinements, and therefore should not be used to turn the exploratory paragraph into a settled assertion. Three smaller source slips remain visible on p.~157: the definition prints $i_Wi^*$ instead of the previously defined $i_Ai^*$, the three-map diagram is called (12) instead of (11), and saturation is called condition b) although it was a). The lower comparison row has primed objects, which the supplied witness row loses. The guide identifies these dependencies without changing the source formulas or replacing the author's subsequent change of mind with a modern theorem.

\subsection*{63. Relative asphericity and terminology under revision (pp.~159--163)}
Grothendieck revisits terminology after allowing general classes $\underline W$. In a category $M$ with a final object $e$, an object with $X\to e$ in $\underline W$ is renamed a $\underline W$-object, and $M(\underline W)$ denotes the full subcategory of such objects; a map is $\underline W$-aspheric when every base change belongs to $\underline W$. The change responds to a real ambiguity: in a topos, relative weak equivalence to its final object need not make the slice topos absolutely aspheric unless the original topos is itself aspheric. For categories the induced morphism of presheaf topoi leads to the weaker condition that every $A_{/b}\to B_{/b}$ belongs to $\underline W$, called weak $\underline W$-asphericity; it does not require arbitrary base change in $(Cat)$. Page 162 compares these notions through the category-of-elements functor: under the slice criterion f), a map $u:F\to G$ of presheaves is universally in $W_A$ exactly when $i_A(u)$ is weakly $\underline W$-aspheric. This makes section 61's frequently used criterion into an explicit standing hypothesis and prepares section 64's localization axiom. The section ends by distinguishing the interval homotopy structure $h_W$ from the canonical $h_M$ of p.~123: their equality is expected in the intended examples, but the author expressly allows a setting in which additional assumptions would be unnecessary and the chosen structure must instead be specified. The terminological decision is itself qualified in the margin of p.~160 and again on p.~164, where the author keeps saying $\underline W$-aspheric for a category whose map to $e$ lies in $\underline W$. The witness's $G'=A\to G$ on p.~162 must be read against the scan's lowercase representable $a$, and its $A_F\to A_G$ omits the visible slice slashes. The same page prints the class $W_{(Cat)}$, not merely $(Cat)$. At p.~163 the overwritten verb is have, not the earlier transcription's doubtful had; the homotopy-structure discussion is not a retraction of the earlier definition. These distinctions explain the subsequent axioms without imposing a uniform modern meaning on every occurrence of aspheric.

\subsection*{64. Basic localizers, optional axioms and a degeneracy claim\newline (pp.~163--168)}
This section fixes the general framework prepared in sections 61--63. A saturated class satisfies identities, two-out-of-three and the two-sided composite condition; strong saturation means that it is exactly the inverse image of the isomorphisms of its localization, and is sought as a consequence rather than imposed initially. The provisional list consists of L1 saturation, L2 interval homotopy invariance, L3 the final-object condition, its special case L3-prime for the interval, L4 the slice localization criterion and L5 the fibered-category criterion. A basic localizer is then defined by L1,L3,L4, which supply the homotopy axiom and recover the previous test-category theory, initially with final objects in the values of test functors. Dropping that final-object restriction motivates L5 and the next section's promised comparison theorems. The oblique note on p.~165 qualifies the implication from L5 to localization: it invokes $W=W^o$, and otherwise asks for L5 plus a dual axiom concerning cofibrations. This is materially different from the earlier transcription's supposed reference to p.~30, and the witness omits the qualification. Grothendieck next considers optional requirements: 0-admissibility and its weaker connectedness condition, preservation of sums and products by localization, invariance under passage to opposite categories, and a stronger relative slice criterion L4-prime. A French note on p.~167 says that the duality assertion is true for fundamental localizers and refers to the duals of L5 and L4-prime; the overlapped words and final label remain doubtful in the checked source. The section concludes by extending all test terminology relative to a fixed $\underline W$, while reserving $(Hot)$ for ordinary homotopy types and using $(Hot)_{\underline W}$ in general. A substantive gap occurs in the proposition on p.~166: equality of all parallel arrows after localization does not imply that the localization is terminal. Let $q:(Cat)\to[1]$ record whether a category is empty, and let $W=q^{-1}(\operatorname{Iso}[1])$. This class consists of all functors between nonempty categories together with the identity of the empty category, is saturated, and satisfies L2 because product with the nonempty interval preserves emptiness. The natural maps from a nonempty category to the one-object category exhibit $W^{-1}(Cat)\simeq[1]$, with distinct empty and nonempty objects. The discrete two-point category is universally $W$-aspheric and makes any parallel pair homotopic for $h_W$, whereas its two points are not connected by the canonical homotopy relation. Thus $h_W\ne h_{(Cat)}$ although the localization is not terminal. The source argument correctly collapses parallel arrows but overlooks the remaining empty-to-nonempty distinction; the readers retain the statement. Smaller cautions are p.~163's composition $gf$, reversed to $fg$ by the witness; the source's basic modelizer in the propositions on pp.~166--167, where the class under discussion is a basic localizer; and p.~167's literal for any $A$ in $C$, where the product criterion requires arbitrary $A$ in $(Cat)$. The clearly typed so si on that page alone is corrected as an explicitly recorded letter transposition. These limitations belong alongside the evolving hypotheses, before the cleaner test-functor review beginning on p.~168.

\subsection*{65. Model-preserving functors and the general test-category theorem\newline (pp.~168--179)}
After the long investigation of homotopy structures, Grothendieck returns to the comparison problem announced around p.~88. A model-preserving functor both reflects and preserves the chosen weak equivalences and induces an equivalence of localizations; neither category is assumed beforehand to be a modelizer. With a fixed basic localizer $\underline W$, he reviews the category-of-elements adjunction $i_A\dashv i_A^*$, the class $\underline W_A=i_A^{-1}(\underline W)$, and the difference between absolute asphericity of $A_{/F}$ and universal asphericity of the projection $F\to e_{\widehat A}$. The latter is tested on products with representables, while general universally weakly equivalent maps are tested on representable base changes. Part A characterizes closure of aspheric presheaves under products; the handwritten addition in Definition 1 also requires $A$ itself to be aspheric before calling its presheaf topos totally aspheric. This is a substantive added hypothesis, retained in the readers. The subsequent discussion explains how contraction of representables and asphericity of interval products can verify the criterion, and proposes the Mayer--Vietoris condition L e) for a cover of a category by two full subcategories closed under passage to sources. The marginal note calls these open subobjects and identifies their presheaf topoi as open subtopoi, not merely open sets. Whether the extra axiom follows from the preceding axioms remains a question here. Parts B--D characterize weak, local and ordinary test categories through the adjunction, the Lawvere contractor and products with an interval; the right marginal addition on p.~173 records a further local criterion, preservation by $i_A^*$ of weak equivalences as universally $\underline W_A$-aspheric maps, and its restricted final-object version. On p.~174 Grothendieck explicitly notices that a hypothesis on $i_A^*$ may have been omitted in the theorem on p.~46; the guide records this reconsideration without silently adding the hypothesis to the earlier text. Part E extends weak and ordinary test functors to $i:A\to M$ for a category endowed with a saturated class $W$: the three basic requirements are model preservation by $i^*$, the weak-test property of $A$, and its asphericity. The induced composite $f_i=i_Ai^*$ provides an equivalent comparison. Theorem 1 then specializes to $M=(Cat)$, assumes L5 and contractibility of every $i(a)$, and characterizes the test property by universal asphericity of $i^*(\Delta_1)$ together with asphericity of $A$. The proof uses nonemptiness to separate the interval endpoints and the key lemma on p.~156 for model preservation; it does not need a final object in every $i(a)$. The intervening contractibility criterion is weaker than the false final-object converse on p.~151 discussed in section 60. The local variant and corollary on p.~178 separate the interval condition from the additional asphericity of $A$; a boxed marginal claim about a model-preserving functor from $(Cat)_{/A}$ is explicitly marked to check and remains provisional. Part F assumes total asphericity and proves Theorem 2: the test property is equivalent successively to model preservation, preservation of weak equivalences, preservation of aspheric objects, preservation of contractible objects, and the single test on $\Delta_1$. The revised implication diagram and the generating-contractor argument give the equivalence of the last two tests; universal asphericity and contractibility are related using the total-asphericity assumption. This completes the general comparison whose significance section 66 reassesses. Several literal formulas need care: p.~169 omits the prime on the target localization and later prints $\underline W_{A^\wedge}$ where the defined class is $\underline W_A$; its displayed target $A_{/G}$ is a typed correction over $A_{/F}$ that the witness misses. On p.~172 the incomplete correction of the to a symbol with subscript $A$, and the doubtful nearby two-$A$ formula, are not reconstructed from mathematical expectations. In the remark on p.~177, condition a) literally repeats $(i_{/a})^*:M\to(A_{/a})^\wedge$, also used in b), although the explanation that $A$ is a local test category calls for $i_{A_{/a}}^*:(Cat)\to(A_{/a})^\wedge$; the source expression is preserved and this likely index/domain slip is confined to the guide. Page 178 omits the subscript 1 from a displayed $\Delta$, and p.~179's condition (iv) literally says $\underline W_A$-objects, with its explanatory parenthesis supplying the asphericity meaning. The strict-total-asphericity qualification on p.~173 should also be read with section 64's caution about the degeneracy argument. These distinctions preserve the developing hypotheses and the author's corrections without converting the review into a retrospectively uniform theorem statement.

\subsection*{66. Asphericity, universal equivalences\newline and the revision of interval terminology (pp. 180--187)}
The entry of 9.4 begins by reassessing the long route to the ``key result'' of section~65. Grothendieck wants to compare a category with a presheaf category and to construct canonical modelizers; the restriction to functors taking values in $(\mathrm{Cat})$ has made that route indirect. Its psychological gain has been the realization that contractible objects and contractibility might recover the weak equivalences. He returns to the program of section~56, specifically part d) on p.~138 after the work on b), while acknowledging that the preceding development has not followed his original priorities (pp.~180--181). The next section will immediately reconsider the emphasis again: aspheric maps organize comparison between modelizers, whereas the contractibility program addresses a particular canonical construction. The terminology on these pages should therefore be read as an active revision, with its dated qualifications, rather than as a settled list of interchangeable definitions. The geometrical uses of asphericity for spaces, topoi and categories are to govern its extension to new settings; the earlier discussion of topoi on p.~37 supplies a reference point. For categories, ordinary asphericity of a functor is tested after the maps $Y/y\to Y$, whereas universal asphericity asks for the corresponding weak equivalence after every base change $Y'\to Y$. These conditions agree for an object viewed over the final category. Their analogues relative to a basic localizer remain subject to the localizer hypotheses in use (pp.~181--182). A presheaf $F$ on $A$ is called $\underline W$-aspheric when its category of elements $A/F$ is aspheric; weak equivalences of presheaves are likewise detected through the categories of elements. The localization axiom implies that an aspheric map in $A^\wedge$ is universally a $\underline W_A$-equivalence with respect to base changes in $A^\wedge$. This quantifier ranges over the presheaf category; it does not automatically assert the result for every arbitrary base change of the corresponding functor in $(\mathrm{Cat})$ (pp.~182--183). Three tests must consequently be kept distinct: absolute asphericity of $A/F\to e$, relative asphericity of $A/F\to A$, and the weaker assertion that this latter functor belongs to $\underline W$. For arbitrary $A$, the first two do not imply one another. With $A$ the discrete category on two objects, a presheaf supported by a singleton at just one object has an absolutely aspheric category of elements, but its projection fails the relative test at the other object. Conversely, the terminal presheaf has identity projection to $A$, hence is relatively aspheric, although its category of elements is the nonaspheric two-point category. If $A$ is aspheric, the first and third tests agree by the two-out-of-three property. The terminal presheaf and the terminal category cannot be interchanged in this argument (pp.~183--184). In a general pair $(M,W)$ Grothendieck now reserves a separate name for the maps universally in $W$, writing $UW\subset W$, and withholds the word aspheric until an appropriate class of localizing maps has been specified. Universal equivalences contain isomorphisms and are stable under composition, but need not themselves be saturated. The handwritten revision from ``is'nt'' to ``need not be'' changes an asserted universal failure into a possibility, and must be applied when reading the final claim (p.~184). The related terminology of contractors is also reconsidered. He proposes reserving contractor for a stronger notion suited to strict test categories, and renaming the earlier contractors of p.~120 multiplicative intervals with a left unit and a left zero. The source's phrase ``left zero object'' is retained. The proposed alternative intersector does not adequately describe an internal-Hom example. These changes of name do not withdraw the contraction argument of section~51 (p.~185). The discussion then returns to the \v Cech condition L e) of p.~171 and reduces its elementary two-piece form to the category $C_o=(b\leftarrow a\to c)$. A cover by two cribles supplies a functor to $C_o$; the asphericity conditions on the cribles and their intersection make that functor aspheric, and the relevant localization axiom makes it a weak equivalence. The source leaves an empty slot between $C'$ and $C''$ in one statement where the comparison witness supplies an intersection sign. The mathematical role is that of an intersection, but the readers preserve the actual blank. The functor $C_o\to\Delta_1$ carrying $a,b$ to $0$ and $c$ to $1$ is fibering, with fibers $\Delta_1$ and $\Delta_o$. The fibration axiom L 5), together with the interval asphericity requirement, therefore gives this Mayer--Vietoris form of L e). The suggestion that L 5) should imply all reasonable \v Cech criteria is an expectation beyond the argument just supplied, not a proved general theorem (pp.~185--186). The connectedness condition previously discussed on p.~165 is next promoted to L 6): $f\in\underline W$ must induce a bijection on $\pi_o$. Equivalently, the component functor factors through $(\mathrm{Hot}_{\underline W})$. Its limited role in the earlier key result no longer seems a reason to treat it as accessory; Grothendieck explicitly distinguishes this requirement from the weaker statement that the associated homotopy structure is the canonical one on $(\mathrm{Cat})$ (p.~186). The entry of 11.4 finally abandons the qualifier weak in ``weak homotopy interval''. A homotopy interval for a specified structure $h$ is henceforth simply an interval whose two endpoint sections are $h$-homotopic. When $h=h_W$, this is broader than the earlier p.~132 requirement that $I\to e$ lie in $UW$. The stronger condition remains available and useful, but it no longer defines the unqualified term. The requirement $\mathrm{Ker}(\delta_o,\delta_1)=\emptyset_M$ is renamed separating rather than separated, to avoid confusion with a separate notion of separation for objects of $M$ (pp.~186--187). The source evidence is consequential here: ``from now on'', ``page 132'' and ``till this very minute'' are typed interlinear additions; their pen placing curves do not cancel the adjoining live words. In the opening line of p.~187 only the n of an is canceled, leaving a live, and the two quotation pairs around homotopy interval have separate machine cancellations. The readers retain doubtful overprint layers, the clipped canceled cell on p.~184, the unresolved oval/parenthesis patch on p.~185, and the recurring graphic marks around h without assigning them an unsupported pen provenance. The isolated typed slips gontext and notes on p.~184 are treated in the corrected forms only, using the same words correctly typed on other reviewed pages; neither those corrections nor the explanatory distinctions above introduce new mathematical hypotheses into the source.

\subsection*{\S67. Aspheric maps replace contractibility as the organizing notion (pp.\ 188--191).}
Grothendieck reorients the modelizer program: the crucial objects are aspheric objects and weak equivalences, and its natural comparison maps are aspheric functors. He explicitly judges the earlier restriction to contractibility and special test functors to have been an unnecessarily long detour (p.\ 188), while retaining its possible relevance for ``canonical modelizers.'' The discussion now allows arbitrary small categories, rather than only test categories. For a basic localizer $\underline W$ he introduces $(\mathrm{Hot}^{W}_{A})$, also written $(\mathrm{Hot}_{A})$, as the localization of the presheaf category by $\underline W_A$; the construction is then organized by inverse image $i^*$ and the comparison morphism $\lambda_i(F):A_{/i^*F}\longrightarrow B_{/F}$, interpreted through the corresponding induced topoi (pp.\ 189--190). This prepares the promised description of asphericity structures and testing functors at the end of the section (p.\ 191); the next section interrupts that development. The dated marginal correction of 3 May on p.\ 189 expressly retracts the claim that Quillen's weak equivalences are precisely $\underline W_A$, saying that Quillen's are stronger. It is therefore essential to read the original assertion together with the margin. The witness changes handwritten ``or'' into equality in (1), drops the subscript $A$ on $\underline W_A$, omits the handwritten domain and codomain after the action on $(a,p)$ on p.\ 190, and prints the struck word ``real'' on p.\ 188. These differences concern notation, functorial context, or the author's revision; none licenses an unrecorded mathematical repair of the typescript.

\subsection*{\S68. Mathematical work, the Teichmüller tower, and the necessity of groupoids (pp.\ 191--195).}
Dated 27 May, the section explains the six-week interruption in the asphericity discussion through a retrospective of the author's changing mathematical interests. Its mathematical center is the ``Lego-Teichmüller construction game'': cutting and gluing surfaces, the tower of groupoids $T_{g,\nu}$, and the arithmetic action on their profinite completions (pp.\ 191--192, 194--195). The account of \emph{La Longue Marche}, tame topology and ``unfolding'' of stratified structures supplies the context for the anabelian project; the intervening reflection on creative interest, ambition and alternating periods of work belongs to the argument's historical presentation and must not be omitted as an aside (p.\ 193). The final page gives the technical lesson: reducing several natural base points to one can destroy visible symmetries and obscure a presentation by generators and relations, whereas groupoids preserve the geometry. Van Kampen situations supply the comparison, and the modular ``multiplicities'' are identified as topoi whose points have finite automorphism groups (p.\ 195). This is a return to the text's groupoid motivation, rather than a proof of the asphericity claims left pending in \S67; \S69 resumes a retrospective of the technical scratchwork. The section explicitly revises an earlier preference for groups or ``exterior groups'' in favor of groupoids (p.\ 195). The witness has several ordinary wording errors here---including ``case'' for ``care'' on p.\ 195---while the underlining and handwritten exclamation in ``topoi!'' preserve the emphasis of the typescript.

\subsection*{\S69. Derivators and integration of homotopy types (pp.\ 196--202).}
The technical retrospective explains why an ordinary derived or homotopy category retains too little structure: traces, determinants and the compatibility of iterated constructions motivate the system of localized diagram categories $D_A(I)$, contravariant in the indexing category $I$, with inverse-image functors expected to admit both adjoints (pp.\ 196--198). This develops the modelizer program recalled in \S67 and connects it with the geometric motivations of \S68. For categories, a covariant diagram can be integrated by forming its associated cofibered category; the comparison with localized slice categories yields the model for ``integration of homotopy types'' (pp.\ 198--199). The name ``derivator'' is introduced on p.\ 200, together with the internal stability expressed by $D(I\times J)$ and the aspiration to a corresponding homotopy theory of derivators. The geometric example is the system of tubes around strata of Mumford--Deligne compactifications: their $K(\pi,1)$ homotopy types, with Teichmüller-type groups, suggest an algebraic description from a diagram of fundamental groupoids (p.\ 201). A diagram in $(\mathrm{Hot})$ alone is explicitly said to be insufficient; a diagram in a model category is required. Grothendieck retracts his recent expectation that finite limits exist in $(\mathrm{Hot})$ (p.\ 201), and \S70 immediately examines the extra factorization condition proposed for the two adjoints (pp.\ 200, 202). Significant witness differences include $T$ replaced by $\mathscr D$ in the statement about projective limits (p.\ 198), $q=f\circ p$ replaced by an equation beginning with $g$, and ``basic modelizer'' replaced by ``basic localizer'' (p.\ 199). These readings are clear on the images and are preserved rather than repaired from the witness.

\subsection*{\S70. Base change, duality, and the limitations of $(\mathrm{Cat})$ (pp.\ 202--206).}
Starting from the desired two-stage factorization $X\to Z\to Y$, Grothendieck asks for preservation of weak equivalences by base change along the second map and co-base change along the first, with either factor optionally a weak equivalence (p.\ 202). The goal supplies the extra condition proposed for derivator operations in \S69. The investigation of the base-change side leads to cohomological smoothness and properness of functors, characterized through asphericity; fibered categories are smooth, cofibered categories proper, and passage to opposite categories interchanges the two conditions (p.\ 203). The pairing of presheaves and copresheaves suggests a broader duality of topoi, but the claim that cosheaves on an arbitrary topos form a topos remains a question (pp.\ 203--204). The co-base-change side is less tractable: open or closed immersions in $(\mathrm{Cat})$ fail the hoped-for general assertion, whereas pushouts along monomorphisms in $A^{\wedge}$ have the required homotopy meaning through the integral of the associated span (pp.\ 204--205). The stated smooth/proper exceptions and the possible special case of an immersion that is also a weak equivalence must be distinguished from that rejected general assertion. This forces a return to elementary modelizers and Quillen factorization questions; the author finally accepts the Kan condition after describing his resistance to it, while retaining dissatisfaction with the reduction to semisimplicial methods (p.\ 206). These questions become part B of the program in \S71. The witness omits the monomorphism hook in the p.\ 204 square and replaces the actual objects $Y,X,Y'$ of the p.\ 205 integral diagram by categories of elements. The scan-confirmed diagrams preserve the distinction between these two levels of the construction.

\subsection*{\S71. A provisional program: modelizers, derivators, and topoi (pp.\ 206--212).}
The program dated 30 May organizes the preceding scratchwork into four tasks. Part A returns to the asphericity structures and canonical modelizers suspended in \S67; part B develops the base-change, co-base-change and factorization questions of \S70, with homotopy integrals, amalgamated sums and Mayer--Vietoris sequences as the connecting constructions (pp.\ 206--207). A handwritten transposition on p.\ 207 makes strong saturation one of the properties expected to follow from the first three basic-localizer axioms. Part C develops the derivator proposal of \S69 while asking what additional structure is needed for external Hom's and basic invariants. It then distinguishes group objects modeling pointed connected homotopy types from the truncated output of group objects in $(\mathrm{Cat})$, and asks which modelizers satisfy a Dold--Puppe theorem (pp.\ 207--209). The proposed trilogy consists of semisimplicial, cubical and hemispherical complexes, the last also called globular complexes by Brown; their boundary and degeneracy operations are related to regular cellular subdivisions of spheres. Multicomplexes give multiple chain complexes, so the simple Dold--Puppe statement does not extend unchanged. The classification of the three series and the need for enough degeneracy maps are presented as questions to investigate, not proved theorems (pp.\ 209--210). Part D, ``Back to topoi,'' seeks relative homotopy categories $D(T)$ and their possible opposite counterparts, extending the $f^*,f_!,f_*$ formalism from indexed categories to topoi. The discussion keeps existence of derived adjoints tentative and moves from objects to proobjects when a topos fails to determine an essentially constant homotopy type (pp.\ 210--211). Local asphericity then reconnects the program with \S35; the revised distinction between modelizing and strictly modelizing topoi is recalled from p.\ 68, and the section ends with the Lawvere-element condition for a finitely triangulated contractible space (p.\ 212). The next section begins the promised return to the asphericity game. The witness changes $f_*$ to $f^*$ and $Lf_!$ to a subscript-$1$ expression on p.\ 210, and applies a handwritten transposition there without marking it. The duplicated target-only display on p.\ 211 was an error of the draft transcription, absent from the scan; its removal restores the actual three-display sequence.

\subsection*{\S72. Criteria for aspheric functors and the first stability results (pp.\ 213--219).}
Returning to \S67 after the program of \S71, Grothendieck isolates the localizer assumptions needed for asphericity: mild saturation, the condition on categories with final object, and the locality condition Loc~3 (pp.\ 213--214). Proposition~1 identifies the comparison maps on categories of elements, their restrictions over objects, and the corresponding slice criteria; it then distinguishes these equivalent asphericity conditions from preservation of weak equivalences and equivalence after localization (pp.\ 214--216). The extra hypotheses in the implication diagram matter: the discrete-target example separates aspheric fibers from merely nonempty fibers, and the product-projection example shows that inducing the localized equivalence need not determine an aspheric map without the stated assumptions (p.\ 216). Under asphericity, the localized triangle commutes up to the canonical comparison isomorphism; this reconnects the discussion with the earlier ``inspiring assumption'' of \S28 without using that assumption as a proof (pp.\ 216--217). Proposition~2 transfers total asphericity along an aspheric functor and relates the inverse image of a separating homotopy interval to asphericity of the functor. Its handwritten addition makes the source category aspheric, and the long interlinear insertion specifies the homotopy condition (TH 1) of p.\ 50; the corollary deduces strict test categories from a totally aspheric source and a local test target (pp.\ 217--218). Proposition~3 supplies composition and right cancellation, base change along a fibering map, and products (pp.\ 218--219). The base-change proof uses the fibration condition L5, whose derivation from the stated localizer assumptions is expressly postponed; Grothendieck explains that only part~(a) is needed for the immediate asphericity-structure formalism, preventing the announced application from relying on that deferred argument (p.\ 219). Remark~4 denies the opposite cancellation assertion and leads directly to the examples in \S73. The visual collation restores the proof-diagram arrow from (i'') to (iii), rather than to (i'), on p.\ 215; it also identifies the apparent deletion of the t in Hot on p.\ 214 as pen retracing. On p.\ 218 the comparison sign is a handwritten tilde over a bar, with no arrowhead, contrary to the draft transcription and its diff description. These are distinctions between what the image shows and what a mathematically plausible reconstruction might supply.

\subsection*{\S73. Adjoints, coasphericity, and cofinality (pp.\ 220--221).}
The continuation dated 11 June begins by suppressing the qualifying $\underline W$ when one basic localizer is understood. It sharpens Remark~4 of \S72: a functor from the one-point category into an aspheric category is aspheric exactly when its chosen object is initial, illustrating how much stronger a map condition is than asphericity of its source and target (p.\ 220). Proposition~4 enlarges the preceding isomorphism example: every equivalence is aspheric, and more generally a left adjoint is aspheric because its slice categories identify with slices having final objects. Passage to opposite categories then motivates the name ``coaspheric'' for the dual condition, satisfied by the right adjoint. The corresponding undercategories and their initial objects connect coasphericity with a promised all-dimensional cofinality statement, whose dimension-zero consequence is invariance of inductive limits (pp.\ 220--221). The analogy from \S70 returns: asphericity resembles properness with aspheric fibers, and coasphericity resembles smoothness with aspheric fibers. The comparison is explicitly weakened because arbitrary base change does not preserve either condition; the applicable base changes are fibering/smooth or cofibering/proper, respectively. Even an equivalence of categories can have an empty fiber, so its simultaneous asphericity and coasphericity do not supply arbitrary base-change stability. Universal weak equivalences, called ``trivial Serre fibrations'' here, are proposed as a stronger common condition and assigned to a later systematic investigation (p.\ 221). The section then turns to the overdue product discussion of \S74. On p.\ 220 the crossed handwritten heading above Indeed is Remark~5, not the draft's tentative Proof; it is written again at the subsequent paragraph. Whether the placement strokes also delete its typed initial The remains unresolved and is marked as such. The p.\ 221 witness changes singular condition to plural conditions and supplies a title for \S74 absent from the typescript.

\subsection*{\S74. Products, diagonal functors, and canonical comparison of models (pp.\ 221--230).}
The overdue product discussion supplies concrete constructions for the abstract stability results of \S72. The external product of presheaves has a category of elements identified with the product of the two categories of elements; it preserves aspheric objects and weak equivalences in the stated variables, and yields the product assertions for locally aspheric and totally aspheric categories (pp.\ 221--222). The corollaries distinguish local test, test, and strict test categories: the hypothesis on the other factor changes from arbitrary to aspheric to totally aspheric. Mixed multicomplexes illustrate the construction (p.\ 223), reconnecting it with the trilogy in \S71. The comparison functor associated with an aspheric map between modelizing categories leads to a functor from $\mathrm{Asph}(A,B)$ into a category of comparison functors and factors through its fundamental groupoid. Nonemptiness and connectedness suffice for the existence of an isomorphism class of comparison, while a canonical transitive system requires the stronger simple-connectedness conclusion; the text carefully separates these claims (pp.\ 224--225). The diagrams on p.\ 226 distinguish categories, presheaf modelizers, and localized homotopy categories, including the middle presheaf triangle omitted from the witness. The diagonal criterion then identifies total asphericity of $A$ with asphericity of $A\to A\times A$, and permits simultaneous comparison through products of strict test categories (p.\ 227). Proposition~4 treats pointwise products of functors into a category $B$ with binary products and aspheric presheaves $\underline{\mathrm{Hom}}(b,y)$: the total-asphericity assumption on $A$, moved by a handwritten transposition into the hypotheses, is essential to the stated construction (p.\ 228). The external-product corollary is sharpened by revision to an equivalence: when $i$ is aspheric, $i\boxtimes i'$ is aspheric if and only if $A'$ is aspheric. Its proof uses a lemma on a presheaf over $A\times A'$ whose restrictions to $A$ are aspheric. The resulting category-of-elements projection is a weak equivalence by the fibration condition L5, explicitly recalling the claim and deferred justification on p.\ 219; this dependency should remain visible (p.\ 229). Contractibility of the objects of $B$ is discussed as a stronger practical hypothesis. Finally, closure of $\mathrm{Asph}(A,B)$ under binary products makes its nonempty case connected with the needed homotopy strength, giving the canonical transitive system of isomorphisms between the localized functors and the canonical functor announced on p.\ 225 (p.\ 230). A handwritten margin points forward to the generalization in \S82, p.\ 272. The image-based edition retains the unusual $X\longrightarrow B_{/X}$ display on p.\ 224, the missing slash in the final category-of-elements notation on p.\ 226, and the visibly typed product x and lowercase subscript o on p.\ 228. The hemisphere-like symbol on pp.\ 223 and 229 remains a doubtful approximate glyph; neither mathematical expectation nor the witness settles its exact drawn form.

\subsection*{\S75. Asphericity structures, testing functors, and weak equivalences (pp.\ 230--237).}
The generalities of \S\S72--74 now become a structure on a possibly large $\underline U$-category $M$. After making the universe convention explicit, Grothendieck defines a class $M_{as}$ of aspheric objects through a small category $A$ and a functor $i:A\to M$ that lands in $M_{as}$ and detects it by the presheaf $i^*(x)=\mathrm{Hom}(i(-),x)$ (pp.\ 230--231). Such an $i$ is a “testing functor.” Every small full subcategory supplies an asphericity structure, and Proposition~1 plus its corollary shows that any structure can be generated this way; enlarging the generating full subcategory within $M_{as}$ preserves the structure (p.\ 232). The following terminology discussion distinguishes testing functors from the earlier test functors of pp.\ 175--176. The essential condition depends on the possibly large category $M_{as}$ itself, rather than on its embedding in $M$, so the author tentatively adopts $M_{as}$-aspheric, or simply aspheric, as a synonym. His use of both “aspherical” and “aspheric” in this discussion is preserved (p.\ 233). Proposition~2 defines $W=(i^*)^{-1}(\underline W_A)$ and proves independence of the testing functor by passing to a common small full subcategory. Mild saturation is immediate; strong saturation remains conditional on the corresponding result for the basic localizer (p.\ 234). When $M$ has a final object, calling the structure aspheric means that this final object belongs to $M_{as}$; the resulting criterion recovers aspheric objects from the maps to the final object. The small-generator descriptions then characterize both the admissible weak-equivalence sets and the asphericity classes. Without the aspheric final-object hypothesis, a weak-equivalence set need not determine a unique asphericity structure. The analogy with contractibility and homotopism structures in \S\S51--52 motivates a promised construction from contractibility to asphericity and a uniqueness assertion for the intended canonical examples (p.\ 235). Proposition~4 recasts the criteria of p.\ 214 as six conditions: three about aspheric objects and three about weak equivalences. The object conditions are equivalent and imply the others; the revised converse requires both $A$ and the asphericity structure to be aspheric. All six conditions become equivalent when the specified generating category $B$ contains the final object. If no $B$ is fixed, enlarging one to contain $i(A)$ and the final object yields equivalence of the first five conditions (pp.\ 236--237). These added hypotheses, some typed interlinearly and some handwritten, are central to the proposition and were partly concealed by malformed insertion macros in the draft. Section~76 immediately tests the formalism on examples. Two visible mathematical slips are retained for the reader to assess: the universe paragraph says “two objects of $U$” where the witness substitutes $M$ (p.\ 230), and Proposition~2 says “Then $M$ is” a mildly saturated subset of $\mathrm{Fl}(M)$ after defining $W$, where the witness prints $W$ (p.\ 234). Proposition~3 and its corollary also call the relevant map “aspheric,” while the explicit revisions of Proposition~4 replace that word by “a weak equivalence” (pp.\ 234, 236); the edition does not silently harmonize the earlier wording. Several correction layers remain uncertain, including the usual-localizer subscript at the preceding section boundary and overtyped letters in these definitions and proofs.

\subsection*{\S76. Basic examples and total asphericity (pp.\ 237--239).}
Three examples test the definitions of \S75. For $M=(\mathrm{Cat})$, the $\underline W$-aspheric categories form the basic asphericity structure, generated through $a\mapsto A_{/a}$ for a weak test category $A$. The shift from the basic modelizer to this structure puts aspheric objects before weak-equivalence arrows. The distinction matters for functors into $(\mathrm{Cat})$: asphericity supplies the induced localized functor, whereas the weak-test requirement asks for a homotopy-model comparison. Theorem~1 of p.\ 176 suggests equivalence of these conditions when the values are contractible, and the author identifies a comprehensive answer to the “silly question” of \S46, p.\ 95, as the motivation for the present formalism (pp.\ 237--238). The second example takes topological spaces with the weakly contractible spaces as aspheric objects; the geometric realization of simplices supplies the testing functor, and the text appeals to Gabriel--Zisman for the stronger modelizer statement. Grothendieck adds the previously omitted assertion that the first two structures are aspheric and later notes that the resulting weak equivalences are the usual ones (p.\ 238). The third example is $A^\wedge$ with its usual aspheric presheaves, generated by the canonical fully faithful inclusion of any small $A$. Its weak equivalences are $\underline W_A$, while its asphericity is equivalent to asphericity of $A$; the modelizing property is not asserted for arbitrary $A$. A functor $A'\to A$ is aspheric precisely when its composite into this presheaf structure is $M_{as}$-aspheric (pp.\ 238--239). This leads to “totally aspheric” structures: $M$ admits finite products and finite products of aspheric objects, including the empty product, are aspheric. Proposition~5 identifies this condition with the existence of a small generating subcategory stable under finite products, and with a testing functor from a category whose presheaf category is totally aspheric. It reconnects the product arguments of \S74 with the general formalism and prepares the localized canonical functor of \S77 (p.\ 239). The visual check restores a missing typed a before the repaired l in “realization,” retains the literal thatnthis and appheric typing, and distinguishes typewritten interlinear additions from pen text. The upward margin arrow at the end of the examples had been misrepresented by three brackets in the draft; its direction is restored, without attributing a mathematical meaning to the layout mark.

\subsection*{\S77. The canonical homotopy functor and the meaning of modelizing (pp.\ 239--241).}
An asphericity structure now determines more than its weak-equivalence class: it supplies a canonical functor from $\mathrm{Hot}_M=W^{-1}M$ to $(\mathrm{Hot})$. Choosing a testing functor $i:A\to M$, Grothendieck uses the composite through $A^\wedge$ and its category-of-elements functor; the inverse-image descriptions of the weak equivalences permit localization, while conservativity is stated conditionally on strong saturation (pp.\ 239--240). To compare two choices, he places their images in one small full subcategory $B$ of $M_{as}$, obtains the comparison isomorphisms from aspheric functors, and then enlarges $B$ to prove independence. A common $B$ for three choices gives transitivity, completing the canonical construction promised in \S74. Formula (13) on p.\ 240 actually names its composite $i_M$, while the following localized functor carries the extra choice index; the witness's harmonization to $i_{M,i}$ in (13) is not adopted. The two sloping comparison signs there are hand-drawn bars with tildes, without arrowheads. The new definition calls the structure modelizing when this particular canonical functor is an equivalence; this is stronger than merely requiring an abstract equivalence between $W^{-1}M$ and $(\mathrm{Hot})$. In the presheaf example it says exactly that $A$ is a pseudo-test category (p.\ 241). The author immediately hesitates over the terminology and proposes a parallel to the hierarchy of pseudo-test, weak test, test, strict test categories and contractors, but leaves the final qualification as question marks. His deletion of “progressively stronger” in that discussion remains visible; the witness prints the deleted words as live text. Section~78 therefore begins with the modelizing name expressly provisional, and uses the canonical construction to revisit the observation and silly question of \S46 rather than treating the proposed hierarchy as already settled.

\subsection*{\S78. Test functors separated into asphericity and modelizing conditions (pp.\ 241--246).}
The provisional definition of modelizing from \S77 supplies the setting for a reformulation of the observation and silly question of \S46. Theorem~1 assumes a modelizing structure $(M,M_{as})$, a pseudo-test category $A$, a functor $i:A\to M$ with aspheric values, and the existence of a final object of $M$; the last hypothesis is a typed interlinear addition, not an original underlined parenthetical sentence. It equates the asphericity and weak-equivalence criteria of Proposition~4, p.\ 236, with the condition that $i^*$ induce an equivalence between the localized categories (p.\ 242). A lemma says that a final object in a modelizing structure is aspheric. Its supporting corollary characterizes $M_{as}$ as the objects whose images under the canonical map to $(\mathrm{Hot})$ are final. The proof uses strong saturation of the basic localizer, announced as a consequence of Loc~1--3 but deferred to part~V; the author explicitly offers the additional asphericity assumptions on the two final objects as a way of avoiding the feared circularity before that proof is supplied (pp.\ 242--243). The canonical triangle from \S77 then proves the localized equivalence. The proof's literal reference to conditions through (ii') is retained even though the statement speaks of six conditions through (ii''). The theorem prompts a revision of the terminology inherited from \S65. The evolving pseudo-test-functor definition first describes model-preserving composites, then separates the intended conditions into: $A$ is pseudo-test, $(M,M_{as})$ is modelizing, and $i$ is $M_{as}$-aspheric. The discussion revises both its hypotheses and its formulation on the page; these are stages of the argument rather than interchangeable completed definitions (pp.\ 243--244). Passing to weak test functors strengthens the condition on $A$ to weak test. Passing to test functors requires a test category and asphericity both of $i$ and of all its slice-induced functors. Definition~2 calls the latter family of slice conditions locally aspheric, and the conjunction with asphericity of $i$ totally aspheric. The local criterion additionally requires that $i$ factor through $M_{as}$ and that $i^*(x)$ be aspheric over the final presheaf, with $x$ restricted to a generating subcategory containing $i(A)$ when desired (p.\ 245). The Remarks paragraph draws the conceptual conclusion: the independent asphericity conditions have wider relevance than the combined test-functor terminology. Yet the earlier theorem on $(\mathrm{Cat})$ still has one advantage, namely its criterion involving only $i^*(\Delta_1)$ under the contractibility hypothesis on the values of $i$. Recovering such a minimal criterion requires relating asphericity structures to the contractibility structures of \S51, which becomes the subject of \S79 (p.\ 246). The typed minimal is inserted before criterion without deleting generating in the preceding line. Most of the final sentence on p.\ 246 is physically covered by a laid-over sheet; its tentative transcription remains doubtful and is not treated as fully recovered evidence.

\subsection*{\S79. From contractibility to asphericity and the general test-functor criterion (pp.\ 247--255).}
The entry dated 20 June supplies the missing bridge identified at the end of \S78. Proposition~1 removes a smallness assumption on a generating full subcategory $N$: the existence of a small generating part, a small inclusion detecting asphericity on $N$, and the assertion that $(N,N)$ is an asphericity structure are equivalent. The name “aspherator” emphasizes that the condition belongs to $N$ itself, while its dependence on the chosen universe and basic localizer remains explicit (p.\ 247). The author then strengthens the standing convention for contractibility structures to require a small generating full subcategory and recalls the finite-product assumption on $M$. Proposition~2 makes $M_c$ generate a totally aspheric structure $M_{as}$; its clause about a chosen generating category $C$ acquires the additional handwritten hypothesis that $C^\wedge$ be totally aspheric (p.\ 248). The proof uses finite-product closure, contractibility of the induced presheaves, and Proposition~5 of p.\ 239. The following remark distinguishes this construction from arbitrary asphericity structures and asks about a canonical largest contractibility structure generating a given one. Its claim about every previously encountered modelizing structure is immediately retracted using nonstrict test categories, then narrowed to modelizing totally aspheric structures (pp.\ 249--250). The second part recalls locally and totally aspheric presheaves and the homotopy interval structure generated by the locally aspheric ones. The source's reference to part IV, \S\S54--55, is retained as written. Six equivalent compatibility conditions H1--H4' express preservation of homotopies, homotopisms, intervals, or contractible generators by a product-preserving functor; Remark~2 explains exactly what generation by $C$ means (pp.\ 250--252). Theorem~1 applies this to $i^*:M\to A^\wedge$ when $i:A\to M$ lands in $M_c$: compatibility with the homotopy structures, local $M_{as}$-asphericity, and the same presheaf criterion restricted to $C$ are equivalent. Its corollaries add asphericity of $A$ for total asphericity, and the test-category and modelizing requirements for a test functor. A separating multiplicative interval supplies a sufficient criterion forcing $A$ to be a local test category; later margin notes question smallness of $C$ and say that multiplicativity is unnecessary, pointing to p.\ 348, Remark~1 (pp.\ 252--253). The proof abandons a boxed localization argument in favor of a lemma and corollary about testing on generators; the blank relation positions in the deleted draft are not completed editorially (p.\ 254). Remark~3 specializes to totally aspheric $A$ and obtains a strict test category under the interval hypothesis, closing the promised restatement of the earlier result. Its corollary reference is corrected by hand from 2 to 3, and a separate margin note points to a stronger version on p.\ 347 (p.\ 255). Section~80 then reverses the direction of inquiry, asking what homotopy structure is recovered from the associated weak equivalences.

\subsection*{\S80. Recovering homotopy from weak equivalences and the canonical-modelizer proposal (pp.\ 255--258).}
The construction $M_c\mapsto M_{as}$ from \S79 is set beside the reverse-looking association $W_a\mapsto h_{W_a}$ of \S54. The first direction distinguishes contractibility from the weaker homotopy interval and homotopism structures; it also makes clear why the associated weak equivalences $W_a$ should not be confused with the generally smaller class $W_c$ of homotopisms. The author recalls how the aspheric-object viewpoint emerged after the earlier “awkward approach,” and corrects and supplements the references to that development (p.\ 255). For an asphericity structure with aspheric final object, $W_a$ and $M_{as}$ determine each other. The intervals generating $h_{W_a}$ are those for which every projection $x\times I\to x$ lies in $W_a$. Under total asphericity this is equivalent to asphericity of $I$, an observation promoted from a tentative expectation to an immediate statement. Two questions then ask whether passing from $M_c$ to $M_{as}$ and back recovers the same homotopy structure, and which extra conditions on $M_{as}$ ensure that the recovered structure comes from a contractibility structure generating it (p.\ 256). A third question, now redundant with the preceding observation, is boxed and deleted; its formula reads $UW_a$, with both letters visible (p.\ 257). Before answering, Grothendieck resumes the “idyllic picture” of canonical modelizers from \S50, incorporating \S\S57 and 59. A finite-product category $M$ with a functor $\pi_o:M\to(\mathrm{Sets})$ receives an interval structure from projections inducing bijections under $\pi_o$. If this structure comes from $M_c$, then $M_{as}$, $W_a$, and the canonical functor $W_a^{-1}M\to(\mathrm{Hot})$ follow. The further compatibility between $\pi_o$ and the composite through $(\mathrm{Hot})$ remains to be imposed; the mild localizer condition L~a) on p.\ 165 is needed to define the final connected-components functor (pp.\ 257--258). When $M$ is totally 0-connected, $\pi_o$ is its intrinsic connected-components functor, and the resulting asphericity structure is modelizing, the construction deserves the name canonical modelizer because it depends on the category itself. The closing Remark calls the definition essentially complete while explicitly retaining the compatibility question. This unresolved point, and the two recovery questions, lead into \S81. The visual collation distinguishes the typed interlinear hypotheses from their pen locators and preserves the literal “homotopisms into homotopisms” on p.\ 257, where the witness substitutes “isomorphisms.”

\subsection*{\S81. Recovering contractibility and comparing the three families of notions (pp.\ 259--266).}
The entry dated 25 June announces the positive answer to the first recovery question of \S80: a contractibility structure $M_c$ on a category with finite products is recovered from its associated asphericity structure by taking the objects contractible for the homotopy interval structure generated by all aspheric intervals. Thus $\mathrm{Cont}(M)\to\underline{W}\text{-}\mathrm{Asph}(M)$ is injective when the basic localizer satisfies Loc~4; the proof is deferred to \S82. The four homotopy notions recalled in (2)--(4) place this assertion beside the earlier inclusions from contractibility to interval structures and homotopism structures, with the latter equivalent to homotopy relations (p.\ 259). Changing the basic localizer from $\underline{W}$ to a coarser $\underline{W}'$ induces a canonical change of asphericity structure, and the triangle on p.\ 260 makes this construction compatible with the inclusion of contractibility structures. Loc~4 requires weak equivalences of small categories to induce bijections on connected components. Among localizers satisfying it, $\underline{W}_o$ is the coarsest; the finest basic localizer $\underline{W}_{\mathrm{oo}}$ is introduced as their intersection and its identification with the usual weak equivalences is promised for part V. The strongest recovery statement is therefore the elementary case associated with $\underline{W}_o$ (pp.\ 260--261). Writing $M_o$ for the 0-connected objects and $W_o$ for the corresponding maps, (17)--(20) show how either determines the other and how both are recovered from the connected-components functor $\pi_o:M\to(\mathrm{Sets})$. The first line of (19) literally omits “is” before “a one-point set.” A handwritten marginal note objects to the name 0-connected maps and proposes 0-équivalences, but its crowded continuation remains partly unreadable; the later asterisk on p.\ 266 refers back to this objection (p.\ 262). The retrospective discussion explains how the special canonical connected-components functor delayed recognition of this general construction and why much of part III is superseded by part IV, while its homotopy notions remain essential (pp.\ 263--264). The three families are test notions, homotopy notions, and asphericity notions. Test and asphericity notions depend on the basic localizer; homotopy notions are absolute. Their order in a deductive account would reverse much of their order of discovery. Contractibility becomes the hinge between homotopy and asphericity: it is both the strongest of the four homotopy structures and a specially constrained asphericity structure (pp.\ 264--265). The revised displays then enumerate ten mutually determining data, including $M_c,J_c,W_c,R_c$, the 0-connected and $\mathrm{oo}$-aspheric objects and maps, and the two canonical functors to sets and homotopy types. This count is reached through visible revisions from four to three asphericity-flavored data and from eight through nine to ten structures. Two further relations $R_o,R_{\mathrm{oo}}$, and more generally $R_{\underline{W}}$, compare maps having equal images under the corresponding functors; whether such a relation alone recovers $M_c$ is explicitly left unclear. The inclusions and commutative comparison in (26), followed by the localization chain in (27), complete the section (pp.\ 265--266). The checked source preserves the typed “oo” notation, identifies the underlined $\pi_o$ in (26), and distinguishes pen arrows from typed diagram objects. Its revision record also retains the five-dot reinstatement of “being” on p.\ 259, the phrase transposition on p.\ 260, and the canceled draft displays on p.\ 265.

\subsection*{\S82. Proof of recovery and consequences for internal Hom objects (pp.\ 267--272).}
Theorem~1 supplies the proof promised in \S81. Under Loc~4, with finite products and the stipulated small full generating category, the contractible objects for the homotopy interval structure generated by all aspheric intervals are exactly the original $M_c$. The argument compares $M_c$, the recovered contractible objects, and $M_a$, then reduces the essential point to Lemma~1: any two sections of an aspheric object are already $M_c$-homotopic (pp.\ 267--268). Passing to the elementary localizer and a small full generating category $C$ reduces the problem to its presheaf topos. Lemma~2 restates the finite-chain argument from \S60, p.\ 149: in a totally 0-connected topos with the stated section and generator assumptions, two sections of a 0-connected object can be joined through a finite chain of objects of $C$. The proof of this lemma is left to the reader. The displayed inclusion $\underline W_o\subset\underline W$ on p.\ 268 is retained as printed, although the witness reverses it; it should be read alongside the coarseness discussion on pp.\ 260--261. The next application assumes the relevant internal Hom objects exist. For a pointed object $X$, the endomorphism object $\underline{\mathrm{Hom}}(X,X)$ has the identity and the constant map as endpoints and supplies a multiplicative interval. The resulting equivalences connect contractibility of $X$, contractibility or 0-connectedness of internal Hom objects, and homotopy equivalence of the constant-map inclusion $Y\to\underline{\mathrm{Hom}}(X,Y)$ (pp.\ 269--270). These are consequences of the recovered homotopy structure, not new definitions of the asphericity class. The source uses $\pi_o,\pi_1$ for the endpoints in one display; that notation is preserved. Lemma~3 compares homotopy classes of maps from $Y\in C$ with the corresponding classes after applying $i^*$, referring back to p.\ 144. Together with the availability of sections of contractible generators, this gives Proposition~1: homotopy classes of sections map bijectively to connected components. Its internal-Hom corollary identifies homotopy classes of maps $X\to Y$ with the connected components of $\underline{\mathrm{Hom}}(X,Y)$ when that object is representable (pp.\ 270--271). Proposition~2 generalizes the product argument of \S74, p.\ 228, under internal-Hom hypotheses. The revisions add total asphericity of the indexing category and alter several occurrences of the generating category and object variables; the canceled final criterion remains visible in the diplomatic account (p.\ 271). Setting $B=C=M_c$ yields the corollary that the product of an aspheric functor with a functor valued in $M_c$ is aspheric. The author explicitly wonders whether the internal-Hom existence assumption for pairs of contractible objects is superfluous. Hence the category of aspheric functors valued in $M_c$ is stable under binary products. If nonempty, it gives the canonical functor $\mathrm{Hot}_{M,\underline W}\to\mathrm{Hot}_{A,\underline W}$, defined up to unique isomorphism through the transitive system recalled from p.\ 230 (p.\ 272). This closes the recovery proof and its product application; the next section asks for a converse characterization of the asphericity structures to which it applies.

\subsection*{\S83. A tautological characterization and an unresolved functor question (pp.\ 272--273).}
Grothendieck returns to the question omitted from the preceding theorem: which asphericity structures arise from contractibility structures? He first imposes total asphericity, then forms the interval structure $J$ generated by intervals whose objects lie in $M_a$ and calls its contractible objects $M_c$, warning that this set need not itself be a contractibility structure (pp.\ 272--273). The two conditions are that $M_c$ generate $J$, equivalently that two sections of any object of $M_a$ be $M_c$-homotopic, and that $M_c$ generate the given asphericity structure through a small full subcategory with aspheric inclusion. The first condition ensures that $J$ does arise from a contractibility structure, necessarily this $M_c$. The characterization is expressly called a tautology. A parallel request asks when a functor $\phi:M\to\mathrm{Hot}_{\underline W}$ is isomorphic to one associated with a contractibility structure, especially when it reduces to $\pi_o:M\to(\mathrm{Sets})$. The text offers no useful necessary-and-sufficient criterion and leaves open whether suitable exactness and internal-Hom assumptions would produce one. Its final underlined question asks whether such a functor can have nontrivial automorphisms, connecting this inquiry to the “inspiring assumption” of \S28 (p.\ 273). The checked text preserves the visible overtyped ending of “generated,” the parentheses around Sets, and “writing it down”; it does not replace these with the witness's readings.

\subsection*{\S84. A provisional notion of morphism and the direct-image problem (pp.\ 274--280).}
The entry dated 29 June begins with letters to Faltings and Porter, then describes the unexpectedly rapid emergence of a proposed morphism notion after prolonged scratchwork (p.\ 274). For two asphericity structures and a functor $f^*:M\to M'$ with left adjoint $f_!$, the proposition parallels the six conditions of \S75, p.\ 236: preservation and detection of aspheric objects, and corresponding assertions about weak equivalences. Its extra assumption (Awk) requires an aspheric testing functor into $M'$ whose composite with $f_!$ lands in the chosen generating part of $M$. The page begins with a generating category $B$ but uses $C$ in that auxiliary clause; the source wording is retained. A triangle through $A^\wedge$ reduces the equivalences to the earlier proposition. The revised conclusion first equates five conditions under the aspheric-final-object hypotheses, then all six if the final object belongs to the chosen generating category (p.\ 275). The canonical localization functors give a second triangle, and the corollary proves that a morphism between modelizing structures induces an equivalence of localized categories. The author calls this the answer to the “silly question,” with the source reference \S45, p.\ 95 (p.\ 276). The rest of the section examines the restrictions rather than treating the proposed definition as settled. The left-adjoint requirement is compared with preservation of projective limits. Condition (Awk) is reformulated through a small generating subcategory of $M'_a$ taken by $f_!$ into $M_a$, or through the generating property of $(f_!)^{-1}(M_a)\cap M'_a$. Replacing it by the stronger condition $f_!(M'_a)\subset M_a$ would pair it neatly with $f^*(M_a)\subset M'_a$, but its reasonableness is unresolved (pp.\ 276--277). Question~1 asks whether the left adjoint $i_!:A^\wedge\to M$ attached to an aspheric functor preserves aspheric objects. In the presheaf case this becomes Question~1': when does $i_!:A^\wedge\to B^\wedge$ preserve them? Dropping asphericity of $i$ reduces the issue to the inductive limit of the representables $i(a)$ over an aspheric category. A cofibering example with 0-connected fibers shows why asphericity of the indexing category alone does not suffice; asphericity of $B$ must at least enter the discussion (pp.\ 277--278). The proposed counterexample on p.\ 279 concerns an amalgamated sum $a\amalg_c b$. The monomorphism case recalls \S70, but the author explicitly has not completed a general counterexample. Neither the extra initial-object condition nor total asphericity appears convincing to him; imposing a strict test category leads to a further question he does not know how to handle. The checked edition preserves the inward arrow directions of diagram (9), despite the surrounding reference to a dual category, and the capital $B$ in the inline map $c\to B$. Finally, for $M=(\mathrm{Cat})$, the canonical slice functor $a\mapsto A_{/a}$ supplies the positive example when $A$ is a weak test category, with both associated functors modelizing and the left adjoint preserving aspheric objects. The standard inclusion $\Delta\hookrightarrow(\mathrm{Cat})$ is the decisive unresolved test: its left adjoint is not generally modelizing, and preservation of aspheric objects still has to be checked (p.\ 280). At the opening of \S85 on the same page, Grothendieck retracts the provisional morphism notion because stability under composition has not been established; this later objection governs how the preceding proposition and corollary should be read.

\subsection*{\S85. Product preservation and the limits of left exactness (pp.\ 280--285).}
The entry dated 30 June begins with Porter's reprints and the author's surprise at Thomason's announced closed model structure on categories. It immediately retracts the preceding proposed morphism notion: stability under composition has not been secured. Requiring $f_!$ to preserve aspheric objects would supply that stability, but the standard inclusion $\Delta\hookrightarrow(\mathrm{Cat})$ remains a difficult test. The author already knows that its left adjoint does not preserve all weak equivalences; he now asks whether it preserves aspheric or at least contractible objects (pp.\ 280--281). Finite-product preservation would prove the latter because the images of the generating simplices are contractible. Proposition~1 supplies the needed criterion for a cocomplete $M$ and $f:A\to M$: under commutation of the product functors with all small inductive limits, the binary-product comparison for $f_!:A^\wedge\to M$ need only be checked on representables. The hypothesis is visibly strengthened from distribution over arbitrary sums. Under the corresponding base-change hypothesis, the fibered-product criterion reduces the two factors to representables, while retaining a general presheaf as base; preservation of a final object then gives the stated left-exactness corollary (p.\ 281). The abandoned proof on p.\ 282 leaves unfinished relation positions inside its cancellation. The replacement proof writes presheaves as inductive limits of representables and uses the displayed interchange of products with two inductive limits. Passing to slice categories handles the fibered-product part. Remark~1 then identifies a limitation: reducing the base itself to an object of $A$ does not yet establish preservation of all fiber products without an additional monomorphism argument; a margin note points forward to \S89 (pp.\ 282--283). Proposition~2 assumes full faithfulness of $f$ and generation of $M$ by strict epimorphisms, expressed by $P\simeq\mathop{\mathrm{colim}}_{A_{/P}}f(x)$. It obtains finite products and the fiber products whose base comes from $A$. The proof uses the full-subcategory identification of the relevant slice categories. Remark~2 refers to SGA~4, volume~1, expos\'e~I, \S7, especially Proposition~7.2 on p.\ 47, for the connection between strict generation and full faithfulness of $f^*$ (pp.\ 283--284). For $(\mathrm{Cat})$, Proposition~3 applies these exactness hypotheses and gives the finite-product conclusion together with the restricted fibered-product assertion. The nerve provides the familiar $\Delta$ example, but the extension to every full subcategory containing $\Delta$ carries the manuscript reservation “hasty!” (p.\ 284). Remark~3 doubts full left exactness even for $\Delta$, appealing to the generators-and-relations description cited to Gabriel--Zisman and the difficulty of preserving monomorphisms. The later margin note explicitly says that $i_!$ is not left exact (p.\ 285). Section~86 therefore proceeds from finite-product preservation and its contractibility consequence, rather than treating the stronger left-exactness claim as proved. The checked source retains literal objedt and the unlabelled pen curve beside hence; neither is silently replaced by the witness's smoother reading.

\subsection*{\S86. Bimorphisms of contractibility structures and the remaining morphism problem (pp.\ 285--290).}
Finite-product preservation from \S85 gives $i_!(A^\wedge_c)\subset(\mathrm{Cat})_c$ for the specified strictly generating subcategories of contractible categories containing $\Delta_1$. The text distinguishes this from the doubtful preservation of all aspheric objects. When both asphericity structures arise from contractibility structures, the condition $f_!(M'_c)\subset M_c$ suffices for the earlier purpose and is stable under composition. The results of \S81 and Loc~4 also allow it to be tested by the apparently weaker requirement that the relevant generators map to 0-connected objects (pp.\ 285--286). This replaces the troublesome (Awk) of \S84 and returns the six equivalent asphericity criteria. The resulting condition can be expressed through the induced functor between the aspherators of contractible objects (p.\ 286). The theorem on p.\ 252 converts the corresponding inverse-image condition into $f^*(M_c)\subset M'_c$, producing a definition independent of the basic localizer. The author calls an adjoint pair satisfying these two inclusions a “bimorphism,” and changes the heading Proposition to the handwritten “Scholie” (p.\ 287). Part (a) reduces the inverse-image inclusion to testing generators for 0-connectedness; the analogous direct-image reduction additionally assumes finite-product preservation. Part (b) says that a bimorphism makes $f^*$ both preserve and detect the aspheric objects and weak equivalences for each basic localizer. Part (c) gives the triangle with the canonical functors of \S77, and hence an equivalence on localizations when both structures are modelizing. Part (d), assuming the direct-image inclusion and Loc~4, equates the bimorphism condition with the indicated inclusions, asphericity of $f_{!c}$, and preservation of an aspheric testing functor valued in $M'_c$ (p.\ 288). The checked page retains its literal $M_{\underline W'}$ in the opening notation and $\underline W_M$ in the final line. The triangle's handwritten $f^*$ has no bar, although the nearby inserted localized functor does. The comments call the Scholie a long tautology in light of \S\S81--82, but emphasize the care required to separate these relationships. A product-preserving $f_!$ also acts on strict homotopy localizations; Grothendieck doubts the geometric relevance of that localized map except in the exceptional situation where both adjoints preserve the weak-equivalence models and induce inverse equivalences. The Scholie replaces the “seducing” proposition of \S84 as his answer to the question of \S46, almost 200 pages earlier. Whether finite-product preservation should enter the definition itself remains a terminological question (p.\ 289). More fundamentally, a satisfactory notion of morphism for arbitrary asphericity structures is still missing, as is a convincing way to dispense with the left adjoint. The closing paragraph explicitly shifts the emphasis toward contractibility structures and entertains the possibility that asphericity structures are principally an important by-product of them (p.\ 290). These reservations qualify the confident language about the bimorphism definition without cancelling it.

\subsection*{\S87. Thomason's Dwyer maps and the independence of homotopy limits (pp.\ 291--292).}
The entry dated 1 July reports an overnight reading of Thomason's preprint. Grothendieck welcomes Dwyer maps as the class of monomorphisms he had been missing: his account describes an open immersion whose map into its closure has a right adjoint, and observes that this implies asphericity. He conjectures that the weaker asphericity condition might itself suffice for preservation of weak equivalences under cobase change, and states the dual proposal for closed immersions and left adjoints. These are proposed extensions, not established results on the page. He hopes they will lead to a more direct model structure on categories and defers the investigation to part~V. The marginal circled interior symbol and the pen repairs in the discussion are retained (p.\ 291). Thomason's remarks on algebraic $K$-theory lead to Anderson's work and to the author's own “integration” and “cointegration” of homotopy types, projected for part~VI on derivators. Grothendieck wants to examine Anderson's indexing-category assumptions, but argues that the notion of a homotopy limit should depend only on weak equivalence. For the direct homotopy limits arising in the cited application, he points to the explicit categorical construction sketched in \S69, pp.\ 198--199, and regards the model structure as unnecessary for that purpose. Cointegration appears less direct in $(\mathrm{Cat})$; he contrasts this with the anticipated geometric behavior of higher groupoid stacks and expects fibrations, characterized intrinsically by preservation of weak equivalence under base change, to enter its categorical construction. The references and evaluations here are the author's account of the papers, not an independent verification of their claims. Section~88 then pauses the development to inventory the questions still outstanding before part~V (p.\ 292).

\subsection*{\S88. An inventory of unfinished questions (pp.\ 292--300).}
Before beginning part~V, Grothendieck lists five subjects still requiring attention. The first is the direct-image problem of \S84: whether an appropriate $f_!$ preserves aspheric objects. His earlier suspicion of counterexamples has not become a result, and he now allows that an unexpected positive statement may exist. The second is the promised systematic treatment of contractors and generators for contractibility, parallel to aspherators for asphericity. Apart from the trivial contractor equivalent to a final category, their strict-test property connects their presheaf homotopy formalism with pure homotopy theory (pp.\ 292--293). The third asks for an autodual homotopy formalism. Interval homotopy degenerates in a pointed category, whereas complexes in an abelian category supply meaningful, autodual notions of homotopy, homotopism, null-homotopy and quasi-isomorphism. The mapping-cylinder comparison belongs within each of these structures; it does not identify the two kinds of equivalence. The analogy between $H^0$ with its shifts and connected components suggests a broader formalism, possibly involving fibers, cofibers and exact sequences. This program is linked to the integration and cointegration of homotopy types, the discussion of derivators in \S69, and program C of \S71, pp.\ 207--210 (pp.\ 293--295). A nearer question concerns the induced structure on $M_{/a}$: the usual aspheric structure on a presheaf slice can fail total asphericity even for a contractor, suggesting that the appropriate equivalences should compare homotopy types fiber by fiber (pp.\ 295--296). The fourth item becomes an extended digression on why semisimplicial homotopy is preserved by an arbitrary functor $A\to B$, although the interval formalism of \S51 assumed preservation of finite products. When products and distributive sums exist, a simplicial set $I$ supplies objects $I_A$, and maps $I_A\times X\to Y$ can be described degreewise by maps from $I_n$ into the sets of morphisms $X_n\to Y_n$, with the corresponding compatibility. The latter description makes sense without those exactness hypotheses. Fully faithful embeddings preserve it; embedding $A$ and $B$ in presheaf categories explains why the induced $f_!$ preserves the particular products $I_A\times X$, since they are sums of copies of $X$, without preserving every finite product. The author therefore contemplates homotopy and contractibility structures, and their morphisms, without global products or final objects. He compares this enlargement with the earlier passage to sites, but does not regard it as urgent (pp.\ 296--299). The checked text retains the differing visible variance marks in the Hom displays rather than harmonizing them. The fifth item asks for intrinsic criteria making categories specified by algebraic structure types into canonical modelizers. Small limits and inductive limits, internal Hom objects in slices, universal and disjoint sums, total 0-connectedness and sections of nonempty objects form the provisional list of properties. The expected examples include elementary presheaf constructions and categories of $n$-stacks and $\mathrm{oo}$-stacks. The desired conclusion is a canonical contractibility structure whose associated usual weak-equivalence structure is modelizing, and then the analogous assertion for every basic localizer. At the end, a separating contractible interval is proposed as a deliberately tentative test conjecture; the author would first seek a counterexample. None of this is presented as an established characterization (pp.\ 299--300). Section~89 interrupts the inventory to repair a specific claim made in \S85 about strict generation.

\subsection*{\S89. Strict generation and the topos obstruction to left exactness (pp.\ 300--302).}
Returning to \S85 after three days, Grothendieck reexamines the claim that a full subcategory of $(\mathrm{Cat})$ containing the standard simplices must be strictly generating. Enlarging a strictly epimorphic family need not preserve strict epimorphicity, so the earlier “a fortiori” was insufficient. He first invokes the stronger universally strict generation furnished by the low-dimensional simplices, citing Giraud's descent article, Proposition~2.5, p.\ 28. He then records a direct argument that enlargement does preserve strict generation for full subcategories, despite the failure for arbitrary families (pp.\ 300--301). The reference to the Glossaire in SGA~4 distinguishes the relevant epimorphism notions. These observations repair the statement on p.\ 284 rather than validate its original justification. The section then settles the earlier doubt about the left adjoint of $\Delta\hookrightarrow(\mathrm{Cat})$. For a cocomplete category $M$ and a small strictly generating full subcategory $A$, the proposition says that $i_!:A^\wedge\to M$ is left exact if and only if $M$ is a topos. The proof identifies $M$ with the full image of its fully faithful $i^*$ and applies Giraud's criterion to the left adjoint of that inclusion. Since $(\mathrm{Cat})$ is not a topos, its $i_!$ is not left exact; the text points to epimorphisms that are not strict or effective (p.\ 301). The corollary locates the failure in fiber products $b\times_F c$ with $b,c$ in $A$ and general $F$ in $A^\wedge$, even while finite products and fiber products over an object of $A$ are preserved in the situation of \S85. This answers the sharper-version question in Remark~1, p.\ 283: additional hypotheses are required (p.\ 302). The explanation on p.\ 302 literally says “if $M$ is strictly generating”; the edition preserves that wording, alongside the proposition's earlier designation of $A$ as the strictly generating subcategory. Section~90 then resumes the unfinished inventory.

\subsection*{\S90. Existence and perfection of test functors,\newline followed by an afterthought on strict generation (pp.\ 302--312).}
The sixth item in the unfinished inventory asks for test functors from a test category $A$ into a modelizing asphericity structure $M$. The usual finest localizer is the strongest case: a functor satisfying the corresponding test conditions also serves for coarser basic localizers. For contractibility structures, the preferred functor lands in $M_c$ and takes the contractible generators of $M$ to contractible presheaves under $i^*$, thereby supplying the bimorphism setting of \S86. This reformulation separates the existence question from the ambient category as far as possible and suggests testing on a small collection of generators, perhaps a single interval (pp.\ 302--303). The image phrase on p.\ 303 literally says “image of A by A”; the checked text does not silently replace its last letter by the witness's $i$. The value of a direct test functor is stronger than the abstract availability of modelizing comparisons. A composite through a presheaf category and $(\mathrm{Cat})$ already transfers homotopy types under the earlier hypotheses, but a direct $i:A\to M$ also gives the external object $\mathrm{Hom}^{(A)}(x,y)$ with value $\mathrm{Hom}_M(i(a)\times x,y)$ at $a$. The discussion links this to Quillen's semisimplicial categories and, through the presheaf enlargement of $M$, avoids assuming all internal Hom objects exist in $M$ itself. Normalizing global sections gives the final-object conditions in (5)--(7). The distinction between a left-exact functor into a slice $A^\wedge/E$ and one into $A^\wedge$ matters here; the earlier composite need not have the latter property (pp.\ 304--305). A further demand is that the left adjoint $i_!$ also modelize homotopy types. The sharper condition that every unit $F\to i^*i_!F$ be a weak equivalence introduces a \emph{perfect test functor}; without the modelizing assumptions, the parallel expression is a perfect aspheric functor. These are intended to give mutually inverse localized adjoints. The ordinary inclusion $\Delta\to(\mathrm{Cat})$ fails this stronger condition, whereas the canonical slice functor associated with a weak test category and Thomason's subdivided construction provide motivating positive examples. The author asks how far these examples extend to arbitrary test categories and contractors, and what special features of $\Delta$ would be responsible if they do not (p.\ 306). This leads to the category $M^A$, its alternative interpretation when $M=B^\wedge$, and the need to distinguish the weak-equivalence structures one can put on it (p.\ 307). The first occurrence of “B is a test functor” is retained literally, while a later occurrence is visibly corrected to “test category.” The dual interpretation of $M^A$ suggests importing Quillen's projective construction for $M^o$ as an injective construction for $M$. Factoring the map from the initial to the final diagram as a cofibration followed by a trivial fibration is proposed as a way to produce a test functor, with the author's explicit qualification “more or less, at any rate imply.” It is a proposed escape from the stalled inductive construction, not a completed proof. Two further questions ask whether Quillen's theorems extend from $\Delta$ to general test categories and whether the localized diagram categories are canonically independent of that choice. The topos case is compared with Illusie's semisimplicial sheaves and projected for part VII. Grothendieck then notices that the first question depends on an asphericity or contractibility structure while the others appear to depend only on the category, and acknowledges that the questions remain insufficiently clear (pp.\ 308--309). The dated comments of 4 July reduce existence questions to elementary presheaf modelizers when a suitable testing functor is already available, with a corresponding reduction in the perfect case. For the standard inclusion $i$, the functors $i_!\mathrm{Sd}^n\alpha$ for $n\geq2$ give the Thomason example; $\alpha:\Delta\to\Delta^\wedge$ is the canonical inclusion. The notation changes from the author's earlier Bar to Thomason's Sd. Grothendieck speculates that analogous subdivision operations, perhaps arising from closed-model factorization, could turn ordinary test functors into perfect ones. This places the existence problem firmly in the projected part V rather than resolving it here (pp.\ 309--310). After the asterism, an unnumbered afterthought returns to \S89. A family $u_i:X_i\to X$ is strictly epimorphic when the induced Hom map is injective and every compatible family, tested against all common-source squares, comes from a map out of $X$. The property depends on the generated sieve, but strictness alone need not survive enlargement. Universal strictness requires stability under arbitrary base change and does survive enlargement; when fiber products exist, the compatibility test reduces to $X_i\times_X X_j$. A inductive-limit cocone is strictly epimorphic (pp.\ 310--311). The proposition then equates four conditions for a small full $A\subset M$: strict epimorphicity of every family of arrows from $A$, universal strict epimorphicity of those families, their presentation of each object as the indicated inductive limit, and full faithfulness of $i^*$. This supplies the definition of a strictly generating subcategory and proves the enlargement corollary by the second criterion. The text cites SGA IV, expos\'e I, Proposition 7.2, p.\ 47, for part of the equivalence and leaves the proof to the reader. Its literal “other three criteria” and the sentence ending the digression are retained. Section~91 begins immediately below (p.\ 312).

\subsection*{\S91. Finite type, boundary operations, and the remaining questions from part II (pp.\ 312--319).}
The seventh item compares finite triangulations of spaces with objects of finite presentation in $\Delta^\wedge$. Grothendieck sketches why the two descriptions should determine the same homotopy types: an ordering of vertices gives one direction, while induction on simplices, piecewise-linear realization and perhaps two barycentric subdivisions suggest the other. The pen additions explicitly require a compact piecewise-linear structure and finite triangulation. For a general category admitting filtered inductive limits, finite presentation means that the covariant Hom functor commutes with those limits. In $A^\wedge$, the finite-presentation objects are finite inductive limits of representables; the author describes this as the 2-universal problem of adjoining finite inductive limits to $A$. The question is when these objects describe exactly the finite-type homotopy types (pp.\ 312--313). The abandoned opening of p.\ 314 is replaced by a proposed counterexample. A finite group $G$ acts freely on an aspheric space $E_G$, with quotient $B_G$. A suitably chosen small full category of aspheric spaces, containing the interval and $E_G$ and closed under finite products, supplies the test category. The finitely presented quotient presheaf $F=E_G/G$ is identified with $i^*(B_G)$; the author argues that its homotopy type is not of finite type when $G\ne1$, using nonvanishing cohomology in arbitrarily high dimensions. This demonstrates why extra restrictions on $A$ are needed in the proposed correspondence. He suggests that quotients of representables remain representable, that each representable have only finitely many sieves, and perhaps that Hom sets be finite. For the converse inclusion, a perfect test functor $\Delta\to A^\wedge$ with finitely presented values would suffice, linking this question back to \S90. These conditions are proposals rather than a proved characterization. The author expects the familiar simplicial, cubical and hemispherical categories, and their finite products, to satisfy the desired finiteness property. He again resists treating the simplicial category as uniquely privileged, while allowing the strict Dold--Puppe theorem as a possible special feature of the familiar three (pp.\ 314--315). The eighth item asks for boundary operations intrinsic to a test category. Computing through a test functor from a familiar category, or through a nerve, does not answer that question. The proposed geometric data are a functor $i:A\to(\mathrm{Ord})$ with finite values, each having a largest element and realizing as a cell, with the complement of that largest element realizing as its boundary sphere. Principal lower ideals must themselves be represented by objects of $A$, through maps $b\to a$; uniqueness in the slice gives the face objects $a_x$ and maps $\partial_x$ for codimension-one elements. Orientations $\omega_a$ then define signs $\varepsilon_a(x)$ by the induced boundary orientation “à la Stokes” (pp.\ 315--317). With the dimension decomposition of $A$, a contravariant additive-valued functor produces the displayed sums $K_n=\coprod_{a\in A_n}K.(a)$ and their signed boundary maps. The manuscript still asks for natural extra conditions ensuring that these complexes calculate the ordinary homology and cohomology of the associated category, also with twisted coefficients, and ensuring that $A$ is a test category. The formulas therefore describe a tentative construction, not a completed comparison theorem (pp.\ 316--317). Degeneracy operations are the next obstacle. The first attempted example is boxed and crossed out on p.\ 317; the composition display inside it has genuinely blank arrow positions. On p.\ 318 Grothendieck explains the failure. He starts with a full category $A_o$ of finite ordered sets satisfying the cell condition, stable under products and passage to principal ideals, and containing the three-point ordered set whose two outer points map to the middle point. He then tries objects consisting of such an ordered set and an orientation. Allowing all underlying maps would preserve the desired product argument but fail invariance of orientations under isomorphisms, an assumption now recognized as necessary for the differential. Enforcing that compatibility destroys the available product and representability argument for the interval. Thus the intended test-functor proof does not go through. The difficulty leads to the opposite possibility: boundary operations might intrinsically single out the simplicial, cubical and hemispherical examples, just as the Dold--Puppe question had suggested earlier. This possibility remains open (pp.\ 317--318). The ninth item returns to \S43. The category of simplices without degeneracies is asserted to be a weak test category but not a test category; the author wants to write down the proof and consider the analogous cubical and hemispherical assertions. He also notes that the hemispherical category itself has not yet been formally introduced, and that its strict-test or stronger contractor property still deserves proof. Six comparison test functors between the three familiar types of complexes would be useful, but he may omit them from a work that is not intended to be an “Elements d'Algèbre Homotopique” (pp.\ 318--319). The other recalled question concerns the ordered set of nonempty finite subsets of an infinite set: a positive weak-test result would support the modelizing role of $(\mathrm{Ord})$. He asks whether that modelizing structure comes from an asphericity structure or even a contractibility structure, and poses the same issue in the general situation of the earlier proposition on p.\ 74. Finally, for basic localizers $\underline W\subset\underline W'$ satisfying Loc~4, he asks whether total $\underline W$-asphericity and total $\underline W'$-asphericity are equivalent, or equivalently whether either is the same as total 0-connectedness. These are the concluding questions of the cote, not claims established by it (p.\ 319).

\subsection*{92. Abelianization, its possible factorization, and universal chain models (pp. 320--331; 134-6/1--12).}
The entry dated 5.7 reorders the unfinished program at the end of part~IV. Questions about contractors, induced asphericity and contractibility, and some remaining material from part~II are to stay in this part; several existence and duality questions are deferred. The organizing issue is whether an arbitrary suitably qualified test category can do the work familiar from the simplicial, cubical, and hemispherical examples. The list distinguishes test-functor existence, closed model structures, independence of derived categories from the choice of indexing category, finite-type conditions, and boundary operations with Dold--Puppe relations (pp.~320--321). The last question motivates the section, but the author now doubts that it should be subordinated to cellular decompositions of spheres. The unexpectedly weak axioms already sufficient for a test category suggest seeking a more flexible construction of chain complexes. The additional ``strict'' qualification and the deletion of ``or contractor'' in the initial comparison are visible revisions to the scope of the claim. The discussion fixes the ordinary localizer of weak homotopy equivalences and calls $\mathrm{D.}(\mathrm{Ab})$ by the name $\mathrm{Hotab}$ (pp.~321--322). Its starting point is the canonical abelianization functor $\mathrm{Hot}\to\mathrm{Hotab}$. Diagram (3) constructs it using a chosen test functor on $\Delta$, the free-abelian-group or ``Whitehead'' functor $\mathrm{Wh}_{\Delta}$, and the Dold--Puppe equivalence. The weak equivalences used on abelian objects are those detected on their underlying semisimplicial sets. Here ``Whitehead's theorem'' names the compatibility needed to descend free abelianization to the localized categories. Different choices of the test functor need not give isomorphic functors before localization; the claimed canonical independence afterward is referred to section~77, replacing a canceled reference to section~82 (pp.~322--323). The corresponding cubical and hemispherical comparisons are proposed with an express qualification about their Whitehead theorem. The following remark also separates the construction of abelianization itself from these stronger equivalence theorems: the ordinary unnormalized chain complex already sends weak equivalences to quasi-isomorphisms and yields the same localized functor. For every small category $A$, formula (6) therefore gives a canonical functor $\mathrm{Hot}_A\to\mathrm{Hotab}$ through $\mathrm{Hot}$; no claim that $A$ is a test category is needed for that composition (p.~323). The new question is whether it factors through the localization $\mathrm{Hotab}_A$ of abelian presheaves, or can be computed by a functor from abelian presheaves to $\mathrm{Ch.}(\mathrm{Ab})$ (p.~324). These are additional requirements. An aspheric functor $j:\Delta\to A^{\wedge}$ gives a route back to the familiar simplicial case, but the text notes that its existence already forces substantial asphericity conditions on $A$. The comparison with a general $j:B\to A^{\wedge}$ leads to diagrams (11)--(14), including an independence question for the resulting factorization. In the contractor case, the argument is referred to the product remarks of section~82, p.~272 (p.~325). This reference has a different role from the corrected section~77 citation above. The live superscript on $j$ in (10') is not fully legible in the scan and remains marked as uncertain in the transcription. Pages~326--328 attempt a construction without assuming such a $j$. Write $u=i^*i_A:A^{\wedge}\to\Delta^{\wedge}$. Although $u$ generally fails to preserve finite products, it preserves fiber products. Passing to the slice over $F=u(e_{A^{\wedge}})=i^*(A)$ gives an exact functor $u_o$ and hence a functor on abelian group objects. The localization morphism determined by $F$ supplies the left adjoint $\alpha^{ab}_!$ used in diagram (17). The lower square commutes by the adjunction argument, but the upper square is explicitly described as failing to look commutative. The proposed remedy is that its comparison map should become a weak equivalence after passage through the lower square. The author immediately revises the claim that this is the exact sufficient hypothesis: the composite functor (18) must also preserve weak equivalences so that it descends to localization. Neither condition is established for arbitrary $A$, and the text admits uncertainty even in the presence of the earlier aspheric functor. Diagram (19) then raises a further question: when $A$ models ordinary homotopy types, does the corresponding abelian comparison also become an equivalence? This is an open prospect in the section, not a conclusion of the diagrammatic construction. The author interrupts this detour on p.~328 and returns to the desired computational description. Diagram (20) asks for a functor $L:A^{\wedge}_{ab}\to\mathrm{Ch.}(\mathrm{Ab})$ whose composite with free abelianization realizes the already defined homotopy abelianization. The distinction matters: the canonical functor to $\mathrm{Hotab}$ is available, while a particular chain-complex model and its factorization through abelian presheaves require further work. Using an aspheric simplicial comparison is called ``cheating'' because it imports the familiar model instead of producing the sought intrinsic construction. General lifting results into a category of models are mentioned as another possible route and then set aside (p.~329). The visible changes from ``abelian'' to ``additive'' for the derived category and from ``chain'' to ``abelian'' for its complexes sharpen the setting without supplying such a theorem. The replacement idea is to vary the coefficient category. For any additive category $M$, a chain construction should give $L_M:\underline{\mathrm{Hom}}(A^o,M)\to\mathrm{Ch.}(M)$ compatibly with additive functors in $M$ (p.~329). The ``envelopping additive category'' $\mathrm{Add}(B)$ is introduced by a universal map from $B$: composition with it induces an equivalence between additive functors out of $\mathrm{Add}(B)$ and arbitrary functors out of $B$. Taking $B=A^o$, the text identifies such a compatible system with a chain complex $L.$ in $\mathrm{Add}(A^o)$. Formula (26), $L(X)=\widetilde X(L.)$, then evaluates this universal complex by the additive extension of an abelian presheaf $X$ (p.~330). The resulting existence problem concerns a complex whose induced square commutes; the associated uniqueness problem asks, for example, about uniqueness up to chain homotopy and about the category of pairs consisting of a complex and a compatibility isomorphism. These remain questions. The universal map is renumbered (23), while the sentence immediately after (24) still prints the reference (22); the source's uncorrected reference is retained. The final revision restricts the apparent example available from a simplicial comparison (pp.~330--331). After first asserting that the construction obtained ``by cheating'' visibly comes from a universal complex, Grothendieck cancels an attempted justification and adds that $j$ must factor through $A$ itself. An aspheric functor $\Delta\to A$ then induces $\mathrm{Add}(j^o):\mathrm{Add}(\Delta^o)\to\mathrm{Add}(A^o)$, and the image of the familiar complex gives the desired candidate $L.$. Landing merely in $A^{\wedge}$ does not justify this particular construction. The section closes at that narrower sufficient situation. The next dated entry begins the promised explicit construction of the additive envelope, rather than resolving the wider existence and uniqueness questions.

\subsection*{93. Additive envelopes, cellular signs, and the proposed standard abelianizator (pp. 331--338; 134-6/12--19).}
The entry dated 6.7 makes the universal category from section~92 explicit. First replace each set of maps in $B$ by the free abelian group on those maps; then add formal finite families of objects, with morphisms given by matrices of such linear combinations (formulas (1)--(4), pp.~331--332). Matrix multiplication extends composition in $B$. The formal sums are deliberately kept distinct from any sums that already exist in $B$, since the desired universal property concerns arbitrary functors from $B$ into additive categories. The text states that the resulting $\mathrm{Add}(B)$ has the required two-dimensional universal property. A coefficient-ring variant replaces $\underline Z$ by a commutative ring $k$ with unit and targets additive categories equipped with $k\to\mathrm{End}(\mathrm{id}_M)$; the word ``commutative'' and the instruction to replace $\underline Z$ by $k$ are typed interlinear additions. This leads to the prospect of $k$-linearization and an associated basic localizer $\underline W_k$, without developing that theory here (p.~332). For $B=A^o$, an \emph{abelianizator} is proposed as a pair $(L.,\lambda)$: a chain complex in $\mathrm{Add}(A^o)$ together with the compatibility isomorphism that makes section~92's square (20) compute homotopy abelianization (p.~332). The induced additive functor on abelian presheaves is available from this universal complex, but preservation of weak equivalences as quasi-isomorphisms is raised as an additional issue. The author considers making that preservation part of the definition if it is not automatic. Only after it holds does localization give the sought functor $\mathrm{Hotab}_A\to\mathrm{Hotab}$. The resulting square (19) of section~92 is then accounted for, while the earlier diagrams of type (15) are still described as less evident. Thus the terminology does not close the existence, uniqueness, or compatibility questions raised in section~92. The cellular construction begins with the three standard categories and their finite products (pp.~333--334). For an object $a$, the ordered set $i(a)$ consists of its subobjects in $A$, explicitly distinguished from sieves on $a$. In the proposed setting it is finite, has a largest element, and gives the dimension of $a$ through the dimension of its geometric realization. The universal term is $L_n=\bigoplus_{\dim a=n}a$, requiring finite dimension fibers because the additive envelope has only formal finite sums. The differential on a summand $a$ is described by a linear combination of maps $b\to a$ from dimension $n-1$; in the standard cases its nonzero terms are precisely the boundary monomorphisms, with coefficients $\varepsilon_{\partial}\in\{\pm1\}$ (formulas (10)--(14)). The revision from ``any'' to ``all'' monomorphisms on p.~334 is therefore substantive. Choosing orientations of the cells $|i(a)|$ determines these coefficients by the stated Stokes rule. The text immediately identifies a difficulty: a subobject can have several representing monomorphisms if isomorphisms are allowed. Its first safe formulation assumes that every isomorphism is an identity. Under those hypotheses, the anticommutativity of successive boundary orientations is invoked to obtain $d_{n-1}d_n=0$ (formula (18), p.~335). If nontrivial automorphisms are admitted, one must choose a representative monomorphism for each boundary subobject and impose the commutativity of the resulting diamonds (19). The page calls this a delicate condition, probably difficult to ensure. An additional factorization of every arrow as an epimorphism with section followed by a monomorphism would connect the construction with section~91 and make the subobject assignment a functor to ordered sets, but the author doubts that the existence of that functor is essential to the abelianizator question. There is also a narrower independence result: for different choices of orientations, the signs comparing them define a transitive system of componentwise isomorphisms of the complexes (22)--(25), p.~336. The circled marginal $\varepsilon_a^{\omega,\omega'}$ clarifies the sign notation. This identifies the orientation-dependent complexes canonically in the safe case. When representative boundary maps also have to be chosen, the text does not establish an analogous independence from those additional choices. The final pages compare the scope of this construction with the test-category program (pp.~336--338). The familiar rigid examples and their finite products satisfy several desirable properties at once: the cellular condition, a canonical abelianizator, and strict test-category behavior. The author wonders whether these ``multistandard'' examples exhaust the possibilities, then contrasts unordered simplices, which retain automorphisms, with $\Delta^f$, which has no degeneracies and is expected here only to be a weak test category. Returning to the full subcategory $A_o$ of ordered sets considered in section~91, p.~318, the text argues that orientations need not be added to its objects or morphisms. Strict inclusions of ordered subsets provide the suggested coherent boundary maps. It accordingly claims that the unchanged category is a strict test category, while qualifying that it is presumably not a contractor. The earlier assumption of closure under finite products is then recognized as too restrictive for the standard examples themselves. A proposed substitute asks for compatible cellular subdivisions of products $|i(a)|\times|i(b)|$ by cells of the form $|i(c)|$. Whether this suffices even to make $A$ a weak test category, and whether the associated complex is an abelianizator, remain open questions in the entry. The concluding remark revises its own terminology while it is being written (p.~338). ``Standard abelianizator'' is acceptable only if the constructed canonical complex actually has the required abelianization property in every case under consideration. ``Dold--Puppe chain complex'' is first offered as another name, with normalization viewed in the universal additive category. The text then corrects the attribution: the name is more appropriate for the \emph{normalized} complex obtained from $L.$, whereas the full complex itself is described as long familiar. Performing that normalization universally also requires enlarging the additive envelope to admit the direct summands associated with projectors. The two opposed marginal bracket arrows mark the end of this remark and the start of section~94; they do not provide a further proof or a resolved replacement for the open conditions above.

\subsection*{94. Recovering homotopy types from abelianization and varying the ground ring (pp. 338--342; 134-6/19--23).}
The emphasis changes from constructing abelianization through different modelizers to asking what extra structure would allow the homotopy type to be recovered from it. For a pointed connected type with fundamental group $\pi$, the first proposed ingredient is the chain complex of the universal covering, considered in $\mathrm{D.}(\underline Z(\pi))$ (formula (26), p.~339). This retains the group action needed for twisted coefficients. The second ingredient is a diagonal $L.\to L.\overset{L}{\otimes}L.$ in $\mathrm{Hotab}$, subject to suitable commutativity and associativity conditions (27). The text suggests combining these structures to include cup products with arbitrary twisted coefficients. Its larger question is whether they suffice to recover the entire homotopy type and therefore the other operations and invariants. That sufficiency is expressly a question which the author says he had not investigated. The recollection of Sullivan's explanation then distinguishes such a diagonal, or the resulting rational cohomology ring under suitable finiteness assumptions, from a differential graded algebra carrying the fuller information (pp.~339--340). The author recalls a rational De Rham construction for spaces or semisimplicial sets and an equivalence after localization and restriction to the appropriate simply connected objects, but repeatedly qualifies his memory. The inserted words ``differential graded'' and ``in the derived category'' sharpen the statement: the relevant comparison is between algebra objects with their differential, up to the stated localized notion of isomorphism. The passage asks why such an algebra need not be recovered from its cohomology algebra with zero differential. It is not asserting that the cohomology ring alone classifies rational homotopy types, and the guide should not turn this remembered explanation into an unqualified theorem with unstated hypotheses. Integral coefficients lead to the next proposal (pp.~340--341). The reported obstruction to using an anticommutative differential graded algebra prompts comparison with the failure of the algebraic Poincar\'e lemma and its divided-power remedy. Grothendieck recalls developing a De Rham formalism with divided powers for arbitrary semisimplicial sets, and considers its dual coalgebra form technically more appropriate. He wonders whether a suitable derived category of this structure could recover homotopy types. The later restatement specifies a bigraduation, and for nontrivial fundamental group it again invokes the universal covering and its group action. These are proposed structures and recollections of earlier work; no reconstruction theorem, equivalence relation on all such objects, or proof of sufficiency is supplied here. The discussion of derived categories and binomial coefficients explains what attracted his attention, rather than completing that missing argument. The feature that drives the remainder of the section is the possibility of defining the algebraic structure over an arbitrary commutative ring, or even a scheme (pp.~341--342). This suggests a notion of homotopy types over a ground ring and a change-of-ring operation. Abelianization provides the familiar model through the dependence of $\mathrm{D.}(k)$ on $k$, but the proposed nonabelian version is less settled. The author recalls expressing some simple spaces, including standard $K(\pi,n)$ examples and fibrations between them, through integral points of finite-type affine semisimplicial schemes. A prospective functor from $\mathrm{Hot}$ to a suitable derived category of such schemes would make naive scheme-theoretic base change $\underline Z\to k$ reflect the expected change of coefficients. The text entertains an equivalence, then stresses that the attempt was brief, was not resumed, and may have been unrealistic. These qualifications belong to the mathematical status of the proposal. The final paragraph tries to extend that picture beyond simply connected examples by using nilpotent fundamental groups (p.~342). It first hopes to recover a suitably restricted finitely presented torsion-free group as the integral points of a finite-type affine nilpotent group scheme, then retreats to a rational group scheme and recovery only up to commensurability. The handwritten ``Possibly'' introduces an anticipated equivalence after localizing the discrete groups at monomorphisms of finite index. On the algebraic side the page says connected nilpotent affine group schemes over $\underline Q$, followed by the further description that their underlying schemes are affine spaces. That latter restriction must remain visible when interpreting the proposed class; the paragraph does not establish the stated equivalence of descriptions or a categorical equivalence. The next section leaves these afterthoughts and returns to the contractor program of section~92.

\subsection*{95. Contractors and the comparison of contractibility with asphericity (pp. 343--351; 134-6/24--33).}
The entry dated 7.7 resumes the three topics announced in section~92, beginning with contractors. Its first definition is canceled because it tacitly requires finite products in $A$, an assumption which excludes the standard test categories that the author wants to include (p.~343). The replacement embeds $A$ in $A^{\wedge}$ and asks that the intervals coming from its objects generate a contractibility structure in which those objects are contractible. This is the condition called Bas~4 in section~51, p.~118. A category satisfying it is first called a \emph{precontractor}; further compatibility with asphericity will define a contractor. The inserted qualification that the ambient category $M$ is stable under finite products belongs to the generalized formulation of this condition. The discarded initial definition and the separate abandoned version of p.~345 remain visible stages of the argument, rather than alternative final definitions. Proposition~1 explains how the structure passes from $A$ to another category $M$ stable under finite products (pp.~344--345). The three sufficient situations are: $A$ itself has finite products and the functor preserves them; an extension from $A^{\wedge}$ exists and preserves the final object and binary products of representables; or the original functor is fully faithful and strictly generating. The discussion relates these conditions to the criteria on pp.~281 and 283 and to full faithfulness of the associated $f^*:M\to A^{\wedge}$. In the last case, identifying $M$ with a full subcategory containing $A$ gives a direct argument without separately requiring all small inductive limits or the existence of $f_!$. The proposition's universal validity in the indicated cases also characterizes the precontractor condition. These distinctions explain why the author replaces an intrinsic product assumption on $A$ by a condition on its presheaf category. Proposition~2 supplies the missing asphericity requirement (p.~345). For a precontractor, let $h$ be generated by the intervals coming from $A$, and compare its contractible objects with the canonical asphericity structure for ordinary weak equivalences. The listed equivalent conditions include total asphericity of $A^{\wedge}$, generation of that asphericity structure by the contractibility structure, the corresponding inclusions of contractible objects, compatibility of $h$-homotopisms with weak equivalences, and equality with the homotopy structure associated with those weak equivalences. The pen changes from ``totally'' to ``locally'' in the object-level conditions are consequential: the initial notation $A^{\wedge}_{\mathrm{totas}}$ is not consistently replaced everywhere, but the revised conditions and Definition~1 on p.~346 use local asphericity for the representables. The resulting \emph{contractor} is a precontractor satisfying these equivalent conditions. The page then records trivial examples, generating subcategories of contractibility structures, and the three standard categories with their finite products. The revision of Proposition~3 separates the existence of a separating interval from the additional asphericity needed for a strict test category (pp.~346--347). A nontrivial precontractor contains such an interval; when it is also totally aspheric, hence a contractor in the terminology here, the strict test-category conclusion follows. The proof uses sections of objects to show that an interval is either separating or has equal endpoints. If every interval from $A$ had equal endpoints, the resulting homotopy relation would collapse to equality and every contractible object would be final, contradicting nontriviality. The subsequent corollary characterizes a trivial contractibility structure by several equivalent collapses; a fifth handwritten condition says that any two sections of a contractible object are equal. The supplied scan cuts off faint traces below that addition, and the input transcription supplies no hidden continuation there. The guide therefore does not reconstruct one. For a nontrivial contractibility structure $(M,M_c)$, the next argument transports a contractible interval with distinct endpoints through $i^*:M\to A^{\wedge}$ and obtains a separating interval (p.~348). With the stated total $\underline W$-asphericity assumptions this yields a $\underline W$-test category, and the stronger assumptions give a strict one. The author explicitly replaces corollary~3 on p.~253: the earlier multiplicative-interval and membership assumptions are declared unnecessary, and the separation requirement reduces to distinct endpoints. The following remark corrects a different possible confusion. The structure $h$ generated by intervals from $A$ must be distinguished from $h'=h_{\underline W_A}$ used in section~79; the latter can depend on the localizer. The diagram's lower adjoint index remains visually uncertain between a heavily retraced shriek and a star, so the transcription does not resolve it from the adjunction argument. Under strict total $\underline W$-asphericity, together with Loc~4, the text then compares both structures with the canonical homotopy structure defined by 0-connected intervals (pp.~349--350). The coarsest localizer under discussion is described by bijectivity on $\pi_o$. The proposed new auxiliary proposition on p.~349 is crossed out and replaced by a reference to Lemma~2 on p.~268; its blank topos names are left blank in the source. The surviving Proposition~4 states the equality $h=h_{\underline W_A}=h_{A^{\wedge}}$, then equates homotopy compatibility of $i^*$, $\underline W$-asphericity of $i$, and the specified 0-connectedness criterion on a generating subcategory of $M_c$. Nontriviality of the contractibility structure gives the strict $\underline W$-test conclusion. The proof on p.~351 refers back to section~79, p.~252, and Proposition~3's corollary. If $M_c$ is moreover $\underline W$-modelizing, the final corollary makes $i$ a $\underline W$-test functor and gives the equivalence between the displayed localizations. The next dated entry returns to the definitions themselves, particularly the treatment of the empty category.

\subsection*{96. Verticality, horizontal topoi, and the limits of the terminology (pp. 351--354; 134-6/33--36).}
The entry dated 8.7 begins by explaining the repeated abandoned pages in section~95: the author had expected a smooth exposition but found the technical adjustments less settled than anticipated. The first terminological issue separates two properties of a topos: every nonempty sheaf has a global section, and the final object has only the empty and full subobjects (pp.~351--352). The typed replacement of ``as being equivalent'' by ``to imply'' is mathematically significant. The classifying topos $B_G$ is then used to distinguish the properties for a nontrivial discrete group. The first property is discussed as a particularly strong global asphericity condition, initially only as far as $H^1$ is concerned; the inserted ``global'' specifies the intended cohomological scope. Several possible names for the second property are considered before the author settles on ``vertical,'' with the image of a small base and a tower of internal structure. The explicit definition excludes the empty topos and requires its open subtopoi to be only the empty and full ones (p.~352). A horizontal topos is defined by generation of its category of sheaves by subobjects of the final object; the topos of a topological space is the principal example (p.~353). The revised statement says that being both horizontal and vertical gives a final topos, represented by the one-point space and the category of sets. The preceding discussion's accidental inclusion of the empty case is thereby removed. The author also notices that the initial precontractor definition formally allowed the empty category, which he does not want to include. There is a scope problem in the accompanying equivalences which the guide must leave explicit. Pages~351 and 353 identify the vertical condition for a topos of the form $A^{\wedge}$ with the requirement that every object of $A$ have a section, and p.~354 repeats the claimed coincidence. If $A$ is an unrestricted small category, the source's own $B_G$ example supplies a counterexample: take $A$ to be a nontrivial one-object group category. Its presheaf topos has only the empty and full subterminal objects, but the regular $G$-set has no fixed point and hence no global section. Thus the converse implication requires an additional restriction on $A$, or a correction to the assertion. The pages do not state that restriction at these occurrences. The transcription preserves the claims; the guide does not turn them into an unrestricted equivalence. Verticality is nevertheless presented as a separate direction from total asphericity or total 0-connectedness, with Lemma~2 on p.~268 and section~60, p.~149 as motivations (p.~353). Grothendieck therefore proposes replacing the earlier phrase ``strictly totally aspheric'' by ``totally aspheric and vertical,'' rather than treating the extra condition as a degree of asphericity. A second terminological conflict remains unresolved: the notions of locally and totally aspheric objects introduced in section~79 do not automatically agree with the terminology obtained by passing to the associated topoi, as the principle of section~66 would suggest. The author elects to keep this useful but unorthodox language until it causes a concrete difficulty. The section closes by returning to the three structures $h,h',h''$ from section~95: they coincide for a contractor, while only $h'$ depends on the chosen basic localizer. Clarifying their general relationship is announced as the next task, then interrupted before the entry dated 12.7.

\subsection*{97. Projective topoi, the three homotopy structures, and two notions of morphism (pp. 354--357; 134-6/36--39).}
The entry dated 12.7 separates the two conditions which the previous discussion had partially conflated. A topos is called \emph{projective} when every epimorphism $F\to e$ to its final object has a section, equivalently when that final object is projective in the category of sheaves (pp.~354--355). Among nonempty topoi, the property that every nonempty sheaf has a section is then the conjunction of projectivity and verticality. The cohomological discussion uses dimension shifting to pass from vanishing of $H^1$ with arbitrary abelian coefficients to vanishing of all higher cohomology; the pen addition ``except for $H^0$'' and the typed word ``global'' delimit the assertion. A further handwritten qualification requires 0-connectedness in the claimed asphericity consequence. The page again asserts that the two original conditions coincide for topoi $A^{\wedge}$, so its inference for arbitrary vertical categories inherits the scope problem identified in section~96. The distinction between projectivity and verticality itself remains explicit and useful. The promised comparison diagram on p.~355 identifies $h=h_{A^{\wedge},A}$, generated by intervals coming from $A$; $h'=h_{\underline W_A}$, depending on the localizer; and the canonical structure $h''=h_{A^{\wedge}}=h_{\underline W_{oA}}$. It displays three \emph{conditional} inequalities $h\leq h'\leq h''\leq h$. Their separate hypotheses are total $\underline W$-asphericity of $A$, the inclusion $\underline W\subset\underline W_o$, and verticality of $A$. The inserted words ``conditional'' and ``on A or W'' guard against reading this as an unconditional cycle. The first hypothesis also bears a circled marginal ``\'ecarter'' connected to its bracket. That mark does not securely establish a deletion, so the diplomatic text preserves the hypothesis, connector, and instruction together. The revised sentence below the diagram states the coincidence of all three structures under the combined assumptions; taking $\underline W=\underline W_o$ recovers the comparison with Lemma~2 on p.~268. The following sentence refers to (1), although the displayed comparison is labelled with a star. This source cross-reference is left unchanged. For a contractor, the three structures are accordingly said to coincide for every localizer satisfying the present Loc~4 (pp.~355--356). If the contractor is nontrivial, the strict test-category and modelizing results give a \emph{canonical modelizer}: the weak-equivalence structure is determined by the canonical homotopy structure. Alongside $(\mathrm{Cat})$ and $(\mathrm{Spaces})$, these presheaf categories are the principal examples available in the text. Stacks are proposed as another substantial family, with no construction or proof supplied in this entry. The point of the comparison is therefore to identify when the modelizing structure depends only on the canonical homotopy relation, rather than on a separate choice of generating intervals. The first proposed definition of a morphism of contractors $f:A\to B$ requires $f^*:B^{\wedge}\to A^{\wedge}$ to preserve the homotopy structures (p.~356). Conditions (2)--(5) express this through preservation of contractible objects, through the representables in $B$, through their 0-connectedness, or through their $\underline W$-asphericity. The A/B exchanges in these displays are visible pen corrections and are essential to the direction of the conditions. The criterion reduces to $\underline W$-asphericity of $f$, and the stronger-looking equality (5') identifies the entire asphericity structure of $B^{\wedge}$ as the inverse image under $f^*$ of that on $A^{\wedge}$. The author recognizes this as an instance of Proposition~4 of section~95, p.~350, after initially presenting it as a fresh surprise. The final page tests a broader situation in which $f$ lands in $B^{\wedge}_c$, without necessarily factoring through $B$ (p.~357). The associated adjoint pair $f_!\dashv f^*$ suggests the two requirements (8), one on $f_!(A)$ and one on $f^*(B)$. Closure under composition leads to the stronger condition (9), $f_!(A^{\wedge}_c)\subset B^{\wedge}_c$, which follows when $f_!$ preserves finite products by the corrected reference to section~85. It need not follow merely from (8), even when the original functor lands in $B$. The author therefore distinguishes two plausible notions, neither claimed to imply the other: he reserves ``morphism of contractors'' for the first and calls the second a bimorphism of the associated contractibility structures, referring to section~86. The entry closes without deciding which will be more useful. The dated section~98 then postpones the topic of induced structures and turns to simplices without degeneracies.

\subsection*{98. Simplices without degeneracies and the remaining asphericity lemma (pp. 357--360; 134-6/39--42).}
The entry dated 16.7 postpones induced asphericity and contractibility structures on slice categories and returns to the assertion announced in section~43: the category $\Delta^f$ of ordered standard simplices with injective maps should be a weak test category. The author recalls having found a simple proof in March without keeping notes, and reports difficulty recovering it (pp.~357--358). Diagram (1) compares ordered and unordered simplices, with and without degeneracies. The ordinary categories $\Delta$ and $\widetilde\Delta$ are contractors; their injective-map variants are said not to be test categories. The claim pursued here is the weaker one for $\Delta^f$, with the analogous result for $\widetilde\Delta^f$ suggested as a possibility. The diagram's lower arrow is literally labelled $\beta f$, and the blank before an isolated subscript $n$ in the following sentence remains unfilled in the scan. The expected common property of all four comparison functors is asphericity, but only $\alpha$ and $\beta$ are treated as checked for the immediate argument (p.~358). The unordered-simplex functor $\beta$ is moreover said to be a morphism of contractors: choosing a point in a finite nonempty set $E$ gives an elementary homotopy from the identity of $\beta^*(E)$ to the associated constant map. The author does not claim that $\beta$ is a bimorphism of the canonical contractibility structures. That stronger question would follow from preservation of finite products by $\beta_!$, which remains unproved here. This is the first concrete use of the distinction made at the end of section~97. The proof strategy identifies the category of injective simplices over $\Delta_n$ with the ordered set of its nonempty subsets (p.~359). Consequently $i_{\Delta^f}$ is canonically isomorphic to the composite of $\beta\alpha$ with the standard unordered-simplex test functor $\widetilde i$ from section~34. The known asphericity of $\widetilde i$, combined with asphericity of the comparison, would yield the required result. The page explicitly fixes the strongest localizer under discussion, $\underline W_{\infty}$ of ordinary weak equivalences. To prove asphericity of $\alpha$, it adjoins an initial empty simplex to form $\Delta'^f$ and identifies each relevant slice with $(\Delta'^f)^{n+1}\setminus U^{n+1}$, where $U$ is the open subcategory consisting of that initial object. The inserted U is typed; the marginal handwritten formula and the canceled exponent mark show the local revision of this description. Two lemmas remain in that reduction. Lemma~1 asserts the asphericity of $\Delta^f$ itself. Lemma~2 says that, for a finite \emph{nonempty} family of pairs consisting of a category and an open subcategory, asphericity of each ambient category and each closed complement implies asphericity of the complement of the product of the opens (pp.~359--360). The nonemptiness qualification is a handwritten insertion. The proof reduces to two factors and applies a Mayer--Vietoris argument to the two closed pieces and their intersection. The resulting set difference $X\setminus U$ is visible in the source and must not be changed into a product. The text then acknowledges that Lemma~1 is where the proof has stalled (p.~360). A heuristic calculation would use the familiar boundary operations for semisimplicial complexes, assuming that the same cohomological formulas apply to face complexes. Applied to the final face complex, and including noncommutative $H^1$ with constant coefficients, this would yield the required asphericity and even establish that of $\alpha$ directly. The entry does not supply that calculation or justify the extension of the formulas at this point. Instead, trying to understand those boundary operations brings the author back to sections~92--93 and to a stronger notion of abelianizator with a prospective existence and uniqueness theorem. Section~99 begins that reformulation before the unfinished asphericity argument is completed.

\subsection*{99. Integrators, arbitrary coefficients, and the finite-sum restriction (pp. 360--366; 134-6/42--48).}
The uneasiness left by sections~92--93 is now traced to their concentration on constant coefficients. Grothendieck's habitual cohomological viewpoint would instead begin with arbitrary sheaves and regard constant or locally constant coefficients as particular cases (pp.~360--361). He therefore puts $B=A^o$ and asks for another interpretation of the same universal chain complex: it should compute cohomology of $B$, meaning its presheaf topos, with coefficients in an arbitrary abelian presheaf on $B$. The variance is essential. The objects whose homotopy types motivated the earlier construction lie in $A^{\wedge}$; the coefficient objects in this reformulation lie in $B^{\wedge}_{ab}$. The inserted ``abelian'' is typed, while the circled marginal $A^{\wedge}_{ab}$ and the inserted phrase locating the earlier abelianizations keep the two sides of this comparison explicit. The additive envelope is first given a concrete realization (pp.~362--363). The canonical inclusion $B\to B^{\wedge}$ followed by free abelianization factors through an additive functor $\gamma_B:\mathrm{Add}(B)\to B^{\wedge}_{ab}$. The universal property identifies abelian presheaves $F$ with additive functors $\widetilde F:\mathrm{Add}(B)^o\to(\mathrm{Ab})$, and formula (5) identifies evaluation on $L$ with $\mathrm{Hom}(\gamma_B(L),F)$. This makes $\gamma_B$ fully faithful and realizes its image as the finite direct sums of free abelian representables. The source's formula (5) prints only $B^{\wedge}$ in the Hom index; p.~364 explicitly specifies that the Hom is taken in the category of abelian presheaves. Likewise, the generalization (3') omits an additive qualifier on its right-hand Hom even though the prose requires additive functors. The guide follows those explicit descriptions while the transcription preserves the displayed notation. The revision from ``can be generalized into'' to ``is a particular case of,'' and the typed replacement of ``abelian'' by ``additive,'' clarify the intended universal-property argument. The image of $\mathrm{Add}(B)$ consists of projectives of finite presentation, but it need not contain every such projective (p.~363). The exact statement is that every projective of finite presentation is a direct factor of one of these finite sums. The typed replacement of ``type'' by ``presentation'' and the insertion ``full sub'' sharpen this description. Passing to the Karoubi envelope, which formally supplies images of projectors, gives the full subcategory of projectives of finite presentation. Its proposed universal property targets additive categories in which every projector has an image. This distinction between finite sums and their retracts will matter when the author asks whether the components of a resolution can lie in the original additive envelope. For a chain complex $L.$ in that envelope, formula (6) identifies $\widetilde F(L.)$ with the componentwise Hom cochain complex (p.~364). The desired computational property is a natural derived isomorphism $\mathrm{Hom}^{\bullet}(L.,F)\simeq R\Gamma_B(F)$ for every abelian presheaf $F$. The text interprets $R\Gamma_B(F)$ as $R\mathrm{Hom}(\underline Z_B,F)$, where $\underline Z_B$ is the constant presheaf, and uses projectivity of the components to compare ordinary Hom with derived Hom. In the projective-resolution setting this leads to the augmentation $L.\to\underline Z_B$ and the requirement that it be a quasi-isomorphism. The page's argument is the familiar one for resolutions in nonnegative chain degrees; it supplies no separate justification for an arbitrary unbounded complex of projectives. The lower index of the projective limit in (7) is visibly corrected to $B^{\circ}$, with a circled marginal repetition, while the older subscript remains partly unreadable. An \emph{integrator} on $B$ is consequently defined to be a projective resolution of $\underline Z_B$ (p.~365). It is \emph{special} when its components belong to $\mathrm{Add}(B)$. Since $\underline Z_B$ need not itself belong to that category, the augmentation is extra structure external to the additive envelope; formula (11) encodes it by the element $\lambda$ obtained from evaluation on $L_o$. The usual comparison theorem gives chain-homotopy uniqueness of integrators, and fullness gives the corresponding statement within $\mathrm{Add}(B)$ for special ones. General integrators exist because the abelian presheaf category has enough projectives. Existence of a special integrator is a further finiteness question: successive kernels need not be quotients of finite sums of representables. The page later says ``strict integrator'' once where the discussion has concerned the defined term ``special''; that wording is retained in the source. The one-object group category supplies the proposed test case. A special integrator would be a resolution of the trivial $\underline Z[G]$-module $\underline Z$ by free modules of finite rank in each degree. The author doubts that this can happen unless $G$ is finite, and a marginal X stands beside that sentence (p.~365). The mark does not establish a deletion or a revised theorem. The proposed necessity of finite $G$ is nevertheless false: for the infinite cyclic group, write $R=\underline Z[t,t^{-1}]$ and use the exact augmented sequence $0\to R\xrightarrow{t-1}R\to\underline Z\to0$. Multiplication by $t-1$ is injective, and its cokernel is the augmentation module, so both nonzero resolution terms are free of rank one. This example leaves intact the author's broader observation that special integrators express a homological finiteness condition, while preventing the guide from adopting his tentative finite-group criterion. The immediate construction relaxes the finite-sum requirement (pp.~365--366). The category $\mathrm{Addinf}(B)$ is formed by allowing arbitrary small direct sums and has the corresponding universal property for functors into infinitely additive categories. It embeds in $B^{\wedge}_{ab}$, its objects remain projective, and every abelian presheaf is a quotient of one of them. Its retracts therefore give all projectives; the Karoubi-envelope formulation is repeated with the appropriate infinite-additivity condition. An integrator with components in $\mathrm{Addinf}(B)$ is called \emph{quasi-special}. Such resolutions always exist, and the text asserts their uniqueness up to chain homotopy inside this larger category. The typed insertion ``up to equivalence'' qualifies the categorical identification; its long placement bracket passes across the following words without deleting them. The last sentence extends the description of augmentation data in (11) to the quasi-special case, allowing the indexing set $I_o$ to be infinite (p.~366). A notation issue must then be kept in view. For $L_o=\bigoplus_{i\in I_o}\mathrm{Wh}_B(b_i)$, a map $L_o\to\underline Z_B$ is specified by an arbitrary family of integers, hence by an element of $\prod_{i\in I_o}\underline Z$. Formula (11) prints $\underline Z^{(I_o)}$, which agrees with the product for finite $I_o$. If its parenthesized exponent denotes the direct sum, the final extension needs that notation changed or interpreted as a product; an unrestricted augmentation need not have finite support. The transcription leaves the formula and the extension as written. The next entry adopts a canonical notation for the homotopy class of the augmented resolution and develops its homological meaning.

\subsection*{100. Intrinsic homology, relative integration, and the comparison with cohomology (pp. 367--374; 134-6/49--56).}
The entry actually dated 17.6 adopts the notation $L^B_{\bullet}\to\underline Z_B$ for an integrator, unique up to homotopy as an augmented complex; the date is retained despite the surrounding July entries (p.~367). Formula (2) identifies $R\Gamma_B(F)$ with $\mathrm{Hom}^{\bullet}(L^B_{\bullet},F)$, and (3) gives the resulting cohomology groups. A quasi-special integrator has components which are arbitrary small direct sums of free abelian representables, while a special one has finite sums in every degree. The disjointness convention for the index sets lets the author suppress the degree on the objects $b_i$. A substantive qualification is needed in (6), p.~368: the scan has a clear handwritten direct-sum sign in the expression for $\mathrm{Hom}(L_n^B,F)$. In general this Hom is the product $\prod_{i\in I_n}F(b_i)$, because a map from a direct sum is an arbitrary family of maps from its summands. The displayed sum is therefore justified for the finite indexing sets of special integrators, but does not give the unrestricted quasi-special formula. This is the same finite-versus-infinite distinction left implicit in the augmentation notation at the end of section~99. The source sign remains unchanged. Formula (7) extends the construction to bounded-below complexes of coefficients by the simple complex associated with the componentwise double Hom complex. The specialization to a constant complex $K_B^{\bullet}=p_B^*K^{\bullet}$ uses the adjunction $p_{B!}^{ab}\dashv p_B^*$ to obtain (10), where Hom is now taken in abelian groups (pp.~368--369). The source's reference to (2') is retained even though the immediately relevant generalization is labelled (7). The resulting complex $p_{B!}^{ab}(L^B_{\bullet})$, consisting of projective and hence free abelian groups, is proposed as an intrinsic representative of the global homology of $B$ and of its homotopy type. Its covariance and the functoriality of (10) are invoked to descend the construction from categories to homotopy types, producing the abelianization functor (12). This should agree with the functor of section~92, p.~321, obtained by a particular test category and the diagram on p.~322. The proposed comparison would establish the same constant-coefficient formula and recover an object of $D^-(\mathrm{Ab})$ from its Hom functor on $D^+(\mathrm{Ab})$. The author explicitly has not checked the latter assertion, particularly its suggested generalization to arbitrary abelian categories; the guide therefore records the route to the comparison rather than presenting the sketch as a completed uniqueness proof. For a functor $f:B'\to B$, formula (14) introduces $L^{B'/B}_{\bullet}=f_!^{ab}(L^{B'}_{\bullet})$ as a relative homology object with projective components (p.~370). It remains quasi-special or special when the chosen source resolution is so. The text asserts that coasphericity of $f$, meaning asphericity of the opposite functor, makes this object an integrator on $B$; its proposed converse is qualified by the choice of the homological localizer $\underline W_{\infty}^{Z}$ and deferred to later notes. The first symbol and the raised correction in (14) retain local overprint doubts, but the prose and surrounding formulas make the intended relative construction explicit. The long first remark contrasts the author's experience with cohomology on general topoi and the more direct geometric history of homology (pp.~370--373). Universal-coefficient formulas, twisted coefficients on a topological variety, and the older sheaf-theoretic comparison results motivate the question, without supplying a single equivalence that preserves every coefficient system and extra structure. The text attributes the asymmetry in general topoi to the lack of enough projectives, despite the availability of enough injectives for abelian sheaves. Presheaf topoi $B^{\wedge}$ provide the favorable case with enough projectives. For morphisms induced by a functor between the indexing categories, the adjoint triples $f_!\dashv f^*\dashv f_*$ and $f_!^{ab}\dashv f^*_{ab}\dashv f_*^{ab}$ explain the additional left-derived operation (p.~372). The set-valued $f_!$ need not take group objects to group objects, so the abelian left adjoint is constructed by its own adjunction. The author proposes calling $Lf_!^{ab}$ \emph{integration} and $Rf_*$ \emph{cointegration}, revising his former use of the first name for the latter functor (p.~373). Projection from a disjoint family of copies of $B$ gives the motivating direct-sum versus direct-product picture. The inserted ``presumably'' qualifies the assertion that integration is as meaningful as cointegration, while the long pen addition says that the author's continuing personal emphasis on cohomology no longer amounts to confidence in its technical superiority. Possible dualities between homology for $f$ and cohomology for $f^o$ remain suggestions. The second remark allows a commutative coefficient ring $k$, replacing abelian groups and additive envelopes by $k$-modules and their $k$-linear counterparts; the exposition stays over $\underline Z$ only for simplicity (pp.~373--374). The next dated entry returns to the postponed computational description of $f_!^{ab}$ and the relation between integrators and abelianizators.

\subsection*{101. Computing integration and the limits of internal adjunction (pp. 374--377; 134-6/56--59).}
The computational description postponed in section~100 begins with the fact that $f_!^{ab}$ preserves small inductive limits and is right exact (p.~374). Every abelian presheaf on $B'$ has a presentation by two objects of $\mathrm{Addinf}(B')$, so the functor is determined by its restriction to this category of special projectives. Diagram (1) identifies that restriction with the tautological infinitely additive extension $\mathrm{Addinf}(f)$. Formula (2), p.~375, makes the description explicit: apply $f$ to the indexing objects of a direct sum of free abelian representables and take the corresponding direct sum on $B$. The free-abelianization identity (4) and diagram (5) explain why this gives $f_!^{ab}$; the page refers to SGA~4, I,5.8.3 for the analogous construction with other algebraic structures. Restriction to finite sums gives $\mathrm{Add}(f)$ as well. Diagram (6) collects the original categories, their finite and infinite additive envelopes, their presheaf categories and free abelianization into one commuting configuration. The circled instruction to center and spread it out and the unattached outer pencil curve are graphic annotations, not additional functors. The scan also shows that the diagonal from $B$ ends at the lower $B^{\wedge}$, a destination that must be preserved when the diagram is reconstructed. The passage to derived functors uses bounded-above complexes of these projectives (p.~376). Two visible corrections are decisive: ``below'' becomes ``above'' in the degree bound, and ``quasi-isomorphisms'' becomes ``homotopy equivalences'' in the localization description (8). On bounded-above complexes of projectives the two classes coincide, so the corrected formulation describes the usual homotopy model of $D^-(B^{\wedge}_{ab})$. Applying $\mathrm{Addinf}(f)$ componentwise then computes $Lf_!^{ab}$. The displayed right side of (9) literally has $\mathrm{Addinf}(L'_{\bullet})$, without the $f$ argument; the immediately following prose supplies the intended operation $\mathrm{Addinf}(f)$. The transcription leaves that formula as written, and preserves the separate prime and lower dot on the argument, repeated in the circled margin. Formula (10) applies the construction to the relative integrator $L^{B'/B}_{\bullet}$, while projection to the final category gives an explicit global homology complex whose $n$th term is $\underline Z^{(I_n)}$ (11). Here the parenthesized exponent really describes the direct sum: the projection sends each free abelian representable to $\underline Z$. The final comparison returns to cointegration with coefficients pulled back from $B$ and asserts the internal-Hom isomorphisms (12), (12') and (13), p.~377. These are presented as a localization of ordinary derived adjunction, with an arbitrary bounded-below coefficient complex $K^{\bullet}$ and, in (13), an arbitrary bounded-above $F'_{\bullet}$. An additional hypothesis is required for this internalization. A concrete counterexample already occurs in degree zero: let $B=[0\to1]$, let $B'$ be the final category, and let $f$ select $1$. An abelian presheaf on $B$ is a diagram $K_1\to K_0$. Then $f^*$ is evaluation at $1$, $f_!^{ab}(V)=(V\xrightarrow{1}V)$, and $f_*(V)=(V\to0)$; all three functors are exact. Taking the integrator $\underline Z$ on $B'$ gives $L^{B'/B}=f_!^{ab}\underline Z=(\underline Z\xrightarrow{1}\underline Z)$, the tensor unit and the free abelian representable associated with $1$. Internal Hom from that unit is the identity functor. For $K=(0\to\underline Z)$, therefore, the left side of (12) or (12') is $Rf_*f^*K=0$, while the claimed right side is $R\underline{\mathrm{Hom}}(L^{B'/B},K)=K\ne0$. The ordinary external adjunction remains valid: $\mathrm{Hom}(L^{B'/B},K)=K_1=0$. Equivalently, the projection-formula map needed here cannot be an isomorphism, since $f_!^{ab}(\underline Z\otimes f^*K)=0$ whereas $f_!^{ab}\underline Z\otimes K=K$. Thus existence of the three adjoints and the projective-resolution construction does not by itself justify the claimed internal formula; a suitable local or projection-formula compatibility must be added. All three source formulas remain unchanged, including the visibly missing outer closing parenthesis on the left of (13). The last remark calls this the main relation between homology and cohomology, or abelian integration and cointegration, and proposes the same formalism for ringed topoi possessing $f_!$ on sheaves of sets. Its cancellation of ``commutative'' explicitly permits more general sheaves of rings, while ``striking'' replaces ``interesting'' in the closing assessment of possible applications. That proposed generalization inherits the preceding qualification about internal adjunction; the page does not state the missing local condition. Section~102 next turns from these operations to the pairing between the mutually dual indexing categories and their abelian presheaves.

\subsection*{\S102. Abelianizators, integrators and the dual abelian pairing (pp.\ 378--381).}
For a small category $A$ and $B=A^o$, the section returns explicitly to \S93 and rebuilds the relation between integrators and abelianization through the dual pair of presheaf categories. Comparing (2) with \S99, (3), p.\ 362, gives the finite additive-envelope equivalence $\mathrm{Add}(A^o)\simeq\mathrm{Add}(A)^o$. The proposed infinite analogue (3') is entirely canceled, together with the words introducing it; it is therefore not an assertion that $\mathrm{Addinf}$ commutes with passage to the opposite category. The infinite envelope still enters separately in (4), because quasi-special integrators exist even when special ones need not. The revised pairing (*) between $A^\wedge_{ab}$ and $\mathrm{Addinf}(B)$ is biadditive and exact in its first variable; (5) is correspondingly revised from additive to exact functors. Extension to (6), with second argument in all of $B^\wedge_{ab}$, preserves small inductive limits in both variables (p.\ 378). A complex $L_\bullet$ in $\mathrm{Addinf}(B)$ consequently gives the exact complex-valued functor (7). An abelianizator for $A$ is then a pair $(L_\bullet,\lambda)$ supplying the comparison in (8) between the abelianized presheaf construction and absolute abelianization of homotopy types. The latter is interpreted through the integrators of \S100 on an arbitrary modelizing category, avoiding a privileged auxiliary simplicial category; the earlier constructions of \S\S92--93 are thus placed in a more general functorial setting (p.\ 379). The aim is simultaneous chain-level computation for varying $X$, with a formula extending from $\mathrm{Wh}_A(X)$ to arbitrary abelian presheaves. A quasi-special integrator on $B$, namely a resolution $L_\bullet\to\underline Z_B$ in (10), supplies an abelianizator for $A$ and its canonical comparison. The source explicitly says \emph{for A} here, where the witness substitutes B. The converse is proposed as an equivalence, perhaps even an isomorphism, but is not pursued. Enlarging from special and quasi-special resolutions to arbitrary projective resolutions motivates extending the definition of abelianizator through the whole pairing (6), while existence and uniqueness up to chain homotopism are inherited from the resolution construction (p.\ 380). Formula (11) identifies $B^\wedge_{ab}$ with the cocontinuous abelian-group-valued functors on $A^\wedge_{ab}$. Projective objects give exact functors; the converse is explicitly left to be checked. The image of $\underline Z_B$ is the inductive-limit functor in (13), equivalently the integration functor $p^{ab}_{A!}$ for $p_A:A\to\Delta_0$ in (14), so (15) describes standard abelianizators by projective resolutions of that functor. This makes precise an approach that the cellular and test-category questions of \S91 had left uncertain (p.\ 381). The tentative converse on p.\ 381 must be read together with its explicit rejection in \S106, p.\ 410: already for the final category, exactness of tensor product detects flatness, or torsion-freeness over $\underline Z$, rather than projectivity or freeness. Section~103 revisits the proof strategy and changes the terminology from abelianizator toward integrator and cointegrator.

\subsection*{\S103. Cointegrators, integrators and a change of approach (pp.\ 382--383).}
Returning on 11 August after three weeks away from the notes, the author reconsiders the proposed proof that an integrator on the opposite category gives an abelianizator. His earlier route used Quillen's comparison of a category with its opposite: the functor $(A_{/X})^o\to A^o=B$ is a cofibration with discrete fibers, so its exact direct image gives the displayed comparison of derived sections. The source retains an upper star on the first coefficient category in that display, unlike the arXiv witness, and has a handwritten margin VII beside the still-typed reference to part VI. This cohomological route is then rejected as an awkward way to establish a result about homology. The proposed natural proof is postponed until an overall review makes the relations among homology, cohomology and abelianization conceptually clear (pp.\ 382--383). The terminology is revised at the same time, explicitly in light of \S69: a projective resolution of the constant abelian presheaf $\underline Z_A$ computes $R\Gamma(A,-)$, hence cointegration, and should be called a \emph{cointegrator for A}. For the opposite category $B=A^o$, the same complex computes homology and is an \emph{integrator for B}. Such complexes also give the abelianizators under discussion in \S\S93 and 102. A full converse remains only a possibility, but the existence and uniqueness up to homotopism of the resolution construction make other putative abelianizators seem dispensable. The author therefore proposes dropping the term abelianizator in favor of integrator, paired with cointegrator. The paragraph literally calls the object an abelianizator for B while still writing $X\in A^\wedge$ and $\mathrm{Wh}_A(X)$ in the associated computation; this surviving mismatch of base indices is preserved in the text. The new terminology and the decision to rebuild the formalism lead directly to \S104 (p.\ 383).

\subsection*{\S104. Pseudotopoi, abelianization\newline and internal and external operations (pp.\ 384--393).}
The conceptual review promised in \S103 begins with the full subcategories $\underline{\mathrm{Hom}}_!$ and $\underline{\mathrm{Hom}}^!$ of ordinary functor categories, preserving small inductive and projective limits respectively. Proposition~1 assumes a category $\mathcal A$ with small inductive limits and a small full subcategory generating with respect to monomorphisms; citing SGA~4 I, 8.12.7, it identifies $\mathcal A$ with the projective-limit-preserving functors from $\mathcal A^o$ to sets and deduces small projective limits. The defining test literally refers to objects $Z$ of $C$ after naming the generating subcategory $A$; this index mismatch remains in the source (p.\ 384). Such a category is called a pseudo-topos. Corollary~2 gives the corresponding adjoint criterion and the duality between the two Hom categories. The handwritten revision of (4) places an outer opposite on the first Hom category and cancels the one on the second; (6) explicitly expresses the ordinary reversal of natural transformations under passage to opposite categories. This convention matters for the later formula (37), whose missing opposite is discussed in \S105. Both sides of (4) are identified with the separately projective-limit-preserving bifunctors (5). For actual topoi these are valued sheaves, and for presenting sites they become bisheaves; their category is described as the $2$-product of the two topoi (pp.\ 385--386). Part B introduces $\mathcal A_{ab}$ and its forgetful functor, with left adjoint $\mathrm{Wh}_{\mathcal A}$ when it exists. On p.\ 386 the source literally says that this forgetful functor commutes with ``small direct limits,'' although the displayed adjunction argument and the preceding corollary require the projective-limit statement; no mathematical correction is made in the text. For a topos, the abelian category of abelian group objects has exact small filtered inductive limits, a small generating subcategory, enough injectives and a small cogenerating subcategory, so it is both a pseudotopos and a pseudocotopos. More generally, the pseudocotopos condition on $\mathcal A_{ab}$ ensures existence of abelianization. Proposition~2 and its corollary give the universal property (11)--(12): precomposition with $\mathrm{Wh}_{\mathcal A}$ identifies cocontinuous functors from $\mathcal A_{ab}$ to an additive pseudocotopos with those from $\mathcal A$. The dual formula (13), followed by the choice of abelian groups as coefficients, yields (14), interpreting abelian group objects as abelian-group-valued sheaves. The witness drops the essential co in two pseudocotopos hypotheses on p.\ 388 (pp.\ 386--388). Part C defines tensor product by representability of bilinear maps and recalls its construction from abelianization and cokernels; those ingredients can suffice without the full pseudotopos assumption. The canonical comparison (17), associativity, commutativity and the unit $\underline Z_{\mathcal A}=\mathrm{Wh}_{\mathcal A}(e)$ follow. The explanation of the final object on p.\ 389 again literally invokes ``small direct limits,'' where the preceding theory supplies projective limits. The notation (19) distinguishes $\mathcal A^{\mathcal M}=\underline{\mathrm{Hom}}_!(\mathcal A,\mathcal M)$, called cosheaves, from $\mathcal A_{\mathcal N}=\underline{\mathrm{Hom}}^!(\mathcal A^o,\mathcal N)$, called sheaves. The handwritten condition that $\mathrm{Wh}_{\mathcal A}$ exist is retained; the witness omits it at this point. Under $\mathcal N=\mathcal M^o$, (20) identifies the two categories only after taking an opposite. A later change from $A^\wedge$ to the opposite presheaf topos $B^\wedge$, $B=A^o$, is promised to remove this variance reversal while keeping the coefficient category fixed. The sentence describing the base of the sheaves on p.\ 390 actually names $\mathcal N$ as the topos instead of $\mathcal A$; the scan settles that literal letter (pp.\ 389--390). To define the external action on cosheaves, tensoring with a fixed abelian sheaf must preserve small inductive limits. Formula (22) then sets $(F\circledast L)(G)=L(F\otimes G)$, and (23) abbreviates evaluation as $H*L=L(H)$; the circled star and plain star have distinct roles. The required preservation is reformulated through representability of the presheaf-valued internal Hom (24) and the Cartan relation (25), leading to the stated assumptions that tensor products, corresponding internal Hom objects and abelianization exist. The source writes F,H after (24), where the witness has F,G; it also cites (15) for abelianization, whose definition is (10), and briefly says ``sheaves'' and ``presheaves'' where the surrounding discussion distinguishes sheaves and cosheaves. Under these hypotheses the action (26) has associativity (27) and unit (28) (pp.\ 390--392). Duality gives the exterior Hom action (29)--(30), characterized by $\underline{\mathrm{Hom}}(F,K)(G)=K(G\otimes F)$, while the un-underlined Hom of (31) again denotes evaluation. The source uses K in (30) and (30'), whereas the witness substitutes L in the inner Hom. With $\mathcal N=(Ab)$, these operations recover ordinary Hom and internal Hom of abelian sheaves. An analogous interpretation of the cosheaf tensor operation is not established for a general pseudotopos; the last line of p.\ 392 asks for a choice of $\mathcal N$ despite discussing $\mathcal M$-valued cosheaves, and p.\ 393 returns to $\mathcal M$ without explicitly reassigning the symbol. In the presheaf case $\mathcal A=A^\wedge$, the promised duality identifies $\mathcal A^{\mathcal M}$ with sheaves on $B^\wedge$, and with $B^\wedge_{ab}$ for $\mathcal M=(Ab)$. For sets, cocontinuous endofunctors of abelian groups are tensoring with an abelian group, so both exterior operations recover the familiar tensor product. Compatibility when $A=A^o$ literally, as for sums of commutative one-object categories, is left unchecked. The closing distinction between object-valued and sheaf-valued Hom, and between object-valued star and cosheaf-valued circled star, prepares the systematic presheaf duality of \S105 and the computational formulaire of \S106 (pp.\ 392--393).

\subsection*{\S105. Duality of presheaf topoi\newline and tentative additive generalizations (pp.\ 394--406).}
Dated 12 August with a handwritten addition of 13 August, subsection D explicitly limits its heading to topoi of the type $A^\wedge$ and tentative generalizations. Proposition~3 states the free-cocompletion property of the Yoneda inclusion $\varepsilon_A:A\to A^\wedge$: restriction gives (33), an equivalence between cocontinuous functors $A^\wedge\to\mathcal M$ and arbitrary functors $A\to\mathcal M$, when $\mathcal M$ has small inductive limits. The quasi-inverse sends $i$ to $i_!$, computed in (35) by the inductive limit of $i(a)$ over $A_{/F}$, and (36) gives its right adjoint $i^*(x)=(a\mapsto\mathrm{Hom}(i(a),x))$. The following remark literally assumes a pseudocotopos, where the witness says pseudotopos. Formula (37), however, writes $i\mapsto i^*$ as a covariant functor between the ordinary Hom categories without an opposite. Natural transformations reverse under this assignment, and the explicit opposite in (6), p.\ 385, rules out treating the missing variance mark as an unstated convention; the formula is preserved literally rather than mathematically amended (p.\ 394). The universal property adjoins arbitrary inductive limits while disregarding any that may already exist in $A$. Its projective-limit dual is (38), whose restriction map visibly includes $F\circ\varepsilon_A^o$, with the F and composition sign omitted by the witness. Thus $\mathcal M$-valued cosheaves on $A^\wedge$ identify with functors $A\to\mathcal M$, while $\mathcal N$-valued sheaves identify with functors $A^o\to\mathcal N$, equivalently presheaves on the chaotic site. If $\mathcal M$ has both kinds of small limits, combining these descriptions for $B=A^o$ gives (39), identifying cosheaves on $A^\wedge$ with sheaves on $B^\wedge$ with the same values. The final word canonically on p.\ 395 is canceled; diagram (40) then explains the comparison through $\mathrm{Hom}(A,\mathcal M)=\mathrm{Hom}(B^o,\mathcal M)$ and a chosen quasi-inverse supplied by projective limits. Its dual, with $\mathcal N=\mathcal M^o$, is canonically quasi-inverse to the comparison for B, using the ordinary opposite-category identity (6). The source leaves the two inner Hom arguments at the end of p.\ 396 blank, whereas the witness supplies them (pp.\ 395--396). Case~1 takes sets as coefficients. Formula (41) and its symmetric counterpart identify each presheaf topos with the set-valued cosheaves on the other, through the separately cocontinuous pairing (42), denoted $F*G$ in (43). Two applications of Proposition~3 give (44), so the pairing is determined on representables by $a*b^o=\mathrm{Hom}_A(b,a)=\mathrm{Hom}_B(a^o,b^o)$ in (45). Its symmetry compares the pairing for $(A,B)$ with the pairing for $(B,A)$; it is not an undeclared symmetry inside one fixed pair of categories. Formula (47) gives the evaluation identities $a*G\simeq G(a)$ and $F*b^o\simeq F(b)$ (pp.\ 396--398). The remarks postpone further development of this set-valued duality but sketch the opposite-variance bi-sheaf (48), associated with mutually adjoint contravariant functors. The pairings are proposed as the cross-Hom sets in a category $\mathcal E(A)$ combining $A^\wedge$ and $(B^\wedge)^o$ along A, with $\mathcal E(A)^o\simeq\mathcal E(B)$. Their compatibility is to be expressed through composition and associativity; full details are not supplied here. The later discussion on pp.\ 404--405 returns to this construction, as the marginal pointers anticipate (p.\ 398). Case~2, opened at the foot of p.\ 398, takes abelian-group coefficients and uses the abelianization results (12) and (14) of \S104. For $B=A^o$, (49) identifies $B^\wedge_{ab}$ with the category of additive cosheaves on $A^\wedge_{ab}$ valued in $(Ab)$. The corresponding pairing (50) commutes with small inductive limits in each argument, and its symmetry supplies the analogous description of $A^\wedge_{ab}$. The equivalence (51) reduces such pairings to their values on $A\times B$. Formula (52) gives the free abelian group on $\mathrm{Hom}(b,a)$ when the arguments are the abelianizations of representables; (53) extends this to $\underline Z^{(F*G)}$ for presheaves $F,G$. The insertion on p.\ 400 explicitly says that \emph{both members}, rather than only the left side, commute with the relevant inductive limits (pp.\ 399--400). Remark~1 abstracts this situation to two additive categories $\mathcal P,\mathcal Q$ with a pairing to $(Ab)$, cocontinuous in each variable, which identifies each category with the additive cosheaves on the other. The presheaf examples are abelian topoi: abelian categories with exact small filtered inductive limits and a small generating subcategory. The visible revisions add the filtering qualification and clarify that this terminology does not mean a category which is simultaneously an ordinary topos and an abelian category. Replacing abelian groups by modules over a commutative ring gives the corresponding $k$-linear examples. For a noncommutative ring, the two sides use right and left modules; the proposed pairing is the abelian-group pairing followed by quotienting the resulting bimodule by the differences $sx-xs$. Grothendieck explicitly says that he has not checked that this construction gives the desired duality. He does verify the ordinary one-object example through tensor product, evaluation at the regular module, and the associated tensor--Hom equivalences. Their comparison is an additive analogue of (39) when the target category has both kinds of small limits (pp.\ 400--402). Remark~2 singles out the existence of a small generating family of ultraprojective objects. For a module category, one ultraprojective generator suffices; ultraprojectivity is decomposed into projectivity and finite presentation. The proposed general structure theorem describes the resulting “algebraic” abelian topoi as categories of additive functors from $P^o$ to $(Ab)$, for a small additive category $P$. One can take $P$ essentially small and canonically equal to the full subcategory of ultraprojective objects. For $A^\wedge_{ab}$, this is the abelian Karoubi envelope of $A$, or alternatively one may use $\mathrm{Add}(A)$ for more explicit computation; the text refers back to \S\S93 and 99 (p.\ 402). The notation $P^{\&}$ and the inclusion $\varepsilon_P:P\to P^{\&}$ lead to the expected equivalence (**) between cocontinuous additive functors out of $P^{\&}$ and additive functors out of $P$. The suggested proof imitates the nonadditive representable-density argument. Its dual gives the projective-limit version and then (***), comparing sheaves and cosheaves on the dual algebraic abelian topoi. These statements are presented as the expected additive theory, with the proof still to be checked (p.\ 403). The next page calls an abelian topos \emph{reflexive} when it belongs to such a dual pair. Algebraicity appears sufficient, and the author wonders whether it is necessary, at least for the stronger equivalence with arbitrary additive target categories and suitable compatibility as the target varies. Remark~3 records the central limitation: for a general dual pair he does not see how to construct the comparison functor from the pairing alone. The same perplexity occurs in the nonabelian case, so the algebraic examples should not be read as an established general duality theorem (p.\ 404). Remark~4 returns to the auxiliary category $\mathcal E(A)$ from p.\ 398. Inside the functors from $B^\wedge$ to sets, it is represented by the union of the cocontinuous functors and the representable functors. Their intersection is the Karoubi envelope of $A$, equivalently the opposite of the ultraprojective objects of $B^\wedge$. The margin note describes the union by commutation with either type of small limits. These constructions are then said to extend to a small additive $P$ (pp.\ 404--405). Case~3 resumes the intended application of (39). With an additive target $\mathcal M$, initially under the extra pseudotopos assumptions on $\mathcal M$ and its opposite, (12) and (13) yield (54), relating $\mathcal M$-valued cosheaves on $A^\wedge_{ab}$ to sheaves on $B^\wedge_{ab}$. The earlier discussion suggests that those extra assumptions should be unnecessary. The required assertion (55) says that precomposition with abelianization gives the expected equivalence for any additive $\mathcal M$ with small inductive limits. Formula (56) reduces it to additive functors on $\mathrm{Add}(A)$, and (57) is the projective-limit dual. The four-arrow diagram (58) shows that the two left comparisons are equivalent conditions because the two right comparisons are already equivalences. Finally, identifying $A^\wedge_{ab}$ with $P^{\&}$ for $P=\mathrm{Add}(A)$ reduces the remaining task to (**). The author prefers to prove that general statement and ends by saying that the suggested proof should be checked during later scratchwork. The scan's handwritten reference in the insertion is (56), and its final wording is “proof given” followed by the correction from “there” to “here” (pp.\ 405--406).

\subsection*{\S106. Dual operations, a formulaire and explicit computation (pp.\ 407--419).}
The new subsection E), dated 14 and 15 August, organizes the preceding duality into a calculus for the operations $*$ and Hom. Formula (59) relates $\mathcal M$-valued cosheaves on $A^\wedge_{ab}$ to sheaves on the opposite presheaf topos. Formula (60) instead reverses the values by passing from $\mathcal M$ to $\mathcal M^o$ while retaining the same base. Combining them gives (61), a duality expressed entirely in terms of sheaves; the author prefers this presentation for the intended homology and cohomology applications. The enlargement of diagram (58), and its dual (62), compare functors on $\mathrm{Add}(A)$, its completion under small sums, the abelianized topoi, and their opposites. The text observes that the ten displayed occurrences represent nine distinct categories, using the canonical identification $\mathrm{Add}(B)^o\simeq\mathrm{Add}(B^o)=\mathrm{Add}(A)$. The dual diagram and its introducing sentence consistently use the handwritten coefficient $\mathcal N$, an important distinction from the earlier transcription's $\mathcal M$. Formulas (63)--(64) return to the abelian pairing and its interpretation as a cosheaf (pp.\ 407--408). Tensor structures on the two abelian topoi provide additional operations. Formula (65) is the internal Hom action of an abelian sheaf on an $\mathcal M$-valued sheaf. Through the duality, (66)--(67) express the corresponding covariant operation, denoted by a circled star. In the source this sign is repeatedly a typed star surrounded by a hand-drawn circle. The final-category example reduces these constructions to familiar Hom and tensor-product operations. Their commutation with small limits is not automatic: the text distinguishes the finitely presented projective situation from the weaker flatness detected by exactness of tensor product. The proposed converse from exactness to projectivity is explicitly rejected using flat versus free abelian groups (pp.\ 409--410). The next page retracts the initial feeling that the Hom side was already computationally understood. The data (68) are just the functors $F:A^o\to\mathcal M$, $L:A^o\to(Ab)$ and $L':B^o\to(Ab)$. Passing to the opposite coefficient category and interchanging $A,B$ produces (68') and the tautological identifications (69). Thus an adequate computation and formulaire for either operation should yield one for the other by duality; this is a program for the following pages, rather than a claim that the computation was already available (p.\ 411). The author first broadens the Hom calculation to a topos $\mathcal A$, and provisionally to suitable pseudotopoi. The coefficient category is additive and stable under small projective limits; the additional pseudotopos hypothesis on $\mathcal M$ is used when $\mathcal A$ is not a presheaf category. External Hom, with values in $\mathcal M$, is distinguished from internal Hom, with values in $\mathcal A_{\mathcal M}$, in (70). A sheaf $F$ is viewed as the projective-limit-compatible functor (71). The abelianization $X\mapsto\underline Z^{(X)}$, its componentwise form for a presheaf topos, the unit $\underline Z_{\mathcal A}=\underline Z^{(e)}$ in (72), and the sections functor (73)--(74) supply the ingredients for the main list (75) (pp.\ 412--413). Three entries of (75), circled on the page, are basic: a) evaluates external Hom on an abelianized representable; e) gives the relevant projective-limit compatibility in the first argument; and b) is the tensor--Hom relation. The remaining entries express their internal-Hom versions, the unit identities and the relation between internal Hom and sections. The text derives them through (76), the compatibility of abelianization with products, associativity of tensor product, and specialization at the final object. The comments explicitly leave aside coherence among the canonical isomorphisms and postpone change of topos to a later section connected with homology and cohomology invariants of maps in $(Cat)$ (pp.\ 413--414). For computation, a site $A$ presenting $\mathcal A$ allows $F$ and $L$ to be treated as the functors (77)--(79). When representables are sheaves, the canonical functor (78) is fully faithful; the text permits a change of site to reach that situation. The additive envelope and the extension $\varepsilon_{ab}:\mathrm{Add}(A)\to\mathcal A_{ab}$ in (80) turn the evaluation formula into the finite-product expression (81). The handwritten insertion and rewritten formula call the extension $\mathrm{Add}(F)^*$, although the subsequent expression (83) writes $\mathrm{Add}(F)$ without the star. A inductive limit presentation of $L$ by objects of $\mathrm{Add}(A)_{/L}$, compared with p.\ 403, and exactness e) then give (83). Its arrow under lim points \emph{left}: the displayed expression for $\mathrm{Hom}(L,F)$ is an projective limit, in contrast to the arXiv witness's inductive-limit arrow. The source also writes $x_i$ in the live tuple describing objects of the indexing category, then $a_i$ in the explanation; this inconsistency is preserved rather than silently repaired (pp.\ 415--416). The remark after (83) separates the construction from the provisional assumptions. The formula makes sense for an additive $\mathcal M$ with small projective limits, without a pseudotopos hypothesis on $\mathcal M$. This suggests an extension of the equivalence from (13), p.\ 388, when $\mathcal A$ is an actual topos. The proposed quasi-inverse still has to be checked on the abelianized representables. That check is supplied in the presheaf case by full faithfulness of (80), which gives a final object in the relevant indexing category. The author then regards the immediate computational program as complete and turns to its dual, while the proposed general extension remains qualified by the stated verification (p.\ 416). Dualizing gives the list (84), now under the assumption $\mathcal A=A^\wedge$ and with small inductive limits in $\mathcal M$. Its external star and internal circled-star formulas parallel (75), but the source retains a plain star in e'). The asymmetry of the a') formulas comes from the possible absence of binary products or sums in $A$. Formula (85) spells out the Hom calculation using the category over $a\times b$; the dual symbol $a\vee b$ initially denotes an object of $(B^\wedge)^o$ and agrees with the sum in $A$ when that sum exists. Formula (86) describes the dual indexing category concretely by triples $(x,u,v)$ with maps $a\to x$ and $b\to x$ (pp.\ 417--418). Remark~1 treats additive $A$, where finite sum and product agree. The antiequivalence (88) gives the comparison (87) between internal Hom and circled-star for $L\in\mathrm{Add}(A)$, compatible with tensor products. A faint margin note proposes extension to $\mathrm{Kar}\,\mathrm{Add}$; it is not a proof of that extension. Remark~2 begins the comparison with two more elementary operations by introducing componentwise tensor product (89); its dual is taken up on the next page (p.\ 418). The closing page gives the dual componentwise operation (90), taking external Hom at each object and requiring small projective limits in the additive coefficient category. Both elementary operations extend to arbitrary topoi, with abelian cosheaves replacing the opposite presheaf category in the dual construction. The author does not develop them here, but points to the role of tensor product and its derived operation for ringed topoi and anticipates a parallel role for (90). The final paragraph explains the reuse of the Hom notation: its coefficient argument lies in different categories in (65) and (90), while in the final-category case the operations agree canonically. The same reasoning applies to tensor product and the circled-star operation. This closes the notation review before the dated transition, 17 August, to ground-ring extension in \S107 (p.\ 419).

\subsection*{\S107. Ground-ring extension, additive envelopes and resolutions (pp.\ 419--425).}
Subsection F), dated 17 August, extends the preceding constructions from $\underline Z$ to a commutative ring $k$. Formula (91) specifies a $k$-additive category by an action of $k$ on the identity functor of an additive category. The compatible functors form the full subcategory (92) of the additive functors in (93). Lower and upper exclamation indices retain their respective direct- and projective-limit meanings; bilinearity means $k$-additivity in each variable. The most intrinsic possible ground ring for a given additive category is the endomorphism ring of its identity functor, and other scalar structures are obtained by restriction along a ring map. Formula (94) then defines $A^\wedge_k$ as the category of presheaves of $k$-modules on $A$ (pp.\ 419--420). A homomorphism $k\to k'$ induces the componentwise scalar-extension functor (95), $F\mapsto F\otimes_k k'$. The text interprets $A^\wedge_{k'}$ as a ground-ring extension of $A^\wedge_k$ through a universal problem for cocontinuous $k$-linear functors into $k'$-additive targets. The $k$-abelianization (96), $X\mapsto k^{(X)}$, is the composition of ordinary abelianization with extension from $\underline Z$ to $k$. Precomposition gives (97), identifying cocontinuous $k$-linear functors on $A^\wedge_k$ with the corresponding functors on $A^\wedge$, and hence with arbitrary functors from $A$. Formula (98) is the projective-limit dual. Replacing $A$ by $B=A^o$ and assuming both kinds of small limits in the coefficient category produces the duality (99), interpreted through $\mathcal M$-valued sheaves on $B$ (p.\ 421). The two descriptions induce the external star and internal circled-star operations (100), and the external and internal Hom operations (101). Their rules are inherited from (84) and (75), with $k$ subscripts. A typed insertion supplies the limit hypothesis; its coefficient-category position is left blank, while the handwritten replacement specifies inductive limits for the star side and projective limits for the Hom side. Explicit computation requires the $k$-additive envelope $\mathrm{Add}_k(A)$ in (102), consisting of finite sums of the objects $k^{(a)}$, and its extension $\mathrm{Addinf}_k(A)$ in (103), allowing arbitrary small sums. The latter solves the universal problem of mapping $A$ into an infinitely additive $k$-linear category. The source briefly writes $A^\wedge$ without the lower $k$ in the prose description of the first envelope; that omission is preserved in the edition rather than supplied from the surrounding formulas (p.\ 422). Adjoining direct factors gives the projective objects when one starts from $\mathrm{Addinf}_k(A)$, and the ultraprojective objects when one starts from $\mathrm{Add}_k(A)$. Ultraprojectivity here means projectivity together with finite presentation. Formula (104) identifies these categories with the respective Karoubi envelopes; the printed operator for the latter subcategory is $\mathrm{ULProj}$. Thus the five $k$-additive categories are characterized by different combinations of finite or arbitrary sums, images of projectors, and small inductive limits. Adding $A$ and $A^\wedge$ gives the seven-object diagram (105). The handwritten qualification says that the functors among the five middle categories are fully faithful (p.\ 423). Restriction of cocontinuous $\mathcal M$-valued functors gives the seven-term diagram (106). Its two middle horizontal comparisons are drawn as three-stroke equivalence signs without arrowheads; the other five comparisons are arrows. The suffixes ``add'', ``addinf'' and ``kar'' encode the corresponding preservation requirements. For $k=\underline Z$, part of this diagram recovers (58); its projective-limit dual contains (62). When $\mathcal M$ has both kinds of small limits, the last terms coincide, and the author describes the combined system as fourteen mutually equivalent categories, giving alternative interpretations of an $\mathcal M$-valued copresheaf on $A$ (p.\ 424). The closing discussion separates the purposes of the five additive categories. The ambient presheaf category supports general $k$-linear constructions; its projectives support resolutions. The equivalence (*) between the derived category bounded above and the homotopy category of bounded-above projective complexes extends in (107) to complexes in $\mathrm{Addinf}_k(A)$. Direct sums of the $k^{(a)}$ are convenient because their morphisms can be described by matrices with entries in free modules on the sets $\mathrm{Hom}(a,b)$. A cointegrator with coefficients in $k$ is a projective resolution of the constant presheaf $k_A$, viewed up to the stated unique isomorphism in these homotopy categories. Complexes with components in $\mathrm{Addinf}_k(A)$ are called quasi-special; those with components already in $\mathrm{Add}_k(A)$ are special. Formula (108) identifies the homotopy categories formed from $\mathrm{Add}_k(A)$ and its Karoubi envelope $\mathrm{ULProj}(A^\wedge_k)$. The claim that their canonical functor into (107) is fully faithful remains explicitly tentative, subject at least to suitable coherence conditions; it is not presented as an unconditional conclusion (pp.\ 424--425).

\subsection*{\S108. Homology, cohomology\newline and the limits of the projective construction (pp.\ 426--432).}
The entry dated 22 August begins with a personal interruption: the death of the author's granddaughter Ella and the earlier death of Ronnie Brown's son Gabriel. His response to Brown's letter frames mathematical correspondence as a possible human relationship as well as an exchange of ideas. Subsection G), on homology and cohomology in the absolute case, then resumes the review. Its organizing aim is a precise duality between the two theories, requiring coefficients in general abelian categories with the relevant small direct or projective limits. Passing to the opposite coefficient category and to $B=A^o$ should interchange homology and cohomology (p.\ 426). Cohomology supplies the familiar starting point: sheaf cohomology is expressed by Ext from the constant integral sheaf, or from the structure sheaf on a ringed topos. The author fixes a commutative ground ring $k$ for the intended duality and notes, without pursuing it, the possibility of a locally constant sheaf of ground rings. He emphasizes that the total derived object carries more information than its individual Ext groups. Pairings (*) and (**) recall external and internal derived Hom, and a projective resolution of the first argument computes them in the presheaf situation. Taking the first argument to be the structure sheaf gives derived sections in (i); (ii) relates external derived Hom to derived sections of internal derived Hom. These are recalled first in the familiar ringed-topos setting, before the more general coefficient problem is attempted (pp.\ 427--428). The ordinary pairing (109) takes a $k$-module presheaf and an $\mathcal M$-valued presheaf to an object of $\mathcal M$. The proposed derived pairings (110) and (111) have bounded-above first argument and bounded-below second argument; the second pairing takes values in the derived category of $\mathcal M$-valued presheaves. Abelianity of $\mathcal M$ is needed to speak of these derived categories. Formula (112) names the cohomology objects as external and internal Ext. The source's handwritten arguments in its second expression occur as $F_\bullet,L^\bullet$, and this literal order is retained in the edition. The distinction between external and internal refers to the value category, while the two arguments already belong to different presheaf categories (p.\ 428). Because enough injectives in the coefficient category have not been established, (113) attempts to construct both pairings using a projective resolution of the first argument. The resulting double-Hom simple complexes are defined up to homotopy even before the abelian hypothesis. To descend with respect to quasi-isomorphisms in the second argument, Hom from each fixed projective must be exact. For external Hom this follows first for a representable $k^{(a)}$ by evaluation, then for small sums by the hypothesis that small products of epimorphisms are epimorphisms, and finally for their direct factors. Thus the external pairing (110) is obtained under the stated exactness assumption. The analogous internal calculation involves $F(a\times b)$ and immediately exposes a difficulty when $A$ is not stable under binary products (p.\ 429). The next page explicitly retracts the proposed projective computation of internal derived Hom. Projectivity of a sheaf of modules is not generally preserved by localization, in contrast to the injective behavior recalled there. Localizing $k^{(a)}$ over $b$ leads to the object represented by $a\times b$ in the slice presheaf topos; the displayed Ext calculation need not vanish in positive degree. Such vanishing would impose strong acyclicity conditions on the connected components of that product. Consequently (111) has not been constructed by the proposed argument. The possibility of using enough injectives is postponed. External derived Hom is nevertheless sufficient for (114): derived sections are computed by Hom from a projective resolution (115) of $k_A$. Such a resolution is the cointegrator, and (116) is the resulting functor on bounded-below complexes of $\mathcal M$-valued presheaves. The canceled underline beneath Hom in (114) visibly reinforces this distinction between external and internal constructions (p.\ 430). The author acknowledges that the expected comparison (ii) has therefore not been obtained in this generality, then dualizes the construction that survives. The coefficient category $\mathcal N$ is abelian, has small inductive limits, and satisfies the dual exactness condition on small sums of monomorphisms. Formula (117) is the derived external star pairing; (118) computes it by resolving its second argument in the opposite presheaf category. The composite equivalence (119) and the tautological isomorphisms (120), (120') express the interchange with external derived Hom under $\mathcal M=\mathcal N^o$ and $B=A^o$. The two occurrences of $A^\wedge$ in (119) have no coefficient subscript in the scan. Integration is then described, with the inserted qualification ``at least morally'', as total left derivation of inductive limit; (121) and the number corrected by hand to (121') suggest the dual derived-limit notation (p.\ 431). The closing page seeks notation that retains the geometric associations of homology and cohomology while displaying their symmetry. Formula (123) introduces $\mathrm{RH}^{\bullet}$ and $\mathrm{LH}_{\bullet}$, and (124) states their interchange on passing to the dual base and opposite coefficients. The displayed targets of (123) are literally $\mathcal M$ and $\mathcal N$; the edition does not silently supply derived-category symbols. The final remarks delimit the achievement. The symmetry uses fairly general abelian coefficient categories satisfying the specified exactness properties, but $\mathrm{RH}^{\bullet}$ has not yet been shown to be the total right derived functor of sections or projective limit. The notation (121), (121') is consequently called misleading. Establishing enough injectives and comparing the injective and projective definitions should also recover a satisfactory internal derived Hom formalism. This unresolved requirement leads directly into the next day's investigation in \S109 (p.\ 432).

\subsection*{\S109. Projectives, injectives and a return to derived Hom\newline (pp.\ 433--442).}
The entry dated 23 August answers the question left open in \S108 by working first with a small $k$-additive category $P$. Formula (125) defines $P^{\&}_{\mathcal N}$ as the category of $k$-additive functors from $P^o$ to $\mathcal N$; its small limits and abelian structure are computed componentwise when those structures exist in $\mathcal N$. Proposition~4 says that small inductive limits and enough projectives in $\mathcal N$ give enough projectives in $P^{\&}_{\mathcal N}$. If those projectives are also $k$-flat, a projective coefficient presheaf makes the relevant star pairing exact in its other argument. The assumptions are necessary already when $P$ is the additive category generated by the final category. The special case $P=\mathrm{Add}_k(A)$ recovers the presheaf categories needed in the preceding section (pp.\ 433--434). The proof uses the external product $L\otimes x$ of (126), computed componentwise. The three Hom expressions in (127) give its tensor--Hom adjunction; the comparison involving the first expression requires small projective limits in $\mathcal N$, while the direct comparison of the other two does not. It follows that the tensor product of projective arguments is projective. Formula (127') specializes to a representable and gives the evaluation adjunction. Choosing projective covers of the values $F(a)$ then produces a projective epimorphism onto $F$ from the sum of the corresponding representables tensored with those covers. Compatibility (129) proves the supplementary flatness assertion. Its necessity is illustrated by $k=\underline Z$ and coefficients in vector spaces over $\mathbf F_p$: every coefficient object is projective, while only the zero object is flat over $\underline Z$. The dual statement gives enough injectives and the corresponding $k$-coflatness assertion. Several coefficient-category positions are literally blank in the scan and remain blank in the transcription (pp.\ 434--435). Corollary~1 returns to $A$ and lists four conditions on $\mathcal M$: abelianity with small projective limits, enough injectives, $k$-coflatness of its injectives, and preservation of epimorphisms by small products. Under these hypotheses, external derived Hom in (130) can be computed by a projective resolution of its first argument, an injective resolution of its second argument, or both; (131) expresses the compatible constructions. The dual statements (130'), (131') concern derived star and projective coefficient resolutions. This settles the availability of the resolutions required for the external pairing and opens the way to the internal problem (p.\ 436). The internal operations (132), (132') still require care. Projective first arguments alone do not generally suffice, as \S108 already showed. Resolving the second argument by injectives does work if tensoring with a representable is exact in the relevant sense. For $P=\mathrm{Add}_k(A)$, componentwise tensor product makes representables, and their finite sums, $k$-flat. Corollary~2 therefore constructs internal derived Hom (133) from an injective resolution of its second argument. Formula (134) gives the desired comparison: external derived Hom is obtained by applying derived cohomology to internal derived Hom. The circled-star dual (135) uses a projective resolution of its coefficient argument, and (134') compares it with derived homology. The arrow into the target category is absent from the displayed (135) in the source; the edition preserves that omission. These results supply the internal formalism that the earlier projective construction had failed to provide (pp.\ 437--438). Remark~1 broadens the setting again. For a general small additive $P$, an internal tensor structure would come from a diagonal $P\to P\otimes_k P$, together with the appropriate coalgebra structure. An ordinary category supplies such a diagonal before linearization. The author proposes that a coalgebra object in the $2$-category of small $k$-additive categories might give a finer model of homotopy types than the De Rham constructions with divided powers considered earlier, referring back to \S94, p.\ 339. Associativity and commutativity should be specified by coherent isomorphisms, rather than imposed as identities; they are additional data with coherence conditions. The references to Saavedra's thesis and to work of Mac Lane, Bénabou and Sinh indicate hoped-for sources of this formalism. The text presents this as an avenue to investigate and explicitly hopes that reversing arrows will yield the parallel theory (pp.\ 438--439). The next page develops some consequences of that proposal. Extension from $k$ to $k'$ should preserve the objects of $P$ and tensor its Hom modules with $k'$, with the stated $2$-categorical universal property. Units $k_P$ and $k_Q$ correspond to augmentation and coaugmentation data. Derived Hom from the former and derived star with the latter should recover cohomology and homology, with constant coefficients and cup and cap operations interpreted through the tensor structures. Coalgebra objects could also be considered at the derived-category level. The author does not expect this alone to recover a general fundamental group or homotopy type: outside the $1$-connected case he anticipates a groupoid action on the universal-cover data. The following page asks whether a category of homotopy types over $k$, compatible with base change, might be approached this way without first constructing the full proposed models (pp.\ 440--441). Remark~2 sharpens the hypotheses of Corollary~1. The product exactness condition is used for the projective-resolution construction, while enough injectives and their coflatness support the injective construction; abelianity with projective limits supplies the ambient categories. When all four conditions hold, the two constructions agree. Cointegrators remain the author's preferred method for external cohomology, but internal derived Hom uses injective coefficient resolutions. In particular, the first three conditions suffice for the internal construction and the comparison (134); the product condition is not required when one works throughout with injectives. The source calls the injective complex $K^\bullet$ in this discussion, and that literal letter is retained (p.\ 441). Remark~3 closes the section with the general diagonal case. Exactness of $L\mapsto L\otimes a$ follows if each $\mathrm{Hom}_P(b,a)$ is $k$-flat; equivalently, the projectives in $P^{\&}$ must be flat for the external tensor action. In the case $P=\mathrm{Add}_k(A)$, the relevant Hom modules are finite sums of free modules on the sets of morphisms in $A$, hence projective. The author would prefer projectivity in the general theory, but flatness suffices for the asserted injectivity of external and internal Hom when the first argument is projective and the second injective. Finally, for a $k$-additive functor $u:\mathcal M\to\mathcal M'$ commuting with small projective limits, (136) compares $Ru$ applied to derived Hom with derived Hom after the componentwise derived functor $Ru^P$. If $u$ is exact, the derivation of $u$ is unnecessary. The analogous internal and dual formulas are explicitly omitted. This completes \S109 and the cote, before the next day's \S110 returns to schematization (p.\ 442).

\subsection*{110. Schematization through Postnikov fibrations, and its obstruction (pp. 443--445; 134-7/1--3).}
Grothendieck resumes the question of homotopy types over a ground ring by comparing an abelian group $M$ with the functor $W(M)(k)=M\otimes_{\underline Z}k$: the discrete object acquires a structure analogous to a vector bundle, suggesting an analogous construction for semisimplicial models and successive Postnikov fibers. The proposed ``schematization above level n'' would leave a lower truncation discrete and require the comparison from schematic cohomology with quasi-coherent coefficients to ordinary cohomology to be bijective. The reference to section 94 situates the hope of replacing $\underline Z$ by $\underline Q$ when nilpotence replaces an abelian fundamental group. On p.~444 he immediately questions this construction: torsion and the universal coefficients formula prevent the homotopy groups after base change from being obtained by tensoring the individual homotopy groups. The alternative is to reconstruct a homotopy type from a suitably enriched ``abelianization'' or ``linearization''; section 111 then seeks a different construction without explicit Postnikov decomposition. The handwritten replacement of ``(= homotopy)'' by ``(or rather, homotopy)'' on p.~444 matters here: the arXiv witness retains the canceled equality wording. Its mapsto in the displayed base-change formula also differs from the scan's ordinary arrow; the scan alternates star subscripts with the typed notation X. Page 445 closes with a methodological reflection on the tension between detailed preparatory work and the desire to move quickly. The rust-obscured words on p.~443 remain marked rather than restored from the mathematical argument.

\subsection*{111. Unipotent bundles, divided powers, and a revised notion of equivalence (pp. 446--459; 134-7/4--17).}
The starting problem is to construct categories $Hot_n(k)$ of schematized $n$-connected homotopy types, with extension and restriction of the ground ring and a comparison with the ordinary homotopy category over $\underline Z$ (pp.~446--449). Grothendieck builds $U(k)$ inductively from affine bundles $W(M)$ and torsors, first in the fpqc setting. He asks whether the successive fibration data belong to the object or merely witness its existence, then proposes semisimplicial objects subject to connectivity and filtration conditions. The proposed weak equivalences are initially detected by taking sections, and the hoped-for $k$-module structures on homotopy groups lead to a search for a linearization into $D.(Ab_k)$. On p.~448 the model-category symbol has a typed subscript $m$, although the surrounding development uses $n$; this apparent index slip is retained in the text. Pages 450--455 test divided powers as the required linearization. The candidate universal map $x\mapsto\sum x^{(i)}$ immediately forces a completion, and the completed divided-power coalgebra is considered through its group-like elements. After the Eilenberg--Maclane test, the discussion returns to pointed objects and uncompleted divided powers. The comparison (w) on p.~454 places ordinary free abelian linearization, $L$, and the pointed version $L_{pt}$ in one diagram. The marginal correction on p.~452 is consequential: ``One has to take k=Z, and components $X_n$ projective.'' Both its initial words and the subscript are legible on the scan. The subsequent distinction between quasi-isomorphisms and chain homotopisms supplies the projectivity restriction. The final reference in the sentence on p.~455 is to (x''), after a preceding reference to (x'''); the arXiv witness repeats (x''') and loses this distinction. The argument is substantially revised on pp.~456--459. The author retracts the unrestricted derived-functor claim of p.~454, requires projective or flat components, and remarks that the Postnikov approach must again be set aside. The plain typed $I_k(M_*')$ on p.~456 is retained as printed, although the neighboring formulas suggest a notation slip. Page 457 then exposes the problem with ordinary homology as a $k$-module: the exterior-power calculation for $K(M,1)$, applied to $k=\underline Z[T]/(T^2+1)$, yields an obstruction. Page 458 accordingly defines total linearized homology $LH.(X_*)=L_{pt}(X_*)$ and considers weak equivalences detected by its homology invariants. The handwritten revision says it is becoming \emph{dubious} whether equivalences should be defined through $X_*(k)$; the arXiv witness reads \emph{obvious}, reversing the author's change of judgment. At the end of the section, Question 2 asks whether the two definitions coincide for projective pointed unipotent bundles satisfying $X_0=X_1=e$ (p.~459). The weaker condition $X_0=e$ raises an additional issue about abelian fundamental groups. These pages formulate questions and revised expectations, rather than establish an unrestricted equivalence theorem; section 112 resumes from the projective $\underline Z$ case.

\subsection*{112. Kan conditions and the scope of the linearization program (pp. 460--465; 134-7/18--23).}
The entry dated 26.8 takes the projective $\underline Z$ case of the comparison maps (w) as its working foothold. Grothendieck sketches a passage from the Eilenberg--Maclane case to finite chain complexes, then through finite and ``way-out'' infinite successive fibrations, using multiplicativity and hoped-for spectral-sequence or derived-category transitivity statements (pp.~460--461). The proposed Kan--Postnikov complexes use shifted projective resolutions of homotopy groups. To transfer the construction from discrete to schematic objects, he returns explicitly to the cohomology obstruction of p.~444 and hopes to dualize the homology comparison. This is an outline requiring substantial verification, as the author himself stresses. Constant complexes then show why $X_0=e$ is necessary, and the monoid discussion turns the issue into comparison of polynomial with arbitrary cochains. The multiplicative-monoid example on p.~462 is offered as a counterexample to the assumptions previously used; the circled marginal cross is part of the page's evidence and is absent from the arXiv witness. The response is to ``drop Postnikow and keep Kan'' (p.~462). Since $U(k)$ is not closed under arbitrary fiber products, a schematic horn condition must ensure both representability of the relevant finite limits and a suitably strong surjectivity property (p.~463). The tentative name ``submersions'' covers alternatives ranging from trivial bundles to formally smooth maps admitting sections. The first nontrivial condition asks for the three maps $X_2\to X_1\times X_1$ to be submersions; for a monoid nerve this is meant to force a group. Page 464 lists the desired consequences: the linearization theorem over $\underline Z$, a homology criterion for homotopisms of connected Kan complexes over a general ring, homotopy fibers and Postnikov decomposition, and eventually equivalence of $Hot_n(\underline Z)$ with ordinary $Hot_n$. Page 465 distinguishes restriction of rings from extension, noting that restriction does not preserve projective components without an additional assumption and suggesting resolutions to remove that restriction. Its last sentence proposes compatibility of ring extension with the individual homotopy invariants; the next entry immediately revises that proposal. This transition is essential to reading the section as a developing argument rather than a list of established properties.

\subsection*{113. Soft Postnikov decomposition and a schematic fundamental group (pp. 466--472; 134-7/24--30).}
The entry dated 27.8 corrects the previous evening's assertion. Formula (A) concerns total homology and derived extension of scalars; the analogous termwise assertion (B) for homotopy modules fails without flatness hypotheses (p.~466). The revised arrow in (B) points left, a direction separately emphasized by a circled marginal arrow omitted by the arXiv witness. The reason is the use of shifted projective resolutions in a Postnikov fiber: after tensoring, higher homology can appear, so the base-changed successive fibration need not remain a Postnikov decomposition (p.~467). This explicitly answers the puzzle of p.~444 and motivates the judgment that ordinary Postnikov decomposition is not ``geometric,'' meaning compatible with extension of the ground ring. The proposed alternative keeps the entire abelianization complex. Using the completed divided-power linearization and the embedding $X\to W(L(X))$, Grothendieck proposes repeatedly taking the homotopy fiber of $X_*\to W(L(X_*))$, without first extracting the first nonzero homology module (pp.~467--468). He names this ``soft Postnikov d\'evissage.'' For a discrete $K(G,1)$ the analogy leads to iterated commutators, and he distinguishes the versions associated with solvability and nilpotence. Formula (D) contains an index problem: the visible inner exponent is $n+1$, like the outer exponent, while the arXiv witness gives $n-1$. The checked text preserves the scan's reinforced plus and does not silently substitute a mathematically expected recurrence. The next page adds the issue of the action of $\pi_1$ on higher homotopy groups; trivial or unipotent action appears relevant to unipotent schematization (p.~469). The second part of the section revisits complexes with $X_0=e$ but possibly nontrivial $X_1$. A group object gives a fundamental group that need not be abelian, and finite-presentation and smoothness assumptions lead to questions about nilpotence (p.~469). The Kan relations then motivate the presheaf $k'\mapsto\pi_1(X_*(k'))$. The desired invariant is a group object of $U(k)$ representing this sheaf, with an underlying parametrization by a module and a polynomial group law (pp.~470--471). This changes what the notation $\pi_1$ denotes: the group of sections alone is no longer the complete object. Formulas (G) and (G') express evaluation and base-change compatibility for the proposed schematic group. The following duplicated equation on p.~471 is nevertheless literally tautological in the scan, with $\pi_1(X_*(k))$ on both sides; the witness edits one side. The blank between $k$ and $k'$ in the preceding prose likewise contains no arrow on the page. To admit arbitrary quotient modules, the notion of submersion is reconsidered once more: maps isomorphic to those induced by epimorphisms of modules are proposed as the most convincing of the three possibilities considered so far (p.~471). The printed ``unless \ldots{} does'nt give rise'' contains an apparent extra negation, retained despite the witness's change to ``gives rise.'' The stronger requirement that the schematic fundamental group come from addition on a module is distinguished from commutativity; the discussion ends with tests involving the linearization theorem, a fibration over $K(G,1)$, and the relationship between nilpotent discrete groups and group objects of $U(k)$ (p.~472, referring back to section 94). The handwritten addition about unipotent action on the higher $\pi_i$ belongs here and is not merely incidental: it qualifies the contemplated range of homotopy types. Section 114 begins on the same page by describing the accumulated picture as provisional and preparing a review of the program.

\subsection*{114. A provisional model category and the completion gap (pp. 472--480; 134-7/30--38).}
This review arranges the preceding experiments into a prospective construction. The category $U(k)$ consists of unipotent bundles built from modules; smooth objects correspond to projective modules, submersions to epimorphisms, and smooth submersions to maps with projective kernel (pp.~472--473). The completed divided-power coalgebra describes the unpointed linearization, while the pointed version removes the completion in the basic construction (pp.~474--475). Completion is a substantive distinction: p.~474 explicitly says that the map (C) of p.~467 ``does'nt quite exist'' with $W$ as written, and introduces the completed version $W^{\wedge}$ instead. Base change for the completed functor consequently uses completed tensor products. The subsequent model categories impose $X_i=e$ for $i\leq n$, take sections to pointed semisimplicial sets, and localize at the proposed weak equivalences. The hoped-for comparison with ordinary homotopy types is restricted to the appropriate connected range over $\underline Z$ (p.~476). On p.~477 the first expression for the category of models visibly uses $\pi_n(k)$, although the surrounding formulas use the handwritten model-category symbol $\mathcal M_n(k)$; this notation anomaly is retained rather than silently regularized. Smooth Kan resolutions are then proposed to derive extension of scalars and the total linearization $LH.(X_*)=L_{pt}(X_*)$ (p.~477). The next two pages ask whether this homology functor is conservative and whether pointed linearization and completed linearization are quasi-isomorphic. They also reconsider comparison with ordinary free-module homology. The second question undergoes an immediate correction: on p.~479 the author says the proposed comparison with $k^{(X_*(k))}$ is generally false away from $k=\underline Z$. In the first nonvanishing degree, the expected $k$-module invariant is identified with integral homology of the section complex; tensoring that homology with $k$ produces a different candidate. The word \emph{integral} is explicit in the scan, whereas the arXiv witness gives \emph{internal}. The correction renews the need to construct homotopy modules intrinsically and to obtain them through a Postnikov argument, rather than transfer all ordinary homology assertions indiscriminately. The final page makes the outstanding construction precise. The exponential map gives $X_*\to W^{\wedge}(L(X_*))$, but the required homotopy-class map is $X_*\to W(L_{pt}(X_*))$ (p.~480). Even a quasi-isomorphism between the two linearizations does not by itself remove the difference between $W^{\wedge}$ and $W$; the handwritten phrase ``flaw in this `argument''' marks exactly this gap. The target of the completed map is a semisimplicial sheaf and need not belong to the proposed model category. Hard or soft Postnikov decomposition therefore still requires this factorization, a suitable notion of homotopy fiber, and homotopy constructions internal to the model category. The failure of $U(k)$ to have fiber products explains why Quillen's standard machinery cannot simply be applied as stated. These are explicit unresolved tasks, and section 115 returns the following morning to the meaning of the coefficient sheaf itself.

\subsection*{115. Quasicoherent coefficients and prorepresentable linearization (pp. 481--489; 134-7/39--47).}
The entry dated 28.8 restarts the formula numbering and identifies the source of the preceding confusion. The appropriate free linearization is over the quasicoherent ring sheaf $\underline O_k$ on the fpqc site, not over a constant ring (pp.~481--482). Grothendieck explains that the obstruction groups in schematic Postnikov decomposition use the quasicoherent sheaf associated with a homotopy module, rather than the same module regarded as a constant coefficient group. The circled marginal clarification of the argument $X(n)_*,\pi_{n+1}$ belongs to this distinction; the overwritten version in the running formula remains marked as doubtful. He connects the earlier mistake to having translated geometric language into commutative algebra too thoroughly, losing sight of which coefficient objects were intended. The universal free Module $\underline O_k^{(X)}$ gives the Hom, Ext and derived-Hom expressions (4)--(5') (p.~483), but is generally not quasicoherent. The proposed replacement would represent the same functor on quasicoherent coefficients. For $X=W(M)$, the continuous-Hom calculation through truncated divided-power modules shows that the correct object is a proobject, rather than the completed module regarded as an ordinary module (pp.~484--485). The obstruction is again scalar extension: projective limits do not generally commute with tensoring by a new ground ring. Formula (13) packages the finite products of divided powers into a strict projective system; ordinary representability would require it to become essentially constant. The added phrase ``but with slightly different meaning'' therefore changes the interpretation of $L$ in (14), now valued in pro-modules. The pointed functor still takes values in ordinary modules, embedded as constant proobjects. Formula (18), on p.~486, supplies the exponential map into $WL(X)$ in the procategory and distinguishes this abelianization from its richer linear structure. The proposed derived comparison (20) and its Ext form are then tested on $X=W(M)$ (pp.~486--487). Formula (21) contains a meaningful handwritten revision from a limit expression to a direct sum of Ext groups of divided powers. Projectivity of $M$ gives the expected vanishing in positive degree, with a passage to finite-type submodules sketched for the infinite case. For nonprojective $M$, the potential nonzero Ext$^1$ term challenges the author's previous expectation that quasicoherent cohomology of these bundles should vanish. His reference to Larry Breen is a question about this expectation, not a proof of the general formula. The text visibly says ``quasiunipotent'' where the witness shortens it to ``unipotent'' (p.~487). The next page restricts attention to the smooth case and postpones the projective-resolution argument for full generality. For a semisimplicial complex, (22) uses Hyperext rather than a degreewise Ext construction; (23) compares the resulting pro-linearization with pointed linearization (p.~488). The comparison map (24) revisits the map from p.~454 with its proobject interpretation made explicit. Its relationship to a quasi-isomorphism, and then to the componentwise projective limit used on p.~478, still requires foundational work. Page 489 plainly admits that no argument has yet been found for the desired assertion under the proposed smoothness, connectivity and possibly Kan hypotheses. The second comparison is reformulated as $\underline O_k^{(X_*)}\to L(X_*)$ in (25). The author does not confidently assert that this is a quasi-isomorphism: quasicoherence of the first homology sheaves and essential constancy of the second remain in doubt. The practical hope is instead that it gives the needed comparisons against quasicoherent coefficients through (20). Section 116 subsequently reports the collapse of the earlier general abelianization theorem, so these provisional claims must be read with that later revision in view.

\subsection*{116. The failed abelianization theorem and the interruption of the notes (pp. 490--491; 134-7/48--49).}
The entry dated 7.9 reports the decisive revision of the preceding program. After finding a simple description of schematic homotopy groups, Grothendieck learned in a telephone conversation with Illusie that the asserted comparison between discrete and schematic homology or cohomology was false. The anticipated equivalence between schematic and ordinary homotopy types therefore fails (p.~490). This is an explicit retraction of a central assumption underlying sections 111--115; it does not cause him to abandon schematic homotopy types, whose formal properties he now wants to study in their own right, possibly through different models. The remainder of the section explains the intervening pause in writing through fatigue and the need to attend to the death of his granddaughter Ella. Page 491 distinguishes work allowed to take place through attention from a compulsive forward movement that leaves perceptions disconnected from one's picture of reality. The revisions matter to the voice of this passage: the hand moves ``who'' after ``finally,'' retains ``his language'' for the body's intervention, and changes ``cause'' to ``caused.'' The witness substitutes ``which'' and ``its'' and loses the last correction. These pages provide both the mathematical warning needed to read the preceding conjectures and the author's own account of why the notebook paused before the technical review resumes in section 117.

\subsection*{117. Size, exactness, and the distinction between schemes and formal schemes (pp. 492--496; 134-7/50--54).}
The entry dated 10.9 begins with two technical qualifications. The affine site used earlier is not automatically small in the working universe; a possible remedy is to bound the sizes of the modules and the rings used as arguments, using the fppf site for finite-type bundles (p.~492). More consequentially, $M\mapsto W(M)$ is right exact but need not preserve kernels. Thus a free resolution of $M$ need not become a resolution of $W(M)$. The derived-Hom comparison used to interpret formula (21) of p.~487 is therefore questioned in general and retained for the flat case (p.~493). This is a specific dependency of the cohomology calculation, not merely a set-theoretic caveat. The finite projective case clarifies the geometry through duality. Polynomial maps can be described by maps of symmetric algebras or by continuous maps of the dual completed coalgebras; preserving the origin translates into augmentation compatibility (pp.~493--494). The uncompleted divided-power coalgebra is instead dual to the adically completed symmetric algebra. Consequently, the two algebraic descriptions lead to schemes and formal completions along their zero sections, respectively. The scan's formulas require care: the prose on p.~493 includes a dual check on $M$, but formulas (2), (3) and (4) omit dual signs that the arXiv witness supplies. Those omissions remain in the checked text. The upward-arrow-shaped completion mark in (3') on p.~494 is likewise recorded as drawn, and a separate unreadable pen mark above the star in (5) remains marked. This geometric interpretation overturns the earlier hoped-for equivalence between the pointed and nonpointed linearizations (pp.~495--496). For the additive group in characteristic $p$, the degree-one comparison distinguishes additive polynomials from formal series $\sum_{i\geq0}c_i t^{p^i}$: the schematic side allows only finitely many nonzero coefficients. Grothendieck acknowledges that this example is neither simply connected nor over $\underline Z$, but states that he also no longer expects the comparison to be a quasi-isomorphism for the basic connected examples previously contemplated. He reports that Illusie rejected the vanishing assumptions behind his attempted ``way-out'' argument. The closing underline therefore warns explicitly against substituting $L_{pt}(W(M))$ for $L(W(M))$ when computing schematic homology or cohomology. Formal invariants have their own meaning and are distinct. Section 118 begins on the same page with a different and deliberately striking claim: homotopy invariants themselves should nevertheless be recoverable from the first-order tangent object.

\subsection*{118. The tangent complex as a proposed homotopy invariant (pp. 496--500; 134-7/54--58).}
Against the distinction just established between schematic and formal cohomology, Grothendieck proposes that homotopy groups should depend only on the first-order tangent object. The Lie functor sends a smooth pointed bundle $W(M)$ to $M$, and formula (9) identifies the homotopy groups of $X_*(k)$ with those of the semisimplicial module $\mathrm{Lie}(X_*)$, hence with homology of its associated chain complex (p.~496). The next page sharply qualifies the claim: he has no direct description of a natural isomorphism under the stated hypotheses. It is tautological for complexes arising from chain complexes and is expected to extend through Postnikov decomposition, if every relevant model can be replaced by one admitting such a decomposition; naturality would still require work (p.~497). The typed revision makes this formula a ``cornerstone'' of a theory of schematic homotopy types and a proposed intrinsic definition of the invariants, rather than subordinating them from the outset to discrete homotopy groups. Weak equivalences would then be detected by the tangent complexes. Both subsequent comparison displays map an $X_*$ object to a $Y_*$ object; the arXiv witness repeats $X_*$ as target in each. The contemplated application is to explicit models of spheres (pp.~498--499). To make the Lie formula useful for ordinary homotopy groups, one must first realize the discrete type schematically over $\underline Z$ and obtain a workable model. Successive suspension from $S^1$ suggests the need for suspension constructions, complementing the earlier need for fibers and loops; the proposed derivator would organize the finite integration and cointegration operations. These are motivations and difficult construction problems, not completed calculations of sphere homotopy groups. The discussion of $S^2$ as a homogeneous space under $S^3$, and of the corresponding group/flag-variety geometry, asks for a model of the entire relationship: a group, a subgroup playing the role of a torus or Borel subgroup, and a quotient. Several overprinted group-name letters remain uncertain in the text, while the star on $\underline C^*$ is visibly an asterisk rather than the witness's multiplication cross. Hemispherical complexes supply a second route (pp.~499--500). Their smaller list of boundary and degeneracy operations, together with a prospective direct computation of homotopy in the discrete setting, is compared with the tangent-complex method. The marginal exclamation and vertical rule emphasize this alternative. The author immediately admits that he has not checked the proposed computation for arbitrary hemispherical sets that do not arise from an infinity-groupoid or infinity-stack. Comparing the two descriptions is therefore offered as a possible route toward understanding formula (9), with that unresolved premise preserved. Section 119 then resumes on 26.9 after another interruption and returns to the failure of the earlier homology comparison.

\subsection*{119. Integer-valued polynomials and a different comparison problem (pp. 500--505; 134-7/58--63).}
The entry dated 26.9 returns to the failed comparison of discrete and schematic linearization. Grothendieck reports Illusie's objection from Breen's work on stable cohomology, while admitting that the precise counterexample has not yet been made explicit to him (p.~500). He then gives a concrete degree-two obstruction using central extensions of the additive group (p.~501). The multiplicative group of truncated series $1+at+bt^2$ defines the extension $W_2$; a polynomial splitting requires $P(a+a')=P(a)+P(a')+aa'$. The solutions displayed over $\underline Q$ involve $t(t-1)/2$, which has no integral-coefficient representative after addition of the displayed linear term. This separates polynomial cochains over $\underline Z$ from arbitrary cochains on the discrete group. The final splitting display literally has $t$, without the square appearing in the preceding general homomorphism display; the checked text preserves that visible slip rather than adopting the arXiv witness's $t^2$. The main argument should therefore be read through the preceding coefficient equation, not by treating the last display as an independently correct formula. The next proposal replaces polynomial algebras by ``binomial algebras,'' built from the integer-valued expressions $P_n(t)=t(t-1)\cdots(t-n+1)/n!$ (p.~502). Grothendieck tentatively requires the ground ring to carry the corresponding binomial operations and compares this enrichment with divided powers in crystalline cohomology. The marginal reference to pp.~506--507 anticipates the next section's reconsideration of these structures and of the restrictions on the ground ring. On p.~503 the finite-type qualification on the free module $M$ is an above-line typed insertion. The algebra $\mathrm{Symbin}_{\underline Z}(M)$ is defined through rational polynomial functions integral-valued on the dual lattice, and the map to all functions on a semisimplicial free abelian group is posed as a possible quasi-isomorphism in the connected range. Reduction to $K(\underline Z,n)$ is proposed by d\'evissage, not proved. The broader question asks whether Eilenberg--Maclane homology or cohomology is derived from a suitable nonadditive functor $B$, leading to ``B-schematic homotopy types'' and unipotent $B$-bundles. The dark superscript on the first Symbin display remains unresolved and is not silently discarded. The author declines to pursue that alternative immediately, because the ordinary schematic program still promises a simple Lie-functor description of homotopy invariants (pp.~503--504). Loops, suspensions, and compatibility with passage to discrete types remain the practical tests. Reading Bousfield--Kan then changes the comparison: their construction is identified as $k$-completion, distinct from a fully developed category of types over $k$. The marginal ``definitely not, in general!'' rejects the tentative identification of completion with extension of the ground ring; the faint underline beneath ``ring extension'' identifies the disputed phrase (p.~504). Page~505 links this reading to the integration and cointegration operations of section~69 and to Don Anderson's existence theorem in closed model categories. A pen transposition moves the closed-model-category qualification directly after ``operations,'' making the setting of that theorem clearer. The section closes by recognizing Bousfield--Kan as foundational reading while postponing it long enough to record the recent schematic reflections. Section~120 starts on 28.9 by abstracting the extra algebraic operations just proposed.

\subsection*{120. Extra algebraic operations and polyhedral models (pp. 506--507; 134-7/64--65).}
The entry dated 28.9 abstracts the previous section's binomial proposal together with divided powers, lambda operations, and Adams-type operations. Grothendieck uses the notion of an ``analyzer'' and asks for an extension of the analyzer of commutative rings in which unary operations determine higher components by tensor powers. The handwritten changes from $n$ to $n+1$ tensor factors and from the older bounds to $n\geq-1$ and $n\geq1$ affect the indexing of the proposal (p.~506). The live corrected bounds are clear, although overwritten earlier digits remain doubtful. The essential datum is described as a ground ring $k_0$ and an algebra $\Omega_0$ equipped with composition. An $\Omega$-structure on a ring is then an action of these operations satisfying compatibility with addition, multiplication, and constants. Page~507 displays the two diagonals that encode the first two conditions. This is an outline of the desired axioms; the author expressly declines to write all of them down. The circled instructions ``\`a la ligne'' and ``id'' on p.~506 are retained as typist instructions, whose precise targets are not sufficiently explicit to justify speculative paragraph restructuring. The second half of p.~507 turns from these operations to David W.~Jones's work on polyhedral cells and poly-T-complexes, supplied by Ronnie Brown. Grothendieck relates it to the questions of sections~91 and~93: when does a category of cells form a weak test category, and when does the proposed standard chain complex compute the homology of its presheaves? He finds that Jones's emphasis is instead on generalizing thin structures and a corresponding equivalence of categories. The original questions therefore remain open for this notebook. Both appearances of the presheaf-category notation have a hand-drawn upward caret in the scan, rather than the downward vee supplied by the witness. Section~121 next returns to the formal role that linearization should play across the various possible schematic approaches.

\subsection*{121. Abelianization, adjunction, and the pointed correction (pp. 508--510; 134-7/66--68).}
Grothendieck begins the formal review by defining $\mathrm{Hotab}=D.(Ab)$ and considering abel\-ian\-iza\-tion $LH.$ together with the Eilenberg--Maclane functor $K_\pi$ (p.~508). He relates their adjunction to the free-abelian-group and forgetful functors on semisimplicial objects, then interprets ordinary cohomology through the representing objects $K(\pi,n)$. The notation is visibly under revision: a blacked-out earlier functor label is replaced by $K_\pi$, a handwritten $K_\pi$ is itself inserted and canceled before the typed phrase ``defined via,'' and several category-name marks remain doubtful. The checked text does not reconstruct those older labels from the mathematical context. Page~509 explains why $K_\pi$ is conservative using $\pi_i(K_\pi(L.))\simeq H_i(L.)$, while explicitly denying full faithfulness and the assertion that the counit $LH.(K_\pi(L.))\to L.$ is an isomorphism. The typed letter $E$ in the preceding adjunction sentence is retained even though the surrounding notation uses $K_\pi$. The unit $X\to K_\pi(LH.(X))$ leads to the Hurewicz maps and a proposed homotopy-fiber interpretation of the Whitehead exact sequence (p.~509). The reduced homology in that sequence immediately calls for care with base points. Diagrams (8) and (10) factor the construction through pointed objects, using free pointing $X\amalg e$, the forgetful functor, and reduced free abelianization (pp.~509--510). The quotient description is replaced by an exact triangle, and the parenthetical typed ``one hopes!'' qualifies the continuing claim of adjunction. At the end of p.~510 the author explicitly corrects the earlier passage: the sequence and its fiber should come from the pointed unit (6'), since the unpointed expression lacks the canonical base points needed to define the indicated homotopy groups and fibers. The scan's(6') uses a tilde on both $K$ and $LH$, with no $\pi$ subscript on $K$; the witness regularizes this differently. Section~122 then revisits reduced abelianization itself, seeking a definition that no longer requires the chosen point.

\subsection*{122. Reduced abelianization without a point, and the failed homotopy-linearization functor (pp. 511--517; 134-7/69--75).}
The entry dated 2.10 revises the pointed construction of section~121. Instead of taking the quotient by the chosen base point, Grothendieck defines $\widetilde W(X)$ as the kernel of the augmentation $\underline Z^{(X)}\to\underline Z$ (p.~511). The fiber $W(X)(1)$ is the associated torsor and receives the universal map from $X$ to an abelian torsor. A point splits the extension, recovering the earlier reduced construction and giving the direct-sum description of homology; without a chosen point the natural exact triangle is arranged differently from the triangle on p.~510 (p.~512). The comparison is more informative over an object $Y$, because a section need not exist and the torsor class $c(X/Y)$ can be nontrivial. This supplies the connection with Postnikov invariants and motivates a more general Kan--Dold--Puppe correspondence for coefficients in the relevant slice topos (p.~513). The generalization to local test categories is announced as an intuition to revisit, not established here. The paper's notation is also simplified: the old $K_\pi$ and its pointed variant become $K$ and $\widetilde K$. A proposed partner $S$ or $\widetilde S$ would realize prescribed reduced homology, in contrast with the homotopy groups prescribed by $\widetilde K$ (pp.~514--515). The coefficient letter in the sphere-like objects $S(H,n)$ is changed in pen from $\pi$ to $H$, and circled marginal $H$ notes emphasize the change. Restriction to connected complexes addresses the immediate obstruction that degree-zero homology of a space is free. The author admits that he has neither constructed the desired functor nor established its existence or uniqueness. In the schematic setting even the relevant low-dimensional sphere objects have not yet been constructed. The possible value of this program depends on the tentative Lie-functor description from section~118, not on an available computation of sphere homotopy groups. The proposed fourth functor $L\pi.$ is more sharply rejected. It would turn homotopy groups into homology of a chain complex and satisfy $L\pi.(\widetilde K(L.))\simeq L.$. The continuation dated 3.10 states that these properties cannot all hold (p.~515). Applying the functor to the abelianization unit would provide a canonical section of the counit. On $L.=\pi[n]$, this would split the universal-coefficient sequence naturally in the groups involved (p.~516). The corrected indices matter: formula (18) has $\mathrm{Ext}^1$, $H^{n+1}$ and $H_{n+1}$, where the witness gives different indices. The text then specializes to degree one over $\underline F_2$ and the sequence $0\to V'\to\mathrm{Sym}^2(V')\to\bigwedge^2V'\to0$ (p.~517). An editorial mathematical qualification is needed at the last step. The manuscript claims equivariant nonsplitting for $\dim V\geq2$; the literal bound is retained. If $\mathrm{Aut}(V)$ means the finite group $\mathrm{Aut}_{\underline F_2}(V)$, as the preceding discussion of automorphisms of the abstract abelian group suggests, dimension two is an exception: $q=x^2+xy+y^2$ is fixed by all six invertible linear substitutions over $\underline F_2$, while its polar form maps to the generator of $\bigwedge^2V'$. It therefore supplies an equivariant section for that finite-group action. The qualification is specific to this interpretation: the calculation does not give equivariance for the algebraic group scheme $\mathrm{GL}(V)$ or for all matrices after scalar extension. This small calculation identifies a limitation of the asserted bound on the finite-group reading and does not repair the manuscript's text. The ensuing remark is conditional in a second way: it assumes that an analogous nonsplitting example can be found for $n\geq2$ before drawing the conclusion for simply connected homotopy types. Grothendieck presents that expected example as an alternative route to the failure of the earlier discrete/schematic comparison, avoiding reliance on the reported Breen results. Section~123 opens by acknowledging that, even after these discarded expectations, the basic schematic homotopy operations and the correct weak equivalences are still not in hand.

\subsection*{123. Sphere complexes and scalar operations as tests of the proposed theory (pp. 518--520; 134-7/76--78).}
The entry dated 4.10 takes stock after the preceding counterarguments. Even the equivalence of the proposed criteria for weak equivalence, and the construction of fibers and cofibers, are not yet settled (p.~518). Grothendieck therefore treats the realization of the ordinary two-sphere over $\underline Z$ as a basic test. Suspension should then produce the higher spheres; a larger category admitting the connected complexes needed for $S(1,k)$ might supply an alternative starting point (pp.~518--519). These are requirements for a workable theory, not constructions already obtained. If schematic sphere objects and the Lie description of their homotopy invariants exist, their homotopy groups should be the homology of canonical complexes $L\pi.(S^n_{\underline Z})$. Maps between spheres would then induce maps of those complexes, compatibly with composition (p.~519). The qualification from section~122 is crucial: the author has just rejected such a functor on arbitrary ordinary homotopy types. He now asks for a special structure on the sphere subcategory and explicitly allows the possibility that this structure is impossible, which would refute the intended schematic theory. The notation is still rough: the defining display for ${}^n\Pi.$ literally writes $L\,.(S^n)$, omitting the $\pi$ present in the surrounding discussion. The checked text retains that omission and the hand-corrected exponent and punctuation in the preceding displays. The first operation to examine is the action of the multiplicative monoid of integers, obtained by taking maps $S^n\to S^n$ (p.~520). The author relates the resulting actions on finite homotopy groups to the profinite monoid $\widehat{\underline Z}^{\mathrm{mult}}$ and proposes its lift to the derived base extension ${}^n\Pi.\otimes^{L}_{\underline Z}\widehat{\underline Z}$. The derived tensor symbol in the scan is not a direct-sum symbol, and no parentheses surround the tensor expression. For odd-dimensional spheres, restriction to the units is interpreted through the cyclotomic Galois action on the profinite homotopy type of punctured affine space. The pen changes from ``maximal'' to ``largest'' subgroup, adds ``(complex),'' and brackets the exponent $2m-1$; the above-line ``profinite'' remains typed. The point of the passage is the additional structure carried by these operations and a hoped-for lift to complexes, rather than a calculation of the sphere groups. Section~124 next separates the distinctive features of this schematic approach from the two alternatives considered earlier.

\subsection*{124. Four functors and a provisional formal calculus (pp. 521--530; 134-7/79--88).}
Grothendieck contrasts the present unipotent-bundle models with the earlier de Rham and diagonal-category approaches of sections~94 and~109. All seek a homology functor, but the present construction also paraphrases the Eilenberg--Maclane functor $\widetilde K$ compatibly with extension of scalars (p.~521). Simply forgetting a chain complex's $k$-module structure and taking a de Rham complex would not supply that compatibility. Moreover, the present schematic types map to ordinary types, whereas the earlier approaches had the reverse direction of comparison. The notation $\mathrm{Hotab}_0(k)$ deliberately leaves its algebraic realization open: section~115 required complexes of proobjects, and taking their projective limits would lose the asserted base-change behavior. Consequently $\widetilde K$ may require pro-unipotent models on the nonlinear side as well. Pointed linearization offers a different formal theory, whose invariants section~117 has already distinguished from schematic ones (p.~522). The continuation dated 5.10 organizes four expected functors between $\mathrm{Hotab}_0(k)$ and $\mathrm{Hot}_0(k)$: total homology $L\widetilde H.$, total homotopy $L\pi.$, and their partners $\widetilde K$ and $\widetilde S$ (p.~523). Even the zero-connected convention is provisional: the author lacks an intrinsic satisfactory formulation of abelian fundamental group with trivial higher action, and briefly considers restricting to simply connected objects. Property~A posits the adjoint pairs $(L\widetilde H.,\widetilde K)$ and $(\widetilde S,L\pi.)$. Their units and counits give linearization and colinearization. The blank after $L$ in formula (5)b remains unfilled in the scan; the checked text does not silently add $\pi$. The Hom formulas interpret $K(\pi,n)$ as representing cohomology and $S(H,n)$ as representing mixed homotopy with a ``co-coefficient'' $H$. Taking $H=k$ recovers the expected ordinary homotopy-module functor (p.~524). This is the stated motivation for the spherical adjunction, whose existence has not been established. The circled marginal $H$ and $\widetilde W(X)$ record changes to the coefficient and linearization notation. Property~B gives the inversion identities $L\pi.(\widetilde K(L.))\simeq L.$ and $L\widetilde H.(\widetilde S(L.))\simeq L.$ (p.~525). Property~C obtains a Hurewicz map by either applying total homotopy to the linearization map or total homology to the colinearization map. The desired compatibility is that these constructions agree. The first nonvanishing homology objects should occur in the same dimension, with an isomorphism there; a pen addition also demands an epimorphism in the next dimension. The lengthy reduction to Eilenberg--Maclane homology in the preceding paragraph is present in the typescript even though the witness omits it. Property~D states preservation of cofibration sequences by the left adjoints and of fibration sequences by the right adjoints. The two exact triangles spell out the homology and homotopy cases (pp.~525--526). In the second triangle the scan's corrected top vertex is $L\Pi.(Y)$ and its bottom row is $L\Pi.(Z)\to L\Pi.(X)$; the witness gives a different cyclic arrangement. Both oblique labels are the numeral~1. The exactness discussion reveals an unwanted asymmetry. Suspension on the connected linear category is a shift and is fully faithful, whereas loops require truncation and lose information (p.~526). The first truncation bound has been corrected from~0 to~1. Remark~1 explains that the loop and fiber constructions therefore mean the neutral or marked connected component of the unrestricted construction. Grothendieck says that he has only just become aware of this restriction while writing, questions the practical value of imposing $X_0=e$, and contemplates enlarging both categories (p.~527). The wished-for trivial action of $\pi_1$ on higher groups is a different condition and is not supplied by $X_0=e$. Remark~2 recalls the Leray spectral sequence for homology of fibrations and asks whether a corresponding mechanism exists for homotopy of cofibers; it does not supply such a mechanism. Remark~3 starts from the evident exchange symmetry of the four-functor formulas, but then distinguishes autodual axioms from an actual autoduality of the categories (pp.~527--529). Full faithfulness of suspension must be excluded from the symmetric axioms. Derived duality on perfect complexes sends chains to cochains, prompting a tentative appeal to stabilization. The subsequent recollection of geometric duality is explicitly qualified by ``as far as I remember'' and ``if I got it right''; the former qualification, the word ``discovered,'' and the shift qualification are typed insertions, placed with pen brackets (p.~529). That discussion pertains to homology and still does not produce a duality exchanging homology and homotopy. The literal sentence ``homotopy type of X'' omits the vee used on the complement elsewhere; the checked text preserves the omission. Properties~E--G return to the operational requirements (pp.~529--530). Homology and homotopy should be conservative in the indicated connected setting, and all four functors together with their structural maps should commute with extension of the ground ring. Restriction of scalars is mentioned without a developed list of formal properties. Comparison with ordinary types combines forgetting $k$ on the linear side with taking sections of schematic models. It should respect $\widetilde K$, but the already disproved homology comparison is explicitly excluded, even for $k=\underline Z$. Compatibility with $\widetilde S$ is still hoped for over $\underline Z$, conditional on constructing the discrete spherical functor. There is no ordinary total-homotopy functor to compare with $L\pi.$; only the separate homotopy invariants are available there. The closing expectation that the comparison respects fiber and cofiber structures is another requirement of the program. Section~125 begins on 13.10 with Baues's proposed axiomatic setting for making those operations precise.

\subsection*{125. From Baues's categories to weak equivalences and derivators (pp. 530--532; 134-7/88--90).}
The entry dated 13.10 records a letter from H.~J.~Baues and a reading of his preprint on the homotopy classification problem. Grothendieck recognizes the relevance of the dual notions of cofibration and fibration category to his integration and cointegration operations (p.~530). The present semisimplicial unipotent models, however, do not admit the pushouts he would need for a cofibration structure, except in the stated trivial case. He has begun to consider a Kan-fibration notion but reports no definite conclusion. He therefore proposes to leave the schematic construction temporarily and return after further work on these categorical methods (p.~531). The more general objection concerns the scope of the axioms. Grothendieck questions whether every meaningful homotopy universe must arise from Quillen's or Baues's extra structures on a chosen category of models. He invokes semisimplicial objects in a topos, Illusie's weak equivalences, and Verdier's derived categories as reasons to seek a wider formulation (p.~531). These are the author's contemporary judgments and questions, not proofs of general nonexistence of model structures. The qualifications ``it seems,'' ``I would think,'' and ``but (it seems)'' are significant: the last is a pen revision to the claim about enough projectives or injectives and bounded derived categories (p.~532). The empty reference brackets after Illusie and the marginal cross remain empty in the text. The circled $C^{-}(A)$ and $D^{-}(A)$ correct the signs used in the discussion of the bounded categories. The proposed replacement starts with a category $M$ and weak equivalences $W$, without treating a selected fibration or cofibration class as the fundamental datum (p.~532). Different compatible choices should yield the same integration and cointegration operations. Section~69 supplies the earlier point of reference; a derivator is now identified as the structure that should express those operations independently of a particular model presentation. Constructing and characterizing the schematic derivator remain hopes. The passage also frames the derivator as the further structure missing when one stops at a localized or triangulated category, in both abelian and nonabelian settings. The typed insertions ``specific,'' ``on it,'' and ``make any mention of models'' sharpen that distinction, while the correction from ``is'' to the typed ``are'' refers to the two operations. Section~126 resumes the schematic discussion by testing the formal symmetry against Postnikov decomposition.

\subsection*{126. Postnikov truncation and the failure of a literal dual filtration (pp. 533--535; 134-7/91--93).}
The section tests the four-functor symmetry on the construction that originally motivated the schematic program. An $n$-co-connected object has no homotopy groups above degree $n$, and the proposed truncation $X\to X_n$ is the universal map to that full subcategory (p.~533). Equivalently, its inclusion should admit a left adjoint. This universal property produces the familiar inverse tower, with $X$ viewed as an appropriate projective limit (p.~534). The discussion is axiomatic: it asks what should follow from the basic functors and their structural maps. The double underline on ``co'' distinguishes the upper-degree vanishing condition from connectedness. A literal dual would be a canonical increasing filtration ${}_nX\to X$, inducing the required homology isomorphisms through degree $n$ and vanishing above that degree. The ordinary skeletal filtration misses the endpoint of the desired isomorphism range. After trying to modify it, Grothendieck reports the feeling that no such canonical filtration exists, while explicitly admitting that he has neither formulated nor proved a nonexistence statement (p.~534). The decisive structural objection is that both proposed truncation subcategories are inverse images of the same linear subcategory under $L\pi.$ or $LH.$. That linear subcategory is stable under loops but not under suspension (p.~535). Even if the doubtful right adjoints existed on the homology side, their successive terms would not have the literally dual Postnikov relationship. A different suspension-stable family is contemplated and then abandoned for lack of a reasonable candidate. The conclusion qualifies the formal symmetry of section~124: homology and homotopy retain materially asymmetric properties. Several visible details prevent the reader from mistaking the formulas for a fully regularized exposition. On p.~533 the object named $Y$ is tested by a formula using $X$; on p.~534 the skeletal vanishing display writes $H_i(X_n)$ with a right subscript despite the surrounding left-subscript filtration. Formula labels are changed by hand, while some references to (5) remain on p.~535 after the relevant subcategory has been renumbered. Those literal features are retained. Section~127 then changes the question from symmetry of operations to the choice of components allowed in a schematic model.

\subsection*{127. Which objects should serve as schematic models? (pp. 535--538; 134-7/93--96).}
The section keeps the usual simplicial indexing category and changes the category of allowed components. The broad ambient setting is the fpqc topos of sheaves over the ground ring, with Illusie's weak equivalences of semisimplicial sheaves (pp.~535--536). Such a weak equivalence need not induce a weak equivalence on ordinary sections. Grothendieck therefore distinguishes the section functor from the canonical cointegration operation that should recover a discrete homotopy type. The above-line words ``canonical'' and ``(= weak equivalence)'' are typed revisions. He hopes that a suitably restricted class of models, especially the earlier unipotent bundles, will make the naive section construction adequate. The handwritten sentence opening p.~537 makes the base-change requirement explicit: a schematic homotopy type over $k$ should give a functor from $k$-algebras to ordinary homotopy types. The relative homotopy sheaves are defined first after choosing a section. The proposed ``pseudo-abelian'' condition combines relative zero-connectedness, abelian fundamental group, and trivial action on the higher homotopy sheaves (p.~537). Under a local formulation of that condition, the abelian sheaves should be independent of the chosen section and exist even without a global section. Their desired quasicoherent module structures are an additional issue. The comparison from the homotopy of sections to sections of the homotopy sheaves should be an isomorphism under mild hypotheses, for example $X_0=X_1=e$; the author expressly says he has not attempted a proof for unipotent bundles. He still calls the module structure a mystery without the Lie functor of section~118. The final remark proposes a spectral sequence from cohomology of the homotopy sheaves to the homotopy of cointegration, suggesting that acyclicity might imply the comparison (pp.~537--538). This remains a suggested argument, marked by ``maybe'' and ``If so,'' rather than a constructed spectral sequence. Section~128 next imposes a separate closure requirement: linearization should itself remain in the chosen category of bundles.

\subsection*{128. Linearization and the competing $W$- and $V$-bundles (pp. 538--541; 134-7/96--99).}
The new requirement is that linearization of an allowed bundle should itself be an allowed bundle, carrying the additional structure of an $\underline O_k$-module (p.~538). A forgetful functor $K:Uab(k)\to U(k)$ should therefore have a left adjoint $L$. The two natural choices have opposite variance: $W(M)(k')=M\otimes_k k'$ gives quasicoherent modules, while $V(M)(k')=\operatorname{Hom}_k(M,k')$ gives the opposite of the category of modules. The typed wording does not introduce an extra $X$ after ``every bundle'' or ``for a given bundle,'' despite the witness's additions, and formula (7) has a semicolon after $K$. On passing to semisimplicial objects, these functors are intended to produce the earlier $\widetilde K$ and $LH.$. The final dot of $LH.$ on p.~539 is baseline punctuation in the source, not a supplied subscript. Section~115 had already shown that the $W$-case gives only a proadjoint. Grothendieck now draws out the cost: one must work with prounipotent bundles and semisimplicial prosheaves outside the sheaf-theoretic foundations he is using (p.~539). A diagonal marginal addition offers working with the projective limit of the corresponding sheaves; its wording is ``we may however,'' and ``corr.'' remains abbreviated. The witness omits this qualification. The $V$-case is more tractable because $V(M)=\operatorname{Spec}(\operatorname{Sym}_kM)$. Pullback--pushforward adjunction and quasicoherence of the direct image for a quasicompact, quasiseparated scheme give $L_V(X)=V(\Gamma(X,\underline O_X))$, retaining only the underlying module of the global function algebra (pp.~539--540). The Hom subscript in the first display on p.~540 is $\underline O_S$, whereas the witness prints $\underline O_X$. Two circled marginal formulas emphasize the global-sections description. The replacement of ``inclusion'' by the typed ``forgetful'' makes the functor under discussion precise. This advantage raises a larger question: why restrict components to schemes isomorphic to vector bundles instead of allowing all affine schemes? In the affine setting, semisimplicial schemes correspond to cosemisimplicial commutative algebras, and linearization simply forgets multiplication (p.~540). The answer returns to the original ambition over $\underline Z$, where the schematic category should remain close to ordinary homotopy types. Grothendieck calls the initial choice of $W$ rather than $V$ almost a matter of chance, but formulates two substantive tests for either choice (p.~541). First, the homotopy-sheaf comparison with sections should hold, possibly under $X_0=e$ or $X_0=X_1=e$. Second, under flatness assumptions, the Lie functor should describe the homotopy sheaves. Postnikov d\'evissage supplies only the tentative special case, without the required naturality. He regards recovery from a first-order neighborhood as plausible for affine spaces and implausible for general affine schemes, citing an elliptic curve minus a point. Thus these tests motivate the restricted components; they are not established by the easier linearization adjunction. The next day's section~129 adds the further problem of descending that adjunction to localized categories.

\subsection*{129. The missing descent of linearization and q.c.h. equivalences (pp. 542--543; 134-7/100--101).}
The entry dated 14.10 corrects an omission in the preceding day's argument: existence of $L:U(k)\to Uab(k)$ does not itself imply a functor on the corresponding localized categories. One must specify weak equivalences on the linear objects and prove that $L$ preserves them (p.~542). Grothendieck expects preservation and reflection for the $W$- and $V$-models under suitable connectedness assumptions, but explicitly says that a proof has not even been attempted in the unipotent case. This is the third test, after the two tests on p.~541. He compares it with the Whitehead problem already considered in section~92 for replacing the simplex category by other indexing categories. The first live initial before Whitehead in the scan is a corrected typed $J$, over a doubtful older impression. The raised qualifications ``maps between'' and ``total'' are typed additions, and the filled versus open round indices remain graphically distinguished. For broader classes of semisimplicial schemes, the section proposes defining ``quasi-coherent homological quasi-isomorphisms'' by their image under linearization (p.~543). Ordinary weak equivalence need not yet be known to imply this condition, and a homological criterion can be weaker without connectedness restrictions. The proposed refinement allows twisted quasicoherent coefficients: a sheaf $F_*$ must be carried isomorphically to itself by the structural pullback maps. With a section and relative zero-connectedness, this is related to a fundamental-group-sheaf action on coefficients from the base. The circled marginal abbreviation q.c.h. is present in the manuscript and absent from the witness. Grothendieck then hopes that localization at these maps could make a mild category of general semisimplicial schemes equivalent to the category obtained from bundles, producing a subtler connection with ordinary homotopy types than the known étale pro-homotopy connection. He immediately suspends the suggestion as possibly out of reach or mistaken. The typed insertions ``very mild'' and ``only'' emphasize the proposed enlargement and comparison, not a proved equivalence. Section~130 returns to a more limited enlargement by allowing disconnected bundles.

\subsection*{130. Disconnected components as a possible aid to integration (pp. 544--545; 134-7/102--103).}
The limited enlargement now proposed allows direct sums of unipotent bundles as the components of a semisimplicial model. Deleting ``finite'' and replacing the following sentence make the role of infinite sums explicit: linearization still makes sense if the linear category is closed under them; otherwise one must retain a finiteness restriction (p.~544). Descent would also call for twisted finite sums, but the author postpones those technicalities. Applying connected components degree by degree produces a discrete semisimplicial object $E_* = \pi_{o/k}(X_*)$ and a map $X_*\to E_*$. Here $\operatorname{Spec}(k)$ is assumed connected and the sums untwisted. The ``strict'' schematic part lies over this discrete homotopy type, whose own invariants are not themselves schematic. To retain the original strict theory, $E_*$ should be aspheric, explicitly meaning isomorphic to a point in $\mathrm{Hot}$. Equivalence of the resulting localized category with the earlier connected-model category is still only presumed. The motivation is operational: the ordinary construction of integration from section~69 introduces many connected components, and standard simplices would become available among the models (p.~545). For an aspheric indexing category, Grothendieck hopes to preserve the asphericity of $E_*$. He proposes performing integration and cointegration first among semisimplicial sheaves, then checking whether the results lie in the restricted bundle class. He reports no clear-cut result. Moreover, disconnected components threaten the tangent-sheaf description of section~118, so the enlargement would be justified only after proving that it gives the same localized category as the connected models; no attempt at that proof has yet been made. The typed insertions ``components of the,'' ``sheaf,'' and ``would'' refine this cautious formulation. The canceled prospective number 131 belongs to the continued warning, while the actual section 131 starts below a literal date 16.9. That date is retained despite the preceding entry dated 14.10. Its opening personal interlude records the kimchi instructions and a letter to Larry Breen before the next page sharpens the mathematical objection to the enlargement.

\subsection*{131. Infinitesimal spheres and the failure of the naive Lie formula (pp. 545--551; 134-7/103--109).}
After the personal interlude opening the section, Grothendieck rejects the disconnected-model proposal as largely a symptom of his inability to construct the required ``left'' homotopy operations within connected unipotent bundles (pp.~545--546). The implicit condition $X_o=e$, or even $X_o=X_1=e$, is essential to the expected Lie description. Without it, tangent invariants at different sections need not be canonically identified, already for a constant semisimplicial bundle. Reimposing it would undermine the ordinary integration construction that motivated disconnected components in the first place. The author therefore distinguishes the restricted models for schematic homotopy types from more general semisimplicial sheaves that are useful auxiliary objects. The added pen $o$ in ``loosing track'' is clear and is retained, although the witness prints ``losing.'' A different enlargement is suggested by the still hypothetical spherical functor $\widetilde S$ of section~124, which should be left adjoint to total homotopy and right inverse to reduced total homology (pp.~547--548). For flat augmented symmetric algebras representing pointed $V$-bundles, the augmentation ideal modulo its square gives the co-Lie module. A module $M$ also defines the square-zero algebra $D(M)=k\oplus M$ and its spectrum $I(M)$. Applied degreewise, this construction gives the formal candidate $\widetilde S(L_*)=I(L^*)$ for $L_*=V(L^*)$. The resulting Hom identity expresses ordinary representability of the tangent functor by first-order infinitesimal schemes. Grothendieck explicitly marks two gaps: it is an identity between model objects, not yet an adjunction of derived categories, and the resulting infinitesimal schemes are outside the existing bundle category. The pen correction of ``adjunction'' and the deleted claim that the first Hom is taken in $\mathrm{Hotab}$ make the first gap especially visible on p.~548. Page~547 also literally introduces ``L'' as the arbitrary module immediately before writing $D(M)$; that variable mismatch is retained. The inversion identity with reduced homology is equally tautological for square-zero algebras, motivating a new definition of bundles as nilpotent thickenings of $V$-bundles (p.~549). The chosen nilpotent ideal is not part of the structure, and further flatness conditions remain possible. The subscheme in the displayed inclusion is $X_o$, whereas the witness prints $X_*$. Three of the functors now have simple model descriptions: augmentation ideals for reduced homology, symmetric algebras for $\widetilde K$, and their first-order truncations for $\widetilde S$. The canonical closed immersion $\widetilde S(L_*)\to\widetilde K(L_*)$ reveals a missing compatibility axiom in section~124 (p.~550). Its two adjuncts should be the inversion isomorphisms, and the maps obtained from those two isomorphisms must agree. Taking this natural transformation as primary data would also give a single Hurewicz map. The literal condition near the foot of p.~550 has right-hand side $k$, with a doubtful round left index; the witness replaces it by $X_0=e$. No such mathematical correction is made in the text. The apparent simplicity breaks down at the interpretation of total homotopy (p.~551). On the algebra side, naive $L\Pi.$ takes an augmentation ideal to its quotient by the square; dualizing gives the inclusion form of the Hurewicz map. But for a first-order thickening the square is already zero, so $L\Pi.(\widetilde S(L_*))\simeq L_*$. If these were the desired spheres, their naive Lie homotopy would therefore equal their homology, which the author emphatically rejects. His remaining hope is restricted to actual flat vector-bundle components, the basepoint condition, and possibly unspecified Kan hypotheses. The triple exclamation and the two marginal corner marks emphasize the failure and the beginning of section~132. That next section replaces the naive tangent functor by the Andr\'e--Quillen--Illusie tangent and cotangent complexes.

\subsection*{132. Derived tangent complexes and the wrong-quadrant obstruction (pp. 551--554; 134-7/109--112).}
The failure of the naive Lie functor in section~131 prompts replacement of ordinary differentials and tangent spaces by the Andr\'e--Quillen--Illusie complexes. Grothendieck emphasizes that the naive objects behave as the desired homological objects only in the smooth case, which here means the flat vector-bundle components under consideration (pp.~551--552). Formula (15) defines the superscript complex as the derived dual of the subscript complex, and formula (16) gives the corresponding co-Lie and Lie objects along the marked section by derived pullback. The prose first lists ``tangent or cotangent'' before those two notations; that visible order is retained, while the displayed duality and the subsequent co-Lie/Lie terminology specify the roles used in the argument. A free semisimplicial model of the co-Lie complex, followed by its dual, leads to the proposed formula (18) for total homotopy. Its handwritten question mark is essential: the formula is a candidate, not a theorem. Applying this construction to each component of a semisimplicial model introduces two indices of opposite variance. Under Kan--Dold--Puppe it becomes a bicomplex in the quadrant $n\leq0$, $p\geq0$ (p.~553). Taking its associated simple complex is tempting, but each total degree involves infinitely many terms. The manuscript therefore leaves open whether totalization should use a direct sum or a direct product. It also observes that the result can have both positive and negative degrees, so it need not lie in the desired chain-complex subcategory of the derived category. Truncation is proposed as a possible further correction if positive cohomology does not vanish. A separate difficulty is the nonquasicoherence introduced by duals without finiteness assumptions. Returning to the cotangent rather than tangent complexes would address that motivation but still leaves a mixed complex in a wrong quadrant (pp.~553--554). The proposed sphere $S(n,k)=\widetilde S(k[n])$ would be a concrete test of (18), especially over $\underline Z$, where its homotopy modules are hoped to recover the ordinary sphere groups. The author has not begun the necessary computation of the Andr\'e--Quillen Lie complex of the square-zero algebra $D(M)$ (p.~554). More generally, even applying the earlier flat-vector-bundle formula requires replacements of arbitrary models by models with flat components. Polynomial resolutions of the associated augmented algebras again introduce a second index of the opposite variance. The closing question is therefore whether these mixed wrong-quadrant objects could themselves form a workable category of models. Grothendieck expressly admits that his formal intuition breaks down and that the suggestion may be meaningless. The stepped double pen stroke before this question is a paragraph-separation cue whose upper arm passes beneath ``spectra''; it is not an independent underline emphasizing that word. The next supplied page is substantially masked, so the visible mathematical argument ends here without a resolution.

\subsection*{133 (witness numbering). A masked transition and the visible seminar interlude (pp. 555--557; 134-8/1--3).}
The supplied images of p.~555 expose the date 22.10. and page number but cover the entire body with an opaque rectangle. The upper body of p.~556 is likewise masked, and two smaller rectangles hide the end of the personal passage on p.~557. The separate arXiv witness supplies a section 133 heading and narrative in these positions, but neither that hidden wording nor the heading can be visually collated from the supplied scans. The checked text therefore retains only visible matter and records the masks as limits of the evidence; it does not silently fill them from the witness. The visible continuation concerns teaching for the concours d'agr\'egation and a small seminar on Teichm\"uller groupoids (p.~556). The typed revisions replace ``get involved in'' by ``start'' and ``great'' by ``special,'' while the pen addition describes the relaxed and friendly atmosphere of the teaching session. Grothendieck anticipates learning mathematics again through teaching, but worries that the renewed fascination of the Teichm\"uller structures will compete with completion of the homotopical-algebra notes. The continuation on p.~557 explicitly compares that fascination with motives; a pencilled marginal qualification reads ``it's very close anyhow.'' The visible text then breaks at the masks. Section~134 begins below them with the literal typed date 26.8., whereas the witness changes it to 26.10. The guide and checked text preserve the source date and make no inference about the hidden transition.

\subsection*{134. Linearization through integrators and two stronger comparison questions (pp. 557--560; 134-8/3--6).}
The return to ordinary homotopy types contrasts the earlier modelizer $(Cat)$ with the schematic discussion, where semisimplicial models and Kan--Dold--Puppe had dominated. For a general category $A$ and its terminal map $p$, the left adjoint $p^{ab}_{!}$ to constant abelian presheaves has a left derived functor $Lp^{ab}_{!}$, interpreted as total homology (pp.~557--558). Restricting it to individual abelian presheaves can be computed by pairing the presheaf with an integrator $L.^B$ for the opposite category $B=A^o$, then passing from chain complexes to the derived category. The text's later description on p.~558 calls this an integrator for $A$, or cointegrator for $B$; the varying terminology is retained, and the explicit pairing and variance supply the reliable description. Abelian weak equivalences are defined by this homology functor, producing the localization $\mathrm{Hotab}_A$. They must not be confused with quasi-isomorphisms of complexes of presheaves, which belong to a different input category. Even for $A=\Delta$, the unnormalized chain-complex functor in diagram (7) is not the Kan--Dold--Puppe equivalence: normalization is still required (p.~559). Grothendieck therefore places the hoped-for general equivalence at the lower, localized level $\mathrm{Hotab}_A\to D.(Ab)$. He suspects this for test categories, possibly even weak test categories. The first additional question asks whether a map of abelian presheaves is an abelian weak equivalence exactly when its underlying map is an ordinary weak equivalence. The proposed route is a functorial identification of the homotopy groups of the pointed category $A_{/F}$ with $H_i(A,F)$. The marginal note weakens the assumed final object in $A$ to 1-connectedness; it is absent from the witness. The author explicitly revises the convention of section~92: if the two classes of maps disagree, the homology-defined localization is now regarded as the correct $\mathrm{Hotab}_A$. Neither their equality for general $A$ nor the proposed homotopy--homology identifications are established here. The second additional question concerns homotopisms (p.~560). In the simplicial abelian case, the Kan property relates ordinary weak equivalence to ordinary homotopy equivalence. For objects with projective module values, Kan--Dold--Puppe also identifies weak equivalences with homotopisms in the stricter abelian sense and identifies equality in the localized category with strict abelian homotopy of maps. The projectivity condition belongs to this stronger statement. Grothendieck recalls the analogous formulations over a ground ring and in an abelian category, then asks for suitable interval structures on a general presheaf category. Sections~51,~54,~79,~95 and~97 provide the dependencies: interval and contractibility structures, the weak-equivalence construction, and contractors. He hopes that the three earlier interval structures coincide for a contractor and support these extra features; for a mere test category he suggests the structure defined from $W_A$. The qualifications ``it would seem,'' ``we may hope,'' and ``possibly'' prevent these proposed extensions from being read as results. The two marginal corner strokes and handwritten 135 separate the next section's attempt to recover an earlier abelianization proof.

\subsection*{135. Why integrators give the absolute Whitehead functor (pp. 560--564; 134-8/6--10).}
Grothendieck sets aside the stronger comparison questions of section~134, suggesting that closed model structures belong in a later chapter, and undertakes the promised result from sections~99--103: integrators are abelianizators (pp.~560--561). He first defines the absolute Whitehead functor without simplicial machinery. For a category $A$, total homology is $LH.(A)=Lp^{ab}_{A!}(\underline Z_A)$, and the desired functor on $\mathrm{Hot}$ follows once this construction is functorial and inverts weak equivalences. The handwritten parenthesis adds the name ``absolute Whitehead functor''; the opening ``First of all'' is instead a typed interlinear revision. For general coefficients the morphism data are a category map $f:A\to A'$ and a coefficient map to the pullback of $F'$. Formula (4) literally prints $u^*(F')$ after introducing $f$ as the category map. That visible variable mismatch is retained; the adjunction argument on the next page uses the pullback corresponding to $f$. Functoriality follows by composing the canonical augmentation $Lf^{ab}_{!}(F)\to f^{ab}_{!}(F)$ with the map to $F'$ adjunct to the coefficient map (p.~562). Composition of left adjoints gives the comparison of total homology, and preservation of projectives permits the derived composition used there. To show invariance under weak equivalence, the author tests the induced map by derived Hom into bounded-below complexes and identifies the result with the corresponding map in cohomology with constant coefficients. He then invokes the earlier definition of weak equivalence in $(Cat)$ through cohomology. The proof is presented as an argument, unlike the speculative extensions immediately preceding it. Several formula underlayers remain unreadable, but the supposed extra comma after $K_{A'}$ is not a separate source mark: the scan has the prime belonging to the subscript, which the initial transcription counted twice. The remaining comparison is the natural isomorphism $LH.(A_{/X},\underline Z)\simeq LH.(A,\underline Z^{(X)})$ (p.~563). With $f:A_{/X}\to A$, it follows from $Lf^{ab}_{!}(\underline Z_{A_{/X}})\simeq\underline Z^{(X)}$. The ordinary direct image of the constant presheaf already gives the right-hand side, and the key strengthening is that this particular $f^{ab}_{!}$ is exact. Page~564 gives the reason: $f$ is a fibration with discrete fibers, and its abelian left adjoint is the direct sum, over $u\in X(a)$, of the corresponding coefficient fibers. The formula is left to the reader to verify, after which exactness supplies the vanishing of all positive derived functors. Thus the square relating presheaf abelianization, the category of elements, total homology and the absolute Whitehead functor commutes canonically. This result does not by itself settle the stronger equivalence and homotopy questions of section~134. The closing reflection explains the delay as the author's habit of recasting homology in cohomological terms; the direct-sum argument finally provides the straightforward homological proof he had postponed.

\subsection*{136. Intrinsic homology and the corrected duality analogy (pp. 565--570; 134-8/11--16).}
The entry dated 27.10 rereads the abelianization arguments of sections~92 and~100--103 from the viewpoint of the unusual direct image. Grothendieck now regards the initial detour through simplicial models as an obstacle to recognizing the intrinsic construction $Lp^{ab}_{!}(\underline Z_A)$ (pp.~565--566). The comparison proved in section~135 supplies the missing compatibility with presheaf abelianization; for $A=\Delta$ it recovers the description through Kan--Dold--Puppe and localization. The new interlinear restriction to $\widehat\Delta_{ab}$ makes clear which domain is being compared with the classical functor. He draws a methodological conclusion: use the description native to the chosen model category when proving formal properties. The discussion of Thomason's homotopy-colimit construction is another example, contrasting the total category of a fibered category with the Bousfield--Kan presentation. This is the author's assessment, accompanied by an explicit promise to check the latter book, rather than a completed comparison proof. The references to page 569 on pp.~565--566 are the literal source numbers; the witness instead supplies 359 and 369 respectively. The compatibility argument extends to $k$-linearization for an arbitrary ring $k$, explicitly without assuming commutativity (p.~567). Formula (1) proposes obtaining it from the integral case by derived extension of scalars, $Lp^k_{!}(k_A)\simeq Lp^{\underline Z}_{!}(\underline Z_A)\otimes^{L}_{\underline Z}k$. This motivates a formulaire for $Lf_{!},f^*,Rf_*$,including a projection formula and the two internal operations derived tensor product and derived Hom. The analogy with six-operation duality is initially pushed too far: p.~567 calls $f^!$ right adjoint to $Rf_*$, while a pencil note corrects that expression to $Rf_!$. Page~568 nevertheless continues the speculation that suitable finite-type maps in $(Cat)$ might give four successive adjoints $Lf_!,f^*,Rf_*,f^!$. The text states uncertainty about whether there is a deeper relation with geometric duality. It does not construct the proposed extra adjoint. The typed interlinear replacement of ``identity'' by ``relationship'' is one visible softening of the claim. Page~569 explicitly diagnoses the confusion. For a proper map in the geometric duality contexts,$Rf_*$ agrees with $Rf_!$; in general the duality adjunction is $Rf_!\dashv f^!$, alongside $f^*\dashv Rf_*$. Even in the proper case,$Rf_!$ is not generally left adjoint to $f^*$,so the earlier chain of four adjoints does not follow. The pen replacement of a doubtful capital by $R$ in the first displayed pair, and the added explanation that $f^*$ is not isomorphic to $f^!$, belong to this correction. Discrete-fiber fibrations $A_{/X}\to A$ are proposed as the categorical analogue of the exceptional étale case. Grothendieck consequently allows that a future categorical duality theory might require a different $Rf_!$ from the $Lf_!$ currently expressing homology. The section therefore preserves an exploratory analogy and its immediate correction, rather than asserting an identification of those operations. The remaining proposal concerns classes of maps and factorization (pp.~569--570). The previously suggested proper and smooth maps are interchanged by passage to opposite categories,while immersions are self-dual. Proper maps include categorical cofibrations and smooth maps include fibrations; the author explicitly distinguishes those cofibrations from the Kan--Quillen notion. He asks how restrictive it is to factor a functor as an immersion followed by a map that is both proper and smooth,perhaps a bifibration. Here an immersion identifies a full subcategory containing every object lying between two of its objects in a composable pair of arrows. No factorization theorem is proved. The closing paragraph sets the larger duality project aside as unnecessary for the present homotopy-model program. Section~137 then immediately revisits the confidence that these definitions of properness and smoothness were the only reasonable ones.

\subsection*{137. Properness tested on stratified realizations (pp. 570--580; 134-8/16--26).}
The immediate reconsideration of section~136 separates two demands that the earlier proper and smooth terminology had conflated. Cohomological base change alone makes every object proper and smooth over the final category, since that projection is a bifibration. Such a weak condition cannot guarantee a right adjoint to direct image: even an infinite discrete category supplies an obstruction (p.~571). Grothendieck therefore asks for compatibility with filtered inductive limits and a finite cohomological-dimension condition. Finite groups do not automatically supply the latter, while infinite groups with Poincar\'e-type duality suggest another source of examples. These are tests for a definition, not a completed characterization. The section's frequent cancellations and replacements record the changing scope of the proposed finiteness hypotheses. The provisional six-operation picture assumes a factorization of a suitable map as an immersion followed by a proper map (p.~572). It proposes defining $Rf_!$ by a composite $Rg_*Li_!$, and defining the exceptional inverse image through the corresponding adjoints. The warning added to the page is essential: the notation $Rf_!$ is not meant to assert that this operation is the right derived functor of the already available $f_!$. For a smooth map the proposed comparison involves a rank-one orientation sheaf $T_f$ and a locally constant dimension $d_f$, with $Rf_!$ expressed through $Lf_!$, the inverse orientation twist and the shift $[-d_f]$. Formula (6) supplies the particular relationship recalled at the start of section~138. The factorization independence, the formalism and the smooth comparison remain proposed work. The typed additions ``partial'' and ``here'' qualify the discussion of the earlier confusion; the duplicated definition of $f^!$ is enclosed in a marginal oval. Finite ordered sets provide the testing ground (pp.~573--575). A presheaf on an ordered set can be interpreted as a sheaf constant along the strata of its simplicial realization, with the order reversed relative to containment of strata. Passing to the opposite ordered set does not change the underlying realized homotopy type, but it changes the stratification data relevant to sheaves. This distinguishes the geometric symmetry of realizations from the proposed categorical interchange of properness and smoothness. Finite ordered sets should give proper objects; smoothness is tentatively related to spherical links. For a map of ordered sets, the first proposed condition asks whether every relation $f(x)<y'$ can be lifted to a relation $x<x'$ with $f(x')=y'$. Requiring a unique lift is recognized as too strong. The stronger replacement is that the ordered set of such lifts, written $I(f,x,y')$ in formula (13), should be aspheric (p.~576). Acyclicity may suffice for a version restricted to commutative coefficient groups, but the text does not conflate that weaker condition with asphericity. Formula (12) is meant to compare the relevant upper cones and thereby explain preservation of sheaves compatible with the stratifications. Its printed target is literally $C_y=|I_{\geq y}|$, although $y$ belongs to $J$; the category letter is retained in the transcription. The pencil note ``namely: aspheric fibers!'' sharpens the intended geometric criterion. Page~577 then distinguishes existence of the six operations and their boundedness restrictions from preservation of constructibility or finiteness. The latter is needed for the intended Grothendieck-group calculations but is not made a condition for every formula in the general duality formulaire. The first remark on p.~578 identifies biduality as a place where finiteness is needed. It asks for a dualizing complex and for its realization as $p^!(k_e)$, with uniqueness qualified by an added dimension shift and an invertible coefficient twist. The later running formula writes $p^!(e_k)$ instead; that index reversal is a visible source feature. For finite ordered sets, or locally finite ones of finite combinatorial dimension, the author expects the relevant dualizing statement to follow from the contemplated formalism. He also suggests an intersection-complex construction, while doubting its relevance when the strata are not even-dimensional and when the outcome depends on the stratification. These are explicit qualifications, not claims of a stratification-independent invariant. Two examples show why merely nonempty lifting sets are insufficient (pp.~578--580). First, adjoining a smallest element $x$ to an arbitrary ordered set $I_1$ gives a cone and a map to the two-element ordered set. With $I_o=f^{-1}(0)=\{x\}$, the lifting condition asks only that $I_1$ be nonempty, while the stronger condition asks for asphericity. The text later names 1 as the exceptional point fiber and writes an unbarred $I_1$ where a realized fiber is discussed; p.~579 repeats the endpoint labels in the complements $|J|\setminus\{1\}$ and $|J|\setminus\{0\}$. Those labels are not harmonized with the earlier definition of $I_o$. The argument's intended contrast remains clear: a nonaspheric general fiber and a point fiber prevent the higher direct images of constant coefficients from remaining constant on the required stratum. The second example avoids relying solely on that two-stratum interval. Its target has three elements $y\to0$ and $y\to1$, giving the ordinary interval with its two endpoints. In a polygon with $2m\geq4$ vertices, $I_o$ and $I_1$ are the alternating vertex sets. The sentence specifying this polygon is present in the typescript and absent from the witness. The nonempty-lift condition holds, yet the relevant fibers are disconnected: at an endpoint there are $m$ points, away from the endpoints and barycenter there are $m$ disjoint segments, and at the barycenter those segments meet and form a contractible union. Thus even $R^0|f|_*$ of a constant sheaf fails the desired constancy on the open one-dimensional stratum. The closing paragraph proposes that the ordered lift sets have exactly the homotopy types of the geometric fibers of the maps between strata. Asphericity would then express the required contractibility and explain preservation of stratification-compatible complexes. The phrasing remains ``should''; no general comparison proof is supplied before section~138 begins on 28.10.

\subsection*{138. Coefficient rings, localization and the existence of $f_!$ (pp. 580--587; 134-8/26--33).}
The entry dated 28.10 first enlarges the preceding proposal from a fixed coefficient ring to categories equipped with sheaves of rings (p.~580). Commutativity is required when the internal tensor and Hom operations are considered. The author also insists that the homological operation $Lf_!$ is an additional operation whose relation with geometric $Rf_!$ must be understood; the smooth comparison on p.~572 is recalled as one such relation. The pencil exclamation and cross emphasize the addition. Yet the claim that $Lf_!$ exists for every map of ringed categories is explicitly withdrawn at the end of the section. The opening recollection on p.~581 itself writes $Rf_!$ before turning to the ordinary $f_!$ and its left derived functor. That visible notation is retained rather than silently made uniform. The first obstruction concerns restriction of operator rings (pp.~581--582). For a map of ringed categories with underlying functor $f_o$, the abelian sheaf $f^{ab}_{o!}(F)$ need not visibly acquire the required action of the target sheaf of rings. The source's test case takes $\underline O_A=\underline{End}_{\underline Z}(F)$ and $\underline O_B=f_*(\underline O_A)$. It supplies an action by global sections of $\underline O_B$, but not the desired sheafwise action. If the target ring is constant, that global action yields an ordinary coefficient action. Derived compatibility imposes a further acyclicity requirement on projective source modules, and the example with $B$ final shows that it can fail for a variable source ring. The square-zero extension $\underline Z_A\oplus M$ makes clear how strong universal vanishing would be: it would force homological dimension zero. The scan has a direct sum here; the witness's tensor-product sign is incorrect. A separate literal mismatch on p.~581 writes $f_!(A)$ after introducing the module $F$. For constant source and target rings, the argument reduces to the generators $M^{(a)}$ (p.~582). Resolving $M$ in abelian groups and applying the exact functor $M\mapsto M^{(a)}$ gives a resolution whose direct image is the corresponding complex $L.^{(b)}$, with $b=f_o(a)$. Exactness of $M\mapsto M^{(b)}$ then proves the vanishing in formula (4). Thus the source presents an argument for the constant-ring comparison and only suggests the locally constant extension. The sentence asserting that one cannot go beyond that assumption is canceled; the interlinear ``rid of such an'' is also separately crossed out. The distinction between a proved calculation and a withdrawn restriction matters for the next discussion. The obstacle is recast as a variance issue (pp.~583--585). Ordinary and derived direct image commute with localization on the base, represented by $B_{/b}\to B$, because the relevant inverse image preserves injectives. The source extends this to discrete-fiber fibrations and then to fibrations more generally. One intermediate display literally reads $B_{/Y}\to Y$; the target letter is retained despite the surrounding discussion of maps to $B$. For $f_!$, the appropriate operation is the opposite construction ${}_{b\backslash}B\to B$, called colocalization. With fixed coefficients, the value of $f_!(F)$ at $b$ is the inductive limit over ${}_{b\backslash}A$, dual to the projective-limit formula for $f_*$. Passing to $Lf_!$ requires preservation of projectives by the corresponding inverse image. The pasted replacement strip on p.~584 supplies the revised proof setup; faint typing beneath it is not reconstructed. Page~585 finishes the argument by expressing the right adjoint along a discrete-fiber cofibration as a product over the fiber. The asserted exactness of those products gives preservation of epimorphisms and hence the required projective statement. The formula is explicitly compared with the direct-sum calculation of section~135. The first remark allows the injective/projective preservation statements with an induced sheaf of rings $\underline O_{A'}=h^*\underline O_A$, while still doubting that this alone suffices for the colocalization comparison (p.~585). The second remark corrects the more basic existence claim. Enough projectives permits deriving an additive functor that already exists; it does not prove that the inverse-image functor has a left adjoint. Even for one-object final categories, extension of scalars need not be left exact or commute with products (p.~586). The source identifies the combined condition with finite projectivity of the relevant coefficient modules, then adopts the simpler induced-ring condition for the intended overview. To recover restriction-of-scalars and colocalization compatibility, it returns to constant, or perhaps locally constant, coefficients. The concluding return to chapter~V therefore has a restricted scope: the author seeks a mutually dual account of $Lf_!,f^*,Rf_*$ and the two internal operations, without yet producing a complete six-operation theory. Section~139 asks whether resolving the integrator instead of the coefficients can actually supply that relative account.

\subsection*{139. The limits of relative integrators and the coalgebra question (pp. 587--591; 134-8/33--37).}
The entry dated 29.10 returns to the origin of chapter~V: standard complexes for test categories, the intermediate notion of an abelianizator, and the integrator construction (p.~587). The pairing between presheaves on $A$ and on $A^o$ had made absolute homology and cohomology computationally attractive. A projective resolution of the constant unit sheaf could be used in the integrator argument instead of first resolving the coefficient complex. The author now questions whether this device carries over to the relative operations $Lf_!$ and $Rf_*$. His criticism is specific to that proposed use. It does not retract the absolute Whitehead construction or the integrator comparison established earlier. The typed corrections from $R$ to $L$ in the opening discussion restore the homological operation under consideration, while the later pen correction on p.~588 literally leaves $L.$ where the input coefficient complex was denoted $F.$; that local notation is not silently harmonized. Formula (1) proposes computing $Rf_*(K^\bullet)$ by applying ordinary $f_*$ to the internal Hom from a cointegrator $L^A.$ (p.~588). The apparent justification would reduce the formula to the known absolute case on the localizations $A_{/b}$. The failure is that restricting a projective resolution to these localizations need not preserve projectivity. The previous section instead obtained the projective preservation statement for colocalization. A concrete counterexample makes the distinction decisive: adjoin a final object to an arbitrary $A_0$ and map the resulting cone to the two-object category $0\to1$. The constant sheaf on the cone is projective, yet the fiber of $R^if_*(K)$ at 0 is $H^i(A_0,K_0)$, which can be nonzero for positive $i$. Thus a final object in the source does not supply the desired acyclicity for an arbitrary relative direct image. The same objection leads to an explicit retraction of formulae (12),(12') and (13) at the end of section~101, p.~377 (pp.~588--589). Their resemblance to familiar duality formulae had encouraged the substitution of $Lf_!$ and $f^*$ for geometric $Rf_!$ and $f^!$. Formula (2) tests this analogy on $Rf_*(f^*K)$. When source and target final objects are preserved, the proposed right-hand side reduces to $K$, incorrectly forcing the positive higher direct images to vanish. The large crossed-out paragraph on p.~589 records an abandoned intermediate justification. It is surrounded by a pen enclosure and several large Xs; the two missing relation signs in its display were never supplied. Those gaps are not completed by mathematical expectation. The surviving counterexample, rather than the rejected calculation, carries the correction. The practical conclusion is to resolve the coefficients in the usual way: projectively for homology and injectively for cohomology (pp.~589--590). Even the absolute internal tensor and Hom operations had already required resolutions of their arguments. The criterion recalled from section~109, proposition 4 on p.~433, supplies enough projectives or injectives in the relevant presheaf categories when the coefficient category has them, together with the limit conditions needed for the ordinary direct-image functors. Grothendieck treats the adjunctions among $Lf_!,f^*,Rf_*$, their transitivity and the expected projection formulae as routine work he does not carry out here. He nevertheless singles out base change, especially for $Lf_!$, as requiring care. Integrators and cointegrators are consequently placed within the familiar operation of resolving the unit sheaves $k_A$ and $k_{A^o}$, rather than presented as a universal shortcut around the relative derived-functor problem. The broader insight he retains from chapter~V is the passage to a $k$-additive envelope equipped with its natural diagonal, recalled from section~109 (pp.~590--591). This is what had prompted the subsequent schematization discussion and the question of homotopy types over arbitrary ground rings. He asks whether the commutative homology and cohomology formalism can be expressed directly on these coalgebra structures. Properness, smoothness and immersion would first need sufficiently clear definitions in $(Cat)$ before their recovery from the linearized structure could be tested. The analogous question for Steenrod operations is also left open; the remark about sheaf coefficients over $\underline F_p$ is presented as a recollection, not proved in this section. General sheaf coefficients move the discussion beyond an ordinary homotopy model (p.~591). A category now determines a whole presheaf topos, carrying information about dimension, compactness and other properties that a bare homotopy type forgets. Grothendieck hopes that the coalgebra structure retains much, perhaps all, of what can be expressed through commutative coefficient sheaves. He is substantially less confident about the noncommutative cohomology that motivates \emph{Pursuing Stacks}. The final possible exception would recover $A$ from $P=\mathrm{Add}_k(A)$ itself, perhaps as a category of ``exponential'' pairs $(x,u)$ with $u:\delta(x)\simeq x\otimes_kx$. The displayed condition is a proposal for further investigation; no reconstruction theorem or sufficiency proof is given. The typed insertion $P=$, the pen addition ``exponential,'' and the retained uncertainty in the heavily corrected involving-like word preserve the exploratory state of this closing page.

\subsection*{140. A personal interruption and the return to scratchwork (pp. 592--593; 134-8/38--39).}
The entry dated 4.11 interrupts the mathematical narrative with two images of personal experience: a stage on which one chooses roles, and a meal whose material must be assimilated (p.~592). Conflict, suffering, joy and tenderness belong to the account of learning about oneself. The typed insertion ``and of joy'' expands the earlier pairing of suffering and violence, while the pen replacement of ``diary'' by ``journal intime'' accompanies the stated wish to return to mathematics instead of narrating recent family events. The guide treats this passage as the author's reflection, without turning its metaphors into psychological or medical claims. The entry dated 12.11 qualifies the preceding self-presentation (pp.~592--593). Grothendieck calls the image of an effortless learner misleading and describes resistance, delay and eventual assimilation, contrasting this with mathematical work on material from which his own person is absent. The revisions make that qualification concrete: ``our'' becomes ``my,'' ``where'' becomes ``which,'' and ``keeping it'' becomes the typed ``kept.'' The final period in that replacement and the separate pen period after ``decennaries'' are source features. The closing paragraph reports that mathematical scratchwork has continued along the lines of abelianization. The supplied typescript therefore ends with renewed work and unresolved questions, rather than a theorem closing the formal program of the preceding sections.

\subsection*{Brown's translation and comments: Overview of the appendix}
The complete appendix preserves Brown's earlier English copy as an independent textual state: an introductory note, eleven comments updating the program in 1982, the February 1975 postscript on local systems and integration, and a July extract moving from weak higher structures to stratified topology and reconstruction. The February claims about successive deloopings must be read with the July withdrawal of the simple identification of iterated loop structures with strict commutativity. Brown's comments then distinguish the accomplishments of crossed complexes, group-groupoids, and multiple groupoids from the still broader homotopy program. The July extract explicitly omits the biduality and homology discussion; the corresponding material, added subtitles, and retrospective notes in the 1983 version are not interpolated. The guide's detailed paragraphs identify consequential translation slips and source formulas, while the supplemental Brown concordance accounts for every supplied witness difference. This copy is particularly useful for following the development of the program: its earlier questions coexist with the later author's retractions only when the two versions are compared, not when they are silently merged.

\subsection*{Brown's introductory note (cote 134-2, site 115)}
Brown identifies the following text as an informal translation of the mathematical part of Grothendieck's letter to Breen dated 17--19 July 1975, prepared with assistance from Breen and Loday, and thanks Grothendieck for copies of that letter and the two related earlier letters. This establishes the attribution and the intended relation to the French letters reproduced earlier in cote 134-2; it is not a claim that Brown authored the mathematical program or that all three letters are translated in this one note. The separate translation of the February postscript and the July letter must therefore be distinguished from the scope stated on this cover sheet. The note contains no mathematical correction or change of position. Its lower-left handwritten 116 is archive pagination, not an original manuscript section or a new part of the letter; the archive's site 115 and PDF page 116 accordingly differ by one.

\subsection*{Brown's comment 1: Subsequent work on duality (site 117)}
Brown opens his November 1982 comments by observing that duality has advanced since the 1975 letter and naming work by Begueri and Bester, with the references attributed to Breen. This places the forthcoming translation in its historical sequence: the speculative account of duality in the July letter is not presented as an inventory of what was known in 1982. The pencilled supplements identify a memoir of the Société mathématique de France and a Michigan thesis, but do not supply complete publication details; the darker marginal request to ask Breen for references, with its grouping line, records that incompleteness. The reader therefore should not infer a precise theorem or bibliographical identification from this comment alone. Its role is to qualify the context of the duality program later translated, not to amend any formula in it; no conjecture is withdrawn here.

\subsection*{Brown's comment 2: Gr-categories and symmetric groupoids (site 117)}
The second comment supplies the reference to Sinh's thesis requested by the letter's discussion of group-like categories, adds Saavedra Rivano's work, and places Takeuchi's study of the Villamayor--Zelinsky sequence beside these notions. Brown's literal notation is ``+-groupoid'', accompanied by the description of a coherent symmetric monoidal groupoid with strict associativity and a strict zero object. These strictness qualifications belong to the definition being cited and must be kept when comparing it with the weak structures envisaged by Grothendieck. The handwriting adds the accent to État and replaces an obscured initial with Saavedra before Rivano; the canceled initial remains doubtful. This is bibliographical supplementation of the earlier program, not an assertion that all its higher coherent structures are strict. It prepares the contrast between the strict models in comments 3--5 and the more flexible multiple-groupoid constructions in comment 6.

\subsection*{Brown's comment 3: Group-groupoids and crossed modules (site 117)}
Brown identifies the letter's description of groups in $(1\text{-cat})$ with the equivalence between group-groupoids, explicitly group objects in the category of groupoids, and crossed modules. The Brown--Spencer reference supplies the precise comparison, and the sentence about Verdier reports a belief about prior knowledge in 1965 rather than a documentary priority claim. This is the first concrete algebraic model in the comments and the starting point for the six-category diagram of comment 4. The meaning of group-groupoid here includes the internal group structure; it must not be replaced by an arbitrary groupoid. Nothing in this equivalence alone establishes the full higher homotopy hypothesis. The later warning about strict infinity-groupoids concerns its attempted higher-dimensional extension rather than a retraction of the dimension-two comparison.

\subsection*{Brown's comment 4: Six equivalent algebraic models (sites 117--118)}
The diagram joins poly-T-complexes, cubical T-complexes, simplicial T-complexes, crossed complexes, omega-groupoids, and infinity-groupoids by seven bidirectional equivalences, with letters referring to the following bibliography. It extends the previous comparison and directs the reader to the cube algebra, inductive limit theorems, Ashley's simplicial construction, and Jones's work in preparation. The two vertical connections, the four horizontal connections, and the diagonal from omega-groupoids to crossed complexes are drawn by hand; their object names and reference labels are typed. The omega/infinity distinction is part of Brown's terminology here and is preserved. The heavily overtyped symbol in reference d remains uncertain rather than being supplied from the expected published title. Comment 5 immediately restricts the homotopy types represented by these models, while comment 8 uses their equivalences to construct a nerve. Thus the diagram is a comparison within this strict algebraic setting, not the sought equivalence with all spaces.

\subsection*{Brown's comment 5: The limitation of crossed complexes (site 118)}
Brown explicitly says that the infinity-groupoids in the preceding diagram are inadequate for Grothendieck's purpose. His deliberately approximate explanation describes spaces fibering over a $K(\pi,1)$ with fiber a product of Eilenberg--Mac Lane spaces. The phrase ``roughly speaking'' matters: the comment identifies a restriction in the higher homotopy information carried by the models, without giving a complete necessary-and-sufficient classification theorem with all actions and hypotheses spelled out. This is the key qualification on comment 4 and the reason that merely citing equivalence with crossed complexes does not settle the program of the translated letters. Comment 9 repeats the limitation after the classifying-space construction. The readers preserve the assertion exactly; the guide distinguishes Brown's approximation from a general claim about all spaces or all meanings of infinity-groupoid.

\subsection*{Brown's comment 6: Multiple groupoids and truncated types (sites 118--119)}
Loday's n-cat-groups enter as n-fold categories in the category of groups and are reported to describe truncated homotopy types. This offers a different route from the strict crossed-complex models just restricted in comment 5, and supplies the background for the handwritten van Kampen update in comment 11. The description is brief and does not state a truncation degree or the connectivity conventions. In comparing it with Grothendieck's general n-types, the group-based construction must be read in its connected, pointed setting; the sentence is not by itself a model for arbitrary disconnected spaces. Brown does not retract the crossed-complex equivalences, but changes the proposed algebraic model when greater nonabelian information is required. The literal wording and the n indexing are preserved rather than expanded into a stronger theorem in the readers.

\subsection*{Brown's comment 7: Length-two nonabelian complexes (site 119)}
Brown cites Conduché's work in preparation on the nonabelian chain-complex structure arising from a simplicial group whose normalized chain complex has length two. This directly addresses the letters' search for noncommutative analogues of complexes and provides another concrete development beyond the abelian behavior diagnosed in comment 5. The comment supplies neither the additional operations nor their identities, so it should be used as a pointer rather than as a definition of a general nonabelian derived category. Grothendieck's later note 27 in cote 134-2 records the same report from Brown, giving a direct connection between this bibliographical update and the author's retrospective annotations. Brown's printed title contains ``longeur'', which remains literal in every form; the guide's account does not silently turn that spelling into evidence of a different mathematical notion.

\subsection*{Brown's comment 8: A nerve and the missing categorical extension (site 119)}
Following the equivalences in comment 4, Brown associates a simplicial T-complex $NG$ to an infinity-groupoid $G$ and describes its n-simplices as maps from the homotopy infinity-groupoid of the decorated standard n-simplex into $G$. He relates that object to the corresponding homotopy crossed complex, then states that the analogous nerve for the infinity-categories of the cited reference is still unknown. The four marks below Delta are each a single wave, not the double-wave approximation signs in the initial transcription; their typed-versus-handwritten provenance remains doubtful. Preserving this decoration matters because the formula names a structured simplex rather than asserting an approximate equality. The direction ``to G from'' is plain typed wording, and the later functor arrow points from infinity-categories to simplicial sets. The limitation is a dated statement about this particular extension, not a claim that no nerve of any higher-category model exists. Comment 9 uses the groupoid/crossed-complex nerve that has been constructed.

\subsection*{Brown's comment 9: Classifying spaces and nonabelian cohomology (sites 119--120)}
The two references connect crossed complexes with nonabelian extensions and with homotopy classification of maps. Brown explains that the equivalence with simplicial T-complexes extends the Dold--Kan comparison and yields a classifying space $BC$; the displayed bijection $[X,BC]\cong H^0(X,C)$ is presented beside the Eilenberg--Mac Lane formula $[X,K(A,n)]\cong H^n(X,A)$. The zero on the left-hand theory's cohomological degree is literal and must not be changed to n by analogy: its coefficients are a crossed complex, not merely the abelian group in the preceding formula. The printed hypothesis is that X is a CW-complex. Brown then returns explicitly to comment 5 and calls the results still ``too abelian'', so this construction is an intermediate achievement within the program rather than its completion. His bibliographical ``Astérixes'' and punctuation remain as typed; no later publication data are substituted into the letter.

\subsection*{Brown's comment 10: Algebraic tubular neighborhoods (site 120)}
Brown points to Cox's two papers on algebraic tubular neighborhoods and a paper on homotopy theory of simplicial schemes. These are intended as later developments relevant to the July letter's geometric program, especially the proposed local models and the organization of stratified spaces. The comment itself makes no claim that Cox's constructions satisfy all the precise requirements Grothendieck subsequently sets for dévissage; the author's note 28 earlier in cote 134-2 explicitly reserves judgment because he has not yet seen the work. This distinction prevents a bibliographical connection from being read as a proof of the proposed stratified theory. The punctuation after Math remains visually uncertain. Brown's signature, affiliation, date November 1982, and handwritten U.K. belong to the commentary's provenance and are kept separate from the 1975 letters.

\subsection*{Brown's comment 11: A van Kampen update (site 120)}
The pencilled final item reports that Brown and Loday have proved a van Kampen theorem for 2-cat-groups. It depends on the multiple-groupoid direction introduced in comment 6 and indicates how such algebraic models can support a local-to-global theorem, a central demand of Grothendieck's program. The short note gives no hypotheses, statement, or reference, so the guide does not infer a van Kampen theorem for arbitrary weak higher groupoids from it. Its placement below the date and lack of a relocation cue make it an added numbered comment, preserved as an anchored marginal note, not a sentence to be silently inserted into the 1975 translation. The expression 2-cat-groups is clear on the scan. The added report changes the status of one proposed line of work while leaving the earlier cautions about strict crossed complexes intact.

\subsection*{Brown's February translation: Local systems and reconstruction (sites 122--123)}
The letter dated 17 February 1975 is an afterthought to the preceding homotopy program. It associates a pro-homotopy type with a topos, passes to an ordinary truncated type under local homotopical triviality, and compares cohomological invariance with conditions A and B on locally constant higher stacks. Condition C proposes their description as higher functors from the fundamental n-groupoid, extending the familiar classification of coverings. The qualification ``locally constant'', explicit throughout the surrounding discussion, is missing from C's opening words and must be understood there; arbitrary higher stacks are not classified by the fundamental groupoid alone. Reconstruction from local systems then leads to higher classifying topoi and exact fiber functors. The arrows under lim correctly distinguish arbitrary inductive limits from finite projective limits. The weaker proposal of an ordinary classifying topos is immediately questioned for n=2; the typed marginal comment ``this seems suspect'', with its clipped first letter still doubtful, records the objection. The proposed obstruction through van Kampen and descent over two closed subsets remains a question, not an established counterexample. Brown writes a low 0 after phi in the displayed action on a morphism; its role is composition with the induced fundamental-groupoid map, but the printed subscript is retained. This portion corresponds to Appendix 6 in the earlier French typescript and supplies the reconstruction problem on which the following integration formulas depend; the marginal topoi and Stack glosses do not replace Brown's running terminology.

\subsection*{Brown's February translation: Integration and cohomology (sites 124--126)}
Integration is described as Hom from a final higher stack, and strict Picard structures permit its homotopy groups to express hypercohomology after iterated application of the classifying-object operation. Groupoid-valued stacks supply the initial degree range; a group structure allows one further degree, while the proposed necessity of strict Picard structure for a second delooping is tentative and must be read with the July reconsideration. Several literal features need care. Site 124 calls the object after extension to degree N an n-champ although the construction raises its categorical degree, and prints $N\leq i,n$ where the immediately following sentence requires $N\geq\sup(i,n)$. Its comparison from the classifying object of sections to sections of the classifying stack is called a morphism of $(n-1)$-groupoids, although the continuation asserts isomorphisms through degree $n+1$. The small Gamma subscripts on sites 124 and 125 are crosses, unlike the capital X elsewhere. These are retained and explained rather than repaired as typography. More substantially, Brown's clearly typed ``one knows it fails'' on site 125 does not translate the French ``on a fait ce qu'il fallait'' at site 53, which says that the construction has done what is needed to obtain the canonical comparison; the English phrase reverses the sense of that transition. Conditions A$'$ and B$'$ then restate the proposed characterization using truncated higher functors. B$'$ has the same truncation problem as the French source: for n=0, a map from a point to the one-object groupoid of a nontrivial group induces an equivalence on zero-truncations, but the section map from the fixed points of a free nonempty group-set to that set is not a bijection. The diagram's right lower term also abbreviates the pulled-back coefficient by F. These qualifications concern the proposed mathematics and the translation, not the accuracy of its now-checked literal transcription.

\subsection*{Brown's February translation: Three approaches to cohomology (site 126)}
The closing paragraphs ask whether integration of higher stacks can furnish cohomology without complexes of abelian sheaves, injective resolutions, or semi-simplicial algebra. Hypercoverings suggest a connection with the Čech approach, but independence from semi-simplicial methods is explicitly left open. The resulting comparison is conditional: commutative homological algebra resolves coefficients, homotopical algebra resolves the base, and the higher-stack approach appears to resolve neither. The three descriptions organize the earlier integration proposal and anticipate the main work's later concern with preserving structured homotopy information rather than only derived-category classes. They do not certify a completed third construction. Brown's source distinguishes a handwritten first hacek from a typed v placed over the second C, and keeps the typed closing salutation and Alexandre signature. This portion corresponds to Appendix 8; unlike the consequential omission in the arXiv witness to the French passage, the translated question about avoiding semi-simplicial algebra is present in full.

\subsection*{Brown's July translation: Axioms and the correction about H-spaces (sites 128--129)}
The July extract begins with Grothendieck's response to Breen's constructive proposal. He values its precision but objects to its technical character and, especially, the absence of the geometrically expected functor from simplicial sets to higher groupoids. A sufficiently complete list of desired notions and relations might lead to an explicit construction or an existence-and-uniqueness theorem. Strict Picard n-categories are required to recover complexes of chains of length n, but on site 129 he explicitly withdraws his earlier belief that arbitrarily high loop-space structures correspond simply to commutative topological groups. This correction must accompany the tentative restriction on repeated delooping in the February letter; strict algebraic commutativity and coherent topological commutativity are not interchangeable assumptions. The passage corresponds to Appendix 9 in the revised 1983 version, which adds subtitles and retrospective notes absent from this Brown copy. Brown's visible heading identifies an extract and its July date; two masked title areas cannot be recovered. His body says ``a theorem'', ``functor nerve'', and ``seems to me to be useless'', where the later revised witness uses different wording. These differences mark distinct textual states and do not authorize replacing Brown's translation by the later author's revisions. The next paragraph returns to a nonabelian Dold--Puppe comparison and supplies the concrete group-category model.

\subsection*{Brown's July translation: Group-categories and nonabelian complexes (sites 129--131)}
The nonabelian question asks what replaces the Dold--Puppe description of strict Picard categories when the additive structure is replaced by a group structure. Groups in 1-categories are described explicitly by a homomorphism $d:L_1\to L_0$ and an action of $L_0$ on $L_1$ satisfying equivariance and the Peiffer identity. The resulting cokernel, central kernel, induced action, and Postnikoff--Sinh class in degree three connect this description with crossed modules and with Brown's comments 3 and 7. The higher-dimensional extension is posed as a problem about noncommutative complexes equipped with additional operations; no general formula for those operations is claimed. The intended homotopy interpretation raises the truncation degree by one on passing from topological groups to connected spaces. Brown's copy omits the word ``pointed'' that appears in the revised witness, so the basepoint convention required by this comparison should be supplied in the reader's interpretation, not inserted into the translation. Site 130 has both closing parentheses in $\theta(d(x_1))$; the witness's misplaced parenthesis is a transcription error, while the outer parenthesis's typed-versus-retouched provenance remains uncertain. Brown also writes ``or commutative formal groups'', where the later version has ``and''. Site 131 then explicitly states that a commentary on biduality and homology is omitted. The corresponding discussion in Appendices 10--14 of the 1983 text is therefore absent by design from this earlier copy, and its later subtitle and discussion of the six operations are not interpolated here. The next surviving portion turns to schemes and stratification.

\subsection*{Brown's July translation: The homotopical triangle (sites 132--133)}
The translation distinguishes the rich algebra of a topos from the information retained by its homotopy type, then arranges spaces modulo n-homotopy, n-groupoids, and special n-categories in a triangle of mutual determination. Special, or n-galois, categories are described through locally constant stacks; exact fiber functors are proposed as the means of recovering the fundamental groupoid. This returns to the February reconstruction problem and supplies the locally constant case from which the stratified proposal departs. Brown's site 133 says local n-systems although its definition uses $(n-1)$-stacks; the revised witness's $(n-1)$-systems exposes the index mismatch. The assertion that the étale topos of a scheme is surely not locally homotopically trivial should be read as a warning about general schemes, not a universal exclusion: the étale topos of an algebraically closed field is the elementary point case. Pro-groupoids and ind-categories are proposed to retain the variation over finer coverings when an ordinary homotopy type is insufficient. The three diagram arrows are handwritten, whereas their vertices are typed. The later subtitle is absent, and the stronger or more cautious formulations of the revised witness do not replace Brown's wording. Only the isolated missing word space in examplethe is corrected in the corrected forms, with an explicit apparatus entry.

\subsection*{Brown's July translation: Constructible coefficients and tubular gluing (sites 134--138)}
The next step retains more than the ordinary homotopy type by fixing an equisingular stratification and studying coefficients locally constant on each stratum. Stability under the usual operations motivates the equisingularity condition. In the two-stratum example, sheaves are reconstructed from restrictions and a comparison through the punctured tubular neighborhood; a fibration from the boundary to the closed stratum and an inclusion into the complement provide a geometric description of the gluing functor. Equivalent spans replace the boundary by the punctured tube and then the closed stratum by its tubular neighborhood. This is the content corresponding to Appendix 16, and it supplies the local model for the diagrammatic treatment that follows. Brown's source has several consequential slips. Site 135 calls both complementary strata closed, points the second inclusion from X to $X_1$, prints an extra superscript star on $i_{1*}$, and says that local constancy is true of $X_1$ where the relevant object is the comparison functor applied to $F_1$. Its wording about stability ``not only for a finite partition ... but for all'' obscures the intended distinction between constructibility after refinement and stability for one fixed partition. Site 136 says ``sheaf'' where the revised version says ``scheme'', writes the bundle as $T/X$, and prints a closed interval in the punctured-cone formula, where deleting the collapsed endpoint requires a half-open interval. A neighborhood of $X_0$ in the disjoint open stratum must be interpreted as a punctured tubular object, not an ordinary inclusion of $X_0$ in $X_1$. The diagram's congruence signs are stronger than the homotopy comparisons established in the prose. All these literal features remain in the readers. The isolated is-for-in slip in ``inclusion of the boundary is a manifold'' on site 137 alone receives an apparatus correction. The later Cox note is absent from Brown's copy; his comment 10 supplies the separate historical update.

\subsection*{Brown's July translation: Diagrammatic and fine homotopy types (sites 138--141)}
The tubular span is generalized to a diagram of homotopy types with locally constant coefficients on its vertices and compatibility maps along its arrows. Tame equisingular spaces are proposed to admit a finite dévissage by manifolds with boundary, locally trivial fibrations, and boundary inclusions; adapted stratifications of a map should then express the direct and inverse image operations. Passage to finer stratifications is intended to retain all constructible information in the fine homotopy type. This corresponds to Appendix 17 and anticipates the main text's insistence on homotopy theories of diagrams. The distinction is essential: one needs to localize the category of diagrams, rather than simply replace each object by its image in the ordinary homotopy category and discard coherence. Brown's site 139 obscures that contrast with ``as based'' and the literal ``category of functions'', where the revised version speaks of fractions; the later note explicitly connecting the construction to derivators is absent. Site 138 says homeomorphisms for the compatibility maps of sheaves, although these are homomorphisms and need not be invertible. Site 140 mixes f and its restriction g in the displayed four operations and prints ``non-étale tubular neighbourhoods'', conflicting with the étale construction just proposed. These remain source-level cautions. The finiteness claim on site 141 is part of an explicitly unfinished program: prime-to-characteristic profinite completions of finite polyhedral or compact manifold models are envisaged, with nonliftable positive-characteristic varieties singled out as mysterious. Its stronger present-tense wording should not be read as a proof of the proposed theory. The next paragraph expressly says that a precise algebraic program has not yet been drawn up and turns to linearization.

\subsection*{Brown's July translation: Linearization and reconstruction (sites 141--143)}
The closing mathematical discussion asks how much cohomological information reconstructs a homotopy type: twisted coefficients, exact-sequence functoriality, products, and other cohomology operations lead to the candidate of a triangulated derived category with derived tensor product and Hom. Grothendieck expresses uncertainty about that abelian candidate and greater confidence in reconstruction from the full higher categories of stacks. The stratified version should similarly recover mixed homotopy information from coefficients subordinate to a fixed stratification and the operations along inclusions of unions of strata. This is the precursor to Appendix 18 and to the later demand for a richer notion than a bare triangulated category. Brown's text must be read with its limitations: site 142 says n-stacks of constructible sheaves where the ordinary reconstruction discussion requires a locally constant qualification; the following paragraph separately imposes local constancy on each stratum. The blank before product is empty and is not filled with the witness's cup. Most significantly, the later author's footnote reporting a negative answer to the proposed abelian recovery is absent from this earlier copy, so Brown's question remains open in its own textual state. Site 143 prints operations of fibrations and cofibrations and retains a full NB about Quillen's early-1970s approach; the revised 1983 version changes these to fibers and cofibers, cancels the NB, and adds the derivator discussion. These are distinct conceptual and chronological layers. The letter closes with the request for the Bloch--Quillen formula, not with a claimed solution to the reconstruction program.


\chapter*{Sources}
Grothendieck, A. (1983). \emph{Pursuing Stacks} [Typescript with handwritten revisions; cotes 134-2--134-8]. Grothendieck archive, Université de Montpellier. The supplied page images and their metadata are the primary witnesses for this edition.

Grothendieck, A. (2021). \emph{Pursuing Stacks} (M. Carmona, with the collaboration of U. Buchholtz, Eds.; arXiv:2111.01000, version 2). Extended from the 2015 edition by ``the scrivener''; CC0 1.0. This edition is consulted as a comparison witness, not reproduced as the source of record.

Brown, R. (Trans.). (n.d.). \emph{English translation of the mathematical part of two letters from Alexander Grothendieck to L. Breen, 1975} [Typescript; prepared with assistance from L. Breen and J-L. Loday]. Grothendieck archive, Université de Montpellier, cote 134-2, sites 115--143.
\end{document}
